\chapter{Additional Results}
\label{chap:additional_results}

This appendix reports the complete results underlying the analysis of Section~\ref{sec:experiments_clustering_quality_of_the_latent_space}.
For each of the six Atari environments studied in Chapter~\ref{chap:experiments}, it gives the evolution of \gls{cu}, \gls{as}, and \gls{acu} for every analyzed layer and for a number of clusters $K \in [2,20]$.
It also gives the three baseline similarities entering the \gls{as}.
These are the similarity to the clustering of the observations ($\Similarity{obs}{}$), to the latent clustering of an untrained network ($\Similarity{untr}{}$), and to the clustering of the action distribution, computed on the policy logits before the softmax ($\Similarity{act}{}$).
All these quantities are computed with the \gls{ari}, consistent with the rest of this thesis.
Images of the clusters obtained for every analyzed layer, which extend the qualitative analysis of Section~\ref{sec:experiments_semantic_interpretation_of_the_clusters}, have been archived on Zenodo \LinkClusterImages.

For each environment, one table is reported per metric and per baseline.
The rows correspond to the analyzed layers and the columns to the values of $K$.
The background shade of a cell is proportional to its value, so that the closer the value is to $1$, the darker the cell.
The tables of \gls{cu} are shown in green, those of \gls{as} and of its baselines in blue, and those of \gls{acu} in red.

A value in \textbf{bold} is the highest of its column, which identifies the best layer for this number of clusters.
An \underline{underlined} value is the highest of its row, which identifies the best number of clusters for this layer.
A value both \textbf{\underline{bold and underlined}} is the best of its row and of its column.
The last column reports the mean $\pm$ standard deviation across $K$ for each layer, and the last row the mean $\pm$ standard deviation across layers for each $K$.
In both, the best mean is shown in bold.
The bottom-right cell reports the mean $\pm$ standard deviation over the whole table.

The baseline tables must be read the other way around.
A high value, shown as a dark cell, indicates that the latent clustering is strongly aligned with the baseline.
It therefore signals an absence of abstraction and lowers the \gls{as}.

\clearpage
\begin{landscape}
% Tighter spacing around the tables (local to this environment) so that the
% six tables of an environment fit on two pages; see the height budget in
% setup/appendix_tables.tex (\ClusterKTableMaxHeight).
\setlength{\intextsep}{4pt plus 1pt minus 1pt}
\setlength{\abovecaptionskip}{3pt}
\setlength{\belowcaptionskip}{0pt}

\section{\BreakoutName}
\label{sec:appendix_additional_results_breakout}
\ClusterKTable{10}{\glsfmtshort{cu} --- \BreakoutName.}{cu_breakout}{
        1 & \SetCell{bg=\HeatCellBg{97}}\CellText{0.97}{0}{1} & \SetCell{bg=\HeatCellBg{97}}\CellText{0.97}{1}{1} & \SetCell{bg=\HeatCellBg{97}}\CellText{0.97}{1}{1} & \SetCell{bg=\HeatCellBg{62}}\CellText{0.62}{0}{0} & \SetCell{bg=\HeatCellBg{84}}\CellText{0.84}{1}{0} & \SetCell{bg=\HeatCellBg{81}}\CellText{0.81}{1}{0} & \SetCell{bg=\HeatCellBg{86}}\CellText{0.86}{1}{0} & \SetCell{bg=\HeatCellBg{84}}\CellText{0.84}{1}{0} & \SetCell{bg=\HeatCellBg{81}}\CellText{0.81}{1}{0} & \SetCell{bg=\HeatCellBg{85}}\CellText{0.85}{1}{0} & \SetCell{bg=\HeatCellBg{85}}\CellText{0.85}{1}{0} & \SetCell{bg=\HeatCellBg{87}}\CellText{0.87}{1}{0} & \SetCell{bg=\HeatCellBg{88}}\CellText{0.88}{1}{0} & \SetCell{bg=\HeatCellBg{84}}\CellText{0.84}{1}{0} & \SetCell{bg=\HeatCellBg{86}}\CellText{0.86}{1}{0} & \SetCell{bg=\HeatCellBg{85}}\CellText{0.85}{1}{0} & \SetCell{bg=\HeatCellBg{85}}\CellText{0.85}{1}{0} & \SetCell{bg=\HeatCellBg{78}}\CellText{0.78}{1}{0} & \SetCell{bg=\HeatCellBg{87}}\CellText{0.87}{1}{0} & \MeanCell{0.85}{0.07}{1} \\
        2 & \SetCell{bg=\HeatCellBg{40}}\CellText{0.40}{0}{0} & \SetCell{bg=\HeatCellBg{91}}\CellText{0.91}{0}{1} & \SetCell{bg=\HeatCellBg{84}}\CellText{0.84}{0}{0} & \SetCell{bg=\HeatCellBg{62}}\CellText{0.62}{0}{0} & \SetCell{bg=\HeatCellBg{79}}\CellText{0.79}{0}{0} & \SetCell{bg=\HeatCellBg{80}}\CellText{0.80}{0}{0} & \SetCell{bg=\HeatCellBg{76}}\CellText{0.76}{0}{0} & \SetCell{bg=\HeatCellBg{76}}\CellText{0.76}{0}{0} & \SetCell{bg=\HeatCellBg{73}}\CellText{0.73}{0}{0} & \SetCell{bg=\HeatCellBg{75}}\CellText{0.75}{0}{0} & \SetCell{bg=\HeatCellBg{67}}\CellText{0.67}{0}{0} & \SetCell{bg=\HeatCellBg{73}}\CellText{0.73}{0}{0} & \SetCell{bg=\HeatCellBg{74}}\CellText{0.74}{0}{0} & \SetCell{bg=\HeatCellBg{76}}\CellText{0.76}{0}{0} & \SetCell{bg=\HeatCellBg{72}}\CellText{0.72}{0}{0} & \SetCell{bg=\HeatCellBg{79}}\CellText{0.79}{0}{0} & \SetCell{bg=\HeatCellBg{66}}\CellText{0.66}{0}{0} & \SetCell{bg=\HeatCellBg{73}}\CellText{0.73}{0}{0} & \SetCell{bg=\HeatCellBg{73}}\CellText{0.73}{0}{0} & \MeanCell{0.73}{0.10}{0} \\
        3 & \SetCell{bg=\HeatCellBg{100}}\CellText{1.00}{1}{1} & \SetCell{bg=\HeatCellBg{85}}\CellText{0.85}{0}{0} & \SetCell{bg=\HeatCellBg{86}}\CellText{0.86}{0}{0} & \SetCell{bg=\HeatCellBg{61}}\CellText{0.61}{0}{0} & \SetCell{bg=\HeatCellBg{46}}\CellText{0.46}{0}{0} & \SetCell{bg=\HeatCellBg{55}}\CellText{0.55}{0}{0} & \SetCell{bg=\HeatCellBg{73}}\CellText{0.73}{0}{0} & \SetCell{bg=\HeatCellBg{52}}\CellText{0.52}{0}{0} & \SetCell{bg=\HeatCellBg{50}}\CellText{0.50}{0}{0} & \SetCell{bg=\HeatCellBg{63}}\CellText{0.63}{0}{0} & \SetCell{bg=\HeatCellBg{49}}\CellText{0.49}{0}{0} & \SetCell{bg=\HeatCellBg{57}}\CellText{0.57}{0}{0} & \SetCell{bg=\HeatCellBg{44}}\CellText{0.44}{0}{0} & \SetCell{bg=\HeatCellBg{38}}\CellText{0.38}{0}{0} & \SetCell{bg=\HeatCellBg{44}}\CellText{0.44}{0}{0} & \SetCell{bg=\HeatCellBg{40}}\CellText{0.40}{0}{0} & \SetCell{bg=\HeatCellBg{40}}\CellText{0.40}{0}{0} & \SetCell{bg=\HeatCellBg{50}}\CellText{0.50}{0}{0} & \SetCell{bg=\HeatCellBg{59}}\CellText{0.59}{0}{0} & \MeanCell{0.57}{0.17}{0} \\
        4 & \SetCell{bg=\HeatCellBg{99}}\CellText{0.99}{0}{1} & \SetCell{bg=\HeatCellBg{85}}\CellText{0.85}{0}{0} & \SetCell{bg=\HeatCellBg{26}}\CellText{0.26}{0}{0} & \SetCell{bg=\HeatCellBg{42}}\CellText{0.42}{0}{0} & \SetCell{bg=\HeatCellBg{38}}\CellText{0.38}{0}{0} & \SetCell{bg=\HeatCellBg{52}}\CellText{0.52}{0}{0} & \SetCell{bg=\HeatCellBg{38}}\CellText{0.38}{0}{0} & \SetCell{bg=\HeatCellBg{49}}\CellText{0.49}{0}{0} & \SetCell{bg=\HeatCellBg{50}}\CellText{0.50}{0}{0} & \SetCell{bg=\HeatCellBg{51}}\CellText{0.51}{0}{0} & \SetCell{bg=\HeatCellBg{53}}\CellText{0.53}{0}{0} & \SetCell{bg=\HeatCellBg{47}}\CellText{0.47}{0}{0} & \SetCell{bg=\HeatCellBg{39}}\CellText{0.39}{0}{0} & \SetCell{bg=\HeatCellBg{47}}\CellText{0.47}{0}{0} & \SetCell{bg=\HeatCellBg{46}}\CellText{0.46}{0}{0} & \SetCell{bg=\HeatCellBg{54}}\CellText{0.54}{0}{0} & \SetCell{bg=\HeatCellBg{47}}\CellText{0.47}{0}{0} & \SetCell{bg=\HeatCellBg{54}}\CellText{0.54}{0}{0} & \SetCell{bg=\HeatCellBg{41}}\CellText{0.41}{0}{0} & \MeanCell{0.50}{0.16}{0} \\
        5 & \SetCell{bg=\HeatCellBg{99}}\CellText{0.99}{0}{1} & \SetCell{bg=\HeatCellBg{85}}\CellText{0.85}{0}{0} & \SetCell{bg=\HeatCellBg{26}}\CellText{0.26}{0}{0} & \SetCell{bg=\HeatCellBg{42}}\CellText{0.42}{0}{0} & \SetCell{bg=\HeatCellBg{38}}\CellText{0.38}{0}{0} & \SetCell{bg=\HeatCellBg{52}}\CellText{0.52}{0}{0} & \SetCell{bg=\HeatCellBg{38}}\CellText{0.38}{0}{0} & \SetCell{bg=\HeatCellBg{49}}\CellText{0.49}{0}{0} & \SetCell{bg=\HeatCellBg{50}}\CellText{0.50}{0}{0} & \SetCell{bg=\HeatCellBg{51}}\CellText{0.51}{0}{0} & \SetCell{bg=\HeatCellBg{53}}\CellText{0.53}{0}{0} & \SetCell{bg=\HeatCellBg{47}}\CellText{0.47}{0}{0} & \SetCell{bg=\HeatCellBg{39}}\CellText{0.39}{0}{0} & \SetCell{bg=\HeatCellBg{47}}\CellText{0.47}{0}{0} & \SetCell{bg=\HeatCellBg{46}}\CellText{0.46}{0}{0} & \SetCell{bg=\HeatCellBg{54}}\CellText{0.54}{0}{0} & \SetCell{bg=\HeatCellBg{47}}\CellText{0.47}{0}{0} & \SetCell{bg=\HeatCellBg{54}}\CellText{0.54}{0}{0} & \SetCell{bg=\HeatCellBg{41}}\CellText{0.41}{0}{0} & \MeanCell{0.50}{0.16}{0} \\
        6b & \SetCell{bg=\HeatCellBg{66}}\CellText{0.66}{0}{0} & \SetCell{bg=\HeatCellBg{57}}\CellText{0.57}{0}{0} & \SetCell{bg=\HeatCellBg{65}}\CellText{0.65}{0}{0} & \SetCell{bg=\HeatCellBg{79}}\CellText{0.79}{1}{1} & \SetCell{bg=\HeatCellBg{59}}\CellText{0.59}{0}{0} & \SetCell{bg=\HeatCellBg{62}}\CellText{0.62}{0}{0} & \SetCell{bg=\HeatCellBg{56}}\CellText{0.56}{0}{0} & \SetCell{bg=\HeatCellBg{56}}\CellText{0.56}{0}{0} & \SetCell{bg=\HeatCellBg{58}}\CellText{0.58}{0}{0} & \SetCell{bg=\HeatCellBg{64}}\CellText{0.64}{0}{0} & \SetCell{bg=\HeatCellBg{59}}\CellText{0.59}{0}{0} & \SetCell{bg=\HeatCellBg{55}}\CellText{0.55}{0}{0} & \SetCell{bg=\HeatCellBg{55}}\CellText{0.55}{0}{0} & \SetCell{bg=\HeatCellBg{54}}\CellText{0.54}{0}{0} & \SetCell{bg=\HeatCellBg{59}}\CellText{0.59}{0}{0} & \SetCell{bg=\HeatCellBg{60}}\CellText{0.60}{0}{0} & \SetCell{bg=\HeatCellBg{53}}\CellText{0.53}{0}{0} & \SetCell{bg=\HeatCellBg{55}}\CellText{0.55}{0}{0} & \SetCell{bg=\HeatCellBg{55}}\CellText{0.55}{0}{0} & \MeanCell{0.59}{0.06}{0} \\
        7b & \SetCell{bg=\HeatCellBg{78}}\CellText{0.78}{0}{1} & \SetCell{bg=\HeatCellBg{41}}\CellText{0.41}{0}{0} & \SetCell{bg=\HeatCellBg{46}}\CellText{0.46}{0}{0} & \SetCell{bg=\HeatCellBg{57}}\CellText{0.57}{0}{0} & \SetCell{bg=\HeatCellBg{52}}\CellText{0.52}{0}{0} & \SetCell{bg=\HeatCellBg{56}}\CellText{0.56}{0}{0} & \SetCell{bg=\HeatCellBg{55}}\CellText{0.55}{0}{0} & \SetCell{bg=\HeatCellBg{41}}\CellText{0.41}{0}{0} & \SetCell{bg=\HeatCellBg{47}}\CellText{0.47}{0}{0} & \SetCell{bg=\HeatCellBg{47}}\CellText{0.47}{0}{0} & \SetCell{bg=\HeatCellBg{47}}\CellText{0.47}{0}{0} & \SetCell{bg=\HeatCellBg{43}}\CellText{0.43}{0}{0} & \SetCell{bg=\HeatCellBg{45}}\CellText{0.45}{0}{0} & \SetCell{bg=\HeatCellBg{47}}\CellText{0.47}{0}{0} & \SetCell{bg=\HeatCellBg{50}}\CellText{0.50}{0}{0} & \SetCell{bg=\HeatCellBg{52}}\CellText{0.52}{0}{0} & \SetCell{bg=\HeatCellBg{48}}\CellText{0.48}{0}{0} & \SetCell{bg=\HeatCellBg{50}}\CellText{0.50}{0}{0} & \SetCell{bg=\HeatCellBg{46}}\CellText{0.46}{0}{0} & \MeanCell{0.50}{0.08}{0} \\
        8b & \SetCell{bg=\HeatCellBg{68}}\CellText{0.68}{0}{1} & \SetCell{bg=\HeatCellBg{50}}\CellText{0.50}{0}{0} & \SetCell{bg=\HeatCellBg{50}}\CellText{0.50}{0}{0} & \SetCell{bg=\HeatCellBg{50}}\CellText{0.50}{0}{0} & \SetCell{bg=\HeatCellBg{50}}\CellText{0.50}{0}{0} & \SetCell{bg=\HeatCellBg{52}}\CellText{0.52}{0}{0} & \SetCell{bg=\HeatCellBg{50}}\CellText{0.50}{0}{0} & \SetCell{bg=\HeatCellBg{44}}\CellText{0.44}{0}{0} & \SetCell{bg=\HeatCellBg{55}}\CellText{0.55}{0}{0} & \SetCell{bg=\HeatCellBg{54}}\CellText{0.54}{0}{0} & \SetCell{bg=\HeatCellBg{55}}\CellText{0.55}{0}{0} & \SetCell{bg=\HeatCellBg{48}}\CellText{0.48}{0}{0} & \SetCell{bg=\HeatCellBg{48}}\CellText{0.48}{0}{0} & \SetCell{bg=\HeatCellBg{49}}\CellText{0.49}{0}{0} & \SetCell{bg=\HeatCellBg{51}}\CellText{0.51}{0}{0} & \SetCell{bg=\HeatCellBg{50}}\CellText{0.50}{0}{0} & \SetCell{bg=\HeatCellBg{51}}\CellText{0.51}{0}{0} & \SetCell{bg=\HeatCellBg{50}}\CellText{0.50}{0}{0} & \SetCell{bg=\HeatCellBg{45}}\CellText{0.45}{0}{0} & \MeanCell{0.51}{0.05}{0} \\
        Mean $\pm$ Std & \MeanCell{0.81}{0.20}{1} & \MeanCell{0.74}{0.20}{0} & \MeanCell{0.60}{0.26}{0} & \MeanCell{0.57}{0.11}{0} & \MeanCell{0.56}{0.16}{0} & \MeanCell{0.61}{0.12}{0} & \MeanCell{0.59}{0.17}{0} & \MeanCell{0.56}{0.14}{0} & \MeanCell{0.58}{0.12}{0} & \MeanCell{0.61}{0.12}{0} & \MeanCell{0.58}{0.12}{0} & \MeanCell{0.57}{0.14}{0} & \MeanCell{0.54}{0.17}{0} & \MeanCell{0.55}{0.15}{0} & \MeanCell{0.57}{0.14}{0} & \MeanCell{0.59}{0.14}{0} & \MeanCell{0.55}{0.13}{0} & \MeanCell{0.58}{0.10}{0} & \MeanCell{0.56}{0.16}{0} & \MeanCell{0.60}{0.17}{0} \\
}

\ClusterKTable{10}{\glsfmtshort{as} --- \BreakoutName.}{as_breakout}{
        1 & \SetCell{bg=\HeatCellBg{5}}\CellText{0.05}{0}{0} & \SetCell{bg=\HeatCellBg{6}}\CellText{0.06}{0}{0} & \SetCell{bg=\HeatCellBg{3}}\CellText{0.03}{0}{0} & \SetCell{bg=\HeatCellBg{16}}\CellText{0.16}{0}{0} & \SetCell{bg=\HeatCellBg{16}}\CellText{0.16}{0}{0} & \SetCell{bg=\HeatCellBg{16}}\CellText{0.16}{0}{0} & \SetCell{bg=\HeatCellBg{16}}\CellText{0.16}{0}{0} & \SetCell{bg=\HeatCellBg{18}}\CellText{0.18}{0}{0} & \SetCell{bg=\HeatCellBg{15}}\CellText{0.15}{0}{0} & \SetCell{bg=\HeatCellBg{15}}\CellText{0.15}{0}{0} & \SetCell{bg=\HeatCellBg{18}}\CellText{0.18}{0}{0} & \SetCell{bg=\HeatCellBg{14}}\CellText{0.14}{0}{0} & \SetCell{bg=\HeatCellBg{13}}\CellText{0.13}{0}{0} & \SetCell{bg=\HeatCellBg{20}}\CellText{0.20}{0}{0} & \SetCell{bg=\HeatCellBg{16}}\CellText{0.16}{0}{0} & \SetCell{bg=\HeatCellBg{16}}\CellText{0.16}{0}{0} & \SetCell{bg=\HeatCellBg{20}}\CellText{0.20}{0}{0} & \SetCell{bg=\HeatCellBg{23}}\CellText{0.23}{0}{1} & \SetCell{bg=\HeatCellBg{15}}\CellText{0.15}{0}{0} & \MeanCell{0.15}{0.05}{0} \\
        2 & \SetCell{bg=\HeatCellBg{60}}\CellText{0.60}{0}{1} & \SetCell{bg=\HeatCellBg{9}}\CellText{0.09}{0}{0} & \SetCell{bg=\HeatCellBg{20}}\CellText{0.20}{0}{0} & \SetCell{bg=\HeatCellBg{29}}\CellText{0.29}{0}{0} & \SetCell{bg=\HeatCellBg{27}}\CellText{0.27}{0}{0} & \SetCell{bg=\HeatCellBg{28}}\CellText{0.28}{0}{0} & \SetCell{bg=\HeatCellBg{24}}\CellText{0.24}{0}{0} & \SetCell{bg=\HeatCellBg{23}}\CellText{0.23}{0}{0} & \SetCell{bg=\HeatCellBg{25}}\CellText{0.25}{0}{0} & \SetCell{bg=\HeatCellBg{28}}\CellText{0.28}{0}{0} & \SetCell{bg=\HeatCellBg{31}}\CellText{0.31}{0}{0} & \SetCell{bg=\HeatCellBg{33}}\CellText{0.33}{0}{0} & \SetCell{bg=\HeatCellBg{28}}\CellText{0.28}{0}{0} & \SetCell{bg=\HeatCellBg{28}}\CellText{0.28}{0}{0} & \SetCell{bg=\HeatCellBg{29}}\CellText{0.29}{0}{0} & \SetCell{bg=\HeatCellBg{29}}\CellText{0.29}{0}{0} & \SetCell{bg=\HeatCellBg{31}}\CellText{0.31}{0}{0} & \SetCell{bg=\HeatCellBg{27}}\CellText{0.27}{0}{0} & \SetCell{bg=\HeatCellBg{27}}\CellText{0.27}{0}{0} & \MeanCell{0.28}{0.09}{0} \\
        3 & \SetCell{bg=\HeatCellBg{85}}\CellText{0.85}{0}{1} & \SetCell{bg=\HeatCellBg{14}}\CellText{0.14}{0}{0} & \SetCell{bg=\HeatCellBg{20}}\CellText{0.20}{0}{0} & \SetCell{bg=\HeatCellBg{36}}\CellText{0.36}{0}{0} & \SetCell{bg=\HeatCellBg{43}}\CellText{0.43}{0}{0} & \SetCell{bg=\HeatCellBg{41}}\CellText{0.41}{0}{0} & \SetCell{bg=\HeatCellBg{29}}\CellText{0.29}{0}{0} & \SetCell{bg=\HeatCellBg{41}}\CellText{0.41}{0}{0} & \SetCell{bg=\HeatCellBg{41}}\CellText{0.41}{0}{0} & \SetCell{bg=\HeatCellBg{44}}\CellText{0.44}{0}{0} & \SetCell{bg=\HeatCellBg{47}}\CellText{0.47}{0}{0} & \SetCell{bg=\HeatCellBg{43}}\CellText{0.43}{0}{0} & \SetCell{bg=\HeatCellBg{50}}\CellText{0.50}{0}{0} & \SetCell{bg=\HeatCellBg{50}}\CellText{0.50}{0}{0} & \SetCell{bg=\HeatCellBg{50}}\CellText{0.50}{0}{0} & \SetCell{bg=\HeatCellBg{47}}\CellText{0.47}{0}{0} & \SetCell{bg=\HeatCellBg{47}}\CellText{0.47}{0}{0} & \SetCell{bg=\HeatCellBg{46}}\CellText{0.46}{0}{0} & \SetCell{bg=\HeatCellBg{57}}\CellText{0.57}{0}{0} & \MeanCell{0.44}{0.14}{0} \\
        4 & \SetCell{bg=\HeatCellBg{88}}\CellText{0.88}{1}{1} & \SetCell{bg=\HeatCellBg{28}}\CellText{0.28}{0}{0} & \SetCell{bg=\HeatCellBg{64}}\CellText{0.64}{0}{0} & \SetCell{bg=\HeatCellBg{72}}\CellText{0.72}{0}{0} & \SetCell{bg=\HeatCellBg{75}}\CellText{0.75}{0}{0} & \SetCell{bg=\HeatCellBg{80}}\CellText{0.80}{0}{0} & \SetCell{bg=\HeatCellBg{77}}\CellText{0.77}{0}{0} & \SetCell{bg=\HeatCellBg{81}}\CellText{0.81}{0}{0} & \SetCell{bg=\HeatCellBg{80}}\CellText{0.80}{0}{0} & \SetCell{bg=\HeatCellBg{81}}\CellText{0.81}{0}{0} & \SetCell{bg=\HeatCellBg{81}}\CellText{0.81}{0}{0} & \SetCell{bg=\HeatCellBg{82}}\CellText{0.82}{0}{0} & \SetCell{bg=\HeatCellBg{80}}\CellText{0.80}{0}{0} & \SetCell{bg=\HeatCellBg{81}}\CellText{0.81}{0}{0} & \SetCell{bg=\HeatCellBg{80}}\CellText{0.80}{0}{0} & \SetCell{bg=\HeatCellBg{79}}\CellText{0.79}{0}{0} & \SetCell{bg=\HeatCellBg{79}}\CellText{0.79}{0}{0} & \SetCell{bg=\HeatCellBg{82}}\CellText{0.82}{0}{0} & \SetCell{bg=\HeatCellBg{80}}\CellText{0.80}{0}{0} & \MeanCell{0.76}{0.12}{0} \\
        5 & \SetCell{bg=\HeatCellBg{88}}\CellText{0.88}{1}{1} & \SetCell{bg=\HeatCellBg{28}}\CellText{0.28}{0}{0} & \SetCell{bg=\HeatCellBg{64}}\CellText{0.64}{0}{0} & \SetCell{bg=\HeatCellBg{72}}\CellText{0.72}{0}{0} & \SetCell{bg=\HeatCellBg{75}}\CellText{0.75}{0}{0} & \SetCell{bg=\HeatCellBg{80}}\CellText{0.80}{0}{0} & \SetCell{bg=\HeatCellBg{77}}\CellText{0.77}{0}{0} & \SetCell{bg=\HeatCellBg{81}}\CellText{0.81}{0}{0} & \SetCell{bg=\HeatCellBg{80}}\CellText{0.80}{0}{0} & \SetCell{bg=\HeatCellBg{81}}\CellText{0.81}{0}{0} & \SetCell{bg=\HeatCellBg{81}}\CellText{0.81}{0}{0} & \SetCell{bg=\HeatCellBg{82}}\CellText{0.82}{0}{0} & \SetCell{bg=\HeatCellBg{80}}\CellText{0.80}{0}{0} & \SetCell{bg=\HeatCellBg{81}}\CellText{0.81}{0}{0} & \SetCell{bg=\HeatCellBg{80}}\CellText{0.80}{0}{0} & \SetCell{bg=\HeatCellBg{79}}\CellText{0.79}{0}{0} & \SetCell{bg=\HeatCellBg{79}}\CellText{0.79}{0}{0} & \SetCell{bg=\HeatCellBg{82}}\CellText{0.82}{0}{0} & \SetCell{bg=\HeatCellBg{80}}\CellText{0.80}{0}{0} & \MeanCell{0.76}{0.12}{0} \\
        6b & \SetCell{bg=\HeatCellBg{62}}\CellText{0.62}{0}{0} & \SetCell{bg=\HeatCellBg{83}}\CellText{0.83}{1}{0} & \SetCell{bg=\HeatCellBg{74}}\CellText{0.74}{1}{0} & \SetCell{bg=\HeatCellBg{82}}\CellText{0.82}{1}{0} & \SetCell{bg=\HeatCellBg{85}}\CellText{0.85}{1}{1} & \SetCell{bg=\HeatCellBg{85}}\CellText{0.85}{1}{1} & \SetCell{bg=\HeatCellBg{85}}\CellText{0.85}{1}{1} & \SetCell{bg=\HeatCellBg{85}}\CellText{0.85}{1}{1} & \SetCell{bg=\HeatCellBg{83}}\CellText{0.83}{1}{0} & \SetCell{bg=\HeatCellBg{82}}\CellText{0.82}{1}{0} & \SetCell{bg=\HeatCellBg{85}}\CellText{0.85}{1}{1} & \SetCell{bg=\HeatCellBg{85}}\CellText{0.85}{1}{1} & \SetCell{bg=\HeatCellBg{84}}\CellText{0.84}{1}{0} & \SetCell{bg=\HeatCellBg{83}}\CellText{0.83}{1}{0} & \SetCell{bg=\HeatCellBg{83}}\CellText{0.83}{1}{0} & \SetCell{bg=\HeatCellBg{83}}\CellText{0.83}{1}{0} & \SetCell{bg=\HeatCellBg{84}}\CellText{0.84}{1}{0} & \SetCell{bg=\HeatCellBg{85}}\CellText{0.85}{1}{1} & \SetCell{bg=\HeatCellBg{83}}\CellText{0.83}{1}{0} & \MeanCell{0.82}{0.05}{1} \\
        7b & \SetCell{bg=\HeatCellBg{64}}\CellText{0.64}{0}{0} & \SetCell{bg=\HeatCellBg{75}}\CellText{0.75}{0}{0} & \SetCell{bg=\HeatCellBg{69}}\CellText{0.69}{0}{0} & \SetCell{bg=\HeatCellBg{67}}\CellText{0.67}{0}{0} & \SetCell{bg=\HeatCellBg{74}}\CellText{0.74}{0}{0} & \SetCell{bg=\HeatCellBg{79}}\CellText{0.79}{0}{0} & \SetCell{bg=\HeatCellBg{79}}\CellText{0.79}{0}{0} & \SetCell{bg=\HeatCellBg{80}}\CellText{0.80}{0}{0} & \SetCell{bg=\HeatCellBg{80}}\CellText{0.80}{0}{0} & \SetCell{bg=\HeatCellBg{82}}\CellText{0.82}{1}{0} & \SetCell{bg=\HeatCellBg{82}}\CellText{0.82}{0}{0} & \SetCell{bg=\HeatCellBg{83}}\CellText{0.83}{0}{0} & \SetCell{bg=\HeatCellBg{83}}\CellText{0.83}{0}{0} & \SetCell{bg=\HeatCellBg{82}}\CellText{0.82}{0}{0} & \SetCell{bg=\HeatCellBg{83}}\CellText{0.83}{1}{0} & \SetCell{bg=\HeatCellBg{83}}\CellText{0.83}{1}{0} & \SetCell{bg=\HeatCellBg{83}}\CellText{0.83}{0}{0} & \SetCell{bg=\HeatCellBg{85}}\CellText{0.85}{1}{1} & \SetCell{bg=\HeatCellBg{83}}\CellText{0.83}{1}{0} & \MeanCell{0.79}{0.06}{0} \\
        8b & \SetCell{bg=\HeatCellBg{71}}\CellText{0.71}{0}{0} & \SetCell{bg=\HeatCellBg{67}}\CellText{0.67}{0}{0} & \SetCell{bg=\HeatCellBg{66}}\CellText{0.66}{0}{0} & \SetCell{bg=\HeatCellBg{64}}\CellText{0.64}{0}{0} & \SetCell{bg=\HeatCellBg{78}}\CellText{0.78}{0}{0} & \SetCell{bg=\HeatCellBg{77}}\CellText{0.77}{0}{0} & \SetCell{bg=\HeatCellBg{75}}\CellText{0.75}{0}{0} & \SetCell{bg=\HeatCellBg{80}}\CellText{0.80}{0}{1} & \SetCell{bg=\HeatCellBg{76}}\CellText{0.76}{0}{0} & \SetCell{bg=\HeatCellBg{78}}\CellText{0.78}{0}{0} & \SetCell{bg=\HeatCellBg{77}}\CellText{0.77}{0}{0} & \SetCell{bg=\HeatCellBg{79}}\CellText{0.79}{0}{0} & \SetCell{bg=\HeatCellBg{77}}\CellText{0.77}{0}{0} & \SetCell{bg=\HeatCellBg{78}}\CellText{0.78}{0}{0} & \SetCell{bg=\HeatCellBg{78}}\CellText{0.78}{0}{0} & \SetCell{bg=\HeatCellBg{78}}\CellText{0.78}{0}{0} & \SetCell{bg=\HeatCellBg{78}}\CellText{0.78}{0}{0} & \SetCell{bg=\HeatCellBg{80}}\CellText{0.80}{0}{1} & \SetCell{bg=\HeatCellBg{78}}\CellText{0.78}{0}{0} & \MeanCell{0.76}{0.05}{0} \\
        Mean $\pm$ Std & \MeanCell{0.65}{0.25}{1} & \MeanCell{0.39}{0.29}{0} & \MeanCell{0.47}{0.26}{0} & \MeanCell{0.55}{0.23}{0} & \MeanCell{0.59}{0.25}{0} & \MeanCell{0.61}{0.26}{0} & \MeanCell{0.58}{0.27}{0} & \MeanCell{0.61}{0.27}{0} & \MeanCell{0.60}{0.26}{0} & \MeanCell{0.61}{0.26}{0} & \MeanCell{0.63}{0.25}{0} & \MeanCell{0.63}{0.26}{0} & \MeanCell{0.62}{0.26}{0} & \MeanCell{0.63}{0.25}{0} & \MeanCell{0.62}{0.25}{0} & \MeanCell{0.62}{0.25}{0} & \MeanCell{0.63}{0.24}{0} & \MeanCell{0.64}{0.25}{0} & \MeanCell{0.63}{0.26}{0} & \MeanCell{0.59}{0.27}{0} \\
}

\ClusterKTable{10}{\glsfmtshort{acu} --- \BreakoutName.}{acu_breakout}{
        1 & \SetCell{bg=\HeatCellBg{1}}\CellText{0.01}{0}{0} & \SetCell{bg=\HeatCellBg{3}}\CellText{0.03}{0}{0} & \SetCell{bg=\HeatCellBg{0}}\CellText{0.00}{0}{0} & \SetCell{bg=\HeatCellBg{-22}}\CellText{-0.22}{0}{0} & \SetCell{bg=\HeatCellBg{-1}}\CellText{-0.01}{0}{0} & \SetCell{bg=\HeatCellBg{-3}}\CellText{-0.03}{0}{0} & \SetCell{bg=\HeatCellBg{3}}\CellText{0.03}{0}{0} & \SetCell{bg=\HeatCellBg{1}}\CellText{0.01}{0}{0} & \SetCell{bg=\HeatCellBg{-4}}\CellText{-0.04}{0}{0} & \SetCell{bg=\HeatCellBg{0}}\CellText{0.00}{0}{0} & \SetCell{bg=\HeatCellBg{2}}\CellText{0.02}{0}{0} & \SetCell{bg=\HeatCellBg{1}}\CellText{0.01}{0}{0} & \SetCell{bg=\HeatCellBg{2}}\CellText{0.02}{0}{0} & \SetCell{bg=\HeatCellBg{4}}\CellText{0.04}{0}{0} & \SetCell{bg=\HeatCellBg{2}}\CellText{0.02}{0}{0} & \SetCell{bg=\HeatCellBg{1}}\CellText{0.01}{0}{0} & \SetCell{bg=\HeatCellBg{5}}\CellText{0.05}{0}{1} & \SetCell{bg=\HeatCellBg{1}}\CellText{0.01}{0}{0} & \SetCell{bg=\HeatCellBg{2}}\CellText{0.02}{0}{0} & \MeanCell{0.00}{0.06}{0} \\
        2 & \SetCell{bg=\HeatCellBg{1}}\CellText{0.01}{0}{0} & \SetCell{bg=\HeatCellBg{0}}\CellText{0.00}{0}{0} & \SetCell{bg=\HeatCellBg{4}}\CellText{0.04}{0}{0} & \SetCell{bg=\HeatCellBg{-9}}\CellText{-0.09}{0}{0} & \SetCell{bg=\HeatCellBg{5}}\CellText{0.05}{0}{0} & \SetCell{bg=\HeatCellBg{8}}\CellText{0.08}{0}{1} & \SetCell{bg=\HeatCellBg{0}}\CellText{0.00}{0}{0} & \SetCell{bg=\HeatCellBg{-1}}\CellText{-0.01}{0}{0} & \SetCell{bg=\HeatCellBg{-2}}\CellText{-0.02}{0}{0} & \SetCell{bg=\HeatCellBg{3}}\CellText{0.03}{0}{0} & \SetCell{bg=\HeatCellBg{-2}}\CellText{-0.02}{0}{0} & \SetCell{bg=\HeatCellBg{6}}\CellText{0.06}{0}{0} & \SetCell{bg=\HeatCellBg{2}}\CellText{0.02}{0}{0} & \SetCell{bg=\HeatCellBg{3}}\CellText{0.03}{0}{0} & \SetCell{bg=\HeatCellBg{1}}\CellText{0.01}{0}{0} & \SetCell{bg=\HeatCellBg{8}}\CellText{0.08}{0}{1} & \SetCell{bg=\HeatCellBg{-3}}\CellText{-0.03}{0}{0} & \SetCell{bg=\HeatCellBg{0}}\CellText{0.00}{0}{0} & \SetCell{bg=\HeatCellBg{0}}\CellText{0.00}{0}{0} & \MeanCell{0.01}{0.04}{0} \\
        3 & \SetCell{bg=\HeatCellBg{85}}\CellText{0.85}{0}{1} & \SetCell{bg=\HeatCellBg{-2}}\CellText{-0.02}{0}{0} & \SetCell{bg=\HeatCellBg{6}}\CellText{0.06}{0}{0} & \SetCell{bg=\HeatCellBg{-3}}\CellText{-0.03}{0}{0} & \SetCell{bg=\HeatCellBg{-10}}\CellText{-0.10}{0}{0} & \SetCell{bg=\HeatCellBg{-5}}\CellText{-0.05}{0}{0} & \SetCell{bg=\HeatCellBg{2}}\CellText{0.02}{0}{0} & \SetCell{bg=\HeatCellBg{-7}}\CellText{-0.07}{0}{0} & \SetCell{bg=\HeatCellBg{-9}}\CellText{-0.09}{0}{0} & \SetCell{bg=\HeatCellBg{7}}\CellText{0.07}{0}{0} & \SetCell{bg=\HeatCellBg{-4}}\CellText{-0.04}{0}{0} & \SetCell{bg=\HeatCellBg{0}}\CellText{0.00}{0}{0} & \SetCell{bg=\HeatCellBg{-6}}\CellText{-0.06}{0}{0} & \SetCell{bg=\HeatCellBg{-12}}\CellText{-0.12}{0}{0} & \SetCell{bg=\HeatCellBg{-5}}\CellText{-0.05}{0}{0} & \SetCell{bg=\HeatCellBg{-13}}\CellText{-0.13}{0}{0} & \SetCell{bg=\HeatCellBg{-12}}\CellText{-0.12}{0}{0} & \SetCell{bg=\HeatCellBg{-5}}\CellText{-0.05}{0}{0} & \SetCell{bg=\HeatCellBg{16}}\CellText{0.16}{0}{0} & \MeanCell{0.01}{0.21}{0} \\
        4 & \SetCell{bg=\HeatCellBg{87}}\CellText{0.87}{1}{1} & \SetCell{bg=\HeatCellBg{13}}\CellText{0.13}{0}{0} & \SetCell{bg=\HeatCellBg{-10}}\CellText{-0.10}{0}{0} & \SetCell{bg=\HeatCellBg{14}}\CellText{0.14}{0}{0} & \SetCell{bg=\HeatCellBg{13}}\CellText{0.13}{0}{0} & \SetCell{bg=\HeatCellBg{31}}\CellText{0.31}{0}{0} & \SetCell{bg=\HeatCellBg{15}}\CellText{0.15}{0}{0} & \SetCell{bg=\HeatCellBg{30}}\CellText{0.30}{0}{0} & \SetCell{bg=\HeatCellBg{30}}\CellText{0.30}{0}{0} & \SetCell{bg=\HeatCellBg{31}}\CellText{0.31}{0}{0} & \SetCell{bg=\HeatCellBg{34}}\CellText{0.34}{0}{0} & \SetCell{bg=\HeatCellBg{28}}\CellText{0.28}{0}{0} & \SetCell{bg=\HeatCellBg{19}}\CellText{0.19}{0}{0} & \SetCell{bg=\HeatCellBg{28}}\CellText{0.28}{0}{0} & \SetCell{bg=\HeatCellBg{26}}\CellText{0.26}{0}{0} & \SetCell{bg=\HeatCellBg{33}}\CellText{0.33}{0}{0} & \SetCell{bg=\HeatCellBg{26}}\CellText{0.26}{0}{0} & \SetCell{bg=\HeatCellBg{36}}\CellText{0.36}{0}{0} & \SetCell{bg=\HeatCellBg{21}}\CellText{0.21}{0}{0} & \MeanCell{0.27}{0.18}{0} \\
        5 & \SetCell{bg=\HeatCellBg{87}}\CellText{0.87}{1}{1} & \SetCell{bg=\HeatCellBg{13}}\CellText{0.13}{0}{0} & \SetCell{bg=\HeatCellBg{-10}}\CellText{-0.10}{0}{0} & \SetCell{bg=\HeatCellBg{14}}\CellText{0.14}{0}{0} & \SetCell{bg=\HeatCellBg{13}}\CellText{0.13}{0}{0} & \SetCell{bg=\HeatCellBg{31}}\CellText{0.31}{0}{0} & \SetCell{bg=\HeatCellBg{15}}\CellText{0.15}{0}{0} & \SetCell{bg=\HeatCellBg{30}}\CellText{0.30}{0}{0} & \SetCell{bg=\HeatCellBg{30}}\CellText{0.30}{0}{0} & \SetCell{bg=\HeatCellBg{31}}\CellText{0.31}{0}{0} & \SetCell{bg=\HeatCellBg{34}}\CellText{0.34}{0}{0} & \SetCell{bg=\HeatCellBg{28}}\CellText{0.28}{0}{0} & \SetCell{bg=\HeatCellBg{19}}\CellText{0.19}{0}{0} & \SetCell{bg=\HeatCellBg{28}}\CellText{0.28}{0}{0} & \SetCell{bg=\HeatCellBg{26}}\CellText{0.26}{0}{0} & \SetCell{bg=\HeatCellBg{33}}\CellText{0.33}{0}{0} & \SetCell{bg=\HeatCellBg{26}}\CellText{0.26}{0}{0} & \SetCell{bg=\HeatCellBg{36}}\CellText{0.36}{0}{0} & \SetCell{bg=\HeatCellBg{21}}\CellText{0.21}{0}{0} & \MeanCell{0.27}{0.18}{0} \\
        6b & \SetCell{bg=\HeatCellBg{28}}\CellText{0.28}{0}{0} & \SetCell{bg=\HeatCellBg{40}}\CellText{0.40}{1}{0} & \SetCell{bg=\HeatCellBg{39}}\CellText{0.39}{1}{0} & \SetCell{bg=\HeatCellBg{61}}\CellText{0.61}{1}{1} & \SetCell{bg=\HeatCellBg{43}}\CellText{0.43}{1}{0} & \SetCell{bg=\HeatCellBg{47}}\CellText{0.47}{1}{0} & \SetCell{bg=\HeatCellBg{41}}\CellText{0.41}{1}{0} & \SetCell{bg=\HeatCellBg{41}}\CellText{0.41}{1}{0} & \SetCell{bg=\HeatCellBg{42}}\CellText{0.42}{1}{0} & \SetCell{bg=\HeatCellBg{46}}\CellText{0.46}{1}{0} & \SetCell{bg=\HeatCellBg{45}}\CellText{0.45}{1}{0} & \SetCell{bg=\HeatCellBg{40}}\CellText{0.40}{1}{0} & \SetCell{bg=\HeatCellBg{39}}\CellText{0.39}{1}{0} & \SetCell{bg=\HeatCellBg{37}}\CellText{0.37}{1}{0} & \SetCell{bg=\HeatCellBg{42}}\CellText{0.42}{1}{0} & \SetCell{bg=\HeatCellBg{43}}\CellText{0.43}{1}{0} & \SetCell{bg=\HeatCellBg{36}}\CellText{0.36}{1}{0} & \SetCell{bg=\HeatCellBg{40}}\CellText{0.40}{1}{0} & \SetCell{bg=\HeatCellBg{38}}\CellText{0.38}{1}{0} & \MeanCell{0.41}{0.06}{1} \\
        7b & \SetCell{bg=\HeatCellBg{42}}\CellText{0.42}{0}{1} & \SetCell{bg=\HeatCellBg{16}}\CellText{0.16}{0}{0} & \SetCell{bg=\HeatCellBg{16}}\CellText{0.16}{0}{0} & \SetCell{bg=\HeatCellBg{24}}\CellText{0.24}{0}{0} & \SetCell{bg=\HeatCellBg{26}}\CellText{0.26}{0}{0} & \SetCell{bg=\HeatCellBg{35}}\CellText{0.35}{0}{0} & \SetCell{bg=\HeatCellBg{35}}\CellText{0.35}{0}{0} & \SetCell{bg=\HeatCellBg{21}}\CellText{0.21}{0}{0} & \SetCell{bg=\HeatCellBg{27}}\CellText{0.27}{0}{0} & \SetCell{bg=\HeatCellBg{29}}\CellText{0.29}{0}{0} & \SetCell{bg=\HeatCellBg{29}}\CellText{0.29}{0}{0} & \SetCell{bg=\HeatCellBg{26}}\CellText{0.26}{0}{0} & \SetCell{bg=\HeatCellBg{28}}\CellText{0.28}{0}{0} & \SetCell{bg=\HeatCellBg{30}}\CellText{0.30}{0}{0} & \SetCell{bg=\HeatCellBg{33}}\CellText{0.33}{0}{0} & \SetCell{bg=\HeatCellBg{35}}\CellText{0.35}{0}{0} & \SetCell{bg=\HeatCellBg{31}}\CellText{0.31}{0}{0} & \SetCell{bg=\HeatCellBg{34}}\CellText{0.34}{0}{0} & \SetCell{bg=\HeatCellBg{29}}\CellText{0.29}{0}{0} & \MeanCell{0.29}{0.06}{0} \\
        8b & \SetCell{bg=\HeatCellBg{40}}\CellText{0.40}{0}{1} & \SetCell{bg=\HeatCellBg{17}}\CellText{0.17}{0}{0} & \SetCell{bg=\HeatCellBg{16}}\CellText{0.16}{0}{0} & \SetCell{bg=\HeatCellBg{14}}\CellText{0.14}{0}{0} & \SetCell{bg=\HeatCellBg{28}}\CellText{0.28}{0}{0} & \SetCell{bg=\HeatCellBg{28}}\CellText{0.28}{0}{0} & \SetCell{bg=\HeatCellBg{25}}\CellText{0.25}{0}{0} & \SetCell{bg=\HeatCellBg{24}}\CellText{0.24}{0}{0} & \SetCell{bg=\HeatCellBg{31}}\CellText{0.31}{0}{0} & \SetCell{bg=\HeatCellBg{32}}\CellText{0.32}{0}{0} & \SetCell{bg=\HeatCellBg{32}}\CellText{0.32}{0}{0} & \SetCell{bg=\HeatCellBg{27}}\CellText{0.27}{0}{0} & \SetCell{bg=\HeatCellBg{26}}\CellText{0.26}{0}{0} & \SetCell{bg=\HeatCellBg{27}}\CellText{0.27}{0}{0} & \SetCell{bg=\HeatCellBg{29}}\CellText{0.29}{0}{0} & \SetCell{bg=\HeatCellBg{28}}\CellText{0.28}{0}{0} & \SetCell{bg=\HeatCellBg{29}}\CellText{0.29}{0}{0} & \SetCell{bg=\HeatCellBg{30}}\CellText{0.30}{0}{0} & \SetCell{bg=\HeatCellBg{23}}\CellText{0.23}{0}{0} & \MeanCell{0.27}{0.06}{0} \\
        Mean $\pm$ Std & \MeanCell{0.46}{0.34}{1} & \MeanCell{0.12}{0.12}{0} & \MeanCell{0.08}{0.15}{0} & \MeanCell{0.12}{0.23}{0} & \MeanCell{0.15}{0.16}{0} & \MeanCell{0.21}{0.18}{0} & \MeanCell{0.17}{0.14}{0} & \MeanCell{0.17}{0.16}{0} & \MeanCell{0.18}{0.18}{0} & \MeanCell{0.22}{0.16}{0} & \MeanCell{0.21}{0.18}{0} & \MeanCell{0.20}{0.14}{0} & \MeanCell{0.16}{0.14}{0} & \MeanCell{0.18}{0.16}{0} & \MeanCell{0.19}{0.16}{0} & \MeanCell{0.21}{0.19}{0} & \MeanCell{0.17}{0.17}{0} & \MeanCell{0.21}{0.18}{0} & \MeanCell{0.19}{0.12}{0} & \MeanCell{0.19}{0.19}{0} \\
}

\ClusterKTable{10}{Observation baseline ($\Similarity{obs}{}$) --- \BreakoutName.}{obs_breakout}{
        1 & \SetCell{bg=\HeatCellBg{66}}\CellText{0.66}{1}{0} & \SetCell{bg=\HeatCellBg{82}}\CellText{0.82}{0}{1} & \SetCell{bg=\HeatCellBg{73}}\CellText{0.73}{1}{0} & \SetCell{bg=\HeatCellBg{55}}\CellText{0.55}{0}{0} & \SetCell{bg=\HeatCellBg{74}}\CellText{0.74}{1}{0} & \SetCell{bg=\HeatCellBg{80}}\CellText{0.80}{1}{0} & \SetCell{bg=\HeatCellBg{80}}\CellText{0.80}{1}{0} & \SetCell{bg=\HeatCellBg{79}}\CellText{0.79}{1}{0} & \SetCell{bg=\HeatCellBg{79}}\CellText{0.79}{1}{0} & \SetCell{bg=\HeatCellBg{73}}\CellText{0.73}{1}{0} & \SetCell{bg=\HeatCellBg{74}}\CellText{0.74}{1}{0} & \SetCell{bg=\HeatCellBg{77}}\CellText{0.77}{1}{0} & \SetCell{bg=\HeatCellBg{74}}\CellText{0.74}{1}{0} & \SetCell{bg=\HeatCellBg{78}}\CellText{0.78}{1}{0} & \SetCell{bg=\HeatCellBg{73}}\CellText{0.73}{1}{0} & \SetCell{bg=\HeatCellBg{76}}\CellText{0.76}{1}{0} & \SetCell{bg=\HeatCellBg{77}}\CellText{0.77}{1}{0} & \SetCell{bg=\HeatCellBg{61}}\CellText{0.61}{0}{0} & \SetCell{bg=\HeatCellBg{74}}\CellText{0.74}{1}{0} & \MeanCell{0.74}{0.07}{1} \\
        2 & \SetCell{bg=\HeatCellBg{32}}\CellText{0.32}{0}{0} & \SetCell{bg=\HeatCellBg{91}}\CellText{0.91}{1}{1} & \SetCell{bg=\HeatCellBg{63}}\CellText{0.63}{0}{0} & \SetCell{bg=\HeatCellBg{61}}\CellText{0.61}{1}{0} & \SetCell{bg=\HeatCellBg{73}}\CellText{0.73}{0}{0} & \SetCell{bg=\HeatCellBg{72}}\CellText{0.72}{0}{0} & \SetCell{bg=\HeatCellBg{76}}\CellText{0.76}{0}{0} & \SetCell{bg=\HeatCellBg{76}}\CellText{0.76}{0}{0} & \SetCell{bg=\HeatCellBg{75}}\CellText{0.75}{0}{0} & \SetCell{bg=\HeatCellBg{65}}\CellText{0.65}{0}{0} & \SetCell{bg=\HeatCellBg{69}}\CellText{0.69}{0}{0} & \SetCell{bg=\HeatCellBg{67}}\CellText{0.67}{0}{0} & \SetCell{bg=\HeatCellBg{71}}\CellText{0.71}{0}{0} & \SetCell{bg=\HeatCellBg{70}}\CellText{0.70}{0}{0} & \SetCell{bg=\HeatCellBg{64}}\CellText{0.64}{0}{0} & \SetCell{bg=\HeatCellBg{71}}\CellText{0.71}{0}{0} & \SetCell{bg=\HeatCellBg{63}}\CellText{0.63}{0}{0} & \SetCell{bg=\HeatCellBg{63}}\CellText{0.63}{1}{0} & \SetCell{bg=\HeatCellBg{67}}\CellText{0.67}{0}{0} & \MeanCell{0.68}{0.11}{0} \\
        3 & \SetCell{bg=\HeatCellBg{9}}\CellText{0.09}{0}{0} & \SetCell{bg=\HeatCellBg{86}}\CellText{0.86}{0}{1} & \SetCell{bg=\HeatCellBg{67}}\CellText{0.67}{0}{0} & \SetCell{bg=\HeatCellBg{59}}\CellText{0.59}{0}{0} & \SetCell{bg=\HeatCellBg{47}}\CellText{0.47}{0}{0} & \SetCell{bg=\HeatCellBg{58}}\CellText{0.58}{0}{0} & \SetCell{bg=\HeatCellBg{71}}\CellText{0.71}{0}{0} & \SetCell{bg=\HeatCellBg{58}}\CellText{0.58}{0}{0} & \SetCell{bg=\HeatCellBg{59}}\CellText{0.59}{0}{0} & \SetCell{bg=\HeatCellBg{56}}\CellText{0.56}{0}{0} & \SetCell{bg=\HeatCellBg{53}}\CellText{0.53}{0}{0} & \SetCell{bg=\HeatCellBg{57}}\CellText{0.57}{0}{0} & \SetCell{bg=\HeatCellBg{50}}\CellText{0.50}{0}{0} & \SetCell{bg=\HeatCellBg{50}}\CellText{0.50}{0}{0} & \SetCell{bg=\HeatCellBg{50}}\CellText{0.50}{0}{0} & \SetCell{bg=\HeatCellBg{48}}\CellText{0.48}{0}{0} & \SetCell{bg=\HeatCellBg{49}}\CellText{0.49}{0}{0} & \SetCell{bg=\HeatCellBg{51}}\CellText{0.51}{0}{0} & \SetCell{bg=\HeatCellBg{42}}\CellText{0.42}{0}{0} & \MeanCell{0.54}{0.14}{0} \\
        4 & \SetCell{bg=\HeatCellBg{9}}\CellText{0.09}{0}{0} & \SetCell{bg=\HeatCellBg{72}}\CellText{0.72}{0}{1} & \SetCell{bg=\HeatCellBg{36}}\CellText{0.36}{0}{0} & \SetCell{bg=\HeatCellBg{28}}\CellText{0.28}{0}{0} & \SetCell{bg=\HeatCellBg{21}}\CellText{0.21}{0}{0} & \SetCell{bg=\HeatCellBg{18}}\CellText{0.18}{0}{0} & \SetCell{bg=\HeatCellBg{21}}\CellText{0.21}{0}{0} & \SetCell{bg=\HeatCellBg{18}}\CellText{0.18}{0}{0} & \SetCell{bg=\HeatCellBg{19}}\CellText{0.19}{0}{0} & \SetCell{bg=\HeatCellBg{19}}\CellText{0.19}{0}{0} & \SetCell{bg=\HeatCellBg{18}}\CellText{0.18}{0}{0} & \SetCell{bg=\HeatCellBg{18}}\CellText{0.18}{0}{0} & \SetCell{bg=\HeatCellBg{20}}\CellText{0.20}{0}{0} & \SetCell{bg=\HeatCellBg{19}}\CellText{0.19}{0}{0} & \SetCell{bg=\HeatCellBg{20}}\CellText{0.20}{0}{0} & \SetCell{bg=\HeatCellBg{21}}\CellText{0.21}{0}{0} & \SetCell{bg=\HeatCellBg{21}}\CellText{0.21}{0}{0} & \SetCell{bg=\HeatCellBg{18}}\CellText{0.18}{0}{0} & \SetCell{bg=\HeatCellBg{20}}\CellText{0.20}{0}{0} & \MeanCell{0.23}{0.13}{0} \\
        5 & \SetCell{bg=\HeatCellBg{9}}\CellText{0.09}{0}{0} & \SetCell{bg=\HeatCellBg{72}}\CellText{0.72}{0}{1} & \SetCell{bg=\HeatCellBg{36}}\CellText{0.36}{0}{0} & \SetCell{bg=\HeatCellBg{28}}\CellText{0.28}{0}{0} & \SetCell{bg=\HeatCellBg{21}}\CellText{0.21}{0}{0} & \SetCell{bg=\HeatCellBg{18}}\CellText{0.18}{0}{0} & \SetCell{bg=\HeatCellBg{21}}\CellText{0.21}{0}{0} & \SetCell{bg=\HeatCellBg{18}}\CellText{0.18}{0}{0} & \SetCell{bg=\HeatCellBg{19}}\CellText{0.19}{0}{0} & \SetCell{bg=\HeatCellBg{19}}\CellText{0.19}{0}{0} & \SetCell{bg=\HeatCellBg{18}}\CellText{0.18}{0}{0} & \SetCell{bg=\HeatCellBg{18}}\CellText{0.18}{0}{0} & \SetCell{bg=\HeatCellBg{20}}\CellText{0.20}{0}{0} & \SetCell{bg=\HeatCellBg{19}}\CellText{0.19}{0}{0} & \SetCell{bg=\HeatCellBg{20}}\CellText{0.20}{0}{0} & \SetCell{bg=\HeatCellBg{21}}\CellText{0.21}{0}{0} & \SetCell{bg=\HeatCellBg{21}}\CellText{0.21}{0}{0} & \SetCell{bg=\HeatCellBg{18}}\CellText{0.18}{0}{0} & \SetCell{bg=\HeatCellBg{20}}\CellText{0.20}{0}{0} & \MeanCell{0.23}{0.13}{0} \\
        6b & \SetCell{bg=\HeatCellBg{38}}\CellText{0.38}{0}{1} & \SetCell{bg=\HeatCellBg{16}}\CellText{0.16}{0}{0} & \SetCell{bg=\HeatCellBg{26}}\CellText{0.26}{0}{0} & \SetCell{bg=\HeatCellBg{18}}\CellText{0.18}{0}{0} & \SetCell{bg=\HeatCellBg{15}}\CellText{0.15}{0}{0} & \SetCell{bg=\HeatCellBg{14}}\CellText{0.14}{0}{0} & \SetCell{bg=\HeatCellBg{15}}\CellText{0.15}{0}{0} & \SetCell{bg=\HeatCellBg{15}}\CellText{0.15}{0}{0} & \SetCell{bg=\HeatCellBg{17}}\CellText{0.17}{0}{0} & \SetCell{bg=\HeatCellBg{18}}\CellText{0.18}{0}{0} & \SetCell{bg=\HeatCellBg{15}}\CellText{0.15}{0}{0} & \SetCell{bg=\HeatCellBg{15}}\CellText{0.15}{0}{0} & \SetCell{bg=\HeatCellBg{16}}\CellText{0.16}{0}{0} & \SetCell{bg=\HeatCellBg{17}}\CellText{0.17}{0}{0} & \SetCell{bg=\HeatCellBg{17}}\CellText{0.17}{0}{0} & \SetCell{bg=\HeatCellBg{17}}\CellText{0.17}{0}{0} & \SetCell{bg=\HeatCellBg{16}}\CellText{0.16}{0}{0} & \SetCell{bg=\HeatCellBg{15}}\CellText{0.15}{0}{0} & \SetCell{bg=\HeatCellBg{17}}\CellText{0.17}{0}{0} & \MeanCell{0.18}{0.05}{0} \\
        7b & \SetCell{bg=\HeatCellBg{36}}\CellText{0.36}{0}{1} & \SetCell{bg=\HeatCellBg{25}}\CellText{0.25}{0}{0} & \SetCell{bg=\HeatCellBg{31}}\CellText{0.31}{0}{0} & \SetCell{bg=\HeatCellBg{33}}\CellText{0.33}{0}{0} & \SetCell{bg=\HeatCellBg{23}}\CellText{0.23}{0}{0} & \SetCell{bg=\HeatCellBg{21}}\CellText{0.21}{0}{0} & \SetCell{bg=\HeatCellBg{21}}\CellText{0.21}{0}{0} & \SetCell{bg=\HeatCellBg{20}}\CellText{0.20}{0}{0} & \SetCell{bg=\HeatCellBg{20}}\CellText{0.20}{0}{0} & \SetCell{bg=\HeatCellBg{18}}\CellText{0.18}{0}{0} & \SetCell{bg=\HeatCellBg{18}}\CellText{0.18}{0}{0} & \SetCell{bg=\HeatCellBg{17}}\CellText{0.17}{0}{0} & \SetCell{bg=\HeatCellBg{17}}\CellText{0.17}{0}{0} & \SetCell{bg=\HeatCellBg{18}}\CellText{0.18}{0}{0} & \SetCell{bg=\HeatCellBg{17}}\CellText{0.17}{0}{0} & \SetCell{bg=\HeatCellBg{17}}\CellText{0.17}{0}{0} & \SetCell{bg=\HeatCellBg{17}}\CellText{0.17}{0}{0} & \SetCell{bg=\HeatCellBg{15}}\CellText{0.15}{0}{0} & \SetCell{bg=\HeatCellBg{17}}\CellText{0.17}{0}{0} & \MeanCell{0.21}{0.06}{0} \\
        8b & \SetCell{bg=\HeatCellBg{29}}\CellText{0.29}{0}{0} & \SetCell{bg=\HeatCellBg{33}}\CellText{0.33}{0}{0} & \SetCell{bg=\HeatCellBg{34}}\CellText{0.34}{0}{0} & \SetCell{bg=\HeatCellBg{36}}\CellText{0.36}{0}{1} & \SetCell{bg=\HeatCellBg{22}}\CellText{0.22}{0}{0} & \SetCell{bg=\HeatCellBg{23}}\CellText{0.23}{0}{0} & \SetCell{bg=\HeatCellBg{23}}\CellText{0.23}{0}{0} & \SetCell{bg=\HeatCellBg{20}}\CellText{0.20}{0}{0} & \SetCell{bg=\HeatCellBg{24}}\CellText{0.24}{0}{0} & \SetCell{bg=\HeatCellBg{22}}\CellText{0.22}{0}{0} & \SetCell{bg=\HeatCellBg{23}}\CellText{0.23}{0}{0} & \SetCell{bg=\HeatCellBg{21}}\CellText{0.21}{0}{0} & \SetCell{bg=\HeatCellBg{21}}\CellText{0.21}{0}{0} & \SetCell{bg=\HeatCellBg{22}}\CellText{0.22}{0}{0} & \SetCell{bg=\HeatCellBg{22}}\CellText{0.22}{0}{0} & \SetCell{bg=\HeatCellBg{22}}\CellText{0.22}{0}{0} & \SetCell{bg=\HeatCellBg{22}}\CellText{0.22}{0}{0} & \SetCell{bg=\HeatCellBg{18}}\CellText{0.18}{0}{0} & \SetCell{bg=\HeatCellBg{22}}\CellText{0.22}{0}{0} & \MeanCell{0.24}{0.05}{0} \\
        Mean $\pm$ Std & \MeanCell{0.28}{0.18}{0} & \MeanCell{0.60}{0.28}{1} & \MeanCell{0.46}{0.17}{0} & \MeanCell{0.40}{0.15}{0} & \MeanCell{0.37}{0.23}{0} & \MeanCell{0.38}{0.26}{0} & \MeanCell{0.41}{0.27}{0} & \MeanCell{0.38}{0.26}{0} & \MeanCell{0.39}{0.25}{0} & \MeanCell{0.36}{0.22}{0} & \MeanCell{0.36}{0.23}{0} & \MeanCell{0.36}{0.24}{0} & \MeanCell{0.36}{0.23}{0} & \MeanCell{0.37}{0.24}{0} & \MeanCell{0.35}{0.22}{0} & \MeanCell{0.37}{0.23}{0} & \MeanCell{0.36}{0.22}{0} & \MeanCell{0.32}{0.20}{0} & \MeanCell{0.35}{0.22}{0} & \MeanCell{0.38}{0.24}{0} \\
}

\ClusterKTable{10}{Untrained-network (latent) baseline ($\Similarity{untr}{}$) --- \BreakoutName.}{untr_breakout}{
        1 & \SetCell{bg=\HeatCellBg{95}}\CellText{0.95}{1}{0} & \SetCell{bg=\HeatCellBg{94}}\CellText{0.94}{1}{0} & \SetCell{bg=\HeatCellBg{97}}\CellText{0.97}{1}{1} & \SetCell{bg=\HeatCellBg{76}}\CellText{0.76}{1}{0} & \SetCell{bg=\HeatCellBg{84}}\CellText{0.84}{1}{0} & \SetCell{bg=\HeatCellBg{83}}\CellText{0.83}{1}{0} & \SetCell{bg=\HeatCellBg{84}}\CellText{0.84}{1}{0} & \SetCell{bg=\HeatCellBg{82}}\CellText{0.82}{1}{0} & \SetCell{bg=\HeatCellBg{85}}\CellText{0.85}{1}{0} & \SetCell{bg=\HeatCellBg{85}}\CellText{0.85}{1}{0} & \SetCell{bg=\HeatCellBg{82}}\CellText{0.82}{1}{0} & \SetCell{bg=\HeatCellBg{86}}\CellText{0.86}{1}{0} & \SetCell{bg=\HeatCellBg{87}}\CellText{0.87}{1}{0} & \SetCell{bg=\HeatCellBg{77}}\CellText{0.77}{1}{0} & \SetCell{bg=\HeatCellBg{84}}\CellText{0.84}{1}{0} & \SetCell{bg=\HeatCellBg{84}}\CellText{0.84}{1}{0} & \SetCell{bg=\HeatCellBg{80}}\CellText{0.80}{1}{0} & \SetCell{bg=\HeatCellBg{77}}\CellText{0.77}{1}{0} & \SetCell{bg=\HeatCellBg{85}}\CellText{0.85}{1}{0} & \MeanCell{0.85}{0.06}{1} \\
        2 & \SetCell{bg=\HeatCellBg{38}}\CellText{0.38}{0}{0} & \SetCell{bg=\HeatCellBg{63}}\CellText{0.63}{0}{0} & \SetCell{bg=\HeatCellBg{80}}\CellText{0.80}{0}{1} & \SetCell{bg=\HeatCellBg{63}}\CellText{0.63}{0}{0} & \SetCell{bg=\HeatCellBg{66}}\CellText{0.66}{0}{0} & \SetCell{bg=\HeatCellBg{67}}\CellText{0.67}{0}{0} & \SetCell{bg=\HeatCellBg{68}}\CellText{0.68}{0}{0} & \SetCell{bg=\HeatCellBg{76}}\CellText{0.76}{0}{0} & \SetCell{bg=\HeatCellBg{67}}\CellText{0.67}{0}{0} & \SetCell{bg=\HeatCellBg{72}}\CellText{0.72}{0}{0} & \SetCell{bg=\HeatCellBg{65}}\CellText{0.65}{0}{0} & \SetCell{bg=\HeatCellBg{66}}\CellText{0.66}{0}{0} & \SetCell{bg=\HeatCellBg{69}}\CellText{0.69}{0}{0} & \SetCell{bg=\HeatCellBg{72}}\CellText{0.72}{0}{0} & \SetCell{bg=\HeatCellBg{71}}\CellText{0.71}{0}{0} & \SetCell{bg=\HeatCellBg{69}}\CellText{0.69}{0}{0} & \SetCell{bg=\HeatCellBg{69}}\CellText{0.69}{0}{0} & \SetCell{bg=\HeatCellBg{73}}\CellText{0.73}{0}{0} & \SetCell{bg=\HeatCellBg{73}}\CellText{0.73}{0}{0} & \MeanCell{0.68}{0.08}{0} \\
        3 & \SetCell{bg=\HeatCellBg{15}}\CellText{0.15}{0}{0} & \SetCell{bg=\HeatCellBg{68}}\CellText{0.68}{0}{0} & \SetCell{bg=\HeatCellBg{80}}\CellText{0.80}{0}{1} & \SetCell{bg=\HeatCellBg{64}}\CellText{0.64}{0}{0} & \SetCell{bg=\HeatCellBg{54}}\CellText{0.54}{0}{0} & \SetCell{bg=\HeatCellBg{47}}\CellText{0.47}{0}{0} & \SetCell{bg=\HeatCellBg{53}}\CellText{0.53}{0}{0} & \SetCell{bg=\HeatCellBg{48}}\CellText{0.48}{0}{0} & \SetCell{bg=\HeatCellBg{45}}\CellText{0.45}{0}{0} & \SetCell{bg=\HeatCellBg{45}}\CellText{0.45}{0}{0} & \SetCell{bg=\HeatCellBg{46}}\CellText{0.46}{0}{0} & \SetCell{bg=\HeatCellBg{50}}\CellText{0.50}{0}{0} & \SetCell{bg=\HeatCellBg{46}}\CellText{0.46}{0}{0} & \SetCell{bg=\HeatCellBg{44}}\CellText{0.44}{0}{0} & \SetCell{bg=\HeatCellBg{46}}\CellText{0.46}{0}{0} & \SetCell{bg=\HeatCellBg{51}}\CellText{0.51}{0}{0} & \SetCell{bg=\HeatCellBg{48}}\CellText{0.48}{0}{0} & \SetCell{bg=\HeatCellBg{54}}\CellText{0.54}{0}{0} & \SetCell{bg=\HeatCellBg{42}}\CellText{0.42}{0}{0} & \MeanCell{0.50}{0.12}{0} \\
        4 & \SetCell{bg=\HeatCellBg{12}}\CellText{0.12}{0}{0} & \SetCell{bg=\HeatCellBg{48}}\CellText{0.48}{0}{1} & \SetCell{bg=\HeatCellBg{32}}\CellText{0.32}{0}{0} & \SetCell{bg=\HeatCellBg{27}}\CellText{0.27}{0}{0} & \SetCell{bg=\HeatCellBg{25}}\CellText{0.25}{0}{0} & \SetCell{bg=\HeatCellBg{20}}\CellText{0.20}{0}{0} & \SetCell{bg=\HeatCellBg{23}}\CellText{0.23}{0}{0} & \SetCell{bg=\HeatCellBg{19}}\CellText{0.19}{0}{0} & \SetCell{bg=\HeatCellBg{20}}\CellText{0.20}{0}{0} & \SetCell{bg=\HeatCellBg{17}}\CellText{0.17}{0}{0} & \SetCell{bg=\HeatCellBg{18}}\CellText{0.18}{0}{0} & \SetCell{bg=\HeatCellBg{18}}\CellText{0.18}{0}{0} & \SetCell{bg=\HeatCellBg{18}}\CellText{0.18}{0}{0} & \SetCell{bg=\HeatCellBg{18}}\CellText{0.18}{0}{0} & \SetCell{bg=\HeatCellBg{18}}\CellText{0.18}{0}{0} & \SetCell{bg=\HeatCellBg{17}}\CellText{0.17}{0}{0} & \SetCell{bg=\HeatCellBg{17}}\CellText{0.17}{0}{0} & \SetCell{bg=\HeatCellBg{18}}\CellText{0.18}{0}{0} & \SetCell{bg=\HeatCellBg{18}}\CellText{0.18}{0}{0} & \MeanCell{0.21}{0.08}{0} \\
        5 & \SetCell{bg=\HeatCellBg{12}}\CellText{0.12}{0}{0} & \SetCell{bg=\HeatCellBg{48}}\CellText{0.48}{0}{1} & \SetCell{bg=\HeatCellBg{32}}\CellText{0.32}{0}{0} & \SetCell{bg=\HeatCellBg{27}}\CellText{0.27}{0}{0} & \SetCell{bg=\HeatCellBg{25}}\CellText{0.25}{0}{0} & \SetCell{bg=\HeatCellBg{20}}\CellText{0.20}{0}{0} & \SetCell{bg=\HeatCellBg{23}}\CellText{0.23}{0}{0} & \SetCell{bg=\HeatCellBg{19}}\CellText{0.19}{0}{0} & \SetCell{bg=\HeatCellBg{20}}\CellText{0.20}{0}{0} & \SetCell{bg=\HeatCellBg{17}}\CellText{0.17}{0}{0} & \SetCell{bg=\HeatCellBg{18}}\CellText{0.18}{0}{0} & \SetCell{bg=\HeatCellBg{18}}\CellText{0.18}{0}{0} & \SetCell{bg=\HeatCellBg{18}}\CellText{0.18}{0}{0} & \SetCell{bg=\HeatCellBg{18}}\CellText{0.18}{0}{0} & \SetCell{bg=\HeatCellBg{18}}\CellText{0.18}{0}{0} & \SetCell{bg=\HeatCellBg{17}}\CellText{0.17}{0}{0} & \SetCell{bg=\HeatCellBg{17}}\CellText{0.17}{0}{0} & \SetCell{bg=\HeatCellBg{18}}\CellText{0.18}{0}{0} & \SetCell{bg=\HeatCellBg{18}}\CellText{0.18}{0}{0} & \MeanCell{0.21}{0.08}{0} \\
        6b & \SetCell{bg=\HeatCellBg{32}}\CellText{0.32}{0}{1} & \SetCell{bg=\HeatCellBg{16}}\CellText{0.16}{0}{0} & \SetCell{bg=\HeatCellBg{22}}\CellText{0.22}{0}{0} & \SetCell{bg=\HeatCellBg{16}}\CellText{0.16}{0}{0} & \SetCell{bg=\HeatCellBg{15}}\CellText{0.15}{0}{0} & \SetCell{bg=\HeatCellBg{15}}\CellText{0.15}{0}{0} & \SetCell{bg=\HeatCellBg{13}}\CellText{0.13}{0}{0} & \SetCell{bg=\HeatCellBg{14}}\CellText{0.14}{0}{0} & \SetCell{bg=\HeatCellBg{14}}\CellText{0.14}{0}{0} & \SetCell{bg=\HeatCellBg{16}}\CellText{0.16}{0}{0} & \SetCell{bg=\HeatCellBg{13}}\CellText{0.13}{0}{0} & \SetCell{bg=\HeatCellBg{13}}\CellText{0.13}{0}{0} & \SetCell{bg=\HeatCellBg{14}}\CellText{0.14}{0}{0} & \SetCell{bg=\HeatCellBg{14}}\CellText{0.14}{0}{0} & \SetCell{bg=\HeatCellBg{14}}\CellText{0.14}{0}{0} & \SetCell{bg=\HeatCellBg{14}}\CellText{0.14}{0}{0} & \SetCell{bg=\HeatCellBg{13}}\CellText{0.13}{0}{0} & \SetCell{bg=\HeatCellBg{13}}\CellText{0.13}{0}{0} & \SetCell{bg=\HeatCellBg{14}}\CellText{0.14}{0}{0} & \MeanCell{0.16}{0.04}{0} \\
        7b & \SetCell{bg=\HeatCellBg{31}}\CellText{0.31}{0}{1} & \SetCell{bg=\HeatCellBg{22}}\CellText{0.22}{0}{0} & \SetCell{bg=\HeatCellBg{26}}\CellText{0.26}{0}{0} & \SetCell{bg=\HeatCellBg{23}}\CellText{0.23}{0}{0} & \SetCell{bg=\HeatCellBg{26}}\CellText{0.26}{0}{0} & \SetCell{bg=\HeatCellBg{16}}\CellText{0.16}{0}{0} & \SetCell{bg=\HeatCellBg{16}}\CellText{0.16}{0}{0} & \SetCell{bg=\HeatCellBg{17}}\CellText{0.17}{0}{0} & \SetCell{bg=\HeatCellBg{17}}\CellText{0.17}{0}{0} & \SetCell{bg=\HeatCellBg{14}}\CellText{0.14}{0}{0} & \SetCell{bg=\HeatCellBg{12}}\CellText{0.12}{0}{0} & \SetCell{bg=\HeatCellBg{14}}\CellText{0.14}{0}{0} & \SetCell{bg=\HeatCellBg{13}}\CellText{0.13}{0}{0} & \SetCell{bg=\HeatCellBg{11}}\CellText{0.11}{0}{0} & \SetCell{bg=\HeatCellBg{12}}\CellText{0.12}{0}{0} & \SetCell{bg=\HeatCellBg{11}}\CellText{0.11}{0}{0} & \SetCell{bg=\HeatCellBg{13}}\CellText{0.13}{0}{0} & \SetCell{bg=\HeatCellBg{11}}\CellText{0.11}{0}{0} & \SetCell{bg=\HeatCellBg{12}}\CellText{0.12}{0}{0} & \MeanCell{0.17}{0.06}{0} \\
        8b & \SetCell{bg=\HeatCellBg{18}}\CellText{0.18}{0}{0} & \SetCell{bg=\HeatCellBg{27}}\CellText{0.27}{0}{1} & \SetCell{bg=\HeatCellBg{17}}\CellText{0.17}{0}{0} & \SetCell{bg=\HeatCellBg{19}}\CellText{0.19}{0}{0} & \SetCell{bg=\HeatCellBg{20}}\CellText{0.20}{0}{0} & \SetCell{bg=\HeatCellBg{16}}\CellText{0.16}{0}{0} & \SetCell{bg=\HeatCellBg{22}}\CellText{0.22}{0}{0} & \SetCell{bg=\HeatCellBg{16}}\CellText{0.16}{0}{0} & \SetCell{bg=\HeatCellBg{17}}\CellText{0.17}{0}{0} & \SetCell{bg=\HeatCellBg{15}}\CellText{0.15}{0}{0} & \SetCell{bg=\HeatCellBg{13}}\CellText{0.13}{0}{0} & \SetCell{bg=\HeatCellBg{11}}\CellText{0.11}{0}{0} & \SetCell{bg=\HeatCellBg{13}}\CellText{0.13}{0}{0} & \SetCell{bg=\HeatCellBg{14}}\CellText{0.14}{0}{0} & \SetCell{bg=\HeatCellBg{14}}\CellText{0.14}{0}{0} & \SetCell{bg=\HeatCellBg{13}}\CellText{0.13}{0}{0} & \SetCell{bg=\HeatCellBg{14}}\CellText{0.14}{0}{0} & \SetCell{bg=\HeatCellBg{13}}\CellText{0.13}{0}{0} & \SetCell{bg=\HeatCellBg{10}}\CellText{0.10}{0}{0} & \MeanCell{0.16}{0.04}{0} \\
        Mean $\pm$ Std & \MeanCell{0.32}{0.26}{0} & \MeanCell{0.48}{0.25}{0} & \MeanCell{0.48}{0.30}{1} & \MeanCell{0.39}{0.22}{0} & \MeanCell{0.39}{0.24}{0} & \MeanCell{0.35}{0.25}{0} & \MeanCell{0.38}{0.25}{0} & \MeanCell{0.36}{0.27}{0} & \MeanCell{0.36}{0.25}{0} & \MeanCell{0.35}{0.27}{0} & \MeanCell{0.33}{0.26}{0} & \MeanCell{0.35}{0.27}{0} & \MeanCell{0.35}{0.27}{0} & \MeanCell{0.34}{0.26}{0} & \MeanCell{0.35}{0.27}{0} & \MeanCell{0.34}{0.27}{0} & \MeanCell{0.34}{0.26}{0} & \MeanCell{0.35}{0.27}{0} & \MeanCell{0.34}{0.28}{0} & \MeanCell{0.37}{0.27}{0} \\
}

\ClusterKTable{10}{Action baseline ($\Similarity{act}{}$) --- \BreakoutName.}{act_breakout}{
        1 & \SetCell{bg=\HeatCellBg{2}}\CellText{0.02}{0}{0} & \SetCell{bg=\HeatCellBg{2}}\CellText{0.02}{0}{0} & \SetCell{bg=\HeatCellBg{4}}\CellText{0.04}{0}{0} & \SetCell{bg=\HeatCellBg{3}}\CellText{0.03}{0}{0} & \SetCell{bg=\HeatCellBg{3}}\CellText{0.03}{0}{0} & \SetCell{bg=\HeatCellBg{5}}\CellText{0.05}{0}{0} & \SetCell{bg=\HeatCellBg{5}}\CellText{0.05}{0}{0} & \SetCell{bg=\HeatCellBg{5}}\CellText{0.05}{0}{0} & \SetCell{bg=\HeatCellBg{3}}\CellText{0.03}{0}{0} & \SetCell{bg=\HeatCellBg{6}}\CellText{0.06}{0}{0} & \SetCell{bg=\HeatCellBg{7}}\CellText{0.07}{0}{0} & \SetCell{bg=\HeatCellBg{8}}\CellText{0.08}{0}{0} & \SetCell{bg=\HeatCellBg{10}}\CellText{0.10}{0}{0} & \SetCell{bg=\HeatCellBg{9}}\CellText{0.09}{0}{0} & \SetCell{bg=\HeatCellBg{9}}\CellText{0.09}{0}{0} & \SetCell{bg=\HeatCellBg{8}}\CellText{0.08}{0}{0} & \SetCell{bg=\HeatCellBg{6}}\CellText{0.06}{0}{0} & \SetCell{bg=\HeatCellBg{9}}\CellText{0.09}{0}{0} & \SetCell{bg=\HeatCellBg{11}}\CellText{0.11}{0}{1} & \MeanCell{0.06}{0.03}{0} \\
        2 & \SetCell{bg=\HeatCellBg{3}}\CellText{0.03}{0}{0} & \SetCell{bg=\HeatCellBg{1}}\CellText{0.01}{0}{0} & \SetCell{bg=\HeatCellBg{3}}\CellText{0.03}{0}{0} & \SetCell{bg=\HeatCellBg{4}}\CellText{0.04}{0}{0} & \SetCell{bg=\HeatCellBg{2}}\CellText{0.02}{0}{0} & \SetCell{bg=\HeatCellBg{4}}\CellText{0.04}{0}{0} & \SetCell{bg=\HeatCellBg{5}}\CellText{0.05}{0}{0} & \SetCell{bg=\HeatCellBg{5}}\CellText{0.05}{0}{0} & \SetCell{bg=\HeatCellBg{3}}\CellText{0.03}{0}{0} & \SetCell{bg=\HeatCellBg{6}}\CellText{0.06}{0}{0} & \SetCell{bg=\HeatCellBg{5}}\CellText{0.05}{0}{0} & \SetCell{bg=\HeatCellBg{6}}\CellText{0.06}{0}{0} & \SetCell{bg=\HeatCellBg{8}}\CellText{0.08}{0}{0} & \SetCell{bg=\HeatCellBg{9}}\CellText{0.09}{0}{1} & \SetCell{bg=\HeatCellBg{8}}\CellText{0.08}{0}{0} & \SetCell{bg=\HeatCellBg{7}}\CellText{0.07}{0}{0} & \SetCell{bg=\HeatCellBg{5}}\CellText{0.05}{0}{0} & \SetCell{bg=\HeatCellBg{9}}\CellText{0.09}{0}{1} & \SetCell{bg=\HeatCellBg{9}}\CellText{0.09}{0}{1} & \MeanCell{0.05}{0.02}{0} \\
        3 & \SetCell{bg=\HeatCellBg{5}}\CellText{0.05}{0}{0} & \SetCell{bg=\HeatCellBg{2}}\CellText{0.02}{0}{0} & \SetCell{bg=\HeatCellBg{4}}\CellText{0.04}{0}{0} & \SetCell{bg=\HeatCellBg{4}}\CellText{0.04}{0}{0} & \SetCell{bg=\HeatCellBg{3}}\CellText{0.03}{0}{0} & \SetCell{bg=\HeatCellBg{5}}\CellText{0.05}{0}{0} & \SetCell{bg=\HeatCellBg{6}}\CellText{0.06}{0}{0} & \SetCell{bg=\HeatCellBg{5}}\CellText{0.05}{0}{0} & \SetCell{bg=\HeatCellBg{4}}\CellText{0.04}{0}{0} & \SetCell{bg=\HeatCellBg{6}}\CellText{0.06}{0}{0} & \SetCell{bg=\HeatCellBg{7}}\CellText{0.07}{0}{0} & \SetCell{bg=\HeatCellBg{6}}\CellText{0.06}{0}{0} & \SetCell{bg=\HeatCellBg{8}}\CellText{0.08}{0}{0} & \SetCell{bg=\HeatCellBg{9}}\CellText{0.09}{0}{1} & \SetCell{bg=\HeatCellBg{7}}\CellText{0.07}{0}{0} & \SetCell{bg=\HeatCellBg{6}}\CellText{0.06}{0}{0} & \SetCell{bg=\HeatCellBg{5}}\CellText{0.05}{0}{0} & \SetCell{bg=\HeatCellBg{8}}\CellText{0.08}{0}{0} & \SetCell{bg=\HeatCellBg{6}}\CellText{0.06}{0}{0} & \MeanCell{0.06}{0.02}{0} \\
        4 & \SetCell{bg=\HeatCellBg{5}}\CellText{0.05}{0}{1} & \SetCell{bg=\HeatCellBg{2}}\CellText{0.02}{0}{0} & \SetCell{bg=\HeatCellBg{2}}\CellText{0.02}{0}{0} & \SetCell{bg=\HeatCellBg{2}}\CellText{0.02}{0}{0} & \SetCell{bg=\HeatCellBg{2}}\CellText{0.02}{0}{0} & \SetCell{bg=\HeatCellBg{3}}\CellText{0.03}{0}{0} & \SetCell{bg=\HeatCellBg{4}}\CellText{0.04}{0}{0} & \SetCell{bg=\HeatCellBg{3}}\CellText{0.03}{0}{0} & \SetCell{bg=\HeatCellBg{2}}\CellText{0.02}{0}{0} & \SetCell{bg=\HeatCellBg{3}}\CellText{0.03}{0}{0} & \SetCell{bg=\HeatCellBg{4}}\CellText{0.04}{0}{0} & \SetCell{bg=\HeatCellBg{3}}\CellText{0.03}{0}{0} & \SetCell{bg=\HeatCellBg{4}}\CellText{0.04}{0}{0} & \SetCell{bg=\HeatCellBg{5}}\CellText{0.05}{0}{1} & \SetCell{bg=\HeatCellBg{4}}\CellText{0.04}{0}{0} & \SetCell{bg=\HeatCellBg{4}}\CellText{0.04}{0}{0} & \SetCell{bg=\HeatCellBg{3}}\CellText{0.03}{0}{0} & \SetCell{bg=\HeatCellBg{4}}\CellText{0.04}{0}{0} & \SetCell{bg=\HeatCellBg{4}}\CellText{0.04}{0}{0} & \MeanCell{0.03}{0.01}{0} \\
        5 & \SetCell{bg=\HeatCellBg{5}}\CellText{0.05}{0}{1} & \SetCell{bg=\HeatCellBg{2}}\CellText{0.02}{0}{0} & \SetCell{bg=\HeatCellBg{2}}\CellText{0.02}{0}{0} & \SetCell{bg=\HeatCellBg{2}}\CellText{0.02}{0}{0} & \SetCell{bg=\HeatCellBg{2}}\CellText{0.02}{0}{0} & \SetCell{bg=\HeatCellBg{3}}\CellText{0.03}{0}{0} & \SetCell{bg=\HeatCellBg{4}}\CellText{0.04}{0}{0} & \SetCell{bg=\HeatCellBg{3}}\CellText{0.03}{0}{0} & \SetCell{bg=\HeatCellBg{2}}\CellText{0.02}{0}{0} & \SetCell{bg=\HeatCellBg{3}}\CellText{0.03}{0}{0} & \SetCell{bg=\HeatCellBg{4}}\CellText{0.04}{0}{0} & \SetCell{bg=\HeatCellBg{3}}\CellText{0.03}{0}{0} & \SetCell{bg=\HeatCellBg{4}}\CellText{0.04}{0}{0} & \SetCell{bg=\HeatCellBg{5}}\CellText{0.05}{0}{1} & \SetCell{bg=\HeatCellBg{4}}\CellText{0.04}{0}{0} & \SetCell{bg=\HeatCellBg{4}}\CellText{0.04}{0}{0} & \SetCell{bg=\HeatCellBg{3}}\CellText{0.03}{0}{0} & \SetCell{bg=\HeatCellBg{4}}\CellText{0.04}{0}{0} & \SetCell{bg=\HeatCellBg{4}}\CellText{0.04}{0}{0} & \MeanCell{0.03}{0.01}{0} \\
        6b & \SetCell{bg=\HeatCellBg{5}}\CellText{0.05}{0}{0} & \SetCell{bg=\HeatCellBg{2}}\CellText{0.02}{0}{0} & \SetCell{bg=\HeatCellBg{3}}\CellText{0.03}{0}{0} & \SetCell{bg=\HeatCellBg{3}}\CellText{0.03}{0}{0} & \SetCell{bg=\HeatCellBg{2}}\CellText{0.02}{0}{0} & \SetCell{bg=\HeatCellBg{3}}\CellText{0.03}{0}{0} & \SetCell{bg=\HeatCellBg{3}}\CellText{0.03}{0}{0} & \SetCell{bg=\HeatCellBg{4}}\CellText{0.04}{0}{0} & \SetCell{bg=\HeatCellBg{3}}\CellText{0.03}{0}{0} & \SetCell{bg=\HeatCellBg{3}}\CellText{0.03}{0}{0} & \SetCell{bg=\HeatCellBg{4}}\CellText{0.04}{0}{0} & \SetCell{bg=\HeatCellBg{4}}\CellText{0.04}{0}{0} & \SetCell{bg=\HeatCellBg{5}}\CellText{0.05}{0}{0} & \SetCell{bg=\HeatCellBg{4}}\CellText{0.04}{0}{0} & \SetCell{bg=\HeatCellBg{5}}\CellText{0.05}{0}{0} & \SetCell{bg=\HeatCellBg{5}}\CellText{0.05}{0}{0} & \SetCell{bg=\HeatCellBg{4}}\CellText{0.04}{0}{0} & \SetCell{bg=\HeatCellBg{5}}\CellText{0.05}{0}{0} & \SetCell{bg=\HeatCellBg{6}}\CellText{0.06}{0}{1} & \MeanCell{0.04}{0.01}{0} \\
        7b & \SetCell{bg=\HeatCellBg{5}}\CellText{0.05}{0}{0} & \SetCell{bg=\HeatCellBg{2}}\CellText{0.02}{0}{0} & \SetCell{bg=\HeatCellBg{4}}\CellText{0.04}{0}{0} & \SetCell{bg=\HeatCellBg{5}}\CellText{0.05}{0}{0} & \SetCell{bg=\HeatCellBg{4}}\CellText{0.04}{0}{0} & \SetCell{bg=\HeatCellBg{6}}\CellText{0.06}{0}{0} & \SetCell{bg=\HeatCellBg{7}}\CellText{0.07}{0}{0} & \SetCell{bg=\HeatCellBg{6}}\CellText{0.06}{0}{0} & \SetCell{bg=\HeatCellBg{5}}\CellText{0.05}{0}{0} & \SetCell{bg=\HeatCellBg{5}}\CellText{0.05}{0}{0} & \SetCell{bg=\HeatCellBg{5}}\CellText{0.05}{0}{0} & \SetCell{bg=\HeatCellBg{6}}\CellText{0.06}{0}{0} & \SetCell{bg=\HeatCellBg{8}}\CellText{0.08}{0}{1} & \SetCell{bg=\HeatCellBg{8}}\CellText{0.08}{0}{1} & \SetCell{bg=\HeatCellBg{7}}\CellText{0.07}{0}{0} & \SetCell{bg=\HeatCellBg{7}}\CellText{0.07}{0}{0} & \SetCell{bg=\HeatCellBg{5}}\CellText{0.05}{0}{0} & \SetCell{bg=\HeatCellBg{8}}\CellText{0.08}{0}{1} & \SetCell{bg=\HeatCellBg{8}}\CellText{0.08}{0}{1} & \MeanCell{0.06}{0.02}{0} \\
        8b & \SetCell{bg=\HeatCellBg{11}}\CellText{0.11}{1}{0} & \SetCell{bg=\HeatCellBg{10}}\CellText{0.10}{1}{0} & \SetCell{bg=\HeatCellBg{13}}\CellText{0.13}{1}{0} & \SetCell{bg=\HeatCellBg{16}}\CellText{0.16}{1}{0} & \SetCell{bg=\HeatCellBg{17}}\CellText{0.17}{1}{0} & \SetCell{bg=\HeatCellBg{21}}\CellText{0.21}{1}{0} & \SetCell{bg=\HeatCellBg{23}}\CellText{0.23}{1}{1} & \SetCell{bg=\HeatCellBg{16}}\CellText{0.16}{1}{0} & \SetCell{bg=\HeatCellBg{13}}\CellText{0.13}{1}{0} & \SetCell{bg=\HeatCellBg{17}}\CellText{0.17}{1}{0} & \SetCell{bg=\HeatCellBg{19}}\CellText{0.19}{1}{0} & \SetCell{bg=\HeatCellBg{18}}\CellText{0.18}{1}{0} & \SetCell{bg=\HeatCellBg{20}}\CellText{0.20}{1}{0} & \SetCell{bg=\HeatCellBg{19}}\CellText{0.19}{1}{0} & \SetCell{bg=\HeatCellBg{18}}\CellText{0.18}{1}{0} & \SetCell{bg=\HeatCellBg{16}}\CellText{0.16}{1}{0} & \SetCell{bg=\HeatCellBg{13}}\CellText{0.13}{1}{0} & \SetCell{bg=\HeatCellBg{18}}\CellText{0.18}{1}{0} & \SetCell{bg=\HeatCellBg{18}}\CellText{0.18}{1}{0} & \MeanCell{0.17}{0.03}{1} \\
        Mean $\pm$ Std & \MeanCell{0.05}{0.02}{0} & \MeanCell{0.03}{0.03}{0} & \MeanCell{0.04}{0.03}{0} & \MeanCell{0.05}{0.04}{0} & \MeanCell{0.04}{0.05}{0} & \MeanCell{0.06}{0.06}{0} & \MeanCell{0.07}{0.06}{0} & \MeanCell{0.06}{0.04}{0} & \MeanCell{0.04}{0.03}{0} & \MeanCell{0.06}{0.04}{0} & \MeanCell{0.07}{0.05}{0} & \MeanCell{0.07}{0.05}{0} & \MeanCell{0.08}{0.05}{0} & \MeanCell{0.09}{0.04}{1} & \MeanCell{0.08}{0.04}{0} & \MeanCell{0.07}{0.04}{0} & \MeanCell{0.06}{0.03}{0} & \MeanCell{0.08}{0.04}{0} & \MeanCell{0.08}{0.04}{0} & \MeanCell{0.06}{0.05}{0} \\
}


\clearpage
\section{\MarioBrosName}
\label{sec:appendix_additional_results_mario_bros}
\ClusterKTable{10}{\gls{cu} --- \MarioBrosName.}{cu_mario_bros}{
        1 & \SetCell{bg=\HeatCellBg{100}}\CellText{1.00}{1}{1} & \SetCell{bg=\HeatCellBg{93}}\CellText{0.93}{1}{0} & \SetCell{bg=\HeatCellBg{82}}\CellText{0.82}{0}{0} & \SetCell{bg=\HeatCellBg{65}}\CellText{0.65}{0}{0} & \SetCell{bg=\HeatCellBg{51}}\CellText{0.51}{0}{0} & \SetCell{bg=\HeatCellBg{53}}\CellText{0.53}{0}{0} & \SetCell{bg=\HeatCellBg{67}}\CellText{0.67}{1}{0} & \SetCell{bg=\HeatCellBg{45}}\CellText{0.45}{0}{0} & \SetCell{bg=\HeatCellBg{58}}\CellText{0.58}{0}{0} & \SetCell{bg=\HeatCellBg{54}}\CellText{0.54}{0}{0} & \SetCell{bg=\HeatCellBg{52}}\CellText{0.52}{0}{0} & \SetCell{bg=\HeatCellBg{54}}\CellText{0.54}{0}{0} & \SetCell{bg=\HeatCellBg{62}}\CellText{0.62}{0}{0} & \SetCell{bg=\HeatCellBg{60}}\CellText{0.60}{0}{0} & \SetCell{bg=\HeatCellBg{54}}\CellText{0.54}{0}{0} & \SetCell{bg=\HeatCellBg{49}}\CellText{0.49}{0}{0} & \SetCell{bg=\HeatCellBg{62}}\CellText{0.62}{0}{0} & \SetCell{bg=\HeatCellBg{55}}\CellText{0.55}{0}{0} & \SetCell{bg=\HeatCellBg{58}}\CellText{0.58}{1}{0} & \MeanCell{0.62}{0.14}{0} \\
        2 & \SetCell{bg=\HeatCellBg{100}}\CellText{1.00}{1}{1} & \SetCell{bg=\HeatCellBg{51}}\CellText{0.51}{0}{0} & \SetCell{bg=\HeatCellBg{91}}\CellText{0.91}{1}{0} & \SetCell{bg=\HeatCellBg{77}}\CellText{0.77}{1}{0} & \SetCell{bg=\HeatCellBg{72}}\CellText{0.72}{1}{0} & \SetCell{bg=\HeatCellBg{70}}\CellText{0.70}{1}{0} & \SetCell{bg=\HeatCellBg{65}}\CellText{0.65}{0}{0} & \SetCell{bg=\HeatCellBg{74}}\CellText{0.74}{1}{0} & \SetCell{bg=\HeatCellBg{79}}\CellText{0.79}{1}{0} & \SetCell{bg=\HeatCellBg{65}}\CellText{0.65}{1}{0} & \SetCell{bg=\HeatCellBg{66}}\CellText{0.66}{1}{0} & \SetCell{bg=\HeatCellBg{62}}\CellText{0.62}{1}{0} & \SetCell{bg=\HeatCellBg{68}}\CellText{0.68}{1}{0} & \SetCell{bg=\HeatCellBg{62}}\CellText{0.62}{1}{0} & \SetCell{bg=\HeatCellBg{57}}\CellText{0.57}{1}{0} & \SetCell{bg=\HeatCellBg{59}}\CellText{0.59}{1}{0} & \SetCell{bg=\HeatCellBg{65}}\CellText{0.65}{1}{0} & \SetCell{bg=\HeatCellBg{59}}\CellText{0.59}{1}{0} & \SetCell{bg=\HeatCellBg{57}}\CellText{0.57}{0}{0} & \MeanCell{0.68}{0.12}{1} \\
        3 & \SetCell{bg=\HeatCellBg{99}}\CellText{0.99}{0}{1} & \SetCell{bg=\HeatCellBg{84}}\CellText{0.84}{0}{0} & \SetCell{bg=\HeatCellBg{58}}\CellText{0.58}{0}{0} & \SetCell{bg=\HeatCellBg{75}}\CellText{0.75}{0}{0} & \SetCell{bg=\HeatCellBg{54}}\CellText{0.54}{0}{0} & \SetCell{bg=\HeatCellBg{57}}\CellText{0.57}{0}{0} & \SetCell{bg=\HeatCellBg{57}}\CellText{0.57}{0}{0} & \SetCell{bg=\HeatCellBg{66}}\CellText{0.66}{0}{0} & \SetCell{bg=\HeatCellBg{62}}\CellText{0.62}{0}{0} & \SetCell{bg=\HeatCellBg{62}}\CellText{0.62}{0}{0} & \SetCell{bg=\HeatCellBg{58}}\CellText{0.58}{0}{0} & \SetCell{bg=\HeatCellBg{62}}\CellText{0.62}{1}{0} & \SetCell{bg=\HeatCellBg{59}}\CellText{0.59}{0}{0} & \SetCell{bg=\HeatCellBg{55}}\CellText{0.55}{0}{0} & \SetCell{bg=\HeatCellBg{50}}\CellText{0.50}{0}{0} & \SetCell{bg=\HeatCellBg{56}}\CellText{0.56}{0}{0} & \SetCell{bg=\HeatCellBg{59}}\CellText{0.59}{0}{0} & \SetCell{bg=\HeatCellBg{54}}\CellText{0.54}{0}{0} & \SetCell{bg=\HeatCellBg{52}}\CellText{0.52}{0}{0} & \MeanCell{0.62}{0.12}{0} \\
        4 & \SetCell{bg=\HeatCellBg{97}}\CellText{0.97}{0}{1} & \SetCell{bg=\HeatCellBg{81}}\CellText{0.81}{0}{0} & \SetCell{bg=\HeatCellBg{50}}\CellText{0.50}{0}{0} & \SetCell{bg=\HeatCellBg{53}}\CellText{0.53}{0}{0} & \SetCell{bg=\HeatCellBg{63}}\CellText{0.63}{0}{0} & \SetCell{bg=\HeatCellBg{63}}\CellText{0.63}{0}{0} & \SetCell{bg=\HeatCellBg{51}}\CellText{0.51}{0}{0} & \SetCell{bg=\HeatCellBg{47}}\CellText{0.47}{0}{0} & \SetCell{bg=\HeatCellBg{48}}\CellText{0.48}{0}{0} & \SetCell{bg=\HeatCellBg{46}}\CellText{0.46}{0}{0} & \SetCell{bg=\HeatCellBg{50}}\CellText{0.50}{0}{0} & \SetCell{bg=\HeatCellBg{47}}\CellText{0.47}{0}{0} & \SetCell{bg=\HeatCellBg{44}}\CellText{0.44}{0}{0} & \SetCell{bg=\HeatCellBg{50}}\CellText{0.50}{0}{0} & \SetCell{bg=\HeatCellBg{48}}\CellText{0.48}{0}{0} & \SetCell{bg=\HeatCellBg{50}}\CellText{0.50}{0}{0} & \SetCell{bg=\HeatCellBg{49}}\CellText{0.49}{0}{0} & \SetCell{bg=\HeatCellBg{49}}\CellText{0.49}{0}{0} & \SetCell{bg=\HeatCellBg{53}}\CellText{0.53}{0}{0} & \MeanCell{0.55}{0.13}{0} \\
        5 & \SetCell{bg=\HeatCellBg{97}}\CellText{0.97}{0}{1} & \SetCell{bg=\HeatCellBg{81}}\CellText{0.81}{0}{0} & \SetCell{bg=\HeatCellBg{50}}\CellText{0.50}{0}{0} & \SetCell{bg=\HeatCellBg{53}}\CellText{0.53}{0}{0} & \SetCell{bg=\HeatCellBg{63}}\CellText{0.63}{0}{0} & \SetCell{bg=\HeatCellBg{63}}\CellText{0.63}{0}{0} & \SetCell{bg=\HeatCellBg{51}}\CellText{0.51}{0}{0} & \SetCell{bg=\HeatCellBg{47}}\CellText{0.47}{0}{0} & \SetCell{bg=\HeatCellBg{48}}\CellText{0.48}{0}{0} & \SetCell{bg=\HeatCellBg{46}}\CellText{0.46}{0}{0} & \SetCell{bg=\HeatCellBg{50}}\CellText{0.50}{0}{0} & \SetCell{bg=\HeatCellBg{47}}\CellText{0.47}{0}{0} & \SetCell{bg=\HeatCellBg{44}}\CellText{0.44}{0}{0} & \SetCell{bg=\HeatCellBg{50}}\CellText{0.50}{0}{0} & \SetCell{bg=\HeatCellBg{48}}\CellText{0.48}{0}{0} & \SetCell{bg=\HeatCellBg{50}}\CellText{0.50}{0}{0} & \SetCell{bg=\HeatCellBg{49}}\CellText{0.49}{0}{0} & \SetCell{bg=\HeatCellBg{49}}\CellText{0.49}{0}{0} & \SetCell{bg=\HeatCellBg{53}}\CellText{0.53}{0}{0} & \MeanCell{0.55}{0.13}{0} \\
        6b & \SetCell{bg=\HeatCellBg{94}}\CellText{0.94}{0}{1} & \SetCell{bg=\HeatCellBg{47}}\CellText{0.47}{0}{0} & \SetCell{bg=\HeatCellBg{39}}\CellText{0.39}{0}{0} & \SetCell{bg=\HeatCellBg{42}}\CellText{0.42}{0}{0} & \SetCell{bg=\HeatCellBg{54}}\CellText{0.54}{0}{0} & \SetCell{bg=\HeatCellBg{52}}\CellText{0.52}{0}{0} & \SetCell{bg=\HeatCellBg{50}}\CellText{0.50}{0}{0} & \SetCell{bg=\HeatCellBg{47}}\CellText{0.47}{0}{0} & \SetCell{bg=\HeatCellBg{47}}\CellText{0.47}{0}{0} & \SetCell{bg=\HeatCellBg{47}}\CellText{0.47}{0}{0} & \SetCell{bg=\HeatCellBg{48}}\CellText{0.48}{0}{0} & \SetCell{bg=\HeatCellBg{44}}\CellText{0.44}{0}{0} & \SetCell{bg=\HeatCellBg{41}}\CellText{0.41}{0}{0} & \SetCell{bg=\HeatCellBg{47}}\CellText{0.47}{0}{0} & \SetCell{bg=\HeatCellBg{43}}\CellText{0.43}{0}{0} & \SetCell{bg=\HeatCellBg{43}}\CellText{0.43}{0}{0} & \SetCell{bg=\HeatCellBg{40}}\CellText{0.40}{0}{0} & \SetCell{bg=\HeatCellBg{42}}\CellText{0.42}{0}{0} & \SetCell{bg=\HeatCellBg{42}}\CellText{0.42}{0}{0} & \MeanCell{0.48}{0.12}{0} \\
        7b & \SetCell{bg=\HeatCellBg{53}}\CellText{0.53}{0}{1} & \SetCell{bg=\HeatCellBg{41}}\CellText{0.41}{0}{0} & \SetCell{bg=\HeatCellBg{34}}\CellText{0.34}{0}{0} & \SetCell{bg=\HeatCellBg{42}}\CellText{0.42}{0}{0} & \SetCell{bg=\HeatCellBg{45}}\CellText{0.45}{0}{0} & \SetCell{bg=\HeatCellBg{38}}\CellText{0.38}{0}{0} & \SetCell{bg=\HeatCellBg{39}}\CellText{0.39}{0}{0} & \SetCell{bg=\HeatCellBg{38}}\CellText{0.38}{0}{0} & \SetCell{bg=\HeatCellBg{35}}\CellText{0.35}{0}{0} & \SetCell{bg=\HeatCellBg{40}}\CellText{0.40}{0}{0} & \SetCell{bg=\HeatCellBg{42}}\CellText{0.42}{0}{0} & \SetCell{bg=\HeatCellBg{40}}\CellText{0.40}{0}{0} & \SetCell{bg=\HeatCellBg{37}}\CellText{0.37}{0}{0} & \SetCell{bg=\HeatCellBg{40}}\CellText{0.40}{0}{0} & \SetCell{bg=\HeatCellBg{40}}\CellText{0.40}{0}{0} & \SetCell{bg=\HeatCellBg{37}}\CellText{0.37}{0}{0} & \SetCell{bg=\HeatCellBg{36}}\CellText{0.36}{0}{0} & \SetCell{bg=\HeatCellBg{38}}\CellText{0.38}{0}{0} & \SetCell{bg=\HeatCellBg{39}}\CellText{0.39}{0}{0} & \MeanCell{0.40}{0.04}{0} \\
        8b & \SetCell{bg=\HeatCellBg{68}}\CellText{0.68}{0}{1} & \SetCell{bg=\HeatCellBg{38}}\CellText{0.38}{0}{0} & \SetCell{bg=\HeatCellBg{49}}\CellText{0.49}{0}{0} & \SetCell{bg=\HeatCellBg{36}}\CellText{0.36}{0}{0} & \SetCell{bg=\HeatCellBg{35}}\CellText{0.35}{0}{0} & \SetCell{bg=\HeatCellBg{33}}\CellText{0.33}{0}{0} & \SetCell{bg=\HeatCellBg{32}}\CellText{0.32}{0}{0} & \SetCell{bg=\HeatCellBg{36}}\CellText{0.36}{0}{0} & \SetCell{bg=\HeatCellBg{36}}\CellText{0.36}{0}{0} & \SetCell{bg=\HeatCellBg{35}}\CellText{0.35}{0}{0} & \SetCell{bg=\HeatCellBg{36}}\CellText{0.36}{0}{0} & \SetCell{bg=\HeatCellBg{33}}\CellText{0.33}{0}{0} & \SetCell{bg=\HeatCellBg{35}}\CellText{0.35}{0}{0} & \SetCell{bg=\HeatCellBg{33}}\CellText{0.33}{0}{0} & \SetCell{bg=\HeatCellBg{36}}\CellText{0.36}{0}{0} & \SetCell{bg=\HeatCellBg{33}}\CellText{0.33}{0}{0} & \SetCell{bg=\HeatCellBg{33}}\CellText{0.33}{0}{0} & \SetCell{bg=\HeatCellBg{33}}\CellText{0.33}{0}{0} & \SetCell{bg=\HeatCellBg{34}}\CellText{0.34}{0}{0} & \MeanCell{0.37}{0.08}{0} \\
        Mean $\pm$ Std & \MeanCell{0.89}{0.17}{1} & \MeanCell{0.65}{0.21}{0} & \MeanCell{0.57}{0.19}{0} & \MeanCell{0.55}{0.15}{0} & \MeanCell{0.55}{0.11}{0} & \MeanCell{0.54}{0.12}{0} & \MeanCell{0.52}{0.11}{0} & \MeanCell{0.50}{0.12}{0} & \MeanCell{0.52}{0.14}{0} & \MeanCell{0.49}{0.10}{0} & \MeanCell{0.50}{0.09}{0} & \MeanCell{0.49}{0.10}{0} & \MeanCell{0.49}{0.12}{0} & \MeanCell{0.50}{0.09}{0} & \MeanCell{0.47}{0.07}{0} & \MeanCell{0.47}{0.08}{0} & \MeanCell{0.49}{0.11}{0} & \MeanCell{0.47}{0.08}{0} & \MeanCell{0.48}{0.08}{0} & \MeanCell{0.53}{0.15}{0} \\
}

\ClusterKTable{10}{\gls{as} --- \MarioBrosName.}{as_mario_bros}{
        1 & \SetCell{bg=\HeatCellBg{1}}\CellText{0.01}{0}{0} & \SetCell{bg=\HeatCellBg{11}}\CellText{0.11}{0}{0} & \SetCell{bg=\HeatCellBg{23}}\CellText{0.23}{0}{0} & \SetCell{bg=\HeatCellBg{50}}\CellText{0.50}{0}{0} & \SetCell{bg=\HeatCellBg{46}}\CellText{0.46}{0}{0} & \SetCell{bg=\HeatCellBg{44}}\CellText{0.44}{0}{0} & \SetCell{bg=\HeatCellBg{49}}\CellText{0.49}{0}{0} & \SetCell{bg=\HeatCellBg{48}}\CellText{0.48}{0}{0} & \SetCell{bg=\HeatCellBg{58}}\CellText{0.58}{0}{1} & \SetCell{bg=\HeatCellBg{51}}\CellText{0.51}{0}{0} & \SetCell{bg=\HeatCellBg{49}}\CellText{0.49}{0}{0} & \SetCell{bg=\HeatCellBg{50}}\CellText{0.50}{0}{0} & \SetCell{bg=\HeatCellBg{44}}\CellText{0.44}{0}{0} & \SetCell{bg=\HeatCellBg{46}}\CellText{0.46}{0}{0} & \SetCell{bg=\HeatCellBg{58}}\CellText{0.58}{0}{1} & \SetCell{bg=\HeatCellBg{54}}\CellText{0.54}{0}{0} & \SetCell{bg=\HeatCellBg{52}}\CellText{0.52}{0}{0} & \SetCell{bg=\HeatCellBg{50}}\CellText{0.50}{0}{0} & \SetCell{bg=\HeatCellBg{48}}\CellText{0.48}{0}{0} & \MeanCell{0.44}{0.15}{0} \\
        2 & \SetCell{bg=\HeatCellBg{0}}\CellText{0.00}{0}{0} & \SetCell{bg=\HeatCellBg{66}}\CellText{0.66}{0}{1} & \SetCell{bg=\HeatCellBg{47}}\CellText{0.47}{0}{0} & \SetCell{bg=\HeatCellBg{55}}\CellText{0.55}{0}{0} & \SetCell{bg=\HeatCellBg{49}}\CellText{0.49}{0}{0} & \SetCell{bg=\HeatCellBg{58}}\CellText{0.58}{0}{0} & \SetCell{bg=\HeatCellBg{53}}\CellText{0.53}{0}{0} & \SetCell{bg=\HeatCellBg{48}}\CellText{0.48}{0}{0} & \SetCell{bg=\HeatCellBg{62}}\CellText{0.62}{0}{0} & \SetCell{bg=\HeatCellBg{62}}\CellText{0.62}{0}{0} & \SetCell{bg=\HeatCellBg{54}}\CellText{0.54}{0}{0} & \SetCell{bg=\HeatCellBg{62}}\CellText{0.62}{0}{0} & \SetCell{bg=\HeatCellBg{56}}\CellText{0.56}{0}{0} & \SetCell{bg=\HeatCellBg{57}}\CellText{0.57}{0}{0} & \SetCell{bg=\HeatCellBg{56}}\CellText{0.56}{0}{0} & \SetCell{bg=\HeatCellBg{59}}\CellText{0.59}{0}{0} & \SetCell{bg=\HeatCellBg{61}}\CellText{0.61}{0}{0} & \SetCell{bg=\HeatCellBg{56}}\CellText{0.56}{0}{0} & \SetCell{bg=\HeatCellBg{57}}\CellText{0.57}{0}{0} & \MeanCell{0.54}{0.14}{0} \\
        3 & \SetCell{bg=\HeatCellBg{0}}\CellText{0.00}{0}{0} & \SetCell{bg=\HeatCellBg{66}}\CellText{0.66}{0}{0} & \SetCell{bg=\HeatCellBg{63}}\CellText{0.63}{0}{0} & \SetCell{bg=\HeatCellBg{65}}\CellText{0.65}{0}{0} & \SetCell{bg=\HeatCellBg{64}}\CellText{0.64}{0}{0} & \SetCell{bg=\HeatCellBg{62}}\CellText{0.62}{0}{0} & \SetCell{bg=\HeatCellBg{63}}\CellText{0.63}{0}{0} & \SetCell{bg=\HeatCellBg{62}}\CellText{0.62}{0}{0} & \SetCell{bg=\HeatCellBg{69}}\CellText{0.69}{0}{1} & \SetCell{bg=\HeatCellBg{69}}\CellText{0.69}{0}{1} & \SetCell{bg=\HeatCellBg{64}}\CellText{0.64}{0}{0} & \SetCell{bg=\HeatCellBg{66}}\CellText{0.66}{0}{0} & \SetCell{bg=\HeatCellBg{62}}\CellText{0.62}{0}{0} & \SetCell{bg=\HeatCellBg{64}}\CellText{0.64}{0}{0} & \SetCell{bg=\HeatCellBg{67}}\CellText{0.67}{0}{0} & \SetCell{bg=\HeatCellBg{65}}\CellText{0.65}{0}{0} & \SetCell{bg=\HeatCellBg{66}}\CellText{0.66}{0}{0} & \SetCell{bg=\HeatCellBg{64}}\CellText{0.64}{0}{0} & \SetCell{bg=\HeatCellBg{69}}\CellText{0.69}{0}{1} & \MeanCell{0.62}{0.15}{0} \\
        4 & \SetCell{bg=\HeatCellBg{2}}\CellText{0.02}{0}{0} & \SetCell{bg=\HeatCellBg{53}}\CellText{0.53}{0}{0} & \SetCell{bg=\HeatCellBg{67}}\CellText{0.67}{0}{0} & \SetCell{bg=\HeatCellBg{71}}\CellText{0.71}{0}{0} & \SetCell{bg=\HeatCellBg{64}}\CellText{0.64}{0}{0} & \SetCell{bg=\HeatCellBg{69}}\CellText{0.69}{0}{0} & \SetCell{bg=\HeatCellBg{65}}\CellText{0.65}{0}{0} & \SetCell{bg=\HeatCellBg{65}}\CellText{0.65}{0}{0} & \SetCell{bg=\HeatCellBg{71}}\CellText{0.71}{0}{0} & \SetCell{bg=\HeatCellBg{72}}\CellText{0.72}{0}{0} & \SetCell{bg=\HeatCellBg{71}}\CellText{0.71}{0}{0} & \SetCell{bg=\HeatCellBg{76}}\CellText{0.76}{0}{1} & \SetCell{bg=\HeatCellBg{72}}\CellText{0.72}{0}{0} & \SetCell{bg=\HeatCellBg{74}}\CellText{0.74}{0}{0} & \SetCell{bg=\HeatCellBg{73}}\CellText{0.73}{0}{0} & \SetCell{bg=\HeatCellBg{76}}\CellText{0.76}{0}{1} & \SetCell{bg=\HeatCellBg{75}}\CellText{0.75}{0}{0} & \SetCell{bg=\HeatCellBg{72}}\CellText{0.72}{0}{0} & \SetCell{bg=\HeatCellBg{75}}\CellText{0.75}{0}{0} & \MeanCell{0.66}{0.16}{0} \\
        5 & \SetCell{bg=\HeatCellBg{2}}\CellText{0.02}{0}{0} & \SetCell{bg=\HeatCellBg{53}}\CellText{0.53}{0}{0} & \SetCell{bg=\HeatCellBg{67}}\CellText{0.67}{0}{0} & \SetCell{bg=\HeatCellBg{71}}\CellText{0.71}{0}{0} & \SetCell{bg=\HeatCellBg{64}}\CellText{0.64}{0}{0} & \SetCell{bg=\HeatCellBg{69}}\CellText{0.69}{0}{0} & \SetCell{bg=\HeatCellBg{65}}\CellText{0.65}{0}{0} & \SetCell{bg=\HeatCellBg{65}}\CellText{0.65}{0}{0} & \SetCell{bg=\HeatCellBg{71}}\CellText{0.71}{0}{0} & \SetCell{bg=\HeatCellBg{72}}\CellText{0.72}{0}{0} & \SetCell{bg=\HeatCellBg{71}}\CellText{0.71}{0}{0} & \SetCell{bg=\HeatCellBg{76}}\CellText{0.76}{0}{1} & \SetCell{bg=\HeatCellBg{72}}\CellText{0.72}{0}{0} & \SetCell{bg=\HeatCellBg{74}}\CellText{0.74}{0}{0} & \SetCell{bg=\HeatCellBg{73}}\CellText{0.73}{0}{0} & \SetCell{bg=\HeatCellBg{76}}\CellText{0.76}{0}{1} & \SetCell{bg=\HeatCellBg{75}}\CellText{0.75}{0}{0} & \SetCell{bg=\HeatCellBg{72}}\CellText{0.72}{0}{0} & \SetCell{bg=\HeatCellBg{75}}\CellText{0.75}{0}{0} & \MeanCell{0.66}{0.16}{0} \\
        6b & \SetCell{bg=\HeatCellBg{5}}\CellText{0.05}{0}{0} & \SetCell{bg=\HeatCellBg{66}}\CellText{0.66}{0}{0} & \SetCell{bg=\HeatCellBg{68}}\CellText{0.68}{0}{0} & \SetCell{bg=\HeatCellBg{72}}\CellText{0.72}{0}{0} & \SetCell{bg=\HeatCellBg{71}}\CellText{0.71}{0}{0} & \SetCell{bg=\HeatCellBg{76}}\CellText{0.76}{0}{0} & \SetCell{bg=\HeatCellBg{72}}\CellText{0.72}{0}{0} & \SetCell{bg=\HeatCellBg{72}}\CellText{0.72}{0}{0} & \SetCell{bg=\HeatCellBg{74}}\CellText{0.74}{0}{0} & \SetCell{bg=\HeatCellBg{76}}\CellText{0.76}{0}{0} & \SetCell{bg=\HeatCellBg{76}}\CellText{0.76}{0}{0} & \SetCell{bg=\HeatCellBg{79}}\CellText{0.79}{0}{1} & \SetCell{bg=\HeatCellBg{76}}\CellText{0.76}{0}{0} & \SetCell{bg=\HeatCellBg{78}}\CellText{0.78}{0}{0} & \SetCell{bg=\HeatCellBg{78}}\CellText{0.78}{0}{0} & \SetCell{bg=\HeatCellBg{79}}\CellText{0.79}{0}{1} & \SetCell{bg=\HeatCellBg{79}}\CellText{0.79}{0}{1} & \SetCell{bg=\HeatCellBg{75}}\CellText{0.75}{0}{0} & \SetCell{bg=\HeatCellBg{78}}\CellText{0.78}{0}{0} & \MeanCell{0.71}{0.16}{0} \\
        7b & \SetCell{bg=\HeatCellBg{76}}\CellText{0.76}{1}{0} & \SetCell{bg=\HeatCellBg{85}}\CellText{0.85}{1}{0} & \SetCell{bg=\HeatCellBg{86}}\CellText{0.86}{1}{1} & \SetCell{bg=\HeatCellBg{83}}\CellText{0.83}{1}{0} & \SetCell{bg=\HeatCellBg{83}}\CellText{0.83}{1}{0} & \SetCell{bg=\HeatCellBg{85}}\CellText{0.85}{1}{0} & \SetCell{bg=\HeatCellBg{83}}\CellText{0.83}{1}{0} & \SetCell{bg=\HeatCellBg{83}}\CellText{0.83}{1}{0} & \SetCell{bg=\HeatCellBg{83}}\CellText{0.83}{1}{0} & \SetCell{bg=\HeatCellBg{82}}\CellText{0.82}{1}{0} & \SetCell{bg=\HeatCellBg{82}}\CellText{0.82}{1}{0} & \SetCell{bg=\HeatCellBg{83}}\CellText{0.83}{1}{0} & \SetCell{bg=\HeatCellBg{82}}\CellText{0.82}{1}{0} & \SetCell{bg=\HeatCellBg{84}}\CellText{0.84}{1}{0} & \SetCell{bg=\HeatCellBg{84}}\CellText{0.84}{1}{0} & \SetCell{bg=\HeatCellBg{86}}\CellText{0.86}{1}{1} & \SetCell{bg=\HeatCellBg{84}}\CellText{0.84}{1}{0} & \SetCell{bg=\HeatCellBg{84}}\CellText{0.84}{1}{0} & \SetCell{bg=\HeatCellBg{84}}\CellText{0.84}{1}{0} & \MeanCell{0.83}{0.02}{1} \\
        8b & \SetCell{bg=\HeatCellBg{39}}\CellText{0.39}{0}{0} & \SetCell{bg=\HeatCellBg{55}}\CellText{0.55}{0}{0} & \SetCell{bg=\HeatCellBg{64}}\CellText{0.64}{0}{0} & \SetCell{bg=\HeatCellBg{68}}\CellText{0.68}{0}{0} & \SetCell{bg=\HeatCellBg{67}}\CellText{0.67}{0}{0} & \SetCell{bg=\HeatCellBg{68}}\CellText{0.68}{0}{0} & \SetCell{bg=\HeatCellBg{69}}\CellText{0.69}{0}{0} & \SetCell{bg=\HeatCellBg{69}}\CellText{0.69}{0}{0} & \SetCell{bg=\HeatCellBg{68}}\CellText{0.68}{0}{0} & \SetCell{bg=\HeatCellBg{71}}\CellText{0.71}{0}{0} & \SetCell{bg=\HeatCellBg{68}}\CellText{0.68}{0}{0} & \SetCell{bg=\HeatCellBg{74}}\CellText{0.74}{0}{0} & \SetCell{bg=\HeatCellBg{74}}\CellText{0.74}{0}{0} & \SetCell{bg=\HeatCellBg{74}}\CellText{0.74}{0}{0} & \SetCell{bg=\HeatCellBg{73}}\CellText{0.73}{0}{0} & \SetCell{bg=\HeatCellBg{74}}\CellText{0.74}{0}{0} & \SetCell{bg=\HeatCellBg{75}}\CellText{0.75}{0}{0} & \SetCell{bg=\HeatCellBg{75}}\CellText{0.75}{0}{0} & \SetCell{bg=\HeatCellBg{76}}\CellText{0.76}{0}{1} & \MeanCell{0.68}{0.08}{0} \\
        Mean $\pm$ Std & \MeanCell{0.16}{0.26}{0} & \MeanCell{0.57}{0.20}{0} & \MeanCell{0.61}{0.17}{0} & \MeanCell{0.67}{0.10}{0} & \MeanCell{0.64}{0.11}{0} & \MeanCell{0.66}{0.11}{0} & \MeanCell{0.65}{0.10}{0} & \MeanCell{0.64}{0.11}{0} & \MeanCell{0.69}{0.07}{0} & \MeanCell{0.69}{0.09}{0} & \MeanCell{0.67}{0.10}{0} & \MeanCell{0.71}{0.10}{0} & \MeanCell{0.67}{0.12}{0} & \MeanCell{0.69}{0.12}{0} & \MeanCell{0.70}{0.09}{0} & \MeanCell{0.71}{0.10}{1} & \MeanCell{0.71}{0.10}{0} & \MeanCell{0.69}{0.10}{0} & \MeanCell{0.70}{0.11}{0} & \MeanCell{0.64}{0.17}{0} \\
}

\ClusterKTable{10}{\gls{acu} --- \MarioBrosName.}{acu_mario_bros}{
        1 & \SetCell{bg=\HeatCellBg{0}}\CellText{0.00}{0}{0} & \SetCell{bg=\HeatCellBg{4}}\CellText{0.04}{0}{0} & \SetCell{bg=\HeatCellBg{5}}\CellText{0.05}{0}{0} & \SetCell{bg=\HeatCellBg{16}}\CellText{0.16}{0}{1} & \SetCell{bg=\HeatCellBg{-3}}\CellText{-0.03}{0}{0} & \SetCell{bg=\HeatCellBg{-3}}\CellText{-0.03}{0}{0} & \SetCell{bg=\HeatCellBg{16}}\CellText{0.16}{0}{1} & \SetCell{bg=\HeatCellBg{-7}}\CellText{-0.07}{0}{0} & \SetCell{bg=\HeatCellBg{16}}\CellText{0.16}{0}{1} & \SetCell{bg=\HeatCellBg{5}}\CellText{0.05}{0}{0} & \SetCell{bg=\HeatCellBg{1}}\CellText{0.01}{0}{0} & \SetCell{bg=\HeatCellBg{4}}\CellText{0.04}{0}{0} & \SetCell{bg=\HeatCellBg{6}}\CellText{0.06}{0}{0} & \SetCell{bg=\HeatCellBg{6}}\CellText{0.06}{0}{0} & \SetCell{bg=\HeatCellBg{12}}\CellText{0.12}{0}{0} & \SetCell{bg=\HeatCellBg{3}}\CellText{0.03}{0}{0} & \SetCell{bg=\HeatCellBg{14}}\CellText{0.14}{0}{0} & \SetCell{bg=\HeatCellBg{5}}\CellText{0.05}{0}{0} & \SetCell{bg=\HeatCellBg{6}}\CellText{0.06}{0}{0} & \MeanCell{0.06}{0.07}{0} \\
        2 & \SetCell{bg=\HeatCellBg{0}}\CellText{0.00}{0}{0} & \SetCell{bg=\HeatCellBg{16}}\CellText{0.16}{0}{0} & \SetCell{bg=\HeatCellBg{38}}\CellText{0.38}{1}{0} & \SetCell{bg=\HeatCellBg{32}}\CellText{0.32}{0}{0} & \SetCell{bg=\HeatCellBg{21}}\CellText{0.21}{0}{0} & \SetCell{bg=\HeatCellBg{29}}\CellText{0.29}{0}{0} & \SetCell{bg=\HeatCellBg{18}}\CellText{0.18}{0}{0} & \SetCell{bg=\HeatCellBg{22}}\CellText{0.22}{0}{0} & \SetCell{bg=\HeatCellBg{41}}\CellText{0.41}{1}{1} & \SetCell{bg=\HeatCellBg{28}}\CellText{0.28}{0}{0} & \SetCell{bg=\HeatCellBg{20}}\CellText{0.20}{0}{0} & \SetCell{bg=\HeatCellBg{24}}\CellText{0.24}{0}{0} & \SetCell{bg=\HeatCellBg{24}}\CellText{0.24}{1}{0} & \SetCell{bg=\HeatCellBg{19}}\CellText{0.19}{0}{0} & \SetCell{bg=\HeatCellBg{13}}\CellText{0.13}{0}{0} & \SetCell{bg=\HeatCellBg{18}}\CellText{0.18}{0}{0} & \SetCell{bg=\HeatCellBg{26}}\CellText{0.26}{1}{0} & \SetCell{bg=\HeatCellBg{16}}\CellText{0.16}{0}{0} & \SetCell{bg=\HeatCellBg{15}}\CellText{0.15}{0}{0} & \MeanCell{0.22}{0.09}{0} \\
        3 & \SetCell{bg=\HeatCellBg{0}}\CellText{0.00}{0}{0} & \SetCell{bg=\HeatCellBg{50}}\CellText{0.50}{1}{1} & \SetCell{bg=\HeatCellBg{21}}\CellText{0.21}{0}{0} & \SetCell{bg=\HeatCellBg{40}}\CellText{0.40}{1}{0} & \SetCell{bg=\HeatCellBg{18}}\CellText{0.18}{0}{0} & \SetCell{bg=\HeatCellBg{20}}\CellText{0.20}{0}{0} & \SetCell{bg=\HeatCellBg{20}}\CellText{0.20}{0}{0} & \SetCell{bg=\HeatCellBg{27}}\CellText{0.27}{1}{0} & \SetCell{bg=\HeatCellBg{31}}\CellText{0.31}{0}{0} & \SetCell{bg=\HeatCellBg{31}}\CellText{0.31}{1}{0} & \SetCell{bg=\HeatCellBg{23}}\CellText{0.23}{0}{0} & \SetCell{bg=\HeatCellBg{28}}\CellText{0.28}{1}{0} & \SetCell{bg=\HeatCellBg{21}}\CellText{0.21}{0}{0} & \SetCell{bg=\HeatCellBg{19}}\CellText{0.19}{0}{0} & \SetCell{bg=\HeatCellBg{16}}\CellText{0.16}{0}{0} & \SetCell{bg=\HeatCellBg{22}}\CellText{0.22}{0}{0} & \SetCell{bg=\HeatCellBg{25}}\CellText{0.25}{0}{0} & \SetCell{bg=\HeatCellBg{18}}\CellText{0.18}{0}{0} & \SetCell{bg=\HeatCellBg{21}}\CellText{0.21}{0}{0} & \MeanCell{0.24}{0.10}{1} \\
        4 & \SetCell{bg=\HeatCellBg{-1}}\CellText{-0.01}{0}{0} & \SetCell{bg=\HeatCellBg{34}}\CellText{0.34}{0}{1} & \SetCell{bg=\HeatCellBg{17}}\CellText{0.17}{0}{0} & \SetCell{bg=\HeatCellBg{24}}\CellText{0.24}{0}{0} & \SetCell{bg=\HeatCellBg{27}}\CellText{0.27}{0}{0} & \SetCell{bg=\HeatCellBg{32}}\CellText{0.32}{1}{0} & \SetCell{bg=\HeatCellBg{16}}\CellText{0.16}{0}{0} & \SetCell{bg=\HeatCellBg{12}}\CellText{0.12}{0}{0} & \SetCell{bg=\HeatCellBg{19}}\CellText{0.19}{0}{0} & \SetCell{bg=\HeatCellBg{18}}\CellText{0.18}{0}{0} & \SetCell{bg=\HeatCellBg{20}}\CellText{0.20}{0}{0} & \SetCell{bg=\HeatCellBg{23}}\CellText{0.23}{0}{0} & \SetCell{bg=\HeatCellBg{17}}\CellText{0.17}{0}{0} & \SetCell{bg=\HeatCellBg{23}}\CellText{0.23}{0}{0} & \SetCell{bg=\HeatCellBg{21}}\CellText{0.21}{0}{0} & \SetCell{bg=\HeatCellBg{26}}\CellText{0.26}{1}{0} & \SetCell{bg=\HeatCellBg{24}}\CellText{0.24}{0}{0} & \SetCell{bg=\HeatCellBg{21}}\CellText{0.21}{0}{0} & \SetCell{bg=\HeatCellBg{28}}\CellText{0.28}{1}{0} & \MeanCell{0.21}{0.07}{0} \\
        5 & \SetCell{bg=\HeatCellBg{-1}}\CellText{-0.01}{0}{0} & \SetCell{bg=\HeatCellBg{34}}\CellText{0.34}{0}{1} & \SetCell{bg=\HeatCellBg{17}}\CellText{0.17}{0}{0} & \SetCell{bg=\HeatCellBg{24}}\CellText{0.24}{0}{0} & \SetCell{bg=\HeatCellBg{27}}\CellText{0.27}{0}{0} & \SetCell{bg=\HeatCellBg{32}}\CellText{0.32}{1}{0} & \SetCell{bg=\HeatCellBg{16}}\CellText{0.16}{0}{0} & \SetCell{bg=\HeatCellBg{12}}\CellText{0.12}{0}{0} & \SetCell{bg=\HeatCellBg{19}}\CellText{0.19}{0}{0} & \SetCell{bg=\HeatCellBg{18}}\CellText{0.18}{0}{0} & \SetCell{bg=\HeatCellBg{20}}\CellText{0.20}{0}{0} & \SetCell{bg=\HeatCellBg{23}}\CellText{0.23}{0}{0} & \SetCell{bg=\HeatCellBg{17}}\CellText{0.17}{0}{0} & \SetCell{bg=\HeatCellBg{23}}\CellText{0.23}{0}{0} & \SetCell{bg=\HeatCellBg{21}}\CellText{0.21}{0}{0} & \SetCell{bg=\HeatCellBg{26}}\CellText{0.26}{1}{0} & \SetCell{bg=\HeatCellBg{24}}\CellText{0.24}{0}{0} & \SetCell{bg=\HeatCellBg{21}}\CellText{0.21}{0}{0} & \SetCell{bg=\HeatCellBg{28}}\CellText{0.28}{1}{0} & \MeanCell{0.21}{0.07}{0} \\
        6b & \SetCell{bg=\HeatCellBg{-2}}\CellText{-0.02}{0}{0} & \SetCell{bg=\HeatCellBg{13}}\CellText{0.13}{0}{0} & \SetCell{bg=\HeatCellBg{7}}\CellText{0.07}{0}{0} & \SetCell{bg=\HeatCellBg{14}}\CellText{0.14}{0}{0} & \SetCell{bg=\HeatCellBg{25}}\CellText{0.25}{0}{0} & \SetCell{bg=\HeatCellBg{28}}\CellText{0.28}{0}{1} & \SetCell{bg=\HeatCellBg{22}}\CellText{0.22}{1}{0} & \SetCell{bg=\HeatCellBg{19}}\CellText{0.19}{0}{0} & \SetCell{bg=\HeatCellBg{21}}\CellText{0.21}{0}{0} & \SetCell{bg=\HeatCellBg{23}}\CellText{0.23}{0}{0} & \SetCell{bg=\HeatCellBg{25}}\CellText{0.25}{1}{0} & \SetCell{bg=\HeatCellBg{23}}\CellText{0.23}{0}{0} & \SetCell{bg=\HeatCellBg{18}}\CellText{0.18}{0}{0} & \SetCell{bg=\HeatCellBg{25}}\CellText{0.25}{1}{0} & \SetCell{bg=\HeatCellBg{21}}\CellText{0.21}{0}{0} & \SetCell{bg=\HeatCellBg{22}}\CellText{0.22}{0}{0} & \SetCell{bg=\HeatCellBg{20}}\CellText{0.20}{0}{0} & \SetCell{bg=\HeatCellBg{17}}\CellText{0.17}{0}{0} & \SetCell{bg=\HeatCellBg{20}}\CellText{0.20}{0}{0} & \MeanCell{0.19}{0.07}{0} \\
        7b & \SetCell{bg=\HeatCellBg{29}}\CellText{0.29}{1}{1} & \SetCell{bg=\HeatCellBg{26}}\CellText{0.26}{0}{0} & \SetCell{bg=\HeatCellBg{19}}\CellText{0.19}{0}{0} & \SetCell{bg=\HeatCellBg{25}}\CellText{0.25}{0}{0} & \SetCell{bg=\HeatCellBg{29}}\CellText{0.29}{1}{1} & \SetCell{bg=\HeatCellBg{22}}\CellText{0.22}{0}{0} & \SetCell{bg=\HeatCellBg{22}}\CellText{0.22}{1}{0} & \SetCell{bg=\HeatCellBg{21}}\CellText{0.21}{0}{0} & \SetCell{bg=\HeatCellBg{17}}\CellText{0.17}{0}{0} & \SetCell{bg=\HeatCellBg{22}}\CellText{0.22}{0}{0} & \SetCell{bg=\HeatCellBg{24}}\CellText{0.24}{0}{0} & \SetCell{bg=\HeatCellBg{23}}\CellText{0.23}{0}{0} & \SetCell{bg=\HeatCellBg{19}}\CellText{0.19}{0}{0} & \SetCell{bg=\HeatCellBg{24}}\CellText{0.24}{0}{0} & \SetCell{bg=\HeatCellBg{25}}\CellText{0.25}{1}{0} & \SetCell{bg=\HeatCellBg{23}}\CellText{0.23}{0}{0} & \SetCell{bg=\HeatCellBg{20}}\CellText{0.20}{0}{0} & \SetCell{bg=\HeatCellBg{22}}\CellText{0.22}{1}{0} & \SetCell{bg=\HeatCellBg{23}}\CellText{0.23}{0}{0} & \MeanCell{0.23}{0.03}{0} \\
        8b & \SetCell{bg=\HeatCellBg{7}}\CellText{0.07}{0}{0} & \SetCell{bg=\HeatCellBg{-7}}\CellText{-0.07}{0}{0} & \SetCell{bg=\HeatCellBg{13}}\CellText{0.13}{0}{1} & \SetCell{bg=\HeatCellBg{4}}\CellText{0.04}{0}{0} & \SetCell{bg=\HeatCellBg{2}}\CellText{0.02}{0}{0} & \SetCell{bg=\HeatCellBg{1}}\CellText{0.01}{0}{0} & \SetCell{bg=\HeatCellBg{1}}\CellText{0.01}{0}{0} & \SetCell{bg=\HeatCellBg{5}}\CellText{0.05}{0}{0} & \SetCell{bg=\HeatCellBg{4}}\CellText{0.04}{0}{0} & \SetCell{bg=\HeatCellBg{6}}\CellText{0.06}{0}{0} & \SetCell{bg=\HeatCellBg{4}}\CellText{0.04}{0}{0} & \SetCell{bg=\HeatCellBg{7}}\CellText{0.07}{0}{0} & \SetCell{bg=\HeatCellBg{8}}\CellText{0.08}{0}{0} & \SetCell{bg=\HeatCellBg{7}}\CellText{0.07}{0}{0} & \SetCell{bg=\HeatCellBg{9}}\CellText{0.09}{0}{0} & \SetCell{bg=\HeatCellBg{7}}\CellText{0.07}{0}{0} & \SetCell{bg=\HeatCellBg{8}}\CellText{0.08}{0}{0} & \SetCell{bg=\HeatCellBg{8}}\CellText{0.08}{0}{0} & \SetCell{bg=\HeatCellBg{10}}\CellText{0.10}{0}{0} & \MeanCell{0.05}{0.04}{0} \\
        Mean $\pm$ Std & \MeanCell{0.04}{0.10}{0} & \MeanCell{0.21}{0.17}{0} & \MeanCell{0.17}{0.09}{0} & \MeanCell{0.22}{0.10}{1} & \MeanCell{0.18}{0.11}{0} & \MeanCell{0.20}{0.13}{0} & \MeanCell{0.16}{0.06}{0} & \MeanCell{0.14}{0.10}{0} & \MeanCell{0.21}{0.10}{0} & \MeanCell{0.19}{0.09}{0} & \MeanCell{0.17}{0.09}{0} & \MeanCell{0.19}{0.08}{0} & \MeanCell{0.16}{0.06}{0} & \MeanCell{0.18}{0.07}{0} & \MeanCell{0.17}{0.05}{0} & \MeanCell{0.18}{0.08}{0} & \MeanCell{0.20}{0.06}{0} & \MeanCell{0.16}{0.06}{0} & \MeanCell{0.19}{0.07}{0} & \MeanCell{0.18}{0.10}{0} \\
}

\ClusterKTable{10}{Observation baseline ($\Similarity{obs}{}$) --- \MarioBrosName.}{obs_mario_bros}{
        1 & \SetCell{bg=\HeatCellBg{99}}\CellText{0.99}{0}{1} & \SetCell{bg=\HeatCellBg{89}}\CellText{0.89}{1}{0} & \SetCell{bg=\HeatCellBg{77}}\CellText{0.77}{1}{0} & \SetCell{bg=\HeatCellBg{37}}\CellText{0.37}{1}{0} & \SetCell{bg=\HeatCellBg{49}}\CellText{0.49}{0}{0} & \SetCell{bg=\HeatCellBg{54}}\CellText{0.54}{1}{0} & \SetCell{bg=\HeatCellBg{47}}\CellText{0.47}{1}{0} & \SetCell{bg=\HeatCellBg{52}}\CellText{0.52}{1}{0} & \SetCell{bg=\HeatCellBg{39}}\CellText{0.39}{1}{0} & \SetCell{bg=\HeatCellBg{47}}\CellText{0.47}{1}{0} & \SetCell{bg=\HeatCellBg{47}}\CellText{0.47}{1}{0} & \SetCell{bg=\HeatCellBg{42}}\CellText{0.42}{1}{0} & \SetCell{bg=\HeatCellBg{56}}\CellText{0.56}{1}{0} & \SetCell{bg=\HeatCellBg{49}}\CellText{0.49}{1}{0} & \SetCell{bg=\HeatCellBg{41}}\CellText{0.41}{0}{0} & \SetCell{bg=\HeatCellBg{39}}\CellText{0.39}{1}{0} & \SetCell{bg=\HeatCellBg{48}}\CellText{0.48}{1}{0} & \SetCell{bg=\HeatCellBg{49}}\CellText{0.49}{1}{0} & \SetCell{bg=\HeatCellBg{52}}\CellText{0.52}{1}{0} & \MeanCell{0.53}{0.16}{1} \\
        2 & \SetCell{bg=\HeatCellBg{100}}\CellText{1.00}{1}{1} & \SetCell{bg=\HeatCellBg{34}}\CellText{0.34}{0}{0} & \SetCell{bg=\HeatCellBg{53}}\CellText{0.53}{0}{0} & \SetCell{bg=\HeatCellBg{33}}\CellText{0.33}{0}{0} & \SetCell{bg=\HeatCellBg{51}}\CellText{0.51}{1}{0} & \SetCell{bg=\HeatCellBg{42}}\CellText{0.42}{0}{0} & \SetCell{bg=\HeatCellBg{47}}\CellText{0.47}{1}{0} & \SetCell{bg=\HeatCellBg{52}}\CellText{0.52}{1}{0} & \SetCell{bg=\HeatCellBg{38}}\CellText{0.38}{0}{0} & \SetCell{bg=\HeatCellBg{37}}\CellText{0.37}{0}{0} & \SetCell{bg=\HeatCellBg{40}}\CellText{0.40}{0}{0} & \SetCell{bg=\HeatCellBg{37}}\CellText{0.37}{0}{0} & \SetCell{bg=\HeatCellBg{42}}\CellText{0.42}{0}{0} & \SetCell{bg=\HeatCellBg{39}}\CellText{0.39}{0}{0} & \SetCell{bg=\HeatCellBg{43}}\CellText{0.43}{1}{0} & \SetCell{bg=\HeatCellBg{32}}\CellText{0.32}{0}{0} & \SetCell{bg=\HeatCellBg{38}}\CellText{0.38}{0}{0} & \SetCell{bg=\HeatCellBg{44}}\CellText{0.44}{0}{0} & \SetCell{bg=\HeatCellBg{38}}\CellText{0.38}{0}{0} & \MeanCell{0.44}{0.14}{0} \\
        3 & \SetCell{bg=\HeatCellBg{99}}\CellText{0.99}{0}{1} & \SetCell{bg=\HeatCellBg{24}}\CellText{0.24}{0}{0} & \SetCell{bg=\HeatCellBg{35}}\CellText{0.35}{0}{0} & \SetCell{bg=\HeatCellBg{29}}\CellText{0.29}{0}{0} & \SetCell{bg=\HeatCellBg{36}}\CellText{0.36}{0}{0} & \SetCell{bg=\HeatCellBg{32}}\CellText{0.32}{0}{0} & \SetCell{bg=\HeatCellBg{37}}\CellText{0.37}{0}{0} & \SetCell{bg=\HeatCellBg{37}}\CellText{0.37}{0}{0} & \SetCell{bg=\HeatCellBg{31}}\CellText{0.31}{0}{0} & \SetCell{bg=\HeatCellBg{28}}\CellText{0.28}{0}{0} & \SetCell{bg=\HeatCellBg{34}}\CellText{0.34}{0}{0} & \SetCell{bg=\HeatCellBg{27}}\CellText{0.27}{0}{0} & \SetCell{bg=\HeatCellBg{31}}\CellText{0.31}{0}{0} & \SetCell{bg=\HeatCellBg{31}}\CellText{0.31}{0}{0} & \SetCell{bg=\HeatCellBg{30}}\CellText{0.30}{0}{0} & \SetCell{bg=\HeatCellBg{26}}\CellText{0.26}{0}{0} & \SetCell{bg=\HeatCellBg{28}}\CellText{0.28}{0}{0} & \SetCell{bg=\HeatCellBg{33}}\CellText{0.33}{0}{0} & \SetCell{bg=\HeatCellBg{28}}\CellText{0.28}{0}{0} & \MeanCell{0.35}{0.16}{0} \\
        4 & \SetCell{bg=\HeatCellBg{98}}\CellText{0.98}{0}{1} & \SetCell{bg=\HeatCellBg{24}}\CellText{0.24}{0}{0} & \SetCell{bg=\HeatCellBg{31}}\CellText{0.31}{0}{0} & \SetCell{bg=\HeatCellBg{25}}\CellText{0.25}{0}{0} & \SetCell{bg=\HeatCellBg{36}}\CellText{0.36}{0}{0} & \SetCell{bg=\HeatCellBg{31}}\CellText{0.31}{0}{0} & \SetCell{bg=\HeatCellBg{35}}\CellText{0.35}{0}{0} & \SetCell{bg=\HeatCellBg{35}}\CellText{0.35}{0}{0} & \SetCell{bg=\HeatCellBg{29}}\CellText{0.29}{0}{0} & \SetCell{bg=\HeatCellBg{28}}\CellText{0.28}{0}{0} & \SetCell{bg=\HeatCellBg{29}}\CellText{0.29}{0}{0} & \SetCell{bg=\HeatCellBg{24}}\CellText{0.24}{0}{0} & \SetCell{bg=\HeatCellBg{28}}\CellText{0.28}{0}{0} & \SetCell{bg=\HeatCellBg{26}}\CellText{0.26}{0}{0} & \SetCell{bg=\HeatCellBg{27}}\CellText{0.27}{0}{0} & \SetCell{bg=\HeatCellBg{21}}\CellText{0.21}{0}{0} & \SetCell{bg=\HeatCellBg{25}}\CellText{0.25}{0}{0} & \SetCell{bg=\HeatCellBg{28}}\CellText{0.28}{0}{0} & \SetCell{bg=\HeatCellBg{25}}\CellText{0.25}{0}{0} & \MeanCell{0.32}{0.16}{0} \\
        5 & \SetCell{bg=\HeatCellBg{98}}\CellText{0.98}{0}{1} & \SetCell{bg=\HeatCellBg{24}}\CellText{0.24}{0}{0} & \SetCell{bg=\HeatCellBg{31}}\CellText{0.31}{0}{0} & \SetCell{bg=\HeatCellBg{25}}\CellText{0.25}{0}{0} & \SetCell{bg=\HeatCellBg{36}}\CellText{0.36}{0}{0} & \SetCell{bg=\HeatCellBg{31}}\CellText{0.31}{0}{0} & \SetCell{bg=\HeatCellBg{35}}\CellText{0.35}{0}{0} & \SetCell{bg=\HeatCellBg{35}}\CellText{0.35}{0}{0} & \SetCell{bg=\HeatCellBg{29}}\CellText{0.29}{0}{0} & \SetCell{bg=\HeatCellBg{28}}\CellText{0.28}{0}{0} & \SetCell{bg=\HeatCellBg{29}}\CellText{0.29}{0}{0} & \SetCell{bg=\HeatCellBg{24}}\CellText{0.24}{0}{0} & \SetCell{bg=\HeatCellBg{28}}\CellText{0.28}{0}{0} & \SetCell{bg=\HeatCellBg{26}}\CellText{0.26}{0}{0} & \SetCell{bg=\HeatCellBg{27}}\CellText{0.27}{0}{0} & \SetCell{bg=\HeatCellBg{21}}\CellText{0.21}{0}{0} & \SetCell{bg=\HeatCellBg{25}}\CellText{0.25}{0}{0} & \SetCell{bg=\HeatCellBg{28}}\CellText{0.28}{0}{0} & \SetCell{bg=\HeatCellBg{25}}\CellText{0.25}{0}{0} & \MeanCell{0.32}{0.16}{0} \\
        6b & \SetCell{bg=\HeatCellBg{95}}\CellText{0.95}{0}{1} & \SetCell{bg=\HeatCellBg{25}}\CellText{0.25}{0}{0} & \SetCell{bg=\HeatCellBg{29}}\CellText{0.29}{0}{0} & \SetCell{bg=\HeatCellBg{22}}\CellText{0.22}{0}{0} & \SetCell{bg=\HeatCellBg{27}}\CellText{0.27}{0}{0} & \SetCell{bg=\HeatCellBg{24}}\CellText{0.24}{0}{0} & \SetCell{bg=\HeatCellBg{28}}\CellText{0.28}{0}{0} & \SetCell{bg=\HeatCellBg{28}}\CellText{0.28}{0}{0} & \SetCell{bg=\HeatCellBg{26}}\CellText{0.26}{0}{0} & \SetCell{bg=\HeatCellBg{24}}\CellText{0.24}{0}{0} & \SetCell{bg=\HeatCellBg{24}}\CellText{0.24}{0}{0} & \SetCell{bg=\HeatCellBg{21}}\CellText{0.21}{0}{0} & \SetCell{bg=\HeatCellBg{24}}\CellText{0.24}{0}{0} & \SetCell{bg=\HeatCellBg{22}}\CellText{0.22}{0}{0} & \SetCell{bg=\HeatCellBg{22}}\CellText{0.22}{0}{0} & \SetCell{bg=\HeatCellBg{21}}\CellText{0.21}{0}{0} & \SetCell{bg=\HeatCellBg{21}}\CellText{0.21}{0}{0} & \SetCell{bg=\HeatCellBg{25}}\CellText{0.25}{0}{0} & \SetCell{bg=\HeatCellBg{22}}\CellText{0.22}{0}{0} & \MeanCell{0.28}{0.16}{0} \\
        7b & \SetCell{bg=\HeatCellBg{-6}}\CellText{-0.06}{0}{0} & \SetCell{bg=\HeatCellBg{0}}\CellText{-0.00}{0}{0} & \SetCell{bg=\HeatCellBg{5}}\CellText{0.05}{0}{0} & \SetCell{bg=\HeatCellBg{7}}\CellText{0.07}{0}{0} & \SetCell{bg=\HeatCellBg{11}}\CellText{0.11}{0}{0} & \SetCell{bg=\HeatCellBg{12}}\CellText{0.12}{0}{0} & \SetCell{bg=\HeatCellBg{14}}\CellText{0.14}{0}{0} & \SetCell{bg=\HeatCellBg{15}}\CellText{0.15}{0}{0} & \SetCell{bg=\HeatCellBg{15}}\CellText{0.15}{0}{0} & \SetCell{bg=\HeatCellBg{18}}\CellText{0.18}{0}{1} & \SetCell{bg=\HeatCellBg{15}}\CellText{0.15}{0}{0} & \SetCell{bg=\HeatCellBg{16}}\CellText{0.16}{0}{0} & \SetCell{bg=\HeatCellBg{17}}\CellText{0.17}{0}{0} & \SetCell{bg=\HeatCellBg{15}}\CellText{0.15}{0}{0} & \SetCell{bg=\HeatCellBg{15}}\CellText{0.15}{0}{0} & \SetCell{bg=\HeatCellBg{14}}\CellText{0.14}{0}{0} & \SetCell{bg=\HeatCellBg{16}}\CellText{0.16}{0}{0} & \SetCell{bg=\HeatCellBg{16}}\CellText{0.16}{0}{0} & \SetCell{bg=\HeatCellBg{16}}\CellText{0.16}{0}{0} & \MeanCell{0.12}{0.06}{0} \\
        8b & \SetCell{bg=\HeatCellBg{-9}}\CellText{-0.09}{0}{0} & \SetCell{bg=\HeatCellBg{-4}}\CellText{-0.04}{0}{0} & \SetCell{bg=\HeatCellBg{2}}\CellText{0.02}{0}{0} & \SetCell{bg=\HeatCellBg{2}}\CellText{0.02}{0}{0} & \SetCell{bg=\HeatCellBg{3}}\CellText{0.03}{0}{0} & \SetCell{bg=\HeatCellBg{6}}\CellText{0.06}{0}{0} & \SetCell{bg=\HeatCellBg{7}}\CellText{0.07}{0}{0} & \SetCell{bg=\HeatCellBg{7}}\CellText{0.07}{0}{0} & \SetCell{bg=\HeatCellBg{6}}\CellText{0.06}{0}{0} & \SetCell{bg=\HeatCellBg{7}}\CellText{0.07}{0}{0} & \SetCell{bg=\HeatCellBg{5}}\CellText{0.05}{0}{0} & \SetCell{bg=\HeatCellBg{6}}\CellText{0.06}{0}{0} & \SetCell{bg=\HeatCellBg{9}}\CellText{0.09}{0}{1} & \SetCell{bg=\HeatCellBg{6}}\CellText{0.06}{0}{0} & \SetCell{bg=\HeatCellBg{7}}\CellText{0.07}{0}{0} & \SetCell{bg=\HeatCellBg{6}}\CellText{0.06}{0}{0} & \SetCell{bg=\HeatCellBg{8}}\CellText{0.08}{0}{0} & \SetCell{bg=\HeatCellBg{8}}\CellText{0.08}{0}{0} & \SetCell{bg=\HeatCellBg{7}}\CellText{0.07}{0}{0} & \MeanCell{0.05}{0.04}{0} \\
        Mean $\pm$ Std & \MeanCell{0.72}{0.46}{1} & \MeanCell{0.27}{0.26}{0} & \MeanCell{0.33}{0.23}{0} & \MeanCell{0.23}{0.11}{0} & \MeanCell{0.31}{0.16}{0} & \MeanCell{0.29}{0.14}{0} & \MeanCell{0.31}{0.13}{0} & \MeanCell{0.33}{0.15}{0} & \MeanCell{0.27}{0.10}{0} & \MeanCell{0.27}{0.11}{0} & \MeanCell{0.28}{0.13}{0} & \MeanCell{0.25}{0.11}{0} & \MeanCell{0.29}{0.14}{0} & \MeanCell{0.27}{0.13}{0} & \MeanCell{0.27}{0.11}{0} & \MeanCell{0.23}{0.10}{0} & \MeanCell{0.26}{0.12}{0} & \MeanCell{0.29}{0.13}{0} & \MeanCell{0.27}{0.13}{0} & \MeanCell{0.30}{0.20}{0} \\
}

\ClusterKTable{10}{Untrained-network (latent) baseline ($\Similarity{untr}{}$) --- \MarioBrosName.}{untr_mario_bros}{
        1 & \SetCell{bg=\HeatCellBg{99}}\CellText{0.99}{0}{1} & \SetCell{bg=\HeatCellBg{70}}\CellText{0.70}{1}{0} & \SetCell{bg=\HeatCellBg{51}}\CellText{0.51}{1}{0} & \SetCell{bg=\HeatCellBg{50}}\CellText{0.50}{1}{0} & \SetCell{bg=\HeatCellBg{45}}\CellText{0.45}{1}{0} & \SetCell{bg=\HeatCellBg{49}}\CellText{0.49}{1}{0} & \SetCell{bg=\HeatCellBg{48}}\CellText{0.48}{1}{0} & \SetCell{bg=\HeatCellBg{46}}\CellText{0.46}{1}{0} & \SetCell{bg=\HeatCellBg{42}}\CellText{0.42}{1}{0} & \SetCell{bg=\HeatCellBg{49}}\CellText{0.49}{1}{0} & \SetCell{bg=\HeatCellBg{50}}\CellText{0.50}{1}{0} & \SetCell{bg=\HeatCellBg{50}}\CellText{0.50}{1}{0} & \SetCell{bg=\HeatCellBg{41}}\CellText{0.41}{0}{0} & \SetCell{bg=\HeatCellBg{54}}\CellText{0.54}{1}{0} & \SetCell{bg=\HeatCellBg{40}}\CellText{0.40}{0}{0} & \SetCell{bg=\HeatCellBg{46}}\CellText{0.46}{1}{0} & \SetCell{bg=\HeatCellBg{46}}\CellText{0.46}{1}{0} & \SetCell{bg=\HeatCellBg{47}}\CellText{0.47}{1}{0} & \SetCell{bg=\HeatCellBg{47}}\CellText{0.47}{1}{0} & \MeanCell{0.51}{0.13}{1} \\
        2 & \SetCell{bg=\HeatCellBg{100}}\CellText{1.00}{1}{1} & \SetCell{bg=\HeatCellBg{32}}\CellText{0.32}{0}{0} & \SetCell{bg=\HeatCellBg{42}}\CellText{0.42}{0}{0} & \SetCell{bg=\HeatCellBg{45}}\CellText{0.45}{0}{0} & \SetCell{bg=\HeatCellBg{45}}\CellText{0.45}{1}{0} & \SetCell{bg=\HeatCellBg{35}}\CellText{0.35}{0}{0} & \SetCell{bg=\HeatCellBg{43}}\CellText{0.43}{0}{0} & \SetCell{bg=\HeatCellBg{42}}\CellText{0.42}{0}{0} & \SetCell{bg=\HeatCellBg{35}}\CellText{0.35}{0}{0} & \SetCell{bg=\HeatCellBg{37}}\CellText{0.37}{0}{0} & \SetCell{bg=\HeatCellBg{46}}\CellText{0.46}{0}{0} & \SetCell{bg=\HeatCellBg{38}}\CellText{0.38}{0}{0} & \SetCell{bg=\HeatCellBg{44}}\CellText{0.44}{1}{0} & \SetCell{bg=\HeatCellBg{43}}\CellText{0.43}{0}{0} & \SetCell{bg=\HeatCellBg{41}}\CellText{0.41}{1}{0} & \SetCell{bg=\HeatCellBg{41}}\CellText{0.41}{0}{0} & \SetCell{bg=\HeatCellBg{38}}\CellText{0.38}{0}{0} & \SetCell{bg=\HeatCellBg{40}}\CellText{0.40}{0}{0} & \SetCell{bg=\HeatCellBg{43}}\CellText{0.43}{0}{0} & \MeanCell{0.44}{0.14}{0} \\
        3 & \SetCell{bg=\HeatCellBg{100}}\CellText{1.00}{1}{1} & \SetCell{bg=\HeatCellBg{34}}\CellText{0.34}{0}{0} & \SetCell{bg=\HeatCellBg{37}}\CellText{0.37}{0}{0} & \SetCell{bg=\HeatCellBg{35}}\CellText{0.35}{0}{0} & \SetCell{bg=\HeatCellBg{30}}\CellText{0.30}{0}{0} & \SetCell{bg=\HeatCellBg{38}}\CellText{0.38}{0}{0} & \SetCell{bg=\HeatCellBg{34}}\CellText{0.34}{0}{0} & \SetCell{bg=\HeatCellBg{37}}\CellText{0.37}{0}{0} & \SetCell{bg=\HeatCellBg{31}}\CellText{0.31}{0}{0} & \SetCell{bg=\HeatCellBg{31}}\CellText{0.31}{0}{0} & \SetCell{bg=\HeatCellBg{35}}\CellText{0.35}{0}{0} & \SetCell{bg=\HeatCellBg{34}}\CellText{0.34}{0}{0} & \SetCell{bg=\HeatCellBg{38}}\CellText{0.38}{0}{0} & \SetCell{bg=\HeatCellBg{36}}\CellText{0.36}{0}{0} & \SetCell{bg=\HeatCellBg{33}}\CellText{0.33}{0}{0} & \SetCell{bg=\HeatCellBg{35}}\CellText{0.35}{0}{0} & \SetCell{bg=\HeatCellBg{34}}\CellText{0.34}{0}{0} & \SetCell{bg=\HeatCellBg{35}}\CellText{0.35}{0}{0} & \SetCell{bg=\HeatCellBg{31}}\CellText{0.31}{0}{0} & \MeanCell{0.38}{0.15}{0} \\
        4 & \SetCell{bg=\HeatCellBg{98}}\CellText{0.98}{0}{1} & \SetCell{bg=\HeatCellBg{47}}\CellText{0.47}{0}{0} & \SetCell{bg=\HeatCellBg{32}}\CellText{0.32}{0}{0} & \SetCell{bg=\HeatCellBg{29}}\CellText{0.29}{0}{0} & \SetCell{bg=\HeatCellBg{26}}\CellText{0.26}{0}{0} & \SetCell{bg=\HeatCellBg{29}}\CellText{0.29}{0}{0} & \SetCell{bg=\HeatCellBg{29}}\CellText{0.29}{0}{0} & \SetCell{bg=\HeatCellBg{23}}\CellText{0.23}{0}{0} & \SetCell{bg=\HeatCellBg{25}}\CellText{0.25}{0}{0} & \SetCell{bg=\HeatCellBg{23}}\CellText{0.23}{0}{0} & \SetCell{bg=\HeatCellBg{23}}\CellText{0.23}{0}{0} & \SetCell{bg=\HeatCellBg{23}}\CellText{0.23}{0}{0} & \SetCell{bg=\HeatCellBg{23}}\CellText{0.23}{0}{0} & \SetCell{bg=\HeatCellBg{22}}\CellText{0.22}{0}{0} & \SetCell{bg=\HeatCellBg{22}}\CellText{0.22}{0}{0} & \SetCell{bg=\HeatCellBg{24}}\CellText{0.24}{0}{0} & \SetCell{bg=\HeatCellBg{23}}\CellText{0.23}{0}{0} & \SetCell{bg=\HeatCellBg{23}}\CellText{0.23}{0}{0} & \SetCell{bg=\HeatCellBg{23}}\CellText{0.23}{0}{0} & \MeanCell{0.30}{0.17}{0} \\
        5 & \SetCell{bg=\HeatCellBg{98}}\CellText{0.98}{0}{1} & \SetCell{bg=\HeatCellBg{47}}\CellText{0.47}{0}{0} & \SetCell{bg=\HeatCellBg{32}}\CellText{0.32}{0}{0} & \SetCell{bg=\HeatCellBg{29}}\CellText{0.29}{0}{0} & \SetCell{bg=\HeatCellBg{26}}\CellText{0.26}{0}{0} & \SetCell{bg=\HeatCellBg{29}}\CellText{0.29}{0}{0} & \SetCell{bg=\HeatCellBg{29}}\CellText{0.29}{0}{0} & \SetCell{bg=\HeatCellBg{23}}\CellText{0.23}{0}{0} & \SetCell{bg=\HeatCellBg{25}}\CellText{0.25}{0}{0} & \SetCell{bg=\HeatCellBg{23}}\CellText{0.23}{0}{0} & \SetCell{bg=\HeatCellBg{23}}\CellText{0.23}{0}{0} & \SetCell{bg=\HeatCellBg{22}}\CellText{0.22}{0}{0} & \SetCell{bg=\HeatCellBg{23}}\CellText{0.23}{0}{0} & \SetCell{bg=\HeatCellBg{22}}\CellText{0.22}{0}{0} & \SetCell{bg=\HeatCellBg{22}}\CellText{0.22}{0}{0} & \SetCell{bg=\HeatCellBg{24}}\CellText{0.24}{0}{0} & \SetCell{bg=\HeatCellBg{23}}\CellText{0.23}{0}{0} & \SetCell{bg=\HeatCellBg{24}}\CellText{0.24}{0}{0} & \SetCell{bg=\HeatCellBg{23}}\CellText{0.23}{0}{0} & \MeanCell{0.30}{0.17}{0} \\
        6b & \SetCell{bg=\HeatCellBg{95}}\CellText{0.95}{0}{1} & \SetCell{bg=\HeatCellBg{34}}\CellText{0.34}{0}{0} & \SetCell{bg=\HeatCellBg{28}}\CellText{0.28}{0}{0} & \SetCell{bg=\HeatCellBg{27}}\CellText{0.27}{0}{0} & \SetCell{bg=\HeatCellBg{27}}\CellText{0.27}{0}{0} & \SetCell{bg=\HeatCellBg{18}}\CellText{0.18}{0}{0} & \SetCell{bg=\HeatCellBg{21}}\CellText{0.21}{0}{0} & \SetCell{bg=\HeatCellBg{21}}\CellText{0.21}{0}{0} & \SetCell{bg=\HeatCellBg{17}}\CellText{0.17}{0}{0} & \SetCell{bg=\HeatCellBg{19}}\CellText{0.19}{0}{0} & \SetCell{bg=\HeatCellBg{21}}\CellText{0.21}{0}{0} & \SetCell{bg=\HeatCellBg{16}}\CellText{0.16}{0}{0} & \SetCell{bg=\HeatCellBg{19}}\CellText{0.19}{0}{0} & \SetCell{bg=\HeatCellBg{18}}\CellText{0.18}{0}{0} & \SetCell{bg=\HeatCellBg{19}}\CellText{0.19}{0}{0} & \SetCell{bg=\HeatCellBg{17}}\CellText{0.17}{0}{0} & \SetCell{bg=\HeatCellBg{17}}\CellText{0.17}{0}{0} & \SetCell{bg=\HeatCellBg{19}}\CellText{0.19}{0}{0} & \SetCell{bg=\HeatCellBg{19}}\CellText{0.19}{0}{0} & \MeanCell{0.25}{0.17}{0} \\
        7b & \SetCell{bg=\HeatCellBg{-6}}\CellText{-0.06}{0}{0} & \SetCell{bg=\HeatCellBg{2}}\CellText{0.02}{0}{0} & \SetCell{bg=\HeatCellBg{6}}\CellText{0.06}{0}{0} & \SetCell{bg=\HeatCellBg{6}}\CellText{0.06}{0}{0} & \SetCell{bg=\HeatCellBg{7}}\CellText{0.07}{0}{0} & \SetCell{bg=\HeatCellBg{9}}\CellText{0.09}{0}{0} & \SetCell{bg=\HeatCellBg{9}}\CellText{0.09}{0}{0} & \SetCell{bg=\HeatCellBg{9}}\CellText{0.09}{0}{0} & \SetCell{bg=\HeatCellBg{11}}\CellText{0.11}{0}{0} & \SetCell{bg=\HeatCellBg{13}}\CellText{0.13}{0}{1} & \SetCell{bg=\HeatCellBg{12}}\CellText{0.12}{0}{0} & \SetCell{bg=\HeatCellBg{12}}\CellText{0.12}{0}{0} & \SetCell{bg=\HeatCellBg{10}}\CellText{0.10}{0}{0} & \SetCell{bg=\HeatCellBg{13}}\CellText{0.13}{0}{1} & \SetCell{bg=\HeatCellBg{13}}\CellText{0.13}{0}{1} & \SetCell{bg=\HeatCellBg{12}}\CellText{0.12}{0}{0} & \SetCell{bg=\HeatCellBg{13}}\CellText{0.13}{0}{1} & \SetCell{bg=\HeatCellBg{13}}\CellText{0.13}{0}{1} & \SetCell{bg=\HeatCellBg{12}}\CellText{0.12}{0}{0} & \MeanCell{0.09}{0.05}{0} \\
        8b & \SetCell{bg=\HeatCellBg{-9}}\CellText{-0.09}{0}{0} & \SetCell{bg=\HeatCellBg{-2}}\CellText{-0.02}{0}{0} & \SetCell{bg=\HeatCellBg{1}}\CellText{0.01}{0}{0} & \SetCell{bg=\HeatCellBg{2}}\CellText{0.02}{0}{0} & \SetCell{bg=\HeatCellBg{0}}\CellText{0.00}{0}{0} & \SetCell{bg=\HeatCellBg{2}}\CellText{0.02}{0}{0} & \SetCell{bg=\HeatCellBg{2}}\CellText{0.02}{0}{0} & \SetCell{bg=\HeatCellBg{2}}\CellText{0.02}{0}{0} & \SetCell{bg=\HeatCellBg{3}}\CellText{0.03}{0}{0} & \SetCell{bg=\HeatCellBg{4}}\CellText{0.04}{0}{0} & \SetCell{bg=\HeatCellBg{4}}\CellText{0.04}{0}{0} & \SetCell{bg=\HeatCellBg{4}}\CellText{0.04}{0}{0} & \SetCell{bg=\HeatCellBg{5}}\CellText{0.05}{0}{1} & \SetCell{bg=\HeatCellBg{3}}\CellText{0.03}{0}{0} & \SetCell{bg=\HeatCellBg{4}}\CellText{0.04}{0}{0} & \SetCell{bg=\HeatCellBg{3}}\CellText{0.03}{0}{0} & \SetCell{bg=\HeatCellBg{4}}\CellText{0.04}{0}{0} & \SetCell{bg=\HeatCellBg{4}}\CellText{0.04}{0}{0} & \SetCell{bg=\HeatCellBg{5}}\CellText{0.05}{0}{1} & \MeanCell{0.02}{0.03}{0} \\
        Mean $\pm$ Std & \MeanCell{0.72}{0.46}{1} & \MeanCell{0.33}{0.22}{0} & \MeanCell{0.29}{0.16}{0} & \MeanCell{0.28}{0.16}{0} & \MeanCell{0.26}{0.15}{0} & \MeanCell{0.26}{0.15}{0} & \MeanCell{0.27}{0.15}{0} & \MeanCell{0.25}{0.14}{0} & \MeanCell{0.24}{0.12}{0} & \MeanCell{0.25}{0.13}{0} & \MeanCell{0.27}{0.15}{0} & \MeanCell{0.25}{0.14}{0} & \MeanCell{0.25}{0.13}{0} & \MeanCell{0.26}{0.16}{0} & \MeanCell{0.24}{0.12}{0} & \MeanCell{0.25}{0.14}{0} & \MeanCell{0.25}{0.13}{0} & \MeanCell{0.26}{0.13}{0} & \MeanCell{0.25}{0.13}{0} & \MeanCell{0.29}{0.21}{0} \\
}

\ClusterKTable{10}{Action baseline ($\Similarity{act}{}$) --- \MarioBrosName.}{act_mario_bros}{
        1 & \SetCell{bg=\HeatCellBg{-9}}\CellText{-0.09}{0}{0} & \SetCell{bg=\HeatCellBg{-7}}\CellText{-0.07}{0}{0} & \SetCell{bg=\HeatCellBg{-1}}\CellText{-0.01}{0}{0} & \SetCell{bg=\HeatCellBg{2}}\CellText{0.02}{0}{0} & \SetCell{bg=\HeatCellBg{3}}\CellText{0.03}{0}{0} & \SetCell{bg=\HeatCellBg{2}}\CellText{0.02}{0}{0} & \SetCell{bg=\HeatCellBg{4}}\CellText{0.04}{0}{1} & \SetCell{bg=\HeatCellBg{3}}\CellText{0.03}{0}{0} & \SetCell{bg=\HeatCellBg{3}}\CellText{0.03}{0}{0} & \SetCell{bg=\HeatCellBg{3}}\CellText{0.03}{0}{0} & \SetCell{bg=\HeatCellBg{3}}\CellText{0.03}{0}{0} & \SetCell{bg=\HeatCellBg{4}}\CellText{0.04}{0}{1} & \SetCell{bg=\HeatCellBg{4}}\CellText{0.04}{0}{1} & \SetCell{bg=\HeatCellBg{3}}\CellText{0.03}{0}{0} & \SetCell{bg=\HeatCellBg{3}}\CellText{0.03}{0}{0} & \SetCell{bg=\HeatCellBg{3}}\CellText{0.03}{0}{0} & \SetCell{bg=\HeatCellBg{4}}\CellText{0.04}{0}{1} & \SetCell{bg=\HeatCellBg{3}}\CellText{0.03}{0}{0} & \SetCell{bg=\HeatCellBg{3}}\CellText{0.03}{0}{0} & \MeanCell{0.02}{0.04}{0} \\
        2 & \SetCell{bg=\HeatCellBg{-9}}\CellText{-0.09}{0}{0} & \SetCell{bg=\HeatCellBg{-2}}\CellText{-0.02}{0}{0} & \SetCell{bg=\HeatCellBg{-2}}\CellText{-0.02}{0}{0} & \SetCell{bg=\HeatCellBg{1}}\CellText{0.01}{0}{0} & \SetCell{bg=\HeatCellBg{1}}\CellText{0.01}{0}{0} & \SetCell{bg=\HeatCellBg{1}}\CellText{0.01}{0}{0} & \SetCell{bg=\HeatCellBg{1}}\CellText{0.01}{0}{0} & \SetCell{bg=\HeatCellBg{1}}\CellText{0.01}{0}{0} & \SetCell{bg=\HeatCellBg{2}}\CellText{0.02}{0}{0} & \SetCell{bg=\HeatCellBg{2}}\CellText{0.02}{0}{0} & \SetCell{bg=\HeatCellBg{1}}\CellText{0.01}{0}{0} & \SetCell{bg=\HeatCellBg{2}}\CellText{0.02}{0}{0} & \SetCell{bg=\HeatCellBg{1}}\CellText{0.01}{0}{0} & \SetCell{bg=\HeatCellBg{2}}\CellText{0.02}{0}{0} & \SetCell{bg=\HeatCellBg{3}}\CellText{0.03}{0}{1} & \SetCell{bg=\HeatCellBg{2}}\CellText{0.02}{0}{0} & \SetCell{bg=\HeatCellBg{2}}\CellText{0.02}{0}{0} & \SetCell{bg=\HeatCellBg{2}}\CellText{0.02}{0}{0} & \SetCell{bg=\HeatCellBg{1}}\CellText{0.01}{0}{0} & \MeanCell{0.01}{0.03}{0} \\
        3 & \SetCell{bg=\HeatCellBg{-9}}\CellText{-0.09}{0}{0} & \SetCell{bg=\HeatCellBg{-2}}\CellText{-0.02}{0}{0} & \SetCell{bg=\HeatCellBg{-1}}\CellText{-0.01}{0}{0} & \SetCell{bg=\HeatCellBg{0}}\CellText{-0.00}{0}{0} & \SetCell{bg=\HeatCellBg{1}}\CellText{0.01}{0}{0} & \SetCell{bg=\HeatCellBg{2}}\CellText{0.02}{0}{1} & \SetCell{bg=\HeatCellBg{1}}\CellText{0.01}{0}{0} & \SetCell{bg=\HeatCellBg{1}}\CellText{0.01}{0}{0} & \SetCell{bg=\HeatCellBg{1}}\CellText{0.01}{0}{0} & \SetCell{bg=\HeatCellBg{1}}\CellText{0.01}{0}{0} & \SetCell{bg=\HeatCellBg{1}}\CellText{0.01}{0}{0} & \SetCell{bg=\HeatCellBg{2}}\CellText{0.02}{0}{1} & \SetCell{bg=\HeatCellBg{2}}\CellText{0.02}{0}{1} & \SetCell{bg=\HeatCellBg{2}}\CellText{0.02}{0}{1} & \SetCell{bg=\HeatCellBg{2}}\CellText{0.02}{0}{1} & \SetCell{bg=\HeatCellBg{2}}\CellText{0.02}{0}{1} & \SetCell{bg=\HeatCellBg{2}}\CellText{0.02}{0}{1} & \SetCell{bg=\HeatCellBg{2}}\CellText{0.02}{0}{1} & \SetCell{bg=\HeatCellBg{2}}\CellText{0.02}{0}{1} & \MeanCell{0.01}{0.03}{0} \\
        4 & \SetCell{bg=\HeatCellBg{-8}}\CellText{-0.08}{0}{0} & \SetCell{bg=\HeatCellBg{1}}\CellText{0.01}{0}{0} & \SetCell{bg=\HeatCellBg{1}}\CellText{0.01}{0}{0} & \SetCell{bg=\HeatCellBg{1}}\CellText{0.01}{0}{0} & \SetCell{bg=\HeatCellBg{3}}\CellText{0.03}{0}{0} & \SetCell{bg=\HeatCellBg{3}}\CellText{0.03}{0}{0} & \SetCell{bg=\HeatCellBg{3}}\CellText{0.03}{0}{0} & \SetCell{bg=\HeatCellBg{3}}\CellText{0.03}{0}{0} & \SetCell{bg=\HeatCellBg{2}}\CellText{0.02}{0}{0} & \SetCell{bg=\HeatCellBg{3}}\CellText{0.03}{0}{0} & \SetCell{bg=\HeatCellBg{4}}\CellText{0.04}{0}{1} & \SetCell{bg=\HeatCellBg{4}}\CellText{0.04}{0}{1} & \SetCell{bg=\HeatCellBg{3}}\CellText{0.03}{0}{0} & \SetCell{bg=\HeatCellBg{3}}\CellText{0.03}{0}{0} & \SetCell{bg=\HeatCellBg{3}}\CellText{0.03}{0}{0} & \SetCell{bg=\HeatCellBg{3}}\CellText{0.03}{0}{0} & \SetCell{bg=\HeatCellBg{3}}\CellText{0.03}{0}{0} & \SetCell{bg=\HeatCellBg{3}}\CellText{0.03}{0}{0} & \SetCell{bg=\HeatCellBg{3}}\CellText{0.03}{0}{0} & \MeanCell{0.02}{0.03}{0} \\
        5 & \SetCell{bg=\HeatCellBg{-8}}\CellText{-0.08}{0}{0} & \SetCell{bg=\HeatCellBg{1}}\CellText{0.01}{0}{0} & \SetCell{bg=\HeatCellBg{1}}\CellText{0.01}{0}{0} & \SetCell{bg=\HeatCellBg{1}}\CellText{0.01}{0}{0} & \SetCell{bg=\HeatCellBg{3}}\CellText{0.03}{0}{0} & \SetCell{bg=\HeatCellBg{3}}\CellText{0.03}{0}{0} & \SetCell{bg=\HeatCellBg{3}}\CellText{0.03}{0}{0} & \SetCell{bg=\HeatCellBg{3}}\CellText{0.03}{0}{0} & \SetCell{bg=\HeatCellBg{2}}\CellText{0.02}{0}{0} & \SetCell{bg=\HeatCellBg{3}}\CellText{0.03}{0}{0} & \SetCell{bg=\HeatCellBg{4}}\CellText{0.04}{0}{1} & \SetCell{bg=\HeatCellBg{4}}\CellText{0.04}{0}{1} & \SetCell{bg=\HeatCellBg{3}}\CellText{0.03}{0}{0} & \SetCell{bg=\HeatCellBg{3}}\CellText{0.03}{0}{0} & \SetCell{bg=\HeatCellBg{3}}\CellText{0.03}{0}{0} & \SetCell{bg=\HeatCellBg{3}}\CellText{0.03}{0}{0} & \SetCell{bg=\HeatCellBg{3}}\CellText{0.03}{0}{0} & \SetCell{bg=\HeatCellBg{3}}\CellText{0.03}{0}{0} & \SetCell{bg=\HeatCellBg{3}}\CellText{0.03}{0}{0} & \MeanCell{0.02}{0.03}{0} \\
        6b & \SetCell{bg=\HeatCellBg{-8}}\CellText{-0.08}{0}{0} & \SetCell{bg=\HeatCellBg{8}}\CellText{0.08}{0}{1} & \SetCell{bg=\HeatCellBg{8}}\CellText{0.08}{0}{1} & \SetCell{bg=\HeatCellBg{6}}\CellText{0.06}{0}{0} & \SetCell{bg=\HeatCellBg{7}}\CellText{0.07}{0}{0} & \SetCell{bg=\HeatCellBg{8}}\CellText{0.08}{0}{1} & \SetCell{bg=\HeatCellBg{7}}\CellText{0.07}{0}{0} & \SetCell{bg=\HeatCellBg{8}}\CellText{0.08}{0}{1} & \SetCell{bg=\HeatCellBg{8}}\CellText{0.08}{0}{1} & \SetCell{bg=\HeatCellBg{7}}\CellText{0.07}{0}{0} & \SetCell{bg=\HeatCellBg{7}}\CellText{0.07}{0}{0} & \SetCell{bg=\HeatCellBg{7}}\CellText{0.07}{0}{0} & \SetCell{bg=\HeatCellBg{7}}\CellText{0.07}{0}{0} & \SetCell{bg=\HeatCellBg{6}}\CellText{0.06}{0}{0} & \SetCell{bg=\HeatCellBg{6}}\CellText{0.06}{0}{0} & \SetCell{bg=\HeatCellBg{6}}\CellText{0.06}{0}{0} & \SetCell{bg=\HeatCellBg{5}}\CellText{0.05}{0}{0} & \SetCell{bg=\HeatCellBg{5}}\CellText{0.05}{0}{0} & \SetCell{bg=\HeatCellBg{5}}\CellText{0.05}{0}{0} & \MeanCell{0.06}{0.03}{0} \\
        7b & \SetCell{bg=\HeatCellBg{24}}\CellText{0.24}{0}{1} & \SetCell{bg=\HeatCellBg{15}}\CellText{0.15}{0}{0} & \SetCell{bg=\HeatCellBg{14}}\CellText{0.14}{0}{0} & \SetCell{bg=\HeatCellBg{16}}\CellText{0.16}{0}{0} & \SetCell{bg=\HeatCellBg{15}}\CellText{0.15}{0}{0} & \SetCell{bg=\HeatCellBg{14}}\CellText{0.14}{0}{0} & \SetCell{bg=\HeatCellBg{15}}\CellText{0.15}{0}{0} & \SetCell{bg=\HeatCellBg{15}}\CellText{0.15}{0}{0} & \SetCell{bg=\HeatCellBg{14}}\CellText{0.14}{0}{0} & \SetCell{bg=\HeatCellBg{14}}\CellText{0.14}{0}{0} & \SetCell{bg=\HeatCellBg{16}}\CellText{0.16}{0}{0} & \SetCell{bg=\HeatCellBg{16}}\CellText{0.16}{0}{0} & \SetCell{bg=\HeatCellBg{16}}\CellText{0.16}{0}{0} & \SetCell{bg=\HeatCellBg{15}}\CellText{0.15}{0}{0} & \SetCell{bg=\HeatCellBg{14}}\CellText{0.14}{0}{0} & \SetCell{bg=\HeatCellBg{13}}\CellText{0.13}{0}{0} & \SetCell{bg=\HeatCellBg{13}}\CellText{0.13}{0}{0} & \SetCell{bg=\HeatCellBg{12}}\CellText{0.12}{0}{0} & \SetCell{bg=\HeatCellBg{12}}\CellText{0.12}{0}{0} & \MeanCell{0.15}{0.02}{0} \\
        8b & \SetCell{bg=\HeatCellBg{61}}\CellText{0.61}{1}{1} & \SetCell{bg=\HeatCellBg{45}}\CellText{0.45}{1}{0} & \SetCell{bg=\HeatCellBg{36}}\CellText{0.36}{1}{0} & \SetCell{bg=\HeatCellBg{32}}\CellText{0.32}{1}{0} & \SetCell{bg=\HeatCellBg{33}}\CellText{0.33}{1}{0} & \SetCell{bg=\HeatCellBg{32}}\CellText{0.32}{1}{0} & \SetCell{bg=\HeatCellBg{31}}\CellText{0.31}{1}{0} & \SetCell{bg=\HeatCellBg{31}}\CellText{0.31}{1}{0} & \SetCell{bg=\HeatCellBg{32}}\CellText{0.32}{1}{0} & \SetCell{bg=\HeatCellBg{29}}\CellText{0.29}{1}{0} & \SetCell{bg=\HeatCellBg{32}}\CellText{0.32}{1}{0} & \SetCell{bg=\HeatCellBg{26}}\CellText{0.26}{1}{0} & \SetCell{bg=\HeatCellBg{26}}\CellText{0.26}{1}{0} & \SetCell{bg=\HeatCellBg{26}}\CellText{0.26}{1}{0} & \SetCell{bg=\HeatCellBg{27}}\CellText{0.27}{1}{0} & \SetCell{bg=\HeatCellBg{26}}\CellText{0.26}{1}{0} & \SetCell{bg=\HeatCellBg{25}}\CellText{0.25}{1}{0} & \SetCell{bg=\HeatCellBg{25}}\CellText{0.25}{1}{0} & \SetCell{bg=\HeatCellBg{24}}\CellText{0.24}{1}{0} & \MeanCell{0.32}{0.08}{1} \\
        Mean $\pm$ Std & \MeanCell{0.04}{0.24}{0} & \MeanCell{0.07}{0.16}{0} & \MeanCell{0.07}{0.12}{0} & \MeanCell{0.07}{0.11}{0} & \MeanCell{0.08}{0.10}{0} & \MeanCell{0.08}{0.10}{0} & \MeanCell{0.08}{0.10}{0} & \MeanCell{0.08}{0.10}{0} & \MeanCell{0.08}{0.10}{0} & \MeanCell{0.08}{0.09}{0} & \MeanCell{0.09}{0.10}{1} & \MeanCell{0.08}{0.08}{0} & \MeanCell{0.08}{0.08}{0} & \MeanCell{0.07}{0.08}{0} & \MeanCell{0.08}{0.08}{0} & \MeanCell{0.07}{0.08}{0} & \MeanCell{0.07}{0.08}{0} & \MeanCell{0.07}{0.08}{0} & \MeanCell{0.07}{0.07}{0} & \MeanCell{0.07}{0.11}{0} \\
}


\clearpage
\section{\MsPacManName}
\label{sec:appendix_additional_results_ms_pacman}
\ClusterKTable{10}{\glsfmtshort{cu} --- \MsPacManName.}{cu_ms_pacman}{
        1 & \SetCell{bg=\HeatCellBg{95}}\CellText{0.95}{1}{1} & \SetCell{bg=\HeatCellBg{77}}\CellText{0.77}{0}{0} & \SetCell{bg=\HeatCellBg{90}}\CellText{0.90}{1}{0} & \SetCell{bg=\HeatCellBg{90}}\CellText{0.90}{1}{0} & \SetCell{bg=\HeatCellBg{74}}\CellText{0.74}{1}{0} & \SetCell{bg=\HeatCellBg{87}}\CellText{0.87}{1}{0} & \SetCell{bg=\HeatCellBg{81}}\CellText{0.81}{0}{0} & \SetCell{bg=\HeatCellBg{79}}\CellText{0.79}{0}{0} & \SetCell{bg=\HeatCellBg{80}}\CellText{0.80}{1}{0} & \SetCell{bg=\HeatCellBg{81}}\CellText{0.81}{1}{0} & \SetCell{bg=\HeatCellBg{75}}\CellText{0.75}{0}{0} & \SetCell{bg=\HeatCellBg{74}}\CellText{0.74}{1}{0} & \SetCell{bg=\HeatCellBg{75}}\CellText{0.75}{1}{0} & \SetCell{bg=\HeatCellBg{75}}\CellText{0.75}{1}{0} & \SetCell{bg=\HeatCellBg{68}}\CellText{0.68}{1}{0} & \SetCell{bg=\HeatCellBg{69}}\CellText{0.69}{0}{0} & \SetCell{bg=\HeatCellBg{77}}\CellText{0.77}{1}{0} & \SetCell{bg=\HeatCellBg{66}}\CellText{0.66}{1}{0} & \SetCell{bg=\HeatCellBg{62}}\CellText{0.62}{0}{0} & \MeanCell{0.78}{0.08}{1} \\
        2 & \SetCell{bg=\HeatCellBg{74}}\CellText{0.74}{0}{0} & \SetCell{bg=\HeatCellBg{80}}\CellText{0.80}{0}{0} & \SetCell{bg=\HeatCellBg{90}}\CellText{0.90}{1}{1} & \SetCell{bg=\HeatCellBg{78}}\CellText{0.78}{0}{0} & \SetCell{bg=\HeatCellBg{66}}\CellText{0.66}{0}{0} & \SetCell{bg=\HeatCellBg{75}}\CellText{0.75}{0}{0} & \SetCell{bg=\HeatCellBg{82}}\CellText{0.82}{1}{0} & \SetCell{bg=\HeatCellBg{80}}\CellText{0.80}{1}{0} & \SetCell{bg=\HeatCellBg{68}}\CellText{0.68}{0}{0} & \SetCell{bg=\HeatCellBg{65}}\CellText{0.65}{0}{0} & \SetCell{bg=\HeatCellBg{76}}\CellText{0.76}{1}{0} & \SetCell{bg=\HeatCellBg{66}}\CellText{0.66}{0}{0} & \SetCell{bg=\HeatCellBg{72}}\CellText{0.72}{0}{0} & \SetCell{bg=\HeatCellBg{68}}\CellText{0.68}{0}{0} & \SetCell{bg=\HeatCellBg{67}}\CellText{0.67}{0}{0} & \SetCell{bg=\HeatCellBg{72}}\CellText{0.72}{1}{0} & \SetCell{bg=\HeatCellBg{63}}\CellText{0.63}{0}{0} & \SetCell{bg=\HeatCellBg{63}}\CellText{0.63}{0}{0} & \SetCell{bg=\HeatCellBg{69}}\CellText{0.69}{1}{0} & \MeanCell{0.72}{0.07}{0} \\
        3 & \SetCell{bg=\HeatCellBg{92}}\CellText{0.92}{0}{1} & \SetCell{bg=\HeatCellBg{58}}\CellText{0.58}{0}{0} & \SetCell{bg=\HeatCellBg{83}}\CellText{0.83}{0}{0} & \SetCell{bg=\HeatCellBg{77}}\CellText{0.77}{0}{0} & \SetCell{bg=\HeatCellBg{74}}\CellText{0.74}{1}{0} & \SetCell{bg=\HeatCellBg{73}}\CellText{0.73}{0}{0} & \SetCell{bg=\HeatCellBg{74}}\CellText{0.74}{0}{0} & \SetCell{bg=\HeatCellBg{74}}\CellText{0.74}{0}{0} & \SetCell{bg=\HeatCellBg{63}}\CellText{0.63}{0}{0} & \SetCell{bg=\HeatCellBg{66}}\CellText{0.66}{0}{0} & \SetCell{bg=\HeatCellBg{62}}\CellText{0.62}{0}{0} & \SetCell{bg=\HeatCellBg{64}}\CellText{0.64}{0}{0} & \SetCell{bg=\HeatCellBg{62}}\CellText{0.62}{0}{0} & \SetCell{bg=\HeatCellBg{64}}\CellText{0.64}{0}{0} & \SetCell{bg=\HeatCellBg{65}}\CellText{0.65}{0}{0} & \SetCell{bg=\HeatCellBg{62}}\CellText{0.62}{0}{0} & \SetCell{bg=\HeatCellBg{63}}\CellText{0.63}{0}{0} & \SetCell{bg=\HeatCellBg{63}}\CellText{0.63}{0}{0} & \SetCell{bg=\HeatCellBg{58}}\CellText{0.58}{0}{0} & \MeanCell{0.68}{0.09}{0} \\
        4 & \SetCell{bg=\HeatCellBg{88}}\CellText{0.88}{0}{1} & \SetCell{bg=\HeatCellBg{76}}\CellText{0.76}{0}{0} & \SetCell{bg=\HeatCellBg{73}}\CellText{0.73}{0}{0} & \SetCell{bg=\HeatCellBg{88}}\CellText{0.88}{0}{1} & \SetCell{bg=\HeatCellBg{60}}\CellText{0.60}{0}{0} & \SetCell{bg=\HeatCellBg{81}}\CellText{0.81}{0}{0} & \SetCell{bg=\HeatCellBg{76}}\CellText{0.76}{0}{0} & \SetCell{bg=\HeatCellBg{74}}\CellText{0.74}{0}{0} & \SetCell{bg=\HeatCellBg{72}}\CellText{0.72}{0}{0} & \SetCell{bg=\HeatCellBg{67}}\CellText{0.67}{0}{0} & \SetCell{bg=\HeatCellBg{66}}\CellText{0.66}{0}{0} & \SetCell{bg=\HeatCellBg{61}}\CellText{0.61}{0}{0} & \SetCell{bg=\HeatCellBg{60}}\CellText{0.60}{0}{0} & \SetCell{bg=\HeatCellBg{61}}\CellText{0.61}{0}{0} & \SetCell{bg=\HeatCellBg{60}}\CellText{0.60}{0}{0} & \SetCell{bg=\HeatCellBg{68}}\CellText{0.68}{0}{0} & \SetCell{bg=\HeatCellBg{67}}\CellText{0.67}{0}{0} & \SetCell{bg=\HeatCellBg{59}}\CellText{0.59}{0}{0} & \SetCell{bg=\HeatCellBg{64}}\CellText{0.64}{0}{0} & \MeanCell{0.70}{0.09}{0} \\
        5 & \SetCell{bg=\HeatCellBg{88}}\CellText{0.88}{0}{1} & \SetCell{bg=\HeatCellBg{76}}\CellText{0.76}{0}{0} & \SetCell{bg=\HeatCellBg{73}}\CellText{0.73}{0}{0} & \SetCell{bg=\HeatCellBg{88}}\CellText{0.88}{0}{1} & \SetCell{bg=\HeatCellBg{60}}\CellText{0.60}{0}{0} & \SetCell{bg=\HeatCellBg{81}}\CellText{0.81}{0}{0} & \SetCell{bg=\HeatCellBg{76}}\CellText{0.76}{0}{0} & \SetCell{bg=\HeatCellBg{74}}\CellText{0.74}{0}{0} & \SetCell{bg=\HeatCellBg{72}}\CellText{0.72}{0}{0} & \SetCell{bg=\HeatCellBg{67}}\CellText{0.67}{0}{0} & \SetCell{bg=\HeatCellBg{66}}\CellText{0.66}{0}{0} & \SetCell{bg=\HeatCellBg{61}}\CellText{0.61}{0}{0} & \SetCell{bg=\HeatCellBg{60}}\CellText{0.60}{0}{0} & \SetCell{bg=\HeatCellBg{61}}\CellText{0.61}{0}{0} & \SetCell{bg=\HeatCellBg{60}}\CellText{0.60}{0}{0} & \SetCell{bg=\HeatCellBg{68}}\CellText{0.68}{0}{0} & \SetCell{bg=\HeatCellBg{67}}\CellText{0.67}{0}{0} & \SetCell{bg=\HeatCellBg{59}}\CellText{0.59}{0}{0} & \SetCell{bg=\HeatCellBg{64}}\CellText{0.64}{0}{0} & \MeanCell{0.70}{0.09}{0} \\
        6b & \SetCell{bg=\HeatCellBg{78}}\CellText{0.78}{0}{0} & \SetCell{bg=\HeatCellBg{88}}\CellText{0.88}{1}{1} & \SetCell{bg=\HeatCellBg{80}}\CellText{0.80}{0}{0} & \SetCell{bg=\HeatCellBg{59}}\CellText{0.59}{0}{0} & \SetCell{bg=\HeatCellBg{61}}\CellText{0.61}{0}{0} & \SetCell{bg=\HeatCellBg{68}}\CellText{0.68}{0}{0} & \SetCell{bg=\HeatCellBg{58}}\CellText{0.58}{0}{0} & \SetCell{bg=\HeatCellBg{52}}\CellText{0.52}{0}{0} & \SetCell{bg=\HeatCellBg{51}}\CellText{0.51}{0}{0} & \SetCell{bg=\HeatCellBg{55}}\CellText{0.55}{0}{0} & \SetCell{bg=\HeatCellBg{51}}\CellText{0.51}{0}{0} & \SetCell{bg=\HeatCellBg{48}}\CellText{0.48}{0}{0} & \SetCell{bg=\HeatCellBg{51}}\CellText{0.51}{0}{0} & \SetCell{bg=\HeatCellBg{48}}\CellText{0.48}{0}{0} & \SetCell{bg=\HeatCellBg{52}}\CellText{0.52}{0}{0} & \SetCell{bg=\HeatCellBg{51}}\CellText{0.51}{0}{0} & \SetCell{bg=\HeatCellBg{53}}\CellText{0.53}{0}{0} & \SetCell{bg=\HeatCellBg{48}}\CellText{0.48}{0}{0} & \SetCell{bg=\HeatCellBg{54}}\CellText{0.54}{0}{0} & \MeanCell{0.58}{0.11}{0} \\
        7b & \SetCell{bg=\HeatCellBg{51}}\CellText{0.51}{0}{0} & \SetCell{bg=\HeatCellBg{78}}\CellText{0.78}{0}{1} & \SetCell{bg=\HeatCellBg{65}}\CellText{0.65}{0}{0} & \SetCell{bg=\HeatCellBg{59}}\CellText{0.59}{0}{0} & \SetCell{bg=\HeatCellBg{56}}\CellText{0.56}{0}{0} & \SetCell{bg=\HeatCellBg{59}}\CellText{0.59}{0}{0} & \SetCell{bg=\HeatCellBg{54}}\CellText{0.54}{0}{0} & \SetCell{bg=\HeatCellBg{54}}\CellText{0.54}{0}{0} & \SetCell{bg=\HeatCellBg{56}}\CellText{0.56}{0}{0} & \SetCell{bg=\HeatCellBg{51}}\CellText{0.51}{0}{0} & \SetCell{bg=\HeatCellBg{52}}\CellText{0.52}{0}{0} & \SetCell{bg=\HeatCellBg{47}}\CellText{0.47}{0}{0} & \SetCell{bg=\HeatCellBg{42}}\CellText{0.42}{0}{0} & \SetCell{bg=\HeatCellBg{43}}\CellText{0.43}{0}{0} & \SetCell{bg=\HeatCellBg{44}}\CellText{0.44}{0}{0} & \SetCell{bg=\HeatCellBg{43}}\CellText{0.43}{0}{0} & \SetCell{bg=\HeatCellBg{41}}\CellText{0.41}{0}{0} & \SetCell{bg=\HeatCellBg{40}}\CellText{0.40}{0}{0} & \SetCell{bg=\HeatCellBg{39}}\CellText{0.39}{0}{0} & \MeanCell{0.51}{0.10}{0} \\
        8b & \SetCell{bg=\HeatCellBg{66}}\CellText{0.66}{0}{0} & \SetCell{bg=\HeatCellBg{78}}\CellText{0.78}{0}{1} & \SetCell{bg=\HeatCellBg{58}}\CellText{0.58}{0}{0} & \SetCell{bg=\HeatCellBg{60}}\CellText{0.60}{0}{0} & \SetCell{bg=\HeatCellBg{67}}\CellText{0.67}{0}{0} & \SetCell{bg=\HeatCellBg{64}}\CellText{0.64}{0}{0} & \SetCell{bg=\HeatCellBg{64}}\CellText{0.64}{0}{0} & \SetCell{bg=\HeatCellBg{62}}\CellText{0.62}{0}{0} & \SetCell{bg=\HeatCellBg{59}}\CellText{0.59}{0}{0} & \SetCell{bg=\HeatCellBg{59}}\CellText{0.59}{0}{0} & \SetCell{bg=\HeatCellBg{53}}\CellText{0.53}{0}{0} & \SetCell{bg=\HeatCellBg{62}}\CellText{0.62}{0}{0} & \SetCell{bg=\HeatCellBg{61}}\CellText{0.61}{0}{0} & \SetCell{bg=\HeatCellBg{48}}\CellText{0.48}{0}{0} & \SetCell{bg=\HeatCellBg{48}}\CellText{0.48}{0}{0} & \SetCell{bg=\HeatCellBg{45}}\CellText{0.45}{0}{0} & \SetCell{bg=\HeatCellBg{49}}\CellText{0.49}{0}{0} & \SetCell{bg=\HeatCellBg{48}}\CellText{0.48}{0}{0} & \SetCell{bg=\HeatCellBg{46}}\CellText{0.46}{0}{0} & \MeanCell{0.58}{0.09}{0} \\
        Mean $\pm$ Std & \MeanCell{0.79}{0.14}{1} & \MeanCell{0.76}{0.08}{0} & \MeanCell{0.77}{0.11}{0} & \MeanCell{0.75}{0.13}{0} & \MeanCell{0.65}{0.06}{0} & \MeanCell{0.73}{0.09}{0} & \MeanCell{0.71}{0.10}{0} & \MeanCell{0.69}{0.10}{0} & \MeanCell{0.65}{0.09}{0} & \MeanCell{0.64}{0.09}{0} & \MeanCell{0.63}{0.09}{0} & \MeanCell{0.60}{0.08}{0} & \MeanCell{0.60}{0.10}{0} & \MeanCell{0.58}{0.10}{0} & \MeanCell{0.58}{0.08}{0} & \MeanCell{0.60}{0.11}{0} & \MeanCell{0.60}{0.11}{0} & \MeanCell{0.56}{0.09}{0} & \MeanCell{0.57}{0.09}{0} & \MeanCell{0.66}{0.12}{0} \\
}

\ClusterKTable{10}{\glsfmtshort{as} --- \MsPacManName.}{as_ms_pacman}{
        1 & \SetCell{bg=\HeatCellBg{4}}\CellText{0.04}{0}{0} & \SetCell{bg=\HeatCellBg{14}}\CellText{0.14}{0}{0} & \SetCell{bg=\HeatCellBg{8}}\CellText{0.08}{0}{0} & \SetCell{bg=\HeatCellBg{17}}\CellText{0.17}{0}{0} & \SetCell{bg=\HeatCellBg{20}}\CellText{0.20}{0}{0} & \SetCell{bg=\HeatCellBg{10}}\CellText{0.10}{0}{0} & \SetCell{bg=\HeatCellBg{15}}\CellText{0.15}{0}{0} & \SetCell{bg=\HeatCellBg{12}}\CellText{0.12}{0}{0} & \SetCell{bg=\HeatCellBg{12}}\CellText{0.12}{0}{0} & \SetCell{bg=\HeatCellBg{18}}\CellText{0.18}{0}{0} & \SetCell{bg=\HeatCellBg{18}}\CellText{0.18}{0}{0} & \SetCell{bg=\HeatCellBg{21}}\CellText{0.21}{0}{0} & \SetCell{bg=\HeatCellBg{26}}\CellText{0.26}{0}{0} & \SetCell{bg=\HeatCellBg{29}}\CellText{0.29}{0}{1} & \SetCell{bg=\HeatCellBg{22}}\CellText{0.22}{0}{0} & \SetCell{bg=\HeatCellBg{29}}\CellText{0.29}{0}{1} & \SetCell{bg=\HeatCellBg{23}}\CellText{0.23}{0}{0} & \SetCell{bg=\HeatCellBg{28}}\CellText{0.28}{0}{0} & \SetCell{bg=\HeatCellBg{28}}\CellText{0.28}{0}{0} & \MeanCell{0.19}{0.07}{0} \\
        2 & \SetCell{bg=\HeatCellBg{28}}\CellText{0.28}{0}{0} & \SetCell{bg=\HeatCellBg{25}}\CellText{0.25}{0}{0} & \SetCell{bg=\HeatCellBg{13}}\CellText{0.13}{0}{0} & \SetCell{bg=\HeatCellBg{31}}\CellText{0.31}{0}{0} & \SetCell{bg=\HeatCellBg{34}}\CellText{0.34}{0}{1} & \SetCell{bg=\HeatCellBg{21}}\CellText{0.21}{0}{0} & \SetCell{bg=\HeatCellBg{16}}\CellText{0.16}{0}{0} & \SetCell{bg=\HeatCellBg{16}}\CellText{0.16}{0}{0} & \SetCell{bg=\HeatCellBg{24}}\CellText{0.24}{0}{0} & \SetCell{bg=\HeatCellBg{27}}\CellText{0.27}{0}{0} & \SetCell{bg=\HeatCellBg{24}}\CellText{0.24}{0}{0} & \SetCell{bg=\HeatCellBg{27}}\CellText{0.27}{0}{0} & \SetCell{bg=\HeatCellBg{29}}\CellText{0.29}{0}{0} & \SetCell{bg=\HeatCellBg{28}}\CellText{0.28}{0}{0} & \SetCell{bg=\HeatCellBg{32}}\CellText{0.32}{0}{0} & \SetCell{bg=\HeatCellBg{31}}\CellText{0.31}{0}{0} & \SetCell{bg=\HeatCellBg{33}}\CellText{0.33}{0}{0} & \SetCell{bg=\HeatCellBg{32}}\CellText{0.32}{0}{0} & \SetCell{bg=\HeatCellBg{34}}\CellText{0.34}{0}{1} & \MeanCell{0.27}{0.06}{0} \\
        3 & \SetCell{bg=\HeatCellBg{18}}\CellText{0.18}{0}{0} & \SetCell{bg=\HeatCellBg{34}}\CellText{0.34}{0}{0} & \SetCell{bg=\HeatCellBg{24}}\CellText{0.24}{0}{0} & \SetCell{bg=\HeatCellBg{33}}\CellText{0.33}{0}{0} & \SetCell{bg=\HeatCellBg{32}}\CellText{0.32}{0}{0} & \SetCell{bg=\HeatCellBg{31}}\CellText{0.31}{0}{0} & \SetCell{bg=\HeatCellBg{31}}\CellText{0.31}{0}{0} & \SetCell{bg=\HeatCellBg{31}}\CellText{0.31}{0}{0} & \SetCell{bg=\HeatCellBg{37}}\CellText{0.37}{0}{0} & \SetCell{bg=\HeatCellBg{37}}\CellText{0.37}{0}{0} & \SetCell{bg=\HeatCellBg{39}}\CellText{0.39}{0}{0} & \SetCell{bg=\HeatCellBg{39}}\CellText{0.39}{0}{0} & \SetCell{bg=\HeatCellBg{40}}\CellText{0.40}{0}{0} & \SetCell{bg=\HeatCellBg{41}}\CellText{0.41}{0}{0} & \SetCell{bg=\HeatCellBg{42}}\CellText{0.42}{0}{1} & \SetCell{bg=\HeatCellBg{42}}\CellText{0.42}{0}{1} & \SetCell{bg=\HeatCellBg{38}}\CellText{0.38}{0}{0} & \SetCell{bg=\HeatCellBg{39}}\CellText{0.39}{0}{0} & \SetCell{bg=\HeatCellBg{42}}\CellText{0.42}{0}{1} & \MeanCell{0.35}{0.06}{0} \\
        4 & \SetCell{bg=\HeatCellBg{12}}\CellText{0.12}{0}{0} & \SetCell{bg=\HeatCellBg{52}}\CellText{0.52}{0}{1} & \SetCell{bg=\HeatCellBg{51}}\CellText{0.51}{0}{0} & \SetCell{bg=\HeatCellBg{47}}\CellText{0.47}{0}{0} & \SetCell{bg=\HeatCellBg{47}}\CellText{0.47}{0}{0} & \SetCell{bg=\HeatCellBg{38}}\CellText{0.38}{0}{0} & \SetCell{bg=\HeatCellBg{32}}\CellText{0.32}{0}{0} & \SetCell{bg=\HeatCellBg{35}}\CellText{0.35}{0}{0} & \SetCell{bg=\HeatCellBg{41}}\CellText{0.41}{0}{0} & \SetCell{bg=\HeatCellBg{38}}\CellText{0.38}{0}{0} & \SetCell{bg=\HeatCellBg{38}}\CellText{0.38}{0}{0} & \SetCell{bg=\HeatCellBg{41}}\CellText{0.41}{0}{0} & \SetCell{bg=\HeatCellBg{41}}\CellText{0.41}{0}{0} & \SetCell{bg=\HeatCellBg{40}}\CellText{0.40}{0}{0} & \SetCell{bg=\HeatCellBg{44}}\CellText{0.44}{0}{0} & \SetCell{bg=\HeatCellBg{45}}\CellText{0.45}{0}{0} & \SetCell{bg=\HeatCellBg{43}}\CellText{0.43}{0}{0} & \SetCell{bg=\HeatCellBg{49}}\CellText{0.49}{0}{0} & \SetCell{bg=\HeatCellBg{42}}\CellText{0.42}{0}{0} & \MeanCell{0.41}{0.08}{0} \\
        5 & \SetCell{bg=\HeatCellBg{12}}\CellText{0.12}{0}{0} & \SetCell{bg=\HeatCellBg{52}}\CellText{0.52}{0}{1} & \SetCell{bg=\HeatCellBg{51}}\CellText{0.51}{0}{0} & \SetCell{bg=\HeatCellBg{47}}\CellText{0.47}{0}{0} & \SetCell{bg=\HeatCellBg{47}}\CellText{0.47}{0}{0} & \SetCell{bg=\HeatCellBg{38}}\CellText{0.38}{0}{0} & \SetCell{bg=\HeatCellBg{32}}\CellText{0.32}{0}{0} & \SetCell{bg=\HeatCellBg{35}}\CellText{0.35}{0}{0} & \SetCell{bg=\HeatCellBg{41}}\CellText{0.41}{0}{0} & \SetCell{bg=\HeatCellBg{38}}\CellText{0.38}{0}{0} & \SetCell{bg=\HeatCellBg{38}}\CellText{0.38}{0}{0} & \SetCell{bg=\HeatCellBg{41}}\CellText{0.41}{0}{0} & \SetCell{bg=\HeatCellBg{41}}\CellText{0.41}{0}{0} & \SetCell{bg=\HeatCellBg{40}}\CellText{0.40}{0}{0} & \SetCell{bg=\HeatCellBg{44}}\CellText{0.44}{0}{0} & \SetCell{bg=\HeatCellBg{45}}\CellText{0.45}{0}{0} & \SetCell{bg=\HeatCellBg{43}}\CellText{0.43}{0}{0} & \SetCell{bg=\HeatCellBg{49}}\CellText{0.49}{0}{0} & \SetCell{bg=\HeatCellBg{42}}\CellText{0.42}{0}{0} & \MeanCell{0.41}{0.08}{0} \\
        6b & \SetCell{bg=\HeatCellBg{19}}\CellText{0.19}{0}{0} & \SetCell{bg=\HeatCellBg{52}}\CellText{0.52}{0}{0} & \SetCell{bg=\HeatCellBg{56}}\CellText{0.56}{0}{0} & \SetCell{bg=\HeatCellBg{60}}\CellText{0.60}{1}{0} & \SetCell{bg=\HeatCellBg{59}}\CellText{0.59}{0}{0} & \SetCell{bg=\HeatCellBg{55}}\CellText{0.55}{0}{0} & \SetCell{bg=\HeatCellBg{55}}\CellText{0.55}{0}{0} & \SetCell{bg=\HeatCellBg{58}}\CellText{0.58}{0}{0} & \SetCell{bg=\HeatCellBg{62}}\CellText{0.62}{0}{1} & \SetCell{bg=\HeatCellBg{58}}\CellText{0.58}{0}{0} & \SetCell{bg=\HeatCellBg{57}}\CellText{0.57}{0}{0} & \SetCell{bg=\HeatCellBg{56}}\CellText{0.56}{0}{0} & \SetCell{bg=\HeatCellBg{60}}\CellText{0.60}{0}{0} & \SetCell{bg=\HeatCellBg{59}}\CellText{0.59}{0}{0} & \SetCell{bg=\HeatCellBg{55}}\CellText{0.55}{0}{0} & \SetCell{bg=\HeatCellBg{57}}\CellText{0.57}{0}{0} & \SetCell{bg=\HeatCellBg{55}}\CellText{0.55}{0}{0} & \SetCell{bg=\HeatCellBg{59}}\CellText{0.59}{0}{0} & \SetCell{bg=\HeatCellBg{58}}\CellText{0.58}{0}{0} & \MeanCell{0.55}{0.09}{0} \\
        7b & \SetCell{bg=\HeatCellBg{42}}\CellText{0.42}{0}{0} & \SetCell{bg=\HeatCellBg{54}}\CellText{0.54}{0}{0} & \SetCell{bg=\HeatCellBg{61}}\CellText{0.61}{1}{0} & \SetCell{bg=\HeatCellBg{59}}\CellText{0.59}{0}{0} & \SetCell{bg=\HeatCellBg{67}}\CellText{0.67}{1}{0} & \SetCell{bg=\HeatCellBg{63}}\CellText{0.63}{1}{0} & \SetCell{bg=\HeatCellBg{65}}\CellText{0.65}{1}{0} & \SetCell{bg=\HeatCellBg{65}}\CellText{0.65}{1}{0} & \SetCell{bg=\HeatCellBg{66}}\CellText{0.66}{1}{0} & \SetCell{bg=\HeatCellBg{65}}\CellText{0.65}{1}{0} & \SetCell{bg=\HeatCellBg{65}}\CellText{0.65}{1}{0} & \SetCell{bg=\HeatCellBg{63}}\CellText{0.63}{1}{0} & \SetCell{bg=\HeatCellBg{69}}\CellText{0.69}{1}{0} & \SetCell{bg=\HeatCellBg{67}}\CellText{0.67}{1}{0} & \SetCell{bg=\HeatCellBg{68}}\CellText{0.68}{1}{0} & \SetCell{bg=\HeatCellBg{67}}\CellText{0.67}{0}{0} & \SetCell{bg=\HeatCellBg{67}}\CellText{0.67}{1}{0} & \SetCell{bg=\HeatCellBg{68}}\CellText{0.68}{1}{0} & \SetCell{bg=\HeatCellBg{70}}\CellText{0.70}{1}{1} & \MeanCell{0.64}{0.06}{1} \\
        8b & \SetCell{bg=\HeatCellBg{56}}\CellText{0.56}{1}{0} & \SetCell{bg=\HeatCellBg{59}}\CellText{0.59}{1}{0} & \SetCell{bg=\HeatCellBg{47}}\CellText{0.47}{0}{0} & \SetCell{bg=\HeatCellBg{57}}\CellText{0.57}{0}{0} & \SetCell{bg=\HeatCellBg{56}}\CellText{0.56}{0}{0} & \SetCell{bg=\HeatCellBg{54}}\CellText{0.54}{0}{0} & \SetCell{bg=\HeatCellBg{55}}\CellText{0.55}{0}{0} & \SetCell{bg=\HeatCellBg{56}}\CellText{0.56}{0}{0} & \SetCell{bg=\HeatCellBg{60}}\CellText{0.60}{0}{0} & \SetCell{bg=\HeatCellBg{58}}\CellText{0.58}{0}{0} & \SetCell{bg=\HeatCellBg{62}}\CellText{0.62}{0}{0} & \SetCell{bg=\HeatCellBg{59}}\CellText{0.59}{0}{0} & \SetCell{bg=\HeatCellBg{59}}\CellText{0.59}{0}{0} & \SetCell{bg=\HeatCellBg{65}}\CellText{0.65}{0}{0} & \SetCell{bg=\HeatCellBg{65}}\CellText{0.65}{0}{0} & \SetCell{bg=\HeatCellBg{68}}\CellText{0.68}{1}{1} & \SetCell{bg=\HeatCellBg{66}}\CellText{0.66}{0}{0} & \SetCell{bg=\HeatCellBg{66}}\CellText{0.66}{0}{0} & \SetCell{bg=\HeatCellBg{67}}\CellText{0.67}{0}{0} & \MeanCell{0.60}{0.05}{0} \\
        Mean $\pm$ Std & \MeanCell{0.24}{0.16}{0} & \MeanCell{0.43}{0.15}{0} & \MeanCell{0.39}{0.19}{0} & \MeanCell{0.44}{0.15}{0} & \MeanCell{0.45}{0.15}{0} & \MeanCell{0.39}{0.17}{0} & \MeanCell{0.38}{0.17}{0} & \MeanCell{0.39}{0.18}{0} & \MeanCell{0.43}{0.18}{0} & \MeanCell{0.42}{0.15}{0} & \MeanCell{0.43}{0.16}{0} & \MeanCell{0.43}{0.14}{0} & \MeanCell{0.46}{0.14}{0} & \MeanCell{0.46}{0.14}{0} & \MeanCell{0.47}{0.15}{0} & \MeanCell{0.48}{0.14}{0} & \MeanCell{0.46}{0.15}{0} & \MeanCell{0.49}{0.14}{1} & \MeanCell{0.48}{0.14}{0} & \MeanCell{0.43}{0.17}{0} \\
}

\ClusterKTable{10}{\glsfmtshort{acu} --- \MsPacManName.}{acu_ms_pacman}{
        1 & \SetCell{bg=\HeatCellBg{-1}}\CellText{-0.01}{0}{0} & \SetCell{bg=\HeatCellBg{-10}}\CellText{-0.10}{0}{0} & \SetCell{bg=\HeatCellBg{-2}}\CellText{-0.02}{0}{0} & \SetCell{bg=\HeatCellBg{8}}\CellText{0.08}{0}{1} & \SetCell{bg=\HeatCellBg{-6}}\CellText{-0.06}{0}{0} & \SetCell{bg=\HeatCellBg{-3}}\CellText{-0.03}{0}{0} & \SetCell{bg=\HeatCellBg{-4}}\CellText{-0.04}{0}{0} & \SetCell{bg=\HeatCellBg{-9}}\CellText{-0.09}{0}{0} & \SetCell{bg=\HeatCellBg{-8}}\CellText{-0.08}{0}{0} & \SetCell{bg=\HeatCellBg{-1}}\CellText{-0.01}{0}{0} & \SetCell{bg=\HeatCellBg{-7}}\CellText{-0.07}{0}{0} & \SetCell{bg=\HeatCellBg{-5}}\CellText{-0.05}{0}{0} & \SetCell{bg=\HeatCellBg{1}}\CellText{0.01}{0}{0} & \SetCell{bg=\HeatCellBg{3}}\CellText{0.03}{0}{0} & \SetCell{bg=\HeatCellBg{-10}}\CellText{-0.10}{0}{0} & \SetCell{bg=\HeatCellBg{-2}}\CellText{-0.02}{0}{0} & \SetCell{bg=\HeatCellBg{0}}\CellText{0.00}{0}{0} & \SetCell{bg=\HeatCellBg{-6}}\CellText{-0.06}{0}{0} & \SetCell{bg=\HeatCellBg{-10}}\CellText{-0.10}{0}{0} & \MeanCell{-0.04}{0.05}{0} \\
        2 & \SetCell{bg=\HeatCellBg{3}}\CellText{0.03}{0}{0} & \SetCell{bg=\HeatCellBg{5}}\CellText{0.05}{0}{0} & \SetCell{bg=\HeatCellBg{3}}\CellText{0.03}{0}{0} & \SetCell{bg=\HeatCellBg{9}}\CellText{0.09}{0}{1} & \SetCell{bg=\HeatCellBg{1}}\CellText{0.01}{0}{0} & \SetCell{bg=\HeatCellBg{-4}}\CellText{-0.04}{0}{0} & \SetCell{bg=\HeatCellBg{-2}}\CellText{-0.02}{0}{0} & \SetCell{bg=\HeatCellBg{-4}}\CellText{-0.04}{0}{0} & \SetCell{bg=\HeatCellBg{-9}}\CellText{-0.09}{0}{0} & \SetCell{bg=\HeatCellBg{-9}}\CellText{-0.09}{0}{0} & \SetCell{bg=\HeatCellBg{1}}\CellText{0.01}{0}{0} & \SetCell{bg=\HeatCellBg{-7}}\CellText{-0.07}{0}{0} & \SetCell{bg=\HeatCellBg{1}}\CellText{0.01}{0}{0} & \SetCell{bg=\HeatCellBg{-4}}\CellText{-0.04}{0}{0} & \SetCell{bg=\HeatCellBg{-1}}\CellText{-0.01}{0}{0} & \SetCell{bg=\HeatCellBg{4}}\CellText{0.04}{0}{0} & \SetCell{bg=\HeatCellBg{-5}}\CellText{-0.05}{0}{0} & \SetCell{bg=\HeatCellBg{-6}}\CellText{-0.06}{0}{0} & \SetCell{bg=\HeatCellBg{3}}\CellText{0.03}{0}{0} & \MeanCell{-0.01}{0.05}{0} \\
        3 & \SetCell{bg=\HeatCellBg{10}}\CellText{0.10}{0}{1} & \SetCell{bg=\HeatCellBg{-8}}\CellText{-0.08}{0}{0} & \SetCell{bg=\HeatCellBg{7}}\CellText{0.07}{0}{0} & \SetCell{bg=\HeatCellBg{10}}\CellText{0.10}{0}{1} & \SetCell{bg=\HeatCellBg{5}}\CellText{0.05}{0}{0} & \SetCell{bg=\HeatCellBg{5}}\CellText{0.05}{0}{0} & \SetCell{bg=\HeatCellBg{5}}\CellText{0.05}{0}{0} & \SetCell{bg=\HeatCellBg{5}}\CellText{0.05}{0}{0} & \SetCell{bg=\HeatCellBg{0}}\CellText{0.00}{0}{0} & \SetCell{bg=\HeatCellBg{3}}\CellText{0.03}{0}{0} & \SetCell{bg=\HeatCellBg{1}}\CellText{0.01}{0}{0} & \SetCell{bg=\HeatCellBg{3}}\CellText{0.03}{0}{0} & \SetCell{bg=\HeatCellBg{2}}\CellText{0.02}{0}{0} & \SetCell{bg=\HeatCellBg{5}}\CellText{0.05}{0}{0} & \SetCell{bg=\HeatCellBg{7}}\CellText{0.07}{0}{0} & \SetCell{bg=\HeatCellBg{4}}\CellText{0.04}{0}{0} & \SetCell{bg=\HeatCellBg{1}}\CellText{0.01}{0}{0} & \SetCell{bg=\HeatCellBg{3}}\CellText{0.03}{0}{0} & \SetCell{bg=\HeatCellBg{0}}\CellText{0.00}{0}{0} & \MeanCell{0.04}{0.04}{0} \\
        4 & \SetCell{bg=\HeatCellBg{0}}\CellText{0.00}{0}{0} & \SetCell{bg=\HeatCellBg{28}}\CellText{0.28}{0}{0} & \SetCell{bg=\HeatCellBg{25}}\CellText{0.25}{0}{0} & \SetCell{bg=\HeatCellBg{35}}\CellText{0.35}{1}{1} & \SetCell{bg=\HeatCellBg{7}}\CellText{0.07}{0}{0} & \SetCell{bg=\HeatCellBg{19}}\CellText{0.19}{0}{0} & \SetCell{bg=\HeatCellBg{8}}\CellText{0.08}{0}{0} & \SetCell{bg=\HeatCellBg{10}}\CellText{0.10}{0}{0} & \SetCell{bg=\HeatCellBg{13}}\CellText{0.13}{0}{0} & \SetCell{bg=\HeatCellBg{6}}\CellText{0.06}{0}{0} & \SetCell{bg=\HeatCellBg{4}}\CellText{0.04}{0}{0} & \SetCell{bg=\HeatCellBg{1}}\CellText{0.01}{0}{0} & \SetCell{bg=\HeatCellBg{1}}\CellText{0.01}{0}{0} & \SetCell{bg=\HeatCellBg{1}}\CellText{0.01}{0}{0} & \SetCell{bg=\HeatCellBg{5}}\CellText{0.05}{0}{0} & \SetCell{bg=\HeatCellBg{12}}\CellText{0.12}{0}{0} & \SetCell{bg=\HeatCellBg{9}}\CellText{0.09}{0}{0} & \SetCell{bg=\HeatCellBg{8}}\CellText{0.08}{0}{0} & \SetCell{bg=\HeatCellBg{6}}\CellText{0.06}{0}{0} & \MeanCell{0.10}{0.10}{0} \\
        5 & \SetCell{bg=\HeatCellBg{0}}\CellText{0.00}{0}{0} & \SetCell{bg=\HeatCellBg{28}}\CellText{0.28}{0}{0} & \SetCell{bg=\HeatCellBg{25}}\CellText{0.25}{0}{0} & \SetCell{bg=\HeatCellBg{35}}\CellText{0.35}{1}{1} & \SetCell{bg=\HeatCellBg{7}}\CellText{0.07}{0}{0} & \SetCell{bg=\HeatCellBg{19}}\CellText{0.19}{0}{0} & \SetCell{bg=\HeatCellBg{8}}\CellText{0.08}{0}{0} & \SetCell{bg=\HeatCellBg{10}}\CellText{0.10}{0}{0} & \SetCell{bg=\HeatCellBg{13}}\CellText{0.13}{0}{0} & \SetCell{bg=\HeatCellBg{6}}\CellText{0.06}{0}{0} & \SetCell{bg=\HeatCellBg{4}}\CellText{0.04}{0}{0} & \SetCell{bg=\HeatCellBg{1}}\CellText{0.01}{0}{0} & \SetCell{bg=\HeatCellBg{1}}\CellText{0.01}{0}{0} & \SetCell{bg=\HeatCellBg{1}}\CellText{0.01}{0}{0} & \SetCell{bg=\HeatCellBg{5}}\CellText{0.05}{0}{0} & \SetCell{bg=\HeatCellBg{12}}\CellText{0.12}{0}{0} & \SetCell{bg=\HeatCellBg{9}}\CellText{0.09}{0}{0} & \SetCell{bg=\HeatCellBg{8}}\CellText{0.08}{0}{0} & \SetCell{bg=\HeatCellBg{6}}\CellText{0.06}{0}{0} & \MeanCell{0.10}{0.10}{0} \\
        6b & \SetCell{bg=\HeatCellBg{-3}}\CellText{-0.03}{0}{0} & \SetCell{bg=\HeatCellBg{40}}\CellText{0.40}{1}{1} & \SetCell{bg=\HeatCellBg{35}}\CellText{0.35}{1}{0} & \SetCell{bg=\HeatCellBg{20}}\CellText{0.20}{0}{0} & \SetCell{bg=\HeatCellBg{20}}\CellText{0.20}{0}{0} & \SetCell{bg=\HeatCellBg{23}}\CellText{0.23}{1}{0} & \SetCell{bg=\HeatCellBg{13}}\CellText{0.13}{0}{0} & \SetCell{bg=\HeatCellBg{11}}\CellText{0.11}{0}{0} & \SetCell{bg=\HeatCellBg{14}}\CellText{0.14}{0}{0} & \SetCell{bg=\HeatCellBg{13}}\CellText{0.13}{0}{0} & \SetCell{bg=\HeatCellBg{8}}\CellText{0.08}{0}{0} & \SetCell{bg=\HeatCellBg{4}}\CellText{0.04}{0}{0} & \SetCell{bg=\HeatCellBg{10}}\CellText{0.10}{0}{0} & \SetCell{bg=\HeatCellBg{7}}\CellText{0.07}{0}{0} & \SetCell{bg=\HeatCellBg{7}}\CellText{0.07}{0}{0} & \SetCell{bg=\HeatCellBg{8}}\CellText{0.08}{0}{0} & \SetCell{bg=\HeatCellBg{8}}\CellText{0.08}{0}{0} & \SetCell{bg=\HeatCellBg{7}}\CellText{0.07}{0}{0} & \SetCell{bg=\HeatCellBg{12}}\CellText{0.12}{0}{0} & \MeanCell{0.14}{0.10}{0} \\
        7b & \SetCell{bg=\HeatCellBg{-7}}\CellText{-0.07}{0}{0} & \SetCell{bg=\HeatCellBg{32}}\CellText{0.32}{0}{1} & \SetCell{bg=\HeatCellBg{25}}\CellText{0.25}{0}{0} & \SetCell{bg=\HeatCellBg{18}}\CellText{0.18}{0}{0} & \SetCell{bg=\HeatCellBg{23}}\CellText{0.23}{1}{0} & \SetCell{bg=\HeatCellBg{22}}\CellText{0.22}{0}{0} & \SetCell{bg=\HeatCellBg{18}}\CellText{0.18}{0}{0} & \SetCell{bg=\HeatCellBg{19}}\CellText{0.19}{1}{0} & \SetCell{bg=\HeatCellBg{22}}\CellText{0.22}{1}{0} & \SetCell{bg=\HeatCellBg{17}}\CellText{0.17}{1}{0} & \SetCell{bg=\HeatCellBg{17}}\CellText{0.17}{1}{0} & \SetCell{bg=\HeatCellBg{10}}\CellText{0.10}{0}{0} & \SetCell{bg=\HeatCellBg{11}}\CellText{0.11}{0}{0} & \SetCell{bg=\HeatCellBg{10}}\CellText{0.10}{0}{0} & \SetCell{bg=\HeatCellBg{11}}\CellText{0.11}{0}{0} & \SetCell{bg=\HeatCellBg{10}}\CellText{0.10}{0}{0} & \SetCell{bg=\HeatCellBg{8}}\CellText{0.08}{0}{0} & \SetCell{bg=\HeatCellBg{9}}\CellText{0.09}{0}{0} & \SetCell{bg=\HeatCellBg{9}}\CellText{0.09}{0}{0} & \MeanCell{0.15}{0.08}{0} \\
        8b & \SetCell{bg=\HeatCellBg{22}}\CellText{0.22}{1}{0} & \SetCell{bg=\HeatCellBg{37}}\CellText{0.37}{0}{1} & \SetCell{bg=\HeatCellBg{5}}\CellText{0.05}{0}{0} & \SetCell{bg=\HeatCellBg{16}}\CellText{0.16}{0}{0} & \SetCell{bg=\HeatCellBg{23}}\CellText{0.23}{1}{0} & \SetCell{bg=\HeatCellBg{18}}\CellText{0.18}{0}{0} & \SetCell{bg=\HeatCellBg{19}}\CellText{0.19}{1}{0} & \SetCell{bg=\HeatCellBg{19}}\CellText{0.19}{1}{0} & \SetCell{bg=\HeatCellBg{19}}\CellText{0.19}{0}{0} & \SetCell{bg=\HeatCellBg{17}}\CellText{0.17}{1}{0} & \SetCell{bg=\HeatCellBg{15}}\CellText{0.15}{0}{0} & \SetCell{bg=\HeatCellBg{21}}\CellText{0.21}{1}{0} & \SetCell{bg=\HeatCellBg{20}}\CellText{0.20}{1}{0} & \SetCell{bg=\HeatCellBg{13}}\CellText{0.13}{1}{0} & \SetCell{bg=\HeatCellBg{12}}\CellText{0.12}{1}{0} & \SetCell{bg=\HeatCellBg{13}}\CellText{0.13}{1}{0} & \SetCell{bg=\HeatCellBg{15}}\CellText{0.15}{1}{0} & \SetCell{bg=\HeatCellBg{14}}\CellText{0.14}{1}{0} & \SetCell{bg=\HeatCellBg{13}}\CellText{0.13}{1}{0} & \MeanCell{0.17}{0.06}{1} \\
        Mean $\pm$ Std & \MeanCell{0.03}{0.09}{0} & \MeanCell{0.19}{0.19}{1} & \MeanCell{0.15}{0.13}{0} & \MeanCell{0.19}{0.10}{0} & \MeanCell{0.10}{0.10}{0} & \MeanCell{0.12}{0.11}{0} & \MeanCell{0.08}{0.08}{0} & \MeanCell{0.08}{0.09}{0} & \MeanCell{0.08}{0.11}{0} & \MeanCell{0.07}{0.08}{0} & \MeanCell{0.05}{0.07}{0} & \MeanCell{0.03}{0.08}{0} & \MeanCell{0.06}{0.07}{0} & \MeanCell{0.05}{0.05}{0} & \MeanCell{0.04}{0.07}{0} & \MeanCell{0.08}{0.05}{0} & \MeanCell{0.06}{0.06}{0} & \MeanCell{0.05}{0.07}{0} & \MeanCell{0.05}{0.07}{0} & \MeanCell{0.08}{0.10}{0} \\
}

\ClusterKTable{10}[Observation baseline --- \MsPacManName.]{Observation baseline ($\Similarity{obs}{}$) --- \MsPacManName.}{obs_ms_pacman}{
        1 & \SetCell{bg=\HeatCellBg{55}}\CellText{0.55}{0}{0} & \SetCell{bg=\HeatCellBg{84}}\CellText{0.84}{1}{0} & \SetCell{bg=\HeatCellBg{48}}\CellText{0.48}{0}{0} & \SetCell{bg=\HeatCellBg{65}}\CellText{0.65}{1}{0} & \SetCell{bg=\HeatCellBg{75}}\CellText{0.75}{1}{0} & \SetCell{bg=\HeatCellBg{73}}\CellText{0.73}{1}{0} & \SetCell{bg=\HeatCellBg{85}}\CellText{0.85}{1}{0} & \SetCell{bg=\HeatCellBg{86}}\CellText{0.86}{1}{1} & \SetCell{bg=\HeatCellBg{86}}\CellText{0.86}{1}{1} & \SetCell{bg=\HeatCellBg{81}}\CellText{0.81}{1}{0} & \SetCell{bg=\HeatCellBg{82}}\CellText{0.82}{1}{0} & \SetCell{bg=\HeatCellBg{75}}\CellText{0.75}{1}{0} & \SetCell{bg=\HeatCellBg{63}}\CellText{0.63}{0}{0} & \SetCell{bg=\HeatCellBg{64}}\CellText{0.64}{0}{0} & \SetCell{bg=\HeatCellBg{73}}\CellText{0.73}{1}{0} & \SetCell{bg=\HeatCellBg{70}}\CellText{0.70}{1}{0} & \SetCell{bg=\HeatCellBg{77}}\CellText{0.77}{1}{0} & \SetCell{bg=\HeatCellBg{72}}\CellText{0.72}{1}{0} & \SetCell{bg=\HeatCellBg{67}}\CellText{0.67}{1}{0} & \MeanCell{0.73}{0.10}{1} \\
        2 & \SetCell{bg=\HeatCellBg{61}}\CellText{0.61}{0}{0} & \SetCell{bg=\HeatCellBg{68}}\CellText{0.68}{0}{0} & \SetCell{bg=\HeatCellBg{50}}\CellText{0.50}{0}{0} & \SetCell{bg=\HeatCellBg{60}}\CellText{0.60}{0}{0} & \SetCell{bg=\HeatCellBg{55}}\CellText{0.55}{0}{0} & \SetCell{bg=\HeatCellBg{70}}\CellText{0.70}{0}{0} & \SetCell{bg=\HeatCellBg{83}}\CellText{0.83}{0}{0} & \SetCell{bg=\HeatCellBg{84}}\CellText{0.84}{0}{1} & \SetCell{bg=\HeatCellBg{73}}\CellText{0.73}{0}{0} & \SetCell{bg=\HeatCellBg{73}}\CellText{0.73}{0}{0} & \SetCell{bg=\HeatCellBg{71}}\CellText{0.71}{0}{0} & \SetCell{bg=\HeatCellBg{69}}\CellText{0.69}{0}{0} & \SetCell{bg=\HeatCellBg{66}}\CellText{0.66}{1}{0} & \SetCell{bg=\HeatCellBg{66}}\CellText{0.66}{1}{0} & \SetCell{bg=\HeatCellBg{63}}\CellText{0.63}{0}{0} & \SetCell{bg=\HeatCellBg{68}}\CellText{0.68}{0}{0} & \SetCell{bg=\HeatCellBg{67}}\CellText{0.67}{0}{0} & \SetCell{bg=\HeatCellBg{65}}\CellText{0.65}{0}{0} & \SetCell{bg=\HeatCellBg{66}}\CellText{0.66}{0}{0} & \MeanCell{0.67}{0.08}{0} \\
        3 & \SetCell{bg=\HeatCellBg{82}}\CellText{0.82}{0}{1} & \SetCell{bg=\HeatCellBg{53}}\CellText{0.53}{0}{0} & \SetCell{bg=\HeatCellBg{55}}\CellText{0.55}{1}{0} & \SetCell{bg=\HeatCellBg{54}}\CellText{0.54}{0}{0} & \SetCell{bg=\HeatCellBg{54}}\CellText{0.54}{0}{0} & \SetCell{bg=\HeatCellBg{60}}\CellText{0.60}{0}{0} & \SetCell{bg=\HeatCellBg{65}}\CellText{0.65}{0}{0} & \SetCell{bg=\HeatCellBg{65}}\CellText{0.65}{0}{0} & \SetCell{bg=\HeatCellBg{55}}\CellText{0.55}{0}{0} & \SetCell{bg=\HeatCellBg{60}}\CellText{0.60}{0}{0} & \SetCell{bg=\HeatCellBg{58}}\CellText{0.58}{0}{0} & \SetCell{bg=\HeatCellBg{59}}\CellText{0.59}{0}{0} & \SetCell{bg=\HeatCellBg{57}}\CellText{0.57}{0}{0} & \SetCell{bg=\HeatCellBg{56}}\CellText{0.56}{0}{0} & \SetCell{bg=\HeatCellBg{57}}\CellText{0.57}{0}{0} & \SetCell{bg=\HeatCellBg{55}}\CellText{0.55}{0}{0} & \SetCell{bg=\HeatCellBg{62}}\CellText{0.62}{0}{0} & \SetCell{bg=\HeatCellBg{57}}\CellText{0.57}{0}{0} & \SetCell{bg=\HeatCellBg{56}}\CellText{0.56}{0}{0} & \MeanCell{0.59}{0.06}{0} \\
        4 & \SetCell{bg=\HeatCellBg{88}}\CellText{0.88}{1}{1} & \SetCell{bg=\HeatCellBg{42}}\CellText{0.42}{0}{0} & \SetCell{bg=\HeatCellBg{39}}\CellText{0.39}{0}{0} & \SetCell{bg=\HeatCellBg{45}}\CellText{0.45}{0}{0} & \SetCell{bg=\HeatCellBg{50}}\CellText{0.50}{0}{0} & \SetCell{bg=\HeatCellBg{62}}\CellText{0.62}{0}{0} & \SetCell{bg=\HeatCellBg{68}}\CellText{0.68}{0}{0} & \SetCell{bg=\HeatCellBg{63}}\CellText{0.63}{0}{0} & \SetCell{bg=\HeatCellBg{56}}\CellText{0.56}{0}{0} & \SetCell{bg=\HeatCellBg{60}}\CellText{0.60}{0}{0} & \SetCell{bg=\HeatCellBg{59}}\CellText{0.59}{0}{0} & \SetCell{bg=\HeatCellBg{57}}\CellText{0.57}{0}{0} & \SetCell{bg=\HeatCellBg{59}}\CellText{0.59}{0}{0} & \SetCell{bg=\HeatCellBg{54}}\CellText{0.54}{0}{0} & \SetCell{bg=\HeatCellBg{54}}\CellText{0.54}{0}{0} & \SetCell{bg=\HeatCellBg{55}}\CellText{0.55}{0}{0} & \SetCell{bg=\HeatCellBg{57}}\CellText{0.57}{0}{0} & \SetCell{bg=\HeatCellBg{50}}\CellText{0.50}{0}{0} & \SetCell{bg=\HeatCellBg{58}}\CellText{0.58}{0}{0} & \MeanCell{0.57}{0.10}{0} \\
        5 & \SetCell{bg=\HeatCellBg{88}}\CellText{0.88}{1}{1} & \SetCell{bg=\HeatCellBg{42}}\CellText{0.42}{0}{0} & \SetCell{bg=\HeatCellBg{39}}\CellText{0.39}{0}{0} & \SetCell{bg=\HeatCellBg{45}}\CellText{0.45}{0}{0} & \SetCell{bg=\HeatCellBg{50}}\CellText{0.50}{0}{0} & \SetCell{bg=\HeatCellBg{62}}\CellText{0.62}{0}{0} & \SetCell{bg=\HeatCellBg{68}}\CellText{0.68}{0}{0} & \SetCell{bg=\HeatCellBg{63}}\CellText{0.63}{0}{0} & \SetCell{bg=\HeatCellBg{56}}\CellText{0.56}{0}{0} & \SetCell{bg=\HeatCellBg{60}}\CellText{0.60}{0}{0} & \SetCell{bg=\HeatCellBg{59}}\CellText{0.59}{0}{0} & \SetCell{bg=\HeatCellBg{57}}\CellText{0.57}{0}{0} & \SetCell{bg=\HeatCellBg{59}}\CellText{0.59}{0}{0} & \SetCell{bg=\HeatCellBg{54}}\CellText{0.54}{0}{0} & \SetCell{bg=\HeatCellBg{54}}\CellText{0.54}{0}{0} & \SetCell{bg=\HeatCellBg{55}}\CellText{0.55}{0}{0} & \SetCell{bg=\HeatCellBg{57}}\CellText{0.57}{0}{0} & \SetCell{bg=\HeatCellBg{50}}\CellText{0.50}{0}{0} & \SetCell{bg=\HeatCellBg{58}}\CellText{0.58}{0}{0} & \MeanCell{0.57}{0.10}{0} \\
        6b & \SetCell{bg=\HeatCellBg{81}}\CellText{0.81}{0}{1} & \SetCell{bg=\HeatCellBg{40}}\CellText{0.40}{0}{0} & \SetCell{bg=\HeatCellBg{34}}\CellText{0.34}{0}{0} & \SetCell{bg=\HeatCellBg{38}}\CellText{0.38}{0}{0} & \SetCell{bg=\HeatCellBg{37}}\CellText{0.37}{0}{0} & \SetCell{bg=\HeatCellBg{43}}\CellText{0.43}{0}{0} & \SetCell{bg=\HeatCellBg{45}}\CellText{0.45}{0}{0} & \SetCell{bg=\HeatCellBg{42}}\CellText{0.42}{0}{0} & \SetCell{bg=\HeatCellBg{37}}\CellText{0.37}{0}{0} & \SetCell{bg=\HeatCellBg{42}}\CellText{0.42}{0}{0} & \SetCell{bg=\HeatCellBg{43}}\CellText{0.43}{0}{0} & \SetCell{bg=\HeatCellBg{44}}\CellText{0.44}{0}{0} & \SetCell{bg=\HeatCellBg{40}}\CellText{0.40}{0}{0} & \SetCell{bg=\HeatCellBg{41}}\CellText{0.41}{0}{0} & \SetCell{bg=\HeatCellBg{45}}\CellText{0.45}{0}{0} & \SetCell{bg=\HeatCellBg{43}}\CellText{0.43}{0}{0} & \SetCell{bg=\HeatCellBg{45}}\CellText{0.45}{0}{0} & \SetCell{bg=\HeatCellBg{41}}\CellText{0.41}{0}{0} & \SetCell{bg=\HeatCellBg{42}}\CellText{0.42}{0}{0} & \MeanCell{0.43}{0.09}{0} \\
        7b & \SetCell{bg=\HeatCellBg{58}}\CellText{0.58}{0}{1} & \SetCell{bg=\HeatCellBg{39}}\CellText{0.39}{0}{0} & \SetCell{bg=\HeatCellBg{30}}\CellText{0.30}{0}{0} & \SetCell{bg=\HeatCellBg{36}}\CellText{0.36}{0}{0} & \SetCell{bg=\HeatCellBg{31}}\CellText{0.31}{0}{0} & \SetCell{bg=\HeatCellBg{37}}\CellText{0.37}{0}{0} & \SetCell{bg=\HeatCellBg{35}}\CellText{0.35}{0}{0} & \SetCell{bg=\HeatCellBg{35}}\CellText{0.35}{0}{0} & \SetCell{bg=\HeatCellBg{34}}\CellText{0.34}{0}{0} & \SetCell{bg=\HeatCellBg{35}}\CellText{0.35}{0}{0} & \SetCell{bg=\HeatCellBg{35}}\CellText{0.35}{0}{0} & \SetCell{bg=\HeatCellBg{37}}\CellText{0.37}{0}{0} & \SetCell{bg=\HeatCellBg{31}}\CellText{0.31}{0}{0} & \SetCell{bg=\HeatCellBg{33}}\CellText{0.33}{0}{0} & \SetCell{bg=\HeatCellBg{32}}\CellText{0.32}{0}{0} & \SetCell{bg=\HeatCellBg{33}}\CellText{0.33}{0}{0} & \SetCell{bg=\HeatCellBg{33}}\CellText{0.33}{0}{0} & \SetCell{bg=\HeatCellBg{32}}\CellText{0.32}{0}{0} & \SetCell{bg=\HeatCellBg{30}}\CellText{0.30}{0}{0} & \MeanCell{0.35}{0.06}{0} \\
        8b & \SetCell{bg=\HeatCellBg{38}}\CellText{0.38}{0}{1} & \SetCell{bg=\HeatCellBg{14}}\CellText{0.14}{0}{0} & \SetCell{bg=\HeatCellBg{23}}\CellText{0.23}{0}{0} & \SetCell{bg=\HeatCellBg{31}}\CellText{0.31}{0}{0} & \SetCell{bg=\HeatCellBg{27}}\CellText{0.27}{0}{0} & \SetCell{bg=\HeatCellBg{30}}\CellText{0.30}{0}{0} & \SetCell{bg=\HeatCellBg{23}}\CellText{0.23}{0}{0} & \SetCell{bg=\HeatCellBg{22}}\CellText{0.22}{0}{0} & \SetCell{bg=\HeatCellBg{25}}\CellText{0.25}{0}{0} & \SetCell{bg=\HeatCellBg{25}}\CellText{0.25}{0}{0} & \SetCell{bg=\HeatCellBg{26}}\CellText{0.26}{0}{0} & \SetCell{bg=\HeatCellBg{24}}\CellText{0.24}{0}{0} & \SetCell{bg=\HeatCellBg{21}}\CellText{0.21}{0}{0} & \SetCell{bg=\HeatCellBg{23}}\CellText{0.23}{0}{0} & \SetCell{bg=\HeatCellBg{19}}\CellText{0.19}{0}{0} & \SetCell{bg=\HeatCellBg{25}}\CellText{0.25}{0}{0} & \SetCell{bg=\HeatCellBg{24}}\CellText{0.24}{0}{0} & \SetCell{bg=\HeatCellBg{19}}\CellText{0.19}{0}{0} & \SetCell{bg=\HeatCellBg{19}}\CellText{0.19}{0}{0} & \MeanCell{0.24}{0.05}{0} \\
        Mean $\pm$ Std & \MeanCell{0.69}{0.17}{1} & \MeanCell{0.48}{0.20}{0} & \MeanCell{0.40}{0.10}{0} & \MeanCell{0.47}{0.11}{0} & \MeanCell{0.47}{0.14}{0} & \MeanCell{0.55}{0.15}{0} & \MeanCell{0.59}{0.21}{0} & \MeanCell{0.57}{0.21}{0} & \MeanCell{0.53}{0.19}{0} & \MeanCell{0.55}{0.18}{0} & \MeanCell{0.54}{0.17}{0} & \MeanCell{0.53}{0.16}{0} & \MeanCell{0.49}{0.16}{0} & \MeanCell{0.49}{0.14}{0} & \MeanCell{0.50}{0.16}{0} & \MeanCell{0.51}{0.15}{0} & \MeanCell{0.53}{0.17}{0} & \MeanCell{0.48}{0.16}{0} & \MeanCell{0.49}{0.16}{0} & \MeanCell{0.52}{0.18}{0} \\
}

\ClusterKTable{10}[Untrained-network baseline --- \MsPacManName.]{Untrained-network baseline ($\Similarity{untr}{}$) --- \MsPacManName.}{untr_ms_pacman}{
        1 & \SetCell{bg=\HeatCellBg{96}}\CellText{0.96}{1}{1} & \SetCell{bg=\HeatCellBg{86}}\CellText{0.86}{1}{0} & \SetCell{bg=\HeatCellBg{92}}\CellText{0.92}{1}{0} & \SetCell{bg=\HeatCellBg{83}}\CellText{0.83}{1}{0} & \SetCell{bg=\HeatCellBg{77}}\CellText{0.77}{1}{0} & \SetCell{bg=\HeatCellBg{90}}\CellText{0.90}{1}{0} & \SetCell{bg=\HeatCellBg{80}}\CellText{0.80}{0}{0} & \SetCell{bg=\HeatCellBg{78}}\CellText{0.78}{1}{0} & \SetCell{bg=\HeatCellBg{79}}\CellText{0.79}{1}{0} & \SetCell{bg=\HeatCellBg{78}}\CellText{0.78}{1}{0} & \SetCell{bg=\HeatCellBg{74}}\CellText{0.74}{0}{0} & \SetCell{bg=\HeatCellBg{78}}\CellText{0.78}{1}{0} & \SetCell{bg=\HeatCellBg{74}}\CellText{0.74}{1}{0} & \SetCell{bg=\HeatCellBg{71}}\CellText{0.71}{1}{0} & \SetCell{bg=\HeatCellBg{73}}\CellText{0.73}{1}{0} & \SetCell{bg=\HeatCellBg{71}}\CellText{0.71}{1}{0} & \SetCell{bg=\HeatCellBg{67}}\CellText{0.67}{1}{0} & \SetCell{bg=\HeatCellBg{66}}\CellText{0.66}{1}{0} & \SetCell{bg=\HeatCellBg{68}}\CellText{0.68}{1}{0} & \MeanCell{0.78}{0.08}{1} \\
        2 & \SetCell{bg=\HeatCellBg{72}}\CellText{0.72}{0}{0} & \SetCell{bg=\HeatCellBg{75}}\CellText{0.75}{0}{0} & \SetCell{bg=\HeatCellBg{87}}\CellText{0.87}{0}{1} & \SetCell{bg=\HeatCellBg{69}}\CellText{0.69}{0}{0} & \SetCell{bg=\HeatCellBg{66}}\CellText{0.66}{0}{0} & \SetCell{bg=\HeatCellBg{79}}\CellText{0.79}{0}{0} & \SetCell{bg=\HeatCellBg{81}}\CellText{0.81}{1}{0} & \SetCell{bg=\HeatCellBg{77}}\CellText{0.77}{0}{0} & \SetCell{bg=\HeatCellBg{76}}\CellText{0.76}{0}{0} & \SetCell{bg=\HeatCellBg{69}}\CellText{0.69}{0}{0} & \SetCell{bg=\HeatCellBg{76}}\CellText{0.76}{1}{0} & \SetCell{bg=\HeatCellBg{70}}\CellText{0.70}{0}{0} & \SetCell{bg=\HeatCellBg{71}}\CellText{0.71}{0}{0} & \SetCell{bg=\HeatCellBg{70}}\CellText{0.70}{0}{0} & \SetCell{bg=\HeatCellBg{68}}\CellText{0.68}{0}{0} & \SetCell{bg=\HeatCellBg{66}}\CellText{0.66}{0}{0} & \SetCell{bg=\HeatCellBg{63}}\CellText{0.63}{0}{0} & \SetCell{bg=\HeatCellBg{65}}\CellText{0.65}{0}{0} & \SetCell{bg=\HeatCellBg{65}}\CellText{0.65}{0}{0} & \MeanCell{0.72}{0.06}{0} \\
        3 & \SetCell{bg=\HeatCellBg{65}}\CellText{0.65}{0}{0} & \SetCell{bg=\HeatCellBg{65}}\CellText{0.65}{0}{0} & \SetCell{bg=\HeatCellBg{76}}\CellText{0.76}{0}{1} & \SetCell{bg=\HeatCellBg{67}}\CellText{0.67}{0}{0} & \SetCell{bg=\HeatCellBg{68}}\CellText{0.68}{0}{0} & \SetCell{bg=\HeatCellBg{69}}\CellText{0.69}{0}{0} & \SetCell{bg=\HeatCellBg{69}}\CellText{0.69}{0}{0} & \SetCell{bg=\HeatCellBg{69}}\CellText{0.69}{0}{0} & \SetCell{bg=\HeatCellBg{63}}\CellText{0.63}{0}{0} & \SetCell{bg=\HeatCellBg{61}}\CellText{0.61}{0}{0} & \SetCell{bg=\HeatCellBg{61}}\CellText{0.61}{0}{0} & \SetCell{bg=\HeatCellBg{60}}\CellText{0.60}{0}{0} & \SetCell{bg=\HeatCellBg{60}}\CellText{0.60}{0}{0} & \SetCell{bg=\HeatCellBg{59}}\CellText{0.59}{0}{0} & \SetCell{bg=\HeatCellBg{55}}\CellText{0.55}{0}{0} & \SetCell{bg=\HeatCellBg{58}}\CellText{0.58}{0}{0} & \SetCell{bg=\HeatCellBg{56}}\CellText{0.56}{0}{0} & \SetCell{bg=\HeatCellBg{60}}\CellText{0.60}{0}{0} & \SetCell{bg=\HeatCellBg{54}}\CellText{0.54}{0}{0} & \MeanCell{0.63}{0.06}{0} \\
        4 & \SetCell{bg=\HeatCellBg{64}}\CellText{0.64}{0}{1} & \SetCell{bg=\HeatCellBg{48}}\CellText{0.48}{0}{0} & \SetCell{bg=\HeatCellBg{49}}\CellText{0.49}{0}{0} & \SetCell{bg=\HeatCellBg{53}}\CellText{0.53}{0}{0} & \SetCell{bg=\HeatCellBg{52}}\CellText{0.52}{0}{0} & \SetCell{bg=\HeatCellBg{56}}\CellText{0.56}{0}{0} & \SetCell{bg=\HeatCellBg{64}}\CellText{0.64}{0}{1} & \SetCell{bg=\HeatCellBg{64}}\CellText{0.64}{0}{1} & \SetCell{bg=\HeatCellBg{59}}\CellText{0.59}{0}{0} & \SetCell{bg=\HeatCellBg{62}}\CellText{0.62}{0}{0} & \SetCell{bg=\HeatCellBg{60}}\CellText{0.60}{0}{0} & \SetCell{bg=\HeatCellBg{57}}\CellText{0.57}{0}{0} & \SetCell{bg=\HeatCellBg{56}}\CellText{0.56}{0}{0} & \SetCell{bg=\HeatCellBg{57}}\CellText{0.57}{0}{0} & \SetCell{bg=\HeatCellBg{55}}\CellText{0.55}{0}{0} & \SetCell{bg=\HeatCellBg{50}}\CellText{0.50}{0}{0} & \SetCell{bg=\HeatCellBg{53}}\CellText{0.53}{0}{0} & \SetCell{bg=\HeatCellBg{49}}\CellText{0.49}{0}{0} & \SetCell{bg=\HeatCellBg{49}}\CellText{0.49}{0}{0} & \MeanCell{0.56}{0.05}{0} \\
        5 & \SetCell{bg=\HeatCellBg{64}}\CellText{0.64}{0}{1} & \SetCell{bg=\HeatCellBg{48}}\CellText{0.48}{0}{0} & \SetCell{bg=\HeatCellBg{49}}\CellText{0.49}{0}{0} & \SetCell{bg=\HeatCellBg{53}}\CellText{0.53}{0}{0} & \SetCell{bg=\HeatCellBg{52}}\CellText{0.52}{0}{0} & \SetCell{bg=\HeatCellBg{56}}\CellText{0.56}{0}{0} & \SetCell{bg=\HeatCellBg{64}}\CellText{0.64}{0}{1} & \SetCell{bg=\HeatCellBg{64}}\CellText{0.64}{0}{1} & \SetCell{bg=\HeatCellBg{59}}\CellText{0.59}{0}{0} & \SetCell{bg=\HeatCellBg{62}}\CellText{0.62}{0}{0} & \SetCell{bg=\HeatCellBg{60}}\CellText{0.60}{0}{0} & \SetCell{bg=\HeatCellBg{57}}\CellText{0.57}{0}{0} & \SetCell{bg=\HeatCellBg{56}}\CellText{0.56}{0}{0} & \SetCell{bg=\HeatCellBg{57}}\CellText{0.57}{0}{0} & \SetCell{bg=\HeatCellBg{55}}\CellText{0.55}{0}{0} & \SetCell{bg=\HeatCellBg{50}}\CellText{0.50}{0}{0} & \SetCell{bg=\HeatCellBg{53}}\CellText{0.53}{0}{0} & \SetCell{bg=\HeatCellBg{51}}\CellText{0.51}{0}{0} & \SetCell{bg=\HeatCellBg{49}}\CellText{0.49}{0}{0} & \MeanCell{0.56}{0.05}{0} \\
        6b & \SetCell{bg=\HeatCellBg{51}}\CellText{0.51}{0}{1} & \SetCell{bg=\HeatCellBg{48}}\CellText{0.48}{0}{0} & \SetCell{bg=\HeatCellBg{44}}\CellText{0.44}{0}{0} & \SetCell{bg=\HeatCellBg{37}}\CellText{0.37}{0}{0} & \SetCell{bg=\HeatCellBg{41}}\CellText{0.41}{0}{0} & \SetCell{bg=\HeatCellBg{44}}\CellText{0.44}{0}{0} & \SetCell{bg=\HeatCellBg{42}}\CellText{0.42}{0}{0} & \SetCell{bg=\HeatCellBg{38}}\CellText{0.38}{0}{0} & \SetCell{bg=\HeatCellBg{37}}\CellText{0.37}{0}{0} & \SetCell{bg=\HeatCellBg{40}}\CellText{0.40}{0}{0} & \SetCell{bg=\HeatCellBg{42}}\CellText{0.42}{0}{0} & \SetCell{bg=\HeatCellBg{40}}\CellText{0.40}{0}{0} & \SetCell{bg=\HeatCellBg{37}}\CellText{0.37}{0}{0} & \SetCell{bg=\HeatCellBg{39}}\CellText{0.39}{0}{0} & \SetCell{bg=\HeatCellBg{40}}\CellText{0.40}{0}{0} & \SetCell{bg=\HeatCellBg{38}}\CellText{0.38}{0}{0} & \SetCell{bg=\HeatCellBg{39}}\CellText{0.39}{0}{0} & \SetCell{bg=\HeatCellBg{37}}\CellText{0.37}{0}{0} & \SetCell{bg=\HeatCellBg{39}}\CellText{0.39}{0}{0} & \MeanCell{0.41}{0.04}{0} \\
        7b & \SetCell{bg=\HeatCellBg{42}}\CellText{0.42}{0}{0} & \SetCell{bg=\HeatCellBg{46}}\CellText{0.46}{0}{1} & \SetCell{bg=\HeatCellBg{37}}\CellText{0.37}{0}{0} & \SetCell{bg=\HeatCellBg{39}}\CellText{0.39}{0}{0} & \SetCell{bg=\HeatCellBg{31}}\CellText{0.31}{0}{0} & \SetCell{bg=\HeatCellBg{36}}\CellText{0.36}{0}{0} & \SetCell{bg=\HeatCellBg{29}}\CellText{0.29}{0}{0} & \SetCell{bg=\HeatCellBg{29}}\CellText{0.29}{0}{0} & \SetCell{bg=\HeatCellBg{29}}\CellText{0.29}{0}{0} & \SetCell{bg=\HeatCellBg{30}}\CellText{0.30}{0}{0} & \SetCell{bg=\HeatCellBg{29}}\CellText{0.29}{0}{0} & \SetCell{bg=\HeatCellBg{29}}\CellText{0.29}{0}{0} & \SetCell{bg=\HeatCellBg{27}}\CellText{0.27}{0}{0} & \SetCell{bg=\HeatCellBg{26}}\CellText{0.26}{0}{0} & \SetCell{bg=\HeatCellBg{27}}\CellText{0.27}{0}{0} & \SetCell{bg=\HeatCellBg{27}}\CellText{0.27}{0}{0} & \SetCell{bg=\HeatCellBg{26}}\CellText{0.26}{0}{0} & \SetCell{bg=\HeatCellBg{26}}\CellText{0.26}{0}{0} & \SetCell{bg=\HeatCellBg{24}}\CellText{0.24}{0}{0} & \MeanCell{0.31}{0.06}{0} \\
        8b & \SetCell{bg=\HeatCellBg{39}}\CellText{0.39}{0}{1} & \SetCell{bg=\HeatCellBg{27}}\CellText{0.27}{0}{0} & \SetCell{bg=\HeatCellBg{29}}\CellText{0.29}{0}{0} & \SetCell{bg=\HeatCellBg{27}}\CellText{0.27}{0}{0} & \SetCell{bg=\HeatCellBg{31}}\CellText{0.31}{0}{0} & \SetCell{bg=\HeatCellBg{22}}\CellText{0.22}{0}{0} & \SetCell{bg=\HeatCellBg{19}}\CellText{0.19}{0}{0} & \SetCell{bg=\HeatCellBg{20}}\CellText{0.20}{0}{0} & \SetCell{bg=\HeatCellBg{19}}\CellText{0.19}{0}{0} & \SetCell{bg=\HeatCellBg{16}}\CellText{0.16}{0}{0} & \SetCell{bg=\HeatCellBg{16}}\CellText{0.16}{0}{0} & \SetCell{bg=\HeatCellBg{18}}\CellText{0.18}{0}{0} & \SetCell{bg=\HeatCellBg{15}}\CellText{0.15}{0}{0} & \SetCell{bg=\HeatCellBg{18}}\CellText{0.18}{0}{0} & \SetCell{bg=\HeatCellBg{17}}\CellText{0.17}{0}{0} & \SetCell{bg=\HeatCellBg{16}}\CellText{0.16}{0}{0} & \SetCell{bg=\HeatCellBg{17}}\CellText{0.17}{0}{0} & \SetCell{bg=\HeatCellBg{20}}\CellText{0.20}{0}{0} & \SetCell{bg=\HeatCellBg{18}}\CellText{0.18}{0}{0} & \MeanCell{0.21}{0.06}{0} \\
        Mean $\pm$ Std & \MeanCell{0.62}{0.17}{1} & \MeanCell{0.55}{0.18}{0} & \MeanCell{0.58}{0.22}{0} & \MeanCell{0.54}{0.18}{0} & \MeanCell{0.52}{0.16}{0} & \MeanCell{0.57}{0.21}{0} & \MeanCell{0.56}{0.22}{0} & \MeanCell{0.55}{0.21}{0} & \MeanCell{0.53}{0.21}{0} & \MeanCell{0.52}{0.20}{0} & \MeanCell{0.52}{0.20}{0} & \MeanCell{0.51}{0.19}{0} & \MeanCell{0.49}{0.20}{0} & \MeanCell{0.50}{0.19}{0} & \MeanCell{0.49}{0.18}{0} & \MeanCell{0.47}{0.18}{0} & \MeanCell{0.47}{0.17}{0} & \MeanCell{0.47}{0.16}{0} & \MeanCell{0.46}{0.17}{0} & \MeanCell{0.52}{0.19}{0} \\
}

\ClusterKTable{10}[Action baseline --- \MsPacManName.]{Action baseline ($\Similarity{act}{}$) --- \MsPacManName.}{act_ms_pacman}{
        1 & \SetCell{bg=\HeatCellBg{13}}\CellText{0.13}{0}{0} & \SetCell{bg=\HeatCellBg{7}}\CellText{0.07}{0}{0} & \SetCell{bg=\HeatCellBg{21}}\CellText{0.21}{0}{0} & \SetCell{bg=\HeatCellBg{25}}\CellText{0.25}{0}{1} & \SetCell{bg=\HeatCellBg{12}}\CellText{0.12}{0}{0} & \SetCell{bg=\HeatCellBg{12}}\CellText{0.12}{0}{0} & \SetCell{bg=\HeatCellBg{9}}\CellText{0.09}{0}{0} & \SetCell{bg=\HeatCellBg{9}}\CellText{0.09}{0}{0} & \SetCell{bg=\HeatCellBg{9}}\CellText{0.09}{0}{0} & \SetCell{bg=\HeatCellBg{10}}\CellText{0.10}{0}{0} & \SetCell{bg=\HeatCellBg{9}}\CellText{0.09}{0}{0} & \SetCell{bg=\HeatCellBg{8}}\CellText{0.08}{0}{0} & \SetCell{bg=\HeatCellBg{10}}\CellText{0.10}{0}{0} & \SetCell{bg=\HeatCellBg{7}}\CellText{0.07}{0}{0} & \SetCell{bg=\HeatCellBg{8}}\CellText{0.08}{0}{0} & \SetCell{bg=\HeatCellBg{11}}\CellText{0.11}{0}{0} & \SetCell{bg=\HeatCellBg{10}}\CellText{0.10}{0}{0} & \SetCell{bg=\HeatCellBg{10}}\CellText{0.10}{0}{0} & \SetCell{bg=\HeatCellBg{10}}\CellText{0.10}{0}{0} & \MeanCell{0.11}{0.04}{0} \\
        2 & \SetCell{bg=\HeatCellBg{10}}\CellText{0.10}{0}{0} & \SetCell{bg=\HeatCellBg{11}}\CellText{0.11}{0}{0} & \SetCell{bg=\HeatCellBg{22}}\CellText{0.22}{0}{0} & \SetCell{bg=\HeatCellBg{24}}\CellText{0.24}{0}{1} & \SetCell{bg=\HeatCellBg{16}}\CellText{0.16}{0}{0} & \SetCell{bg=\HeatCellBg{13}}\CellText{0.13}{0}{0} & \SetCell{bg=\HeatCellBg{8}}\CellText{0.08}{0}{0} & \SetCell{bg=\HeatCellBg{7}}\CellText{0.07}{0}{0} & \SetCell{bg=\HeatCellBg{9}}\CellText{0.09}{0}{0} & \SetCell{bg=\HeatCellBg{11}}\CellText{0.11}{0}{0} & \SetCell{bg=\HeatCellBg{8}}\CellText{0.08}{0}{0} & \SetCell{bg=\HeatCellBg{6}}\CellText{0.06}{0}{0} & \SetCell{bg=\HeatCellBg{9}}\CellText{0.09}{0}{0} & \SetCell{bg=\HeatCellBg{11}}\CellText{0.11}{0}{0} & \SetCell{bg=\HeatCellBg{9}}\CellText{0.09}{0}{0} & \SetCell{bg=\HeatCellBg{9}}\CellText{0.09}{0}{0} & \SetCell{bg=\HeatCellBg{10}}\CellText{0.10}{0}{0} & \SetCell{bg=\HeatCellBg{11}}\CellText{0.11}{0}{0} & \SetCell{bg=\HeatCellBg{9}}\CellText{0.09}{0}{0} & \MeanCell{0.11}{0.05}{0} \\
        3 & \SetCell{bg=\HeatCellBg{5}}\CellText{0.05}{0}{0} & \SetCell{bg=\HeatCellBg{12}}\CellText{0.12}{0}{0} & \SetCell{bg=\HeatCellBg{20}}\CellText{0.20}{0}{0} & \SetCell{bg=\HeatCellBg{21}}\CellText{0.21}{0}{1} & \SetCell{bg=\HeatCellBg{18}}\CellText{0.18}{0}{0} & \SetCell{bg=\HeatCellBg{15}}\CellText{0.15}{0}{0} & \SetCell{bg=\HeatCellBg{18}}\CellText{0.18}{0}{0} & \SetCell{bg=\HeatCellBg{12}}\CellText{0.12}{0}{0} & \SetCell{bg=\HeatCellBg{12}}\CellText{0.12}{0}{0} & \SetCell{bg=\HeatCellBg{11}}\CellText{0.11}{0}{0} & \SetCell{bg=\HeatCellBg{11}}\CellText{0.11}{0}{0} & \SetCell{bg=\HeatCellBg{10}}\CellText{0.10}{0}{0} & \SetCell{bg=\HeatCellBg{9}}\CellText{0.09}{0}{0} & \SetCell{bg=\HeatCellBg{10}}\CellText{0.10}{0}{0} & \SetCell{bg=\HeatCellBg{10}}\CellText{0.10}{0}{0} & \SetCell{bg=\HeatCellBg{10}}\CellText{0.10}{0}{0} & \SetCell{bg=\HeatCellBg{10}}\CellText{0.10}{0}{0} & \SetCell{bg=\HeatCellBg{10}}\CellText{0.10}{0}{0} & \SetCell{bg=\HeatCellBg{12}}\CellText{0.12}{0}{0} & \MeanCell{0.12}{0.04}{0} \\
        4 & \SetCell{bg=\HeatCellBg{0}}\CellText{-0.00}{0}{0} & \SetCell{bg=\HeatCellBg{-1}}\CellText{-0.01}{0}{0} & \SetCell{bg=\HeatCellBg{12}}\CellText{0.12}{0}{0} & \SetCell{bg=\HeatCellBg{16}}\CellText{0.16}{0}{1} & \SetCell{bg=\HeatCellBg{10}}\CellText{0.10}{0}{0} & \SetCell{bg=\HeatCellBg{14}}\CellText{0.14}{0}{0} & \SetCell{bg=\HeatCellBg{10}}\CellText{0.10}{0}{0} & \SetCell{bg=\HeatCellBg{11}}\CellText{0.11}{0}{0} & \SetCell{bg=\HeatCellBg{6}}\CellText{0.06}{0}{0} & \SetCell{bg=\HeatCellBg{6}}\CellText{0.06}{0}{0} & \SetCell{bg=\HeatCellBg{6}}\CellText{0.06}{0}{0} & \SetCell{bg=\HeatCellBg{6}}\CellText{0.06}{0}{0} & \SetCell{bg=\HeatCellBg{4}}\CellText{0.04}{0}{0} & \SetCell{bg=\HeatCellBg{6}}\CellText{0.06}{0}{0} & \SetCell{bg=\HeatCellBg{6}}\CellText{0.06}{0}{0} & \SetCell{bg=\HeatCellBg{8}}\CellText{0.08}{0}{0} & \SetCell{bg=\HeatCellBg{8}}\CellText{0.08}{0}{0} & \SetCell{bg=\HeatCellBg{9}}\CellText{0.09}{0}{0} & \SetCell{bg=\HeatCellBg{9}}\CellText{0.09}{0}{0} & \MeanCell{0.08}{0.04}{0} \\
        5 & \SetCell{bg=\HeatCellBg{0}}\CellText{-0.00}{0}{0} & \SetCell{bg=\HeatCellBg{-1}}\CellText{-0.01}{0}{0} & \SetCell{bg=\HeatCellBg{12}}\CellText{0.12}{0}{0} & \SetCell{bg=\HeatCellBg{16}}\CellText{0.16}{0}{1} & \SetCell{bg=\HeatCellBg{10}}\CellText{0.10}{0}{0} & \SetCell{bg=\HeatCellBg{14}}\CellText{0.14}{0}{0} & \SetCell{bg=\HeatCellBg{10}}\CellText{0.10}{0}{0} & \SetCell{bg=\HeatCellBg{11}}\CellText{0.11}{0}{0} & \SetCell{bg=\HeatCellBg{6}}\CellText{0.06}{0}{0} & \SetCell{bg=\HeatCellBg{6}}\CellText{0.06}{0}{0} & \SetCell{bg=\HeatCellBg{6}}\CellText{0.06}{0}{0} & \SetCell{bg=\HeatCellBg{6}}\CellText{0.06}{0}{0} & \SetCell{bg=\HeatCellBg{4}}\CellText{0.04}{0}{0} & \SetCell{bg=\HeatCellBg{6}}\CellText{0.06}{0}{0} & \SetCell{bg=\HeatCellBg{6}}\CellText{0.06}{0}{0} & \SetCell{bg=\HeatCellBg{8}}\CellText{0.08}{0}{0} & \SetCell{bg=\HeatCellBg{8}}\CellText{0.08}{0}{0} & \SetCell{bg=\HeatCellBg{9}}\CellText{0.09}{0}{0} & \SetCell{bg=\HeatCellBg{9}}\CellText{0.09}{0}{0} & \MeanCell{0.08}{0.04}{0} \\
        6b & \SetCell{bg=\HeatCellBg{-2}}\CellText{-0.02}{0}{0} & \SetCell{bg=\HeatCellBg{5}}\CellText{0.05}{0}{0} & \SetCell{bg=\HeatCellBg{19}}\CellText{0.19}{0}{1} & \SetCell{bg=\HeatCellBg{16}}\CellText{0.16}{0}{0} & \SetCell{bg=\HeatCellBg{19}}\CellText{0.19}{0}{1} & \SetCell{bg=\HeatCellBg{17}}\CellText{0.17}{0}{0} & \SetCell{bg=\HeatCellBg{15}}\CellText{0.15}{0}{0} & \SetCell{bg=\HeatCellBg{16}}\CellText{0.16}{0}{0} & \SetCell{bg=\HeatCellBg{10}}\CellText{0.10}{0}{0} & \SetCell{bg=\HeatCellBg{14}}\CellText{0.14}{0}{0} & \SetCell{bg=\HeatCellBg{11}}\CellText{0.11}{0}{0} & \SetCell{bg=\HeatCellBg{10}}\CellText{0.10}{0}{0} & \SetCell{bg=\HeatCellBg{11}}\CellText{0.11}{0}{0} & \SetCell{bg=\HeatCellBg{12}}\CellText{0.12}{0}{0} & \SetCell{bg=\HeatCellBg{11}}\CellText{0.11}{0}{0} & \SetCell{bg=\HeatCellBg{12}}\CellText{0.12}{0}{0} & \SetCell{bg=\HeatCellBg{12}}\CellText{0.12}{0}{0} & \SetCell{bg=\HeatCellBg{12}}\CellText{0.12}{0}{0} & \SetCell{bg=\HeatCellBg{13}}\CellText{0.13}{0}{0} & \MeanCell{0.12}{0.05}{0} \\
        7b & \SetCell{bg=\HeatCellBg{13}}\CellText{0.13}{0}{0} & \SetCell{bg=\HeatCellBg{8}}\CellText{0.08}{0}{0} & \SetCell{bg=\HeatCellBg{31}}\CellText{0.31}{0}{0} & \SetCell{bg=\HeatCellBg{32}}\CellText{0.32}{0}{1} & \SetCell{bg=\HeatCellBg{28}}\CellText{0.28}{0}{0} & \SetCell{bg=\HeatCellBg{25}}\CellText{0.25}{0}{0} & \SetCell{bg=\HeatCellBg{22}}\CellText{0.22}{0}{0} & \SetCell{bg=\HeatCellBg{23}}\CellText{0.23}{0}{0} & \SetCell{bg=\HeatCellBg{19}}\CellText{0.19}{0}{0} & \SetCell{bg=\HeatCellBg{19}}\CellText{0.19}{0}{0} & \SetCell{bg=\HeatCellBg{19}}\CellText{0.19}{0}{0} & \SetCell{bg=\HeatCellBg{16}}\CellText{0.16}{0}{0} & \SetCell{bg=\HeatCellBg{14}}\CellText{0.14}{0}{0} & \SetCell{bg=\HeatCellBg{16}}\CellText{0.16}{0}{0} & \SetCell{bg=\HeatCellBg{16}}\CellText{0.16}{0}{0} & \SetCell{bg=\HeatCellBg{19}}\CellText{0.19}{0}{0} & \SetCell{bg=\HeatCellBg{16}}\CellText{0.16}{0}{0} & \SetCell{bg=\HeatCellBg{18}}\CellText{0.18}{0}{0} & \SetCell{bg=\HeatCellBg{18}}\CellText{0.18}{0}{0} & \MeanCell{0.20}{0.06}{0} \\
        8b & \SetCell{bg=\HeatCellBg{26}}\CellText{0.26}{1}{0} & \SetCell{bg=\HeatCellBg{41}}\CellText{0.41}{1}{0} & \SetCell{bg=\HeatCellBg{53}}\CellText{0.53}{1}{1} & \SetCell{bg=\HeatCellBg{43}}\CellText{0.43}{1}{0} & \SetCell{bg=\HeatCellBg{44}}\CellText{0.44}{1}{0} & \SetCell{bg=\HeatCellBg{46}}\CellText{0.46}{1}{0} & \SetCell{bg=\HeatCellBg{45}}\CellText{0.45}{1}{0} & \SetCell{bg=\HeatCellBg{44}}\CellText{0.44}{1}{0} & \SetCell{bg=\HeatCellBg{40}}\CellText{0.40}{1}{0} & \SetCell{bg=\HeatCellBg{42}}\CellText{0.42}{1}{0} & \SetCell{bg=\HeatCellBg{38}}\CellText{0.38}{1}{0} & \SetCell{bg=\HeatCellBg{41}}\CellText{0.41}{1}{0} & \SetCell{bg=\HeatCellBg{41}}\CellText{0.41}{1}{0} & \SetCell{bg=\HeatCellBg{35}}\CellText{0.35}{1}{0} & \SetCell{bg=\HeatCellBg{35}}\CellText{0.35}{1}{0} & \SetCell{bg=\HeatCellBg{32}}\CellText{0.32}{1}{0} & \SetCell{bg=\HeatCellBg{34}}\CellText{0.34}{1}{0} & \SetCell{bg=\HeatCellBg{34}}\CellText{0.34}{1}{0} & \SetCell{bg=\HeatCellBg{33}}\CellText{0.33}{1}{0} & \MeanCell{0.39}{0.06}{1} \\
        Mean $\pm$ Std & \MeanCell{0.08}{0.09}{0} & \MeanCell{0.10}{0.12}{0} & \MeanCell{0.24}{0.12}{0} & \MeanCell{0.24}{0.09}{1} & \MeanCell{0.20}{0.11}{0} & \MeanCell{0.20}{0.11}{0} & \MeanCell{0.17}{0.11}{0} & \MeanCell{0.17}{0.11}{0} & \MeanCell{0.14}{0.11}{0} & \MeanCell{0.15}{0.11}{0} & \MeanCell{0.14}{0.10}{0} & \MeanCell{0.13}{0.11}{0} & \MeanCell{0.13}{0.11}{0} & \MeanCell{0.13}{0.09}{0} & \MeanCell{0.13}{0.09}{0} & \MeanCell{0.14}{0.08}{0} & \MeanCell{0.14}{0.08}{0} & \MeanCell{0.14}{0.08}{0} & \MeanCell{0.14}{0.08}{0} & \MeanCell{0.15}{0.11}{0} \\
}


\clearpage
\section{\PongName}
\label{sec:appendix_additional_results_pong}
\SetHeatColor{mygreen}
\ClusterKTable{8}{\glsfmtshort{cu} --- \PongName.}{cu_pong}{
        1 & \SetCell{bg=\HeatCellBg{33}}\CellText{0.33}{0}{0} & \SetCell{bg=\HeatCellBg{1}}\CellText{0.01}{0}{0} & \SetCell{bg=\HeatCellBg{91}}\CellText{0.91}{1}{1} & \SetCell{bg=\HeatCellBg{80}}\CellText{0.80}{1}{0} & \SetCell{bg=\HeatCellBg{64}}\CellText{0.64}{0}{0} & \SetCell{bg=\HeatCellBg{69}}\CellText{0.69}{1}{0} & \SetCell{bg=\HeatCellBg{63}}\CellText{0.63}{0}{0} & \SetCell{bg=\HeatCellBg{56}}\CellText{0.56}{0}{0} & \SetCell{bg=\HeatCellBg{66}}\CellText{0.66}{1}{0} & \SetCell{bg=\HeatCellBg{68}}\CellText{0.68}{1}{0} & \SetCell{bg=\HeatCellBg{57}}\CellText{0.57}{0}{0} & \SetCell{bg=\HeatCellBg{76}}\CellText{0.76}{1}{0} & \SetCell{bg=\HeatCellBg{70}}\CellText{0.70}{1}{0} & \SetCell{bg=\HeatCellBg{70}}\CellText{0.70}{1}{0} & \SetCell{bg=\HeatCellBg{60}}\CellText{0.60}{1}{0} & \SetCell{bg=\HeatCellBg{76}}\CellText{0.76}{1}{0} & \SetCell{bg=\HeatCellBg{61}}\CellText{0.61}{0}{0} & \SetCell{bg=\HeatCellBg{56}}\CellText{0.56}{1}{0} & \SetCell{bg=\HeatCellBg{64}}\CellText{0.64}{1}{0} & \MeanCell{0.62}{0.18}{0} \\
        2 & \SetCell{bg=\HeatCellBg{100}}\CellText{1.00}{1}{1} & \SetCell{bg=\HeatCellBg{91}}\CellText{0.91}{1}{0} & \SetCell{bg=\HeatCellBg{74}}\CellText{0.74}{0}{0} & \SetCell{bg=\HeatCellBg{63}}\CellText{0.63}{0}{0} & \SetCell{bg=\HeatCellBg{67}}\CellText{0.67}{0}{0} & \SetCell{bg=\HeatCellBg{59}}\CellText{0.59}{0}{0} & \SetCell{bg=\HeatCellBg{65}}\CellText{0.65}{1}{0} & \SetCell{bg=\HeatCellBg{51}}\CellText{0.51}{0}{0} & \SetCell{bg=\HeatCellBg{66}}\CellText{0.66}{1}{0} & \SetCell{bg=\HeatCellBg{65}}\CellText{0.65}{0}{0} & \SetCell{bg=\HeatCellBg{59}}\CellText{0.59}{1}{0} & \SetCell{bg=\HeatCellBg{56}}\CellText{0.56}{0}{0} & \SetCell{bg=\HeatCellBg{66}}\CellText{0.66}{0}{0} & \SetCell{bg=\HeatCellBg{59}}\CellText{0.59}{0}{0} & \SetCell{bg=\HeatCellBg{60}}\CellText{0.60}{1}{0} & \SetCell{bg=\HeatCellBg{61}}\CellText{0.61}{0}{0} & \SetCell{bg=\HeatCellBg{63}}\CellText{0.63}{1}{0} & \SetCell{bg=\HeatCellBg{51}}\CellText{0.51}{0}{0} & \SetCell{bg=\HeatCellBg{55}}\CellText{0.55}{0}{0} & \MeanCell{0.65}{0.12}{1} \\
        3 & \SetCell{bg=\HeatCellBg{31}}\CellText{0.31}{0}{0} & \SetCell{bg=\HeatCellBg{67}}\CellText{0.67}{0}{0} & \SetCell{bg=\HeatCellBg{66}}\CellText{0.66}{0}{0} & \SetCell{bg=\HeatCellBg{55}}\CellText{0.55}{0}{0} & \SetCell{bg=\HeatCellBg{73}}\CellText{0.73}{1}{1} & \SetCell{bg=\HeatCellBg{64}}\CellText{0.64}{0}{0} & \SetCell{bg=\HeatCellBg{63}}\CellText{0.63}{0}{0} & \SetCell{bg=\HeatCellBg{62}}\CellText{0.62}{1}{0} & \SetCell{bg=\HeatCellBg{46}}\CellText{0.46}{0}{0} & \SetCell{bg=\HeatCellBg{47}}\CellText{0.47}{0}{0} & \SetCell{bg=\HeatCellBg{39}}\CellText{0.39}{0}{0} & \SetCell{bg=\HeatCellBg{38}}\CellText{0.38}{0}{0} & \SetCell{bg=\HeatCellBg{39}}\CellText{0.39}{0}{0} & \SetCell{bg=\HeatCellBg{40}}\CellText{0.40}{0}{0} & \SetCell{bg=\HeatCellBg{37}}\CellText{0.37}{0}{0} & \SetCell{bg=\HeatCellBg{35}}\CellText{0.35}{0}{0} & \SetCell{bg=\HeatCellBg{37}}\CellText{0.37}{0}{0} & \SetCell{bg=\HeatCellBg{39}}\CellText{0.39}{0}{0} & \SetCell{bg=\HeatCellBg{38}}\CellText{0.38}{0}{0} & \MeanCell{0.48}{0.13}{0} \\
        4 & \SetCell{bg=\HeatCellBg{60}}\CellText{0.60}{0}{1} & \SetCell{bg=\HeatCellBg{52}}\CellText{0.52}{0}{0} & \SetCell{bg=\HeatCellBg{58}}\CellText{0.58}{0}{0} & \SetCell{bg=\HeatCellBg{49}}\CellText{0.49}{0}{0} & \SetCell{bg=\HeatCellBg{44}}\CellText{0.44}{0}{0} & \SetCell{bg=\HeatCellBg{51}}\CellText{0.51}{0}{0} & \SetCell{bg=\HeatCellBg{58}}\CellText{0.58}{0}{0} & \SetCell{bg=\HeatCellBg{48}}\CellText{0.48}{0}{0} & \SetCell{bg=\HeatCellBg{50}}\CellText{0.50}{0}{0} & \SetCell{bg=\HeatCellBg{52}}\CellText{0.52}{0}{0} & \SetCell{bg=\HeatCellBg{49}}\CellText{0.49}{0}{0} & \SetCell{bg=\HeatCellBg{46}}\CellText{0.46}{0}{0} & \SetCell{bg=\HeatCellBg{50}}\CellText{0.50}{0}{0} & \SetCell{bg=\HeatCellBg{45}}\CellText{0.45}{0}{0} & \SetCell{bg=\HeatCellBg{46}}\CellText{0.46}{0}{0} & \SetCell{bg=\HeatCellBg{47}}\CellText{0.47}{0}{0} & \SetCell{bg=\HeatCellBg{48}}\CellText{0.48}{0}{0} & \SetCell{bg=\HeatCellBg{47}}\CellText{0.47}{0}{0} & \SetCell{bg=\HeatCellBg{53}}\CellText{0.53}{0}{0} & \MeanCell{0.50}{0.04}{0} \\
        5 & \SetCell{bg=\HeatCellBg{60}}\CellText{0.60}{0}{1} & \SetCell{bg=\HeatCellBg{52}}\CellText{0.52}{0}{0} & \SetCell{bg=\HeatCellBg{58}}\CellText{0.58}{0}{0} & \SetCell{bg=\HeatCellBg{49}}\CellText{0.49}{0}{0} & \SetCell{bg=\HeatCellBg{44}}\CellText{0.44}{0}{0} & \SetCell{bg=\HeatCellBg{51}}\CellText{0.51}{0}{0} & \SetCell{bg=\HeatCellBg{58}}\CellText{0.58}{0}{0} & \SetCell{bg=\HeatCellBg{48}}\CellText{0.48}{0}{0} & \SetCell{bg=\HeatCellBg{50}}\CellText{0.50}{0}{0} & \SetCell{bg=\HeatCellBg{52}}\CellText{0.52}{0}{0} & \SetCell{bg=\HeatCellBg{49}}\CellText{0.49}{0}{0} & \SetCell{bg=\HeatCellBg{46}}\CellText{0.46}{0}{0} & \SetCell{bg=\HeatCellBg{50}}\CellText{0.50}{0}{0} & \SetCell{bg=\HeatCellBg{45}}\CellText{0.45}{0}{0} & \SetCell{bg=\HeatCellBg{46}}\CellText{0.46}{0}{0} & \SetCell{bg=\HeatCellBg{47}}\CellText{0.47}{0}{0} & \SetCell{bg=\HeatCellBg{48}}\CellText{0.48}{0}{0} & \SetCell{bg=\HeatCellBg{47}}\CellText{0.47}{0}{0} & \SetCell{bg=\HeatCellBg{53}}\CellText{0.53}{0}{0} & \MeanCell{0.50}{0.04}{0} \\
        6a & \SetCell{bg=\HeatCellBg{37}}\CellText{0.37}{0}{1} & \SetCell{bg=\HeatCellBg{32}}\CellText{0.32}{0}{0} & \SetCell{bg=\HeatCellBg{29}}\CellText{0.29}{0}{0} & \SetCell{bg=\HeatCellBg{30}}\CellText{0.30}{0}{0} & \SetCell{bg=\HeatCellBg{28}}\CellText{0.28}{0}{0} & \SetCell{bg=\HeatCellBg{28}}\CellText{0.28}{0}{0} & \SetCell{bg=\HeatCellBg{29}}\CellText{0.29}{0}{0} & \SetCell{bg=\HeatCellBg{31}}\CellText{0.31}{0}{0} & \SetCell{bg=\HeatCellBg{32}}\CellText{0.32}{0}{0} & \SetCell{bg=\HeatCellBg{33}}\CellText{0.33}{0}{0} & \SetCell{bg=\HeatCellBg{32}}\CellText{0.32}{0}{0} & \SetCell{bg=\HeatCellBg{31}}\CellText{0.31}{0}{0} & \SetCell{bg=\HeatCellBg{33}}\CellText{0.33}{0}{0} & \SetCell{bg=\HeatCellBg{32}}\CellText{0.32}{0}{0} & \SetCell{bg=\HeatCellBg{32}}\CellText{0.32}{0}{0} & \SetCell{bg=\HeatCellBg{32}}\CellText{0.32}{0}{0} & \SetCell{bg=\HeatCellBg{32}}\CellText{0.32}{0}{0} & \SetCell{bg=\HeatCellBg{31}}\CellText{0.31}{0}{0} & \SetCell{bg=\HeatCellBg{30}}\CellText{0.30}{0}{0} & \MeanCell{0.31}{0.02}{0} \\
        Mean $\pm$ Std & \MeanCell{0.54}{0.24}{0} & \MeanCell{0.49}{0.28}{0} & \MeanCell{0.63}{0.19}{1} & \MeanCell{0.54}{0.15}{0} & \MeanCell{0.53}{0.16}{0} & \MeanCell{0.54}{0.13}{0} & \MeanCell{0.56}{0.12}{0} & \MeanCell{0.49}{0.10}{0} & \MeanCell{0.52}{0.12}{0} & \MeanCell{0.53}{0.12}{0} & \MeanCell{0.47}{0.09}{0} & \MeanCell{0.49}{0.14}{0} & \MeanCell{0.51}{0.13}{0} & \MeanCell{0.48}{0.13}{0} & \MeanCell{0.47}{0.11}{0} & \MeanCell{0.50}{0.15}{0} & \MeanCell{0.48}{0.11}{0} & \MeanCell{0.45}{0.08}{0} & \MeanCell{0.49}{0.11}{0} & \MeanCell{0.51}{0.15}{0} \\
}

\SetHeatColor{myblue}
\ClusterKTable{8}{\glsfmtshort{as} --- \PongName.}{as_pong}{
        1 & \SetCell{bg=\HeatCellBg{33}}\CellText{0.33}{0}{0} & \SetCell{bg=\HeatCellBg{58}}\CellText{0.58}{0}{1} & \SetCell{bg=\HeatCellBg{26}}\CellText{0.26}{0}{0} & \SetCell{bg=\HeatCellBg{17}}\CellText{0.17}{0}{0} & \SetCell{bg=\HeatCellBg{17}}\CellText{0.17}{0}{0} & \SetCell{bg=\HeatCellBg{28}}\CellText{0.28}{0}{0} & \SetCell{bg=\HeatCellBg{31}}\CellText{0.31}{0}{0} & \SetCell{bg=\HeatCellBg{37}}\CellText{0.37}{0}{0} & \SetCell{bg=\HeatCellBg{38}}\CellText{0.38}{0}{0} & \SetCell{bg=\HeatCellBg{36}}\CellText{0.36}{0}{0} & \SetCell{bg=\HeatCellBg{36}}\CellText{0.36}{0}{0} & \SetCell{bg=\HeatCellBg{24}}\CellText{0.24}{0}{0} & \SetCell{bg=\HeatCellBg{19}}\CellText{0.19}{0}{0} & \SetCell{bg=\HeatCellBg{25}}\CellText{0.25}{0}{0} & \SetCell{bg=\HeatCellBg{29}}\CellText{0.29}{0}{0} & \SetCell{bg=\HeatCellBg{35}}\CellText{0.35}{0}{0} & \SetCell{bg=\HeatCellBg{40}}\CellText{0.40}{0}{0} & \SetCell{bg=\HeatCellBg{46}}\CellText{0.46}{0}{0} & \SetCell{bg=\HeatCellBg{43}}\CellText{0.43}{0}{0} & \MeanCell{0.33}{0.10}{0} \\
        2 & \SetCell{bg=\HeatCellBg{100}}\CellText{1.00}{1}{1} & \SetCell{bg=\HeatCellBg{64}}\CellText{0.64}{0}{0} & \SetCell{bg=\HeatCellBg{44}}\CellText{0.44}{0}{0} & \SetCell{bg=\HeatCellBg{27}}\CellText{0.27}{0}{0} & \SetCell{bg=\HeatCellBg{38}}\CellText{0.38}{0}{0} & \SetCell{bg=\HeatCellBg{49}}\CellText{0.49}{0}{0} & \SetCell{bg=\HeatCellBg{51}}\CellText{0.51}{0}{0} & \SetCell{bg=\HeatCellBg{53}}\CellText{0.53}{0}{0} & \SetCell{bg=\HeatCellBg{61}}\CellText{0.61}{0}{0} & \SetCell{bg=\HeatCellBg{63}}\CellText{0.63}{0}{0} & \SetCell{bg=\HeatCellBg{62}}\CellText{0.62}{0}{0} & \SetCell{bg=\HeatCellBg{64}}\CellText{0.64}{0}{0} & \SetCell{bg=\HeatCellBg{63}}\CellText{0.63}{0}{0} & \SetCell{bg=\HeatCellBg{64}}\CellText{0.64}{0}{0} & \SetCell{bg=\HeatCellBg{61}}\CellText{0.61}{0}{0} & \SetCell{bg=\HeatCellBg{67}}\CellText{0.67}{0}{0} & \SetCell{bg=\HeatCellBg{63}}\CellText{0.63}{0}{0} & \SetCell{bg=\HeatCellBg{68}}\CellText{0.68}{0}{0} & \SetCell{bg=\HeatCellBg{66}}\CellText{0.66}{0}{0} & \MeanCell{0.59}{0.14}{0} \\
        3 & \SetCell{bg=\HeatCellBg{99}}\CellText{0.99}{0}{1} & \SetCell{bg=\HeatCellBg{98}}\CellText{0.98}{0}{0} & \SetCell{bg=\HeatCellBg{99}}\CellText{0.99}{1}{1} & \SetCell{bg=\HeatCellBg{94}}\CellText{0.94}{0}{0} & \SetCell{bg=\HeatCellBg{98}}\CellText{0.98}{1}{0} & \SetCell{bg=\HeatCellBg{96}}\CellText{0.96}{1}{0} & \SetCell{bg=\HeatCellBg{93}}\CellText{0.93}{0}{0} & \SetCell{bg=\HeatCellBg{93}}\CellText{0.93}{0}{0} & \SetCell{bg=\HeatCellBg{92}}\CellText{0.92}{0}{0} & \SetCell{bg=\HeatCellBg{89}}\CellText{0.89}{0}{0} & \SetCell{bg=\HeatCellBg{89}}\CellText{0.89}{0}{0} & \SetCell{bg=\HeatCellBg{87}}\CellText{0.87}{0}{0} & \SetCell{bg=\HeatCellBg{89}}\CellText{0.89}{0}{0} & \SetCell{bg=\HeatCellBg{87}}\CellText{0.87}{0}{0} & \SetCell{bg=\HeatCellBg{85}}\CellText{0.85}{0}{0} & \SetCell{bg=\HeatCellBg{84}}\CellText{0.84}{0}{0} & \SetCell{bg=\HeatCellBg{87}}\CellText{0.87}{0}{0} & \SetCell{bg=\HeatCellBg{87}}\CellText{0.87}{0}{0} & \SetCell{bg=\HeatCellBg{83}}\CellText{0.83}{0}{0} & \MeanCell{0.91}{0.05}{0} \\
        4 & \SetCell{bg=\HeatCellBg{99}}\CellText{0.99}{0}{1} & \SetCell{bg=\HeatCellBg{99}}\CellText{0.99}{1}{1} & \SetCell{bg=\HeatCellBg{98}}\CellText{0.98}{0}{0} & \SetCell{bg=\HeatCellBg{99}}\CellText{0.99}{1}{1} & \SetCell{bg=\HeatCellBg{97}}\CellText{0.97}{0}{0} & \SetCell{bg=\HeatCellBg{96}}\CellText{0.96}{1}{0} & \SetCell{bg=\HeatCellBg{97}}\CellText{0.97}{1}{0} & \SetCell{bg=\HeatCellBg{96}}\CellText{0.96}{1}{0} & \SetCell{bg=\HeatCellBg{93}}\CellText{0.93}{0}{0} & \SetCell{bg=\HeatCellBg{94}}\CellText{0.94}{1}{0} & \SetCell{bg=\HeatCellBg{96}}\CellText{0.96}{1}{0} & \SetCell{bg=\HeatCellBg{93}}\CellText{0.93}{0}{0} & \SetCell{bg=\HeatCellBg{93}}\CellText{0.93}{0}{0} & \SetCell{bg=\HeatCellBg{93}}\CellText{0.93}{0}{0} & \SetCell{bg=\HeatCellBg{93}}\CellText{0.93}{0}{0} & \SetCell{bg=\HeatCellBg{91}}\CellText{0.91}{0}{0} & \SetCell{bg=\HeatCellBg{90}}\CellText{0.90}{0}{0} & \SetCell{bg=\HeatCellBg{90}}\CellText{0.90}{0}{0} & \SetCell{bg=\HeatCellBg{91}}\CellText{0.91}{0}{0} & \MeanCell{0.95}{0.03}{1} \\
        5 & \SetCell{bg=\HeatCellBg{99}}\CellText{0.99}{0}{1} & \SetCell{bg=\HeatCellBg{99}}\CellText{0.99}{1}{1} & \SetCell{bg=\HeatCellBg{98}}\CellText{0.98}{0}{0} & \SetCell{bg=\HeatCellBg{99}}\CellText{0.99}{1}{1} & \SetCell{bg=\HeatCellBg{97}}\CellText{0.97}{0}{0} & \SetCell{bg=\HeatCellBg{96}}\CellText{0.96}{1}{0} & \SetCell{bg=\HeatCellBg{97}}\CellText{0.97}{1}{0} & \SetCell{bg=\HeatCellBg{96}}\CellText{0.96}{1}{0} & \SetCell{bg=\HeatCellBg{93}}\CellText{0.93}{0}{0} & \SetCell{bg=\HeatCellBg{94}}\CellText{0.94}{1}{0} & \SetCell{bg=\HeatCellBg{96}}\CellText{0.96}{1}{0} & \SetCell{bg=\HeatCellBg{92}}\CellText{0.92}{0}{0} & \SetCell{bg=\HeatCellBg{93}}\CellText{0.93}{0}{0} & \SetCell{bg=\HeatCellBg{93}}\CellText{0.93}{0}{0} & \SetCell{bg=\HeatCellBg{93}}\CellText{0.93}{0}{0} & \SetCell{bg=\HeatCellBg{91}}\CellText{0.91}{0}{0} & \SetCell{bg=\HeatCellBg{90}}\CellText{0.90}{0}{0} & \SetCell{bg=\HeatCellBg{90}}\CellText{0.90}{0}{0} & \SetCell{bg=\HeatCellBg{91}}\CellText{0.91}{0}{0} & \MeanCell{0.95}{0.03}{0} \\
        6a & \SetCell{bg=\HeatCellBg{97}}\CellText{0.97}{0}{1} & \SetCell{bg=\HeatCellBg{94}}\CellText{0.94}{0}{0} & \SetCell{bg=\HeatCellBg{95}}\CellText{0.95}{0}{0} & \SetCell{bg=\HeatCellBg{95}}\CellText{0.95}{0}{0} & \SetCell{bg=\HeatCellBg{93}}\CellText{0.93}{0}{0} & \SetCell{bg=\HeatCellBg{95}}\CellText{0.95}{0}{0} & \SetCell{bg=\HeatCellBg{95}}\CellText{0.95}{0}{0} & \SetCell{bg=\HeatCellBg{95}}\CellText{0.95}{0}{0} & \SetCell{bg=\HeatCellBg{94}}\CellText{0.94}{1}{0} & \SetCell{bg=\HeatCellBg{94}}\CellText{0.94}{1}{0} & \SetCell{bg=\HeatCellBg{94}}\CellText{0.94}{0}{0} & \SetCell{bg=\HeatCellBg{94}}\CellText{0.94}{1}{0} & \SetCell{bg=\HeatCellBg{95}}\CellText{0.95}{1}{0} & \SetCell{bg=\HeatCellBg{94}}\CellText{0.94}{1}{0} & \SetCell{bg=\HeatCellBg{95}}\CellText{0.95}{1}{0} & \SetCell{bg=\HeatCellBg{94}}\CellText{0.94}{1}{0} & \SetCell{bg=\HeatCellBg{94}}\CellText{0.94}{1}{0} & \SetCell{bg=\HeatCellBg{94}}\CellText{0.94}{1}{0} & \SetCell{bg=\HeatCellBg{93}}\CellText{0.93}{1}{0} & \MeanCell{0.94}{0.01}{0} \\
        Mean $\pm$ Std & \MeanCell{0.88}{0.25}{1} & \MeanCell{0.85}{0.17}{0} & \MeanCell{0.77}{0.30}{0} & \MeanCell{0.72}{0.35}{0} & \MeanCell{0.73}{0.33}{0} & \MeanCell{0.77}{0.28}{0} & \MeanCell{0.77}{0.26}{0} & \MeanCell{0.78}{0.24}{0} & \MeanCell{0.79}{0.22}{0} & \MeanCell{0.78}{0.22}{0} & \MeanCell{0.79}{0.23}{0} & \MeanCell{0.76}{0.25}{0} & \MeanCell{0.75}{0.27}{0} & \MeanCell{0.76}{0.25}{0} & \MeanCell{0.76}{0.24}{0} & \MeanCell{0.77}{0.21}{0} & \MeanCell{0.77}{0.20}{0} & \MeanCell{0.79}{0.17}{0} & \MeanCell{0.78}{0.18}{0} & \MeanCell{0.78}{0.25}{0} \\
}

\SetHeatColor{myred}
\ClusterKTable{8}{\glsfmtshort{acu} --- \PongName.}{acu_pong}{
        1 & \SetCell{bg=\HeatCellBg{-34}}\CellText{-0.34}{0}{0} & \SetCell{bg=\HeatCellBg{-41}}\CellText{-0.41}{0}{0} & \SetCell{bg=\HeatCellBg{18}}\CellText{0.18}{0}{1} & \SetCell{bg=\HeatCellBg{-4}}\CellText{-0.04}{0}{0} & \SetCell{bg=\HeatCellBg{-19}}\CellText{-0.19}{0}{0} & \SetCell{bg=\HeatCellBg{-3}}\CellText{-0.03}{0}{0} & \SetCell{bg=\HeatCellBg{-6}}\CellText{-0.06}{0}{0} & \SetCell{bg=\HeatCellBg{-7}}\CellText{-0.07}{0}{0} & \SetCell{bg=\HeatCellBg{3}}\CellText{0.03}{0}{0} & \SetCell{bg=\HeatCellBg{4}}\CellText{0.04}{0}{0} & \SetCell{bg=\HeatCellBg{-7}}\CellText{-0.07}{0}{0} & \SetCell{bg=\HeatCellBg{0}}\CellText{0.00}{0}{0} & \SetCell{bg=\HeatCellBg{-10}}\CellText{-0.10}{0}{0} & \SetCell{bg=\HeatCellBg{-6}}\CellText{-0.06}{0}{0} & \SetCell{bg=\HeatCellBg{-11}}\CellText{-0.11}{0}{0} & \SetCell{bg=\HeatCellBg{11}}\CellText{0.11}{0}{0} & \SetCell{bg=\HeatCellBg{1}}\CellText{0.01}{0}{0} & \SetCell{bg=\HeatCellBg{3}}\CellText{0.03}{0}{0} & \SetCell{bg=\HeatCellBg{7}}\CellText{0.07}{0}{0} & \MeanCell{-0.05}{0.14}{0} \\
        2 & \SetCell{bg=\HeatCellBg{100}}\CellText{1.00}{1}{1} & \SetCell{bg=\HeatCellBg{55}}\CellText{0.55}{0}{0} & \SetCell{bg=\HeatCellBg{18}}\CellText{0.18}{0}{0} & \SetCell{bg=\HeatCellBg{-9}}\CellText{-0.09}{0}{0} & \SetCell{bg=\HeatCellBg{5}}\CellText{0.05}{0}{0} & \SetCell{bg=\HeatCellBg{8}}\CellText{0.08}{0}{0} & \SetCell{bg=\HeatCellBg{16}}\CellText{0.16}{0}{0} & \SetCell{bg=\HeatCellBg{4}}\CellText{0.04}{0}{0} & \SetCell{bg=\HeatCellBg{27}}\CellText{0.27}{0}{0} & \SetCell{bg=\HeatCellBg{28}}\CellText{0.28}{0}{0} & \SetCell{bg=\HeatCellBg{21}}\CellText{0.21}{0}{0} & \SetCell{bg=\HeatCellBg{20}}\CellText{0.20}{0}{0} & \SetCell{bg=\HeatCellBg{29}}\CellText{0.29}{0}{0} & \SetCell{bg=\HeatCellBg{23}}\CellText{0.23}{0}{0} & \SetCell{bg=\HeatCellBg{20}}\CellText{0.20}{0}{0} & \SetCell{bg=\HeatCellBg{27}}\CellText{0.27}{0}{0} & \SetCell{bg=\HeatCellBg{26}}\CellText{0.26}{0}{0} & \SetCell{bg=\HeatCellBg{19}}\CellText{0.19}{0}{0} & \SetCell{bg=\HeatCellBg{21}}\CellText{0.21}{0}{0} & \MeanCell{0.24}{0.22}{0} \\
        3 & \SetCell{bg=\HeatCellBg{30}}\CellText{0.30}{0}{0} & \SetCell{bg=\HeatCellBg{65}}\CellText{0.65}{1}{0} & \SetCell{bg=\HeatCellBg{64}}\CellText{0.64}{1}{0} & \SetCell{bg=\HeatCellBg{49}}\CellText{0.49}{1}{0} & \SetCell{bg=\HeatCellBg{71}}\CellText{0.71}{1}{1} & \SetCell{bg=\HeatCellBg{60}}\CellText{0.60}{1}{0} & \SetCell{bg=\HeatCellBg{56}}\CellText{0.56}{1}{0} & \SetCell{bg=\HeatCellBg{55}}\CellText{0.55}{1}{0} & \SetCell{bg=\HeatCellBg{38}}\CellText{0.38}{0}{0} & \SetCell{bg=\HeatCellBg{36}}\CellText{0.36}{0}{0} & \SetCell{bg=\HeatCellBg{29}}\CellText{0.29}{0}{0} & \SetCell{bg=\HeatCellBg{26}}\CellText{0.26}{0}{0} & \SetCell{bg=\HeatCellBg{28}}\CellText{0.28}{0}{0} & \SetCell{bg=\HeatCellBg{27}}\CellText{0.27}{0}{0} & \SetCell{bg=\HeatCellBg{22}}\CellText{0.22}{0}{0} & \SetCell{bg=\HeatCellBg{19}}\CellText{0.19}{0}{0} & \SetCell{bg=\HeatCellBg{23}}\CellText{0.23}{0}{0} & \SetCell{bg=\HeatCellBg{26}}\CellText{0.26}{0}{0} & \SetCell{bg=\HeatCellBg{20}}\CellText{0.20}{0}{0} & \MeanCell{0.39}{0.17}{0} \\
        4 & \SetCell{bg=\HeatCellBg{59}}\CellText{0.59}{0}{1} & \SetCell{bg=\HeatCellBg{51}}\CellText{0.51}{0}{0} & \SetCell{bg=\HeatCellBg{56}}\CellText{0.56}{0}{0} & \SetCell{bg=\HeatCellBg{48}}\CellText{0.48}{0}{0} & \SetCell{bg=\HeatCellBg{41}}\CellText{0.41}{0}{0} & \SetCell{bg=\HeatCellBg{48}}\CellText{0.48}{0}{0} & \SetCell{bg=\HeatCellBg{56}}\CellText{0.56}{1}{0} & \SetCell{bg=\HeatCellBg{44}}\CellText{0.44}{0}{0} & \SetCell{bg=\HeatCellBg{44}}\CellText{0.44}{1}{0} & \SetCell{bg=\HeatCellBg{46}}\CellText{0.46}{1}{0} & \SetCell{bg=\HeatCellBg{45}}\CellText{0.45}{1}{0} & \SetCell{bg=\HeatCellBg{40}}\CellText{0.40}{1}{0} & \SetCell{bg=\HeatCellBg{43}}\CellText{0.43}{1}{0} & \SetCell{bg=\HeatCellBg{38}}\CellText{0.38}{1}{0} & \SetCell{bg=\HeatCellBg{39}}\CellText{0.39}{1}{0} & \SetCell{bg=\HeatCellBg{38}}\CellText{0.38}{1}{0} & \SetCell{bg=\HeatCellBg{38}}\CellText{0.38}{1}{0} & \SetCell{bg=\HeatCellBg{38}}\CellText{0.38}{1}{0} & \SetCell{bg=\HeatCellBg{44}}\CellText{0.44}{1}{0} & \MeanCell{0.45}{0.06}{1} \\
        5 & \SetCell{bg=\HeatCellBg{59}}\CellText{0.59}{0}{1} & \SetCell{bg=\HeatCellBg{51}}\CellText{0.51}{0}{0} & \SetCell{bg=\HeatCellBg{56}}\CellText{0.56}{0}{0} & \SetCell{bg=\HeatCellBg{48}}\CellText{0.48}{0}{0} & \SetCell{bg=\HeatCellBg{41}}\CellText{0.41}{0}{0} & \SetCell{bg=\HeatCellBg{48}}\CellText{0.48}{0}{0} & \SetCell{bg=\HeatCellBg{56}}\CellText{0.56}{1}{0} & \SetCell{bg=\HeatCellBg{44}}\CellText{0.44}{0}{0} & \SetCell{bg=\HeatCellBg{44}}\CellText{0.44}{1}{0} & \SetCell{bg=\HeatCellBg{46}}\CellText{0.46}{1}{0} & \SetCell{bg=\HeatCellBg{45}}\CellText{0.45}{1}{0} & \SetCell{bg=\HeatCellBg{38}}\CellText{0.38}{0}{0} & \SetCell{bg=\HeatCellBg{43}}\CellText{0.43}{1}{0} & \SetCell{bg=\HeatCellBg{38}}\CellText{0.38}{1}{0} & \SetCell{bg=\HeatCellBg{39}}\CellText{0.39}{1}{0} & \SetCell{bg=\HeatCellBg{38}}\CellText{0.38}{1}{0} & \SetCell{bg=\HeatCellBg{38}}\CellText{0.38}{1}{0} & \SetCell{bg=\HeatCellBg{38}}\CellText{0.38}{1}{0} & \SetCell{bg=\HeatCellBg{44}}\CellText{0.44}{1}{0} & \MeanCell{0.45}{0.06}{0} \\
        6a & \SetCell{bg=\HeatCellBg{33}}\CellText{0.33}{0}{1} & \SetCell{bg=\HeatCellBg{27}}\CellText{0.27}{0}{0} & \SetCell{bg=\HeatCellBg{24}}\CellText{0.24}{0}{0} & \SetCell{bg=\HeatCellBg{24}}\CellText{0.24}{0}{0} & \SetCell{bg=\HeatCellBg{21}}\CellText{0.21}{0}{0} & \SetCell{bg=\HeatCellBg{23}}\CellText{0.23}{0}{0} & \SetCell{bg=\HeatCellBg{24}}\CellText{0.24}{0}{0} & \SetCell{bg=\HeatCellBg{26}}\CellText{0.26}{0}{0} & \SetCell{bg=\HeatCellBg{26}}\CellText{0.26}{0}{0} & \SetCell{bg=\HeatCellBg{27}}\CellText{0.27}{0}{0} & \SetCell{bg=\HeatCellBg{27}}\CellText{0.27}{0}{0} & \SetCell{bg=\HeatCellBg{26}}\CellText{0.26}{0}{0} & \SetCell{bg=\HeatCellBg{27}}\CellText{0.27}{0}{0} & \SetCell{bg=\HeatCellBg{26}}\CellText{0.26}{0}{0} & \SetCell{bg=\HeatCellBg{26}}\CellText{0.26}{0}{0} & \SetCell{bg=\HeatCellBg{26}}\CellText{0.26}{0}{0} & \SetCell{bg=\HeatCellBg{26}}\CellText{0.26}{0}{0} & \SetCell{bg=\HeatCellBg{25}}\CellText{0.25}{0}{0} & \SetCell{bg=\HeatCellBg{23}}\CellText{0.23}{0}{0} & \MeanCell{0.26}{0.02}{0} \\
        Mean $\pm$ Std & \MeanCell{0.41}{0.41}{1} & \MeanCell{0.35}{0.36}{0} & \MeanCell{0.39}{0.20}{0} & \MeanCell{0.26}{0.25}{0} & \MeanCell{0.27}{0.29}{0} & \MeanCell{0.31}{0.23}{0} & \MeanCell{0.34}{0.24}{0} & \MeanCell{0.28}{0.23}{0} & \MeanCell{0.30}{0.14}{0} & \MeanCell{0.31}{0.14}{0} & \MeanCell{0.27}{0.18}{0} & \MeanCell{0.25}{0.13}{0} & \MeanCell{0.27}{0.18}{0} & \MeanCell{0.24}{0.15}{0} & \MeanCell{0.23}{0.17}{0} & \MeanCell{0.27}{0.10}{0} & \MeanCell{0.25}{0.12}{0} & \MeanCell{0.25}{0.12}{0} & \MeanCell{0.27}{0.13}{0} & \MeanCell{0.29}{0.22}{0} \\
}

\SetHeatColor{myblue}
\ClusterKTable{8}[Observation baseline --- \PongName.]{Observation baseline ($\Similarity{obs}{}$) --- \PongName.}{obs_pong}{
        1 & \SetCell{bg=\HeatCellBg{67}}\CellText{0.67}{1}{0} & \SetCell{bg=\HeatCellBg{30}}\CellText{0.30}{1}{0} & \SetCell{bg=\HeatCellBg{74}}\CellText{0.74}{1}{0} & \SetCell{bg=\HeatCellBg{83}}\CellText{0.83}{1}{1} & \SetCell{bg=\HeatCellBg{81}}\CellText{0.81}{1}{0} & \SetCell{bg=\HeatCellBg{72}}\CellText{0.72}{1}{0} & \SetCell{bg=\HeatCellBg{69}}\CellText{0.69}{1}{0} & \SetCell{bg=\HeatCellBg{62}}\CellText{0.62}{1}{0} & \SetCell{bg=\HeatCellBg{60}}\CellText{0.60}{1}{0} & \SetCell{bg=\HeatCellBg{61}}\CellText{0.61}{1}{0} & \SetCell{bg=\HeatCellBg{62}}\CellText{0.62}{1}{0} & \SetCell{bg=\HeatCellBg{74}}\CellText{0.74}{1}{0} & \SetCell{bg=\HeatCellBg{81}}\CellText{0.81}{1}{0} & \SetCell{bg=\HeatCellBg{75}}\CellText{0.75}{1}{0} & \SetCell{bg=\HeatCellBg{71}}\CellText{0.71}{1}{0} & \SetCell{bg=\HeatCellBg{62}}\CellText{0.62}{1}{0} & \SetCell{bg=\HeatCellBg{58}}\CellText{0.58}{1}{0} & \SetCell{bg=\HeatCellBg{54}}\CellText{0.54}{1}{0} & \SetCell{bg=\HeatCellBg{54}}\CellText{0.54}{1}{0} & \MeanCell{0.66}{0.12}{1} \\
        2 & \SetCell{bg=\HeatCellBg{0}}\CellText{-0.00}{0}{0} & \SetCell{bg=\HeatCellBg{0}}\CellText{0.00}{0}{0} & \SetCell{bg=\HeatCellBg{56}}\CellText{0.56}{0}{0} & \SetCell{bg=\HeatCellBg{73}}\CellText{0.73}{0}{1} & \SetCell{bg=\HeatCellBg{60}}\CellText{0.60}{0}{0} & \SetCell{bg=\HeatCellBg{50}}\CellText{0.50}{0}{0} & \SetCell{bg=\HeatCellBg{48}}\CellText{0.48}{0}{0} & \SetCell{bg=\HeatCellBg{47}}\CellText{0.47}{0}{0} & \SetCell{bg=\HeatCellBg{38}}\CellText{0.38}{0}{0} & \SetCell{bg=\HeatCellBg{34}}\CellText{0.34}{0}{0} & \SetCell{bg=\HeatCellBg{32}}\CellText{0.32}{0}{0} & \SetCell{bg=\HeatCellBg{33}}\CellText{0.33}{0}{0} & \SetCell{bg=\HeatCellBg{36}}\CellText{0.36}{0}{0} & \SetCell{bg=\HeatCellBg{33}}\CellText{0.33}{0}{0} & \SetCell{bg=\HeatCellBg{34}}\CellText{0.34}{0}{0} & \SetCell{bg=\HeatCellBg{29}}\CellText{0.29}{0}{0} & \SetCell{bg=\HeatCellBg{26}}\CellText{0.26}{0}{0} & \SetCell{bg=\HeatCellBg{25}}\CellText{0.25}{0}{0} & \SetCell{bg=\HeatCellBg{24}}\CellText{0.24}{0}{0} & \MeanCell{0.36}{0.18}{0} \\
        3 & \SetCell{bg=\HeatCellBg{0}}\CellText{-0.00}{0}{0} & \SetCell{bg=\HeatCellBg{0}}\CellText{0.00}{0}{0} & \SetCell{bg=\HeatCellBg{0}}\CellText{0.00}{0}{0} & \SetCell{bg=\HeatCellBg{5}}\CellText{0.05}{0}{0} & \SetCell{bg=\HeatCellBg{0}}\CellText{0.00}{0}{0} & \SetCell{bg=\HeatCellBg{2}}\CellText{0.02}{0}{0} & \SetCell{bg=\HeatCellBg{6}}\CellText{0.06}{0}{0} & \SetCell{bg=\HeatCellBg{5}}\CellText{0.05}{0}{0} & \SetCell{bg=\HeatCellBg{4}}\CellText{0.04}{0}{0} & \SetCell{bg=\HeatCellBg{6}}\CellText{0.06}{0}{0} & \SetCell{bg=\HeatCellBg{6}}\CellText{0.06}{0}{0} & \SetCell{bg=\HeatCellBg{7}}\CellText{0.07}{0}{0} & \SetCell{bg=\HeatCellBg{5}}\CellText{0.05}{0}{0} & \SetCell{bg=\HeatCellBg{7}}\CellText{0.07}{0}{0} & \SetCell{bg=\HeatCellBg{8}}\CellText{0.08}{0}{1} & \SetCell{bg=\HeatCellBg{7}}\CellText{0.07}{0}{0} & \SetCell{bg=\HeatCellBg{5}}\CellText{0.05}{0}{0} & \SetCell{bg=\HeatCellBg{5}}\CellText{0.05}{0}{0} & \SetCell{bg=\HeatCellBg{6}}\CellText{0.06}{0}{0} & \MeanCell{0.04}{0.03}{0} \\
        4 & \SetCell{bg=\HeatCellBg{0}}\CellText{0.00}{0}{0} & \SetCell{bg=\HeatCellBg{0}}\CellText{-0.00}{0}{0} & \SetCell{bg=\HeatCellBg{0}}\CellText{0.00}{0}{0} & \SetCell{bg=\HeatCellBg{0}}\CellText{0.00}{0}{0} & \SetCell{bg=\HeatCellBg{0}}\CellText{0.00}{0}{0} & \SetCell{bg=\HeatCellBg{0}}\CellText{0.00}{0}{0} & \SetCell{bg=\HeatCellBg{0}}\CellText{0.00}{0}{0} & \SetCell{bg=\HeatCellBg{0}}\CellText{0.00}{0}{0} & \SetCell{bg=\HeatCellBg{0}}\CellText{0.00}{0}{0} & \SetCell{bg=\HeatCellBg{0}}\CellText{0.00}{0}{0} & \SetCell{bg=\HeatCellBg{0}}\CellText{0.00}{0}{0} & \SetCell{bg=\HeatCellBg{1}}\CellText{0.01}{0}{0} & \SetCell{bg=\HeatCellBg{0}}\CellText{0.00}{0}{0} & \SetCell{bg=\HeatCellBg{0}}\CellText{0.00}{0}{0} & \SetCell{bg=\HeatCellBg{2}}\CellText{0.02}{0}{1} & \SetCell{bg=\HeatCellBg{0}}\CellText{0.00}{0}{0} & \SetCell{bg=\HeatCellBg{1}}\CellText{0.01}{0}{0} & \SetCell{bg=\HeatCellBg{1}}\CellText{0.01}{0}{0} & \SetCell{bg=\HeatCellBg{0}}\CellText{0.00}{0}{0} & \MeanCell{0.00}{0.01}{0} \\
        5 & \SetCell{bg=\HeatCellBg{0}}\CellText{0.00}{0}{0} & \SetCell{bg=\HeatCellBg{0}}\CellText{-0.00}{0}{0} & \SetCell{bg=\HeatCellBg{0}}\CellText{0.00}{0}{0} & \SetCell{bg=\HeatCellBg{0}}\CellText{0.00}{0}{0} & \SetCell{bg=\HeatCellBg{0}}\CellText{0.00}{0}{0} & \SetCell{bg=\HeatCellBg{0}}\CellText{0.00}{0}{0} & \SetCell{bg=\HeatCellBg{0}}\CellText{0.00}{0}{0} & \SetCell{bg=\HeatCellBg{0}}\CellText{0.00}{0}{0} & \SetCell{bg=\HeatCellBg{0}}\CellText{0.00}{0}{0} & \SetCell{bg=\HeatCellBg{0}}\CellText{0.00}{0}{0} & \SetCell{bg=\HeatCellBg{0}}\CellText{0.00}{0}{0} & \SetCell{bg=\HeatCellBg{1}}\CellText{0.01}{0}{0} & \SetCell{bg=\HeatCellBg{0}}\CellText{0.00}{0}{0} & \SetCell{bg=\HeatCellBg{0}}\CellText{0.00}{0}{0} & \SetCell{bg=\HeatCellBg{2}}\CellText{0.02}{0}{1} & \SetCell{bg=\HeatCellBg{0}}\CellText{0.00}{0}{0} & \SetCell{bg=\HeatCellBg{1}}\CellText{0.01}{0}{0} & \SetCell{bg=\HeatCellBg{1}}\CellText{0.01}{0}{0} & \SetCell{bg=\HeatCellBg{0}}\CellText{0.00}{0}{0} & \MeanCell{0.00}{0.01}{0} \\
        6a & \SetCell{bg=\HeatCellBg{0}}\CellText{0.00}{0}{0} & \SetCell{bg=\HeatCellBg{0}}\CellText{0.00}{0}{0} & \SetCell{bg=\HeatCellBg{0}}\CellText{0.00}{0}{0} & \SetCell{bg=\HeatCellBg{0}}\CellText{0.00}{0}{0} & \SetCell{bg=\HeatCellBg{0}}\CellText{0.00}{0}{0} & \SetCell{bg=\HeatCellBg{0}}\CellText{0.00}{0}{0} & \SetCell{bg=\HeatCellBg{0}}\CellText{0.00}{0}{0} & \SetCell{bg=\HeatCellBg{0}}\CellText{0.00}{0}{0} & \SetCell{bg=\HeatCellBg{1}}\CellText{0.01}{0}{0} & \SetCell{bg=\HeatCellBg{1}}\CellText{0.01}{0}{0} & \SetCell{bg=\HeatCellBg{1}}\CellText{0.01}{0}{0} & \SetCell{bg=\HeatCellBg{1}}\CellText{0.01}{0}{0} & \SetCell{bg=\HeatCellBg{1}}\CellText{0.01}{0}{0} & \SetCell{bg=\HeatCellBg{1}}\CellText{0.01}{0}{0} & \SetCell{bg=\HeatCellBg{3}}\CellText{0.03}{0}{1} & \SetCell{bg=\HeatCellBg{2}}\CellText{0.02}{0}{0} & \SetCell{bg=\HeatCellBg{2}}\CellText{0.02}{0}{0} & \SetCell{bg=\HeatCellBg{2}}\CellText{0.02}{0}{0} & \SetCell{bg=\HeatCellBg{2}}\CellText{0.02}{0}{0} & \MeanCell{0.01}{0.01}{0} \\
        Mean $\pm$ Std & \MeanCell{0.11}{0.25}{0} & \MeanCell{0.05}{0.11}{0} & \MeanCell{0.22}{0.31}{0} & \MeanCell{0.27}{0.36}{1} & \MeanCell{0.24}{0.34}{0} & \MeanCell{0.21}{0.29}{0} & \MeanCell{0.20}{0.28}{0} & \MeanCell{0.19}{0.26}{0} & \MeanCell{0.17}{0.23}{0} & \MeanCell{0.17}{0.23}{0} & \MeanCell{0.17}{0.23}{0} & \MeanCell{0.20}{0.27}{0} & \MeanCell{0.21}{0.30}{0} & \MeanCell{0.19}{0.27}{0} & \MeanCell{0.20}{0.25}{0} & \MeanCell{0.17}{0.23}{0} & \MeanCell{0.15}{0.21}{0} & \MeanCell{0.15}{0.19}{0} & \MeanCell{0.14}{0.20}{0} & \MeanCell{0.18}{0.26}{0} \\
}

\SetHeatColor{myblue}
\ClusterKTable{8}[Untrained-network baseline --- \PongName.]{Untrained-network baseline ($\Similarity{untr}{}$) --- \PongName.}{untr_pong}{
        1 & \SetCell{bg=\HeatCellBg{45}}\CellText{0.45}{1}{0} & \SetCell{bg=\HeatCellBg{23}}\CellText{0.23}{0}{0} & \SetCell{bg=\HeatCellBg{53}}\CellText{0.53}{1}{0} & \SetCell{bg=\HeatCellBg{60}}\CellText{0.60}{1}{0} & \SetCell{bg=\HeatCellBg{79}}\CellText{0.79}{1}{1} & \SetCell{bg=\HeatCellBg{63}}\CellText{0.63}{1}{0} & \SetCell{bg=\HeatCellBg{61}}\CellText{0.61}{1}{0} & \SetCell{bg=\HeatCellBg{54}}\CellText{0.54}{1}{0} & \SetCell{bg=\HeatCellBg{62}}\CellText{0.62}{1}{0} & \SetCell{bg=\HeatCellBg{63}}\CellText{0.63}{1}{0} & \SetCell{bg=\HeatCellBg{59}}\CellText{0.59}{1}{0} & \SetCell{bg=\HeatCellBg{72}}\CellText{0.72}{1}{0} & \SetCell{bg=\HeatCellBg{67}}\CellText{0.67}{1}{0} & \SetCell{bg=\HeatCellBg{67}}\CellText{0.67}{1}{0} & \SetCell{bg=\HeatCellBg{65}}\CellText{0.65}{1}{0} & \SetCell{bg=\HeatCellBg{64}}\CellText{0.64}{1}{0} & \SetCell{bg=\HeatCellBg{59}}\CellText{0.59}{1}{0} & \SetCell{bg=\HeatCellBg{50}}\CellText{0.50}{1}{0} & \SetCell{bg=\HeatCellBg{57}}\CellText{0.57}{1}{0} & \MeanCell{0.59}{0.11}{1} \\
        2 & \SetCell{bg=\HeatCellBg{0}}\CellText{-0.00}{0}{0} & \SetCell{bg=\HeatCellBg{36}}\CellText{0.36}{1}{0} & \SetCell{bg=\HeatCellBg{43}}\CellText{0.43}{0}{0} & \SetCell{bg=\HeatCellBg{52}}\CellText{0.52}{0}{0} & \SetCell{bg=\HeatCellBg{61}}\CellText{0.61}{0}{1} & \SetCell{bg=\HeatCellBg{47}}\CellText{0.47}{0}{0} & \SetCell{bg=\HeatCellBg{46}}\CellText{0.46}{0}{0} & \SetCell{bg=\HeatCellBg{41}}\CellText{0.41}{0}{0} & \SetCell{bg=\HeatCellBg{38}}\CellText{0.38}{0}{0} & \SetCell{bg=\HeatCellBg{37}}\CellText{0.37}{0}{0} & \SetCell{bg=\HeatCellBg{38}}\CellText{0.38}{0}{0} & \SetCell{bg=\HeatCellBg{34}}\CellText{0.34}{0}{0} & \SetCell{bg=\HeatCellBg{37}}\CellText{0.37}{0}{0} & \SetCell{bg=\HeatCellBg{36}}\CellText{0.36}{0}{0} & \SetCell{bg=\HeatCellBg{39}}\CellText{0.39}{0}{0} & \SetCell{bg=\HeatCellBg{33}}\CellText{0.33}{0}{0} & \SetCell{bg=\HeatCellBg{37}}\CellText{0.37}{0}{0} & \SetCell{bg=\HeatCellBg{32}}\CellText{0.32}{0}{0} & \SetCell{bg=\HeatCellBg{34}}\CellText{0.34}{0}{0} & \MeanCell{0.38}{0.11}{0} \\
        3 & \SetCell{bg=\HeatCellBg{0}}\CellText{-0.00}{0}{0} & \SetCell{bg=\HeatCellBg{2}}\CellText{0.02}{0}{0} & \SetCell{bg=\HeatCellBg{0}}\CellText{0.00}{0}{0} & \SetCell{bg=\HeatCellBg{6}}\CellText{0.06}{0}{0} & \SetCell{bg=\HeatCellBg{2}}\CellText{0.02}{0}{0} & \SetCell{bg=\HeatCellBg{4}}\CellText{0.04}{0}{0} & \SetCell{bg=\HeatCellBg{7}}\CellText{0.07}{0}{0} & \SetCell{bg=\HeatCellBg{7}}\CellText{0.07}{0}{0} & \SetCell{bg=\HeatCellBg{8}}\CellText{0.08}{0}{0} & \SetCell{bg=\HeatCellBg{11}}\CellText{0.11}{0}{0} & \SetCell{bg=\HeatCellBg{11}}\CellText{0.11}{0}{0} & \SetCell{bg=\HeatCellBg{13}}\CellText{0.13}{0}{0} & \SetCell{bg=\HeatCellBg{11}}\CellText{0.11}{0}{0} & \SetCell{bg=\HeatCellBg{13}}\CellText{0.13}{0}{0} & \SetCell{bg=\HeatCellBg{15}}\CellText{0.15}{0}{0} & \SetCell{bg=\HeatCellBg{16}}\CellText{0.16}{0}{0} & \SetCell{bg=\HeatCellBg{13}}\CellText{0.13}{0}{0} & \SetCell{bg=\HeatCellBg{13}}\CellText{0.13}{0}{0} & \SetCell{bg=\HeatCellBg{17}}\CellText{0.17}{0}{1} & \MeanCell{0.09}{0.05}{0} \\
        4 & \SetCell{bg=\HeatCellBg{0}}\CellText{0.00}{0}{0} & \SetCell{bg=\HeatCellBg{0}}\CellText{-0.00}{0}{0} & \SetCell{bg=\HeatCellBg{2}}\CellText{0.02}{0}{0} & \SetCell{bg=\HeatCellBg{1}}\CellText{0.01}{0}{0} & \SetCell{bg=\HeatCellBg{3}}\CellText{0.03}{0}{0} & \SetCell{bg=\HeatCellBg{4}}\CellText{0.04}{0}{0} & \SetCell{bg=\HeatCellBg{3}}\CellText{0.03}{0}{0} & \SetCell{bg=\HeatCellBg{4}}\CellText{0.04}{0}{0} & \SetCell{bg=\HeatCellBg{7}}\CellText{0.07}{0}{0} & \SetCell{bg=\HeatCellBg{6}}\CellText{0.06}{0}{0} & \SetCell{bg=\HeatCellBg{4}}\CellText{0.04}{0}{0} & \SetCell{bg=\HeatCellBg{7}}\CellText{0.07}{0}{0} & \SetCell{bg=\HeatCellBg{7}}\CellText{0.07}{0}{0} & \SetCell{bg=\HeatCellBg{7}}\CellText{0.07}{0}{0} & \SetCell{bg=\HeatCellBg{7}}\CellText{0.07}{0}{0} & \SetCell{bg=\HeatCellBg{9}}\CellText{0.09}{0}{0} & \SetCell{bg=\HeatCellBg{10}}\CellText{0.10}{0}{1} & \SetCell{bg=\HeatCellBg{10}}\CellText{0.10}{0}{1} & \SetCell{bg=\HeatCellBg{9}}\CellText{0.09}{0}{0} & \MeanCell{0.05}{0.03}{0} \\
        5 & \SetCell{bg=\HeatCellBg{0}}\CellText{0.00}{0}{0} & \SetCell{bg=\HeatCellBg{0}}\CellText{-0.00}{0}{0} & \SetCell{bg=\HeatCellBg{2}}\CellText{0.02}{0}{0} & \SetCell{bg=\HeatCellBg{1}}\CellText{0.01}{0}{0} & \SetCell{bg=\HeatCellBg{3}}\CellText{0.03}{0}{0} & \SetCell{bg=\HeatCellBg{4}}\CellText{0.04}{0}{0} & \SetCell{bg=\HeatCellBg{3}}\CellText{0.03}{0}{0} & \SetCell{bg=\HeatCellBg{4}}\CellText{0.04}{0}{0} & \SetCell{bg=\HeatCellBg{7}}\CellText{0.07}{0}{0} & \SetCell{bg=\HeatCellBg{6}}\CellText{0.06}{0}{0} & \SetCell{bg=\HeatCellBg{4}}\CellText{0.04}{0}{0} & \SetCell{bg=\HeatCellBg{8}}\CellText{0.08}{0}{0} & \SetCell{bg=\HeatCellBg{7}}\CellText{0.07}{0}{0} & \SetCell{bg=\HeatCellBg{7}}\CellText{0.07}{0}{0} & \SetCell{bg=\HeatCellBg{7}}\CellText{0.07}{0}{0} & \SetCell{bg=\HeatCellBg{9}}\CellText{0.09}{0}{0} & \SetCell{bg=\HeatCellBg{10}}\CellText{0.10}{0}{1} & \SetCell{bg=\HeatCellBg{10}}\CellText{0.10}{0}{1} & \SetCell{bg=\HeatCellBg{9}}\CellText{0.09}{0}{0} & \MeanCell{0.05}{0.03}{0} \\
        6a & \SetCell{bg=\HeatCellBg{1}}\CellText{0.01}{0}{0} & \SetCell{bg=\HeatCellBg{2}}\CellText{0.02}{0}{0} & \SetCell{bg=\HeatCellBg{2}}\CellText{0.02}{0}{0} & \SetCell{bg=\HeatCellBg{1}}\CellText{0.01}{0}{0} & \SetCell{bg=\HeatCellBg{5}}\CellText{0.05}{0}{0} & \SetCell{bg=\HeatCellBg{3}}\CellText{0.03}{0}{0} & \SetCell{bg=\HeatCellBg{3}}\CellText{0.03}{0}{0} & \SetCell{bg=\HeatCellBg{4}}\CellText{0.04}{0}{0} & \SetCell{bg=\HeatCellBg{6}}\CellText{0.06}{0}{0} & \SetCell{bg=\HeatCellBg{6}}\CellText{0.06}{0}{0} & \SetCell{bg=\HeatCellBg{6}}\CellText{0.06}{0}{0} & \SetCell{bg=\HeatCellBg{6}}\CellText{0.06}{0}{0} & \SetCell{bg=\HeatCellBg{5}}\CellText{0.05}{0}{0} & \SetCell{bg=\HeatCellBg{6}}\CellText{0.06}{0}{0} & \SetCell{bg=\HeatCellBg{5}}\CellText{0.05}{0}{0} & \SetCell{bg=\HeatCellBg{6}}\CellText{0.06}{0}{0} & \SetCell{bg=\HeatCellBg{6}}\CellText{0.06}{0}{0} & \SetCell{bg=\HeatCellBg{6}}\CellText{0.06}{0}{0} & \SetCell{bg=\HeatCellBg{7}}\CellText{0.07}{0}{1} & \MeanCell{0.05}{0.02}{0} \\
        Mean $\pm$ Std & \MeanCell{0.08}{0.17}{0} & \MeanCell{0.10}{0.14}{0} & \MeanCell{0.17}{0.22}{0} & \MeanCell{0.20}{0.26}{0} & \MeanCell{0.26}{0.32}{1} & \MeanCell{0.21}{0.25}{0} & \MeanCell{0.21}{0.24}{0} & \MeanCell{0.19}{0.21}{0} & \MeanCell{0.21}{0.21}{0} & \MeanCell{0.21}{0.22}{0} & \MeanCell{0.20}{0.21}{0} & \MeanCell{0.23}{0.24}{0} & \MeanCell{0.22}{0.23}{0} & \MeanCell{0.23}{0.22}{0} & \MeanCell{0.23}{0.22}{0} & \MeanCell{0.23}{0.20}{0} & \MeanCell{0.23}{0.19}{0} & \MeanCell{0.20}{0.16}{0} & \MeanCell{0.22}{0.18}{0} & \MeanCell{0.20}{0.22}{0} \\
}

\SetHeatColor{myblue}
\ClusterKTable{8}[Action baseline --- \PongName.]{Action baseline ($\Similarity{act}{}$) --- \PongName.}{act_pong}{
        1 & \SetCell{bg=\HeatCellBg{0}}\CellText{-0.00}{0}{0} & \SetCell{bg=\HeatCellBg{0}}\CellText{-0.00}{0}{0} & \SetCell{bg=\HeatCellBg{-1}}\CellText{-0.01}{0}{0} & \SetCell{bg=\HeatCellBg{0}}\CellText{0.00}{0}{0} & \SetCell{bg=\HeatCellBg{0}}\CellText{0.00}{0}{0} & \SetCell{bg=\HeatCellBg{1}}\CellText{0.01}{0}{1} & \SetCell{bg=\HeatCellBg{1}}\CellText{0.01}{0}{1} & \SetCell{bg=\HeatCellBg{1}}\CellText{0.01}{0}{1} & \SetCell{bg=\HeatCellBg{1}}\CellText{0.01}{0}{1} & \SetCell{bg=\HeatCellBg{1}}\CellText{0.01}{0}{1} & \SetCell{bg=\HeatCellBg{1}}\CellText{0.01}{0}{1} & \SetCell{bg=\HeatCellBg{1}}\CellText{0.01}{0}{1} & \SetCell{bg=\HeatCellBg{1}}\CellText{0.01}{0}{1} & \SetCell{bg=\HeatCellBg{1}}\CellText{0.01}{0}{1} & \SetCell{bg=\HeatCellBg{1}}\CellText{0.01}{0}{1} & \SetCell{bg=\HeatCellBg{1}}\CellText{0.01}{0}{1} & \SetCell{bg=\HeatCellBg{1}}\CellText{0.01}{0}{1} & \SetCell{bg=\HeatCellBg{1}}\CellText{0.01}{0}{1} & \SetCell{bg=\HeatCellBg{1}}\CellText{0.01}{0}{1} & \MeanCell{0.01}{0.01}{0} \\
        2 & \SetCell{bg=\HeatCellBg{0}}\CellText{-0.00}{0}{0} & \SetCell{bg=\HeatCellBg{0}}\CellText{-0.00}{0}{0} & \SetCell{bg=\HeatCellBg{1}}\CellText{0.01}{0}{1} & \SetCell{bg=\HeatCellBg{1}}\CellText{0.01}{0}{1} & \SetCell{bg=\HeatCellBg{1}}\CellText{0.01}{0}{1} & \SetCell{bg=\HeatCellBg{1}}\CellText{0.01}{0}{1} & \SetCell{bg=\HeatCellBg{1}}\CellText{0.01}{0}{1} & \SetCell{bg=\HeatCellBg{1}}\CellText{0.01}{0}{1} & \SetCell{bg=\HeatCellBg{1}}\CellText{0.01}{0}{1} & \SetCell{bg=\HeatCellBg{1}}\CellText{0.01}{0}{1} & \SetCell{bg=\HeatCellBg{1}}\CellText{0.01}{0}{1} & \SetCell{bg=\HeatCellBg{1}}\CellText{0.01}{0}{1} & \SetCell{bg=\HeatCellBg{1}}\CellText{0.01}{0}{1} & \SetCell{bg=\HeatCellBg{1}}\CellText{0.01}{0}{1} & \SetCell{bg=\HeatCellBg{1}}\CellText{0.01}{0}{1} & \SetCell{bg=\HeatCellBg{1}}\CellText{0.01}{0}{1} & \SetCell{bg=\HeatCellBg{1}}\CellText{0.01}{0}{1} & \SetCell{bg=\HeatCellBg{1}}\CellText{0.01}{0}{1} & \SetCell{bg=\HeatCellBg{1}}\CellText{0.01}{0}{1} & \MeanCell{0.01}{0.00}{0} \\
        3 & \SetCell{bg=\HeatCellBg{1}}\CellText{0.01}{0}{0} & \SetCell{bg=\HeatCellBg{1}}\CellText{0.01}{0}{0} & \SetCell{bg=\HeatCellBg{0}}\CellText{0.00}{0}{0} & \SetCell{bg=\HeatCellBg{1}}\CellText{0.01}{0}{0} & \SetCell{bg=\HeatCellBg{1}}\CellText{0.01}{0}{0} & \SetCell{bg=\HeatCellBg{1}}\CellText{0.01}{0}{0} & \SetCell{bg=\HeatCellBg{2}}\CellText{0.02}{0}{1} & \SetCell{bg=\HeatCellBg{1}}\CellText{0.01}{0}{0} & \SetCell{bg=\HeatCellBg{2}}\CellText{0.02}{0}{1} & \SetCell{bg=\HeatCellBg{1}}\CellText{0.01}{0}{0} & \SetCell{bg=\HeatCellBg{1}}\CellText{0.01}{0}{0} & \SetCell{bg=\HeatCellBg{1}}\CellText{0.01}{0}{0} & \SetCell{bg=\HeatCellBg{1}}\CellText{0.01}{0}{0} & \SetCell{bg=\HeatCellBg{1}}\CellText{0.01}{0}{0} & \SetCell{bg=\HeatCellBg{1}}\CellText{0.01}{0}{0} & \SetCell{bg=\HeatCellBg{1}}\CellText{0.01}{0}{0} & \SetCell{bg=\HeatCellBg{1}}\CellText{0.01}{0}{0} & \SetCell{bg=\HeatCellBg{1}}\CellText{0.01}{0}{0} & \SetCell{bg=\HeatCellBg{1}}\CellText{0.01}{0}{0} & \MeanCell{0.01}{0.00}{0} \\
        4 & \SetCell{bg=\HeatCellBg{-1}}\CellText{-0.01}{0}{0} & \SetCell{bg=\HeatCellBg{0}}\CellText{0.00}{0}{0} & \SetCell{bg=\HeatCellBg{0}}\CellText{0.00}{0}{0} & \SetCell{bg=\HeatCellBg{1}}\CellText{0.01}{0}{0} & \SetCell{bg=\HeatCellBg{2}}\CellText{0.02}{0}{1} & \SetCell{bg=\HeatCellBg{1}}\CellText{0.01}{0}{0} & \SetCell{bg=\HeatCellBg{2}}\CellText{0.02}{0}{1} & \SetCell{bg=\HeatCellBg{2}}\CellText{0.02}{0}{1} & \SetCell{bg=\HeatCellBg{2}}\CellText{0.02}{0}{1} & \SetCell{bg=\HeatCellBg{2}}\CellText{0.02}{0}{1} & \SetCell{bg=\HeatCellBg{2}}\CellText{0.02}{0}{1} & \SetCell{bg=\HeatCellBg{2}}\CellText{0.02}{0}{1} & \SetCell{bg=\HeatCellBg{2}}\CellText{0.02}{0}{1} & \SetCell{bg=\HeatCellBg{2}}\CellText{0.02}{0}{1} & \SetCell{bg=\HeatCellBg{2}}\CellText{0.02}{0}{1} & \SetCell{bg=\HeatCellBg{2}}\CellText{0.02}{0}{1} & \SetCell{bg=\HeatCellBg{2}}\CellText{0.02}{0}{1} & \SetCell{bg=\HeatCellBg{2}}\CellText{0.02}{0}{1} & \SetCell{bg=\HeatCellBg{2}}\CellText{0.02}{0}{1} & \MeanCell{0.02}{0.01}{0} \\
        5 & \SetCell{bg=\HeatCellBg{-1}}\CellText{-0.01}{0}{0} & \SetCell{bg=\HeatCellBg{0}}\CellText{0.00}{0}{0} & \SetCell{bg=\HeatCellBg{0}}\CellText{0.00}{0}{0} & \SetCell{bg=\HeatCellBg{1}}\CellText{0.01}{0}{0} & \SetCell{bg=\HeatCellBg{2}}\CellText{0.02}{0}{1} & \SetCell{bg=\HeatCellBg{1}}\CellText{0.01}{0}{0} & \SetCell{bg=\HeatCellBg{2}}\CellText{0.02}{0}{1} & \SetCell{bg=\HeatCellBg{2}}\CellText{0.02}{0}{1} & \SetCell{bg=\HeatCellBg{2}}\CellText{0.02}{0}{1} & \SetCell{bg=\HeatCellBg{2}}\CellText{0.02}{0}{1} & \SetCell{bg=\HeatCellBg{2}}\CellText{0.02}{0}{1} & \SetCell{bg=\HeatCellBg{2}}\CellText{0.02}{0}{1} & \SetCell{bg=\HeatCellBg{2}}\CellText{0.02}{0}{1} & \SetCell{bg=\HeatCellBg{2}}\CellText{0.02}{0}{1} & \SetCell{bg=\HeatCellBg{2}}\CellText{0.02}{0}{1} & \SetCell{bg=\HeatCellBg{2}}\CellText{0.02}{0}{1} & \SetCell{bg=\HeatCellBg{2}}\CellText{0.02}{0}{1} & \SetCell{bg=\HeatCellBg{2}}\CellText{0.02}{0}{1} & \SetCell{bg=\HeatCellBg{2}}\CellText{0.02}{0}{1} & \MeanCell{0.02}{0.01}{0} \\
        6a & \SetCell{bg=\HeatCellBg{3}}\CellText{0.03}{1}{0} & \SetCell{bg=\HeatCellBg{6}}\CellText{0.06}{1}{1} & \SetCell{bg=\HeatCellBg{4}}\CellText{0.04}{1}{0} & \SetCell{bg=\HeatCellBg{5}}\CellText{0.05}{1}{0} & \SetCell{bg=\HeatCellBg{5}}\CellText{0.05}{1}{0} & \SetCell{bg=\HeatCellBg{5}}\CellText{0.05}{1}{0} & \SetCell{bg=\HeatCellBg{5}}\CellText{0.05}{1}{0} & \SetCell{bg=\HeatCellBg{4}}\CellText{0.04}{1}{0} & \SetCell{bg=\HeatCellBg{4}}\CellText{0.04}{1}{0} & \SetCell{bg=\HeatCellBg{4}}\CellText{0.04}{1}{0} & \SetCell{bg=\HeatCellBg{4}}\CellText{0.04}{1}{0} & \SetCell{bg=\HeatCellBg{3}}\CellText{0.03}{1}{0} & \SetCell{bg=\HeatCellBg{3}}\CellText{0.03}{1}{0} & \SetCell{bg=\HeatCellBg{3}}\CellText{0.03}{1}{0} & \SetCell{bg=\HeatCellBg{3}}\CellText{0.03}{1}{0} & \SetCell{bg=\HeatCellBg{3}}\CellText{0.03}{1}{0} & \SetCell{bg=\HeatCellBg{3}}\CellText{0.03}{1}{0} & \SetCell{bg=\HeatCellBg{3}}\CellText{0.03}{1}{0} & \SetCell{bg=\HeatCellBg{3}}\CellText{0.03}{1}{0} & \MeanCell{0.04}{0.01}{1} \\
        Mean $\pm$ Std & \MeanCell{0.00}{0.01}{0} & \MeanCell{0.01}{0.02}{0} & \MeanCell{0.01}{0.02}{0} & \MeanCell{0.02}{0.02}{0} & \MeanCell{0.02}{0.02}{0} & \MeanCell{0.02}{0.01}{0} & \MeanCell{0.02}{0.01}{1} & \MeanCell{0.02}{0.01}{0} & \MeanCell{0.02}{0.01}{0} & \MeanCell{0.02}{0.01}{0} & \MeanCell{0.02}{0.01}{0} & \MeanCell{0.02}{0.01}{0} & \MeanCell{0.02}{0.01}{0} & \MeanCell{0.02}{0.01}{0} & \MeanCell{0.02}{0.01}{0} & \MeanCell{0.02}{0.01}{0} & \MeanCell{0.02}{0.01}{0} & \MeanCell{0.02}{0.01}{0} & \MeanCell{0.02}{0.01}{0} & \MeanCell{0.02}{0.01}{0} \\
}


\clearpage
\section{\SpaceInvadersName}
\label{sec:appendix_additional_results_space_invaders}
\ClusterKTable{8}{\glsfmtshort{cu} --- \SpaceInvadersName.}{cu_space_invaders}{
        1 & \SetCell{bg=\HeatCellBg{43}}\CellText{0.43}{0}{0} & \SetCell{bg=\HeatCellBg{100}}\CellText{1.00}{1}{1} & \SetCell{bg=\HeatCellBg{96}}\CellText{0.96}{0}{0} & \SetCell{bg=\HeatCellBg{94}}\CellText{0.94}{1}{0} & \SetCell{bg=\HeatCellBg{77}}\CellText{0.77}{0}{0} & \SetCell{bg=\HeatCellBg{71}}\CellText{0.71}{0}{0} & \SetCell{bg=\HeatCellBg{80}}\CellText{0.80}{0}{0} & \SetCell{bg=\HeatCellBg{67}}\CellText{0.67}{0}{0} & \SetCell{bg=\HeatCellBg{65}}\CellText{0.65}{0}{0} & \SetCell{bg=\HeatCellBg{60}}\CellText{0.60}{0}{0} & \SetCell{bg=\HeatCellBg{63}}\CellText{0.63}{0}{0} & \SetCell{bg=\HeatCellBg{71}}\CellText{0.71}{1}{0} & \SetCell{bg=\HeatCellBg{83}}\CellText{0.83}{1}{0} & \SetCell{bg=\HeatCellBg{76}}\CellText{0.76}{0}{0} & \SetCell{bg=\HeatCellBg{64}}\CellText{0.64}{0}{0} & \SetCell{bg=\HeatCellBg{67}}\CellText{0.67}{0}{0} & \SetCell{bg=\HeatCellBg{83}}\CellText{0.83}{1}{0} & \SetCell{bg=\HeatCellBg{82}}\CellText{0.82}{1}{0} & \SetCell{bg=\HeatCellBg{83}}\CellText{0.83}{1}{0} & \MeanCell{0.75}{0.14}{0} \\
        2 & \SetCell{bg=\HeatCellBg{99}}\CellText{0.99}{1}{1} & \SetCell{bg=\HeatCellBg{99}}\CellText{0.99}{0}{1} & \SetCell{bg=\HeatCellBg{98}}\CellText{0.98}{1}{0} & \SetCell{bg=\HeatCellBg{92}}\CellText{0.92}{0}{0} & \SetCell{bg=\HeatCellBg{81}}\CellText{0.81}{1}{0} & \SetCell{bg=\HeatCellBg{85}}\CellText{0.85}{1}{0} & \SetCell{bg=\HeatCellBg{95}}\CellText{0.95}{1}{0} & \SetCell{bg=\HeatCellBg{85}}\CellText{0.85}{1}{0} & \SetCell{bg=\HeatCellBg{67}}\CellText{0.67}{1}{0} & \SetCell{bg=\HeatCellBg{74}}\CellText{0.74}{1}{0} & \SetCell{bg=\HeatCellBg{75}}\CellText{0.75}{1}{0} & \SetCell{bg=\HeatCellBg{71}}\CellText{0.71}{1}{0} & \SetCell{bg=\HeatCellBg{63}}\CellText{0.63}{0}{0} & \SetCell{bg=\HeatCellBg{82}}\CellText{0.82}{1}{0} & \SetCell{bg=\HeatCellBg{73}}\CellText{0.73}{1}{0} & \SetCell{bg=\HeatCellBg{71}}\CellText{0.71}{1}{0} & \SetCell{bg=\HeatCellBg{67}}\CellText{0.67}{0}{0} & \SetCell{bg=\HeatCellBg{78}}\CellText{0.78}{0}{0} & \SetCell{bg=\HeatCellBg{82}}\CellText{0.82}{0}{0} & \MeanCell{0.81}{0.11}{1} \\
        3 & \SetCell{bg=\HeatCellBg{99}}\CellText{0.99}{1}{1} & \SetCell{bg=\HeatCellBg{98}}\CellText{0.98}{0}{0} & \SetCell{bg=\HeatCellBg{93}}\CellText{0.93}{0}{0} & \SetCell{bg=\HeatCellBg{78}}\CellText{0.78}{0}{0} & \SetCell{bg=\HeatCellBg{74}}\CellText{0.74}{0}{0} & \SetCell{bg=\HeatCellBg{71}}\CellText{0.71}{0}{0} & \SetCell{bg=\HeatCellBg{69}}\CellText{0.69}{0}{0} & \SetCell{bg=\HeatCellBg{67}}\CellText{0.67}{0}{0} & \SetCell{bg=\HeatCellBg{63}}\CellText{0.63}{0}{0} & \SetCell{bg=\HeatCellBg{60}}\CellText{0.60}{0}{0} & \SetCell{bg=\HeatCellBg{62}}\CellText{0.62}{0}{0} & \SetCell{bg=\HeatCellBg{60}}\CellText{0.60}{0}{0} & \SetCell{bg=\HeatCellBg{74}}\CellText{0.74}{0}{0} & \SetCell{bg=\HeatCellBg{72}}\CellText{0.72}{0}{0} & \SetCell{bg=\HeatCellBg{71}}\CellText{0.71}{0}{0} & \SetCell{bg=\HeatCellBg{69}}\CellText{0.69}{0}{0} & \SetCell{bg=\HeatCellBg{77}}\CellText{0.77}{0}{0} & \SetCell{bg=\HeatCellBg{72}}\CellText{0.72}{0}{0} & \SetCell{bg=\HeatCellBg{76}}\CellText{0.76}{0}{0} & \MeanCell{0.74}{0.11}{0} \\
        4 & \SetCell{bg=\HeatCellBg{31}}\CellText{0.31}{0}{0} & \SetCell{bg=\HeatCellBg{73}}\CellText{0.73}{0}{0} & \SetCell{bg=\HeatCellBg{92}}\CellText{0.92}{0}{1} & \SetCell{bg=\HeatCellBg{69}}\CellText{0.69}{0}{0} & \SetCell{bg=\HeatCellBg{61}}\CellText{0.61}{0}{0} & \SetCell{bg=\HeatCellBg{57}}\CellText{0.57}{0}{0} & \SetCell{bg=\HeatCellBg{64}}\CellText{0.64}{0}{0} & \SetCell{bg=\HeatCellBg{61}}\CellText{0.61}{0}{0} & \SetCell{bg=\HeatCellBg{61}}\CellText{0.61}{0}{0} & \SetCell{bg=\HeatCellBg{66}}\CellText{0.66}{0}{0} & \SetCell{bg=\HeatCellBg{58}}\CellText{0.58}{0}{0} & \SetCell{bg=\HeatCellBg{52}}\CellText{0.52}{0}{0} & \SetCell{bg=\HeatCellBg{56}}\CellText{0.56}{0}{0} & \SetCell{bg=\HeatCellBg{63}}\CellText{0.63}{0}{0} & \SetCell{bg=\HeatCellBg{67}}\CellText{0.67}{0}{0} & \SetCell{bg=\HeatCellBg{63}}\CellText{0.63}{0}{0} & \SetCell{bg=\HeatCellBg{63}}\CellText{0.63}{0}{0} & \SetCell{bg=\HeatCellBg{68}}\CellText{0.68}{0}{0} & \SetCell{bg=\HeatCellBg{64}}\CellText{0.64}{0}{0} & \MeanCell{0.63}{0.11}{0} \\
        5 & \SetCell{bg=\HeatCellBg{31}}\CellText{0.31}{0}{0} & \SetCell{bg=\HeatCellBg{73}}\CellText{0.73}{0}{0} & \SetCell{bg=\HeatCellBg{92}}\CellText{0.92}{0}{1} & \SetCell{bg=\HeatCellBg{69}}\CellText{0.69}{0}{0} & \SetCell{bg=\HeatCellBg{61}}\CellText{0.61}{0}{0} & \SetCell{bg=\HeatCellBg{57}}\CellText{0.57}{0}{0} & \SetCell{bg=\HeatCellBg{64}}\CellText{0.64}{0}{0} & \SetCell{bg=\HeatCellBg{61}}\CellText{0.61}{0}{0} & \SetCell{bg=\HeatCellBg{61}}\CellText{0.61}{0}{0} & \SetCell{bg=\HeatCellBg{66}}\CellText{0.66}{0}{0} & \SetCell{bg=\HeatCellBg{58}}\CellText{0.58}{0}{0} & \SetCell{bg=\HeatCellBg{52}}\CellText{0.52}{0}{0} & \SetCell{bg=\HeatCellBg{56}}\CellText{0.56}{0}{0} & \SetCell{bg=\HeatCellBg{63}}\CellText{0.63}{0}{0} & \SetCell{bg=\HeatCellBg{67}}\CellText{0.67}{0}{0} & \SetCell{bg=\HeatCellBg{63}}\CellText{0.63}{0}{0} & \SetCell{bg=\HeatCellBg{63}}\CellText{0.63}{0}{0} & \SetCell{bg=\HeatCellBg{68}}\CellText{0.68}{0}{0} & \SetCell{bg=\HeatCellBg{64}}\CellText{0.64}{0}{0} & \MeanCell{0.63}{0.11}{0} \\
        6a & \SetCell{bg=\HeatCellBg{11}}\CellText{0.11}{0}{0} & \SetCell{bg=\HeatCellBg{57}}\CellText{0.57}{0}{1} & \SetCell{bg=\HeatCellBg{51}}\CellText{0.51}{0}{0} & \SetCell{bg=\HeatCellBg{53}}\CellText{0.53}{0}{0} & \SetCell{bg=\HeatCellBg{49}}\CellText{0.49}{0}{0} & \SetCell{bg=\HeatCellBg{53}}\CellText{0.53}{0}{0} & \SetCell{bg=\HeatCellBg{42}}\CellText{0.42}{0}{0} & \SetCell{bg=\HeatCellBg{51}}\CellText{0.51}{0}{0} & \SetCell{bg=\HeatCellBg{52}}\CellText{0.52}{0}{0} & \SetCell{bg=\HeatCellBg{50}}\CellText{0.50}{0}{0} & \SetCell{bg=\HeatCellBg{48}}\CellText{0.48}{0}{0} & \SetCell{bg=\HeatCellBg{51}}\CellText{0.51}{0}{0} & \SetCell{bg=\HeatCellBg{52}}\CellText{0.52}{0}{0} & \SetCell{bg=\HeatCellBg{53}}\CellText{0.53}{0}{0} & \SetCell{bg=\HeatCellBg{51}}\CellText{0.51}{0}{0} & \SetCell{bg=\HeatCellBg{55}}\CellText{0.55}{0}{0} & \SetCell{bg=\HeatCellBg{51}}\CellText{0.51}{0}{0} & \SetCell{bg=\HeatCellBg{54}}\CellText{0.54}{0}{0} & \SetCell{bg=\HeatCellBg{54}}\CellText{0.54}{0}{0} & \MeanCell{0.49}{0.10}{0} \\
        Mean $\pm$ Std & \MeanCell{0.52}{0.34}{0} & \MeanCell{0.83}{0.17}{0} & \MeanCell{0.87}{0.16}{1} & \MeanCell{0.76}{0.14}{0} & \MeanCell{0.67}{0.11}{0} & \MeanCell{0.66}{0.11}{0} & \MeanCell{0.69}{0.16}{0} & \MeanCell{0.65}{0.10}{0} & \MeanCell{0.61}{0.05}{0} & \MeanCell{0.63}{0.07}{0} & \MeanCell{0.61}{0.08}{0} & \MeanCell{0.59}{0.09}{0} & \MeanCell{0.64}{0.11}{0} & \MeanCell{0.68}{0.10}{0} & \MeanCell{0.66}{0.07}{0} & \MeanCell{0.65}{0.05}{0} & \MeanCell{0.67}{0.10}{0} & \MeanCell{0.70}{0.09}{0} & \MeanCell{0.70}{0.11}{0} & \MeanCell{0.67}{0.15}{0} \\
}

\ClusterKTable{8}{\glsfmtshort{as} --- \SpaceInvadersName.}{as_space_invaders}{
        1 & \SetCell{bg=\HeatCellBg{40}}\CellText{0.40}{0}{0} & \SetCell{bg=\HeatCellBg{14}}\CellText{0.14}{0}{0} & \SetCell{bg=\HeatCellBg{8}}\CellText{0.08}{0}{0} & \SetCell{bg=\HeatCellBg{13}}\CellText{0.13}{0}{0} & \SetCell{bg=\HeatCellBg{14}}\CellText{0.14}{0}{0} & \SetCell{bg=\HeatCellBg{22}}\CellText{0.22}{0}{0} & \SetCell{bg=\HeatCellBg{41}}\CellText{0.41}{0}{1} & \SetCell{bg=\HeatCellBg{36}}\CellText{0.36}{0}{0} & \SetCell{bg=\HeatCellBg{38}}\CellText{0.38}{0}{0} & \SetCell{bg=\HeatCellBg{37}}\CellText{0.37}{0}{0} & \SetCell{bg=\HeatCellBg{27}}\CellText{0.27}{0}{0} & \SetCell{bg=\HeatCellBg{35}}\CellText{0.35}{0}{0} & \SetCell{bg=\HeatCellBg{29}}\CellText{0.29}{0}{0} & \SetCell{bg=\HeatCellBg{30}}\CellText{0.30}{0}{0} & \SetCell{bg=\HeatCellBg{35}}\CellText{0.35}{0}{0} & \SetCell{bg=\HeatCellBg{26}}\CellText{0.26}{0}{0} & \SetCell{bg=\HeatCellBg{32}}\CellText{0.32}{0}{0} & \SetCell{bg=\HeatCellBg{28}}\CellText{0.28}{0}{0} & \SetCell{bg=\HeatCellBg{29}}\CellText{0.29}{0}{0} & \MeanCell{0.28}{0.10}{0} \\
        2 & \SetCell{bg=\HeatCellBg{83}}\CellText{0.83}{0}{1} & \SetCell{bg=\HeatCellBg{20}}\CellText{0.20}{0}{0} & \SetCell{bg=\HeatCellBg{22}}\CellText{0.22}{0}{0} & \SetCell{bg=\HeatCellBg{49}}\CellText{0.49}{0}{0} & \SetCell{bg=\HeatCellBg{50}}\CellText{0.50}{0}{0} & \SetCell{bg=\HeatCellBg{37}}\CellText{0.37}{0}{0} & \SetCell{bg=\HeatCellBg{42}}\CellText{0.42}{0}{0} & \SetCell{bg=\HeatCellBg{39}}\CellText{0.39}{0}{0} & \SetCell{bg=\HeatCellBg{54}}\CellText{0.54}{0}{0} & \SetCell{bg=\HeatCellBg{43}}\CellText{0.43}{0}{0} & \SetCell{bg=\HeatCellBg{46}}\CellText{0.46}{0}{0} & \SetCell{bg=\HeatCellBg{34}}\CellText{0.34}{0}{0} & \SetCell{bg=\HeatCellBg{43}}\CellText{0.43}{0}{0} & \SetCell{bg=\HeatCellBg{49}}\CellText{0.49}{0}{0} & \SetCell{bg=\HeatCellBg{48}}\CellText{0.48}{0}{0} & \SetCell{bg=\HeatCellBg{47}}\CellText{0.47}{0}{0} & \SetCell{bg=\HeatCellBg{47}}\CellText{0.47}{0}{0} & \SetCell{bg=\HeatCellBg{48}}\CellText{0.48}{0}{0} & \SetCell{bg=\HeatCellBg{47}}\CellText{0.47}{0}{0} & \MeanCell{0.45}{0.13}{0} \\
        3 & \SetCell{bg=\HeatCellBg{83}}\CellText{0.83}{0}{1} & \SetCell{bg=\HeatCellBg{27}}\CellText{0.27}{0}{0} & \SetCell{bg=\HeatCellBg{52}}\CellText{0.52}{0}{0} & \SetCell{bg=\HeatCellBg{57}}\CellText{0.57}{0}{0} & \SetCell{bg=\HeatCellBg{48}}\CellText{0.48}{0}{0} & \SetCell{bg=\HeatCellBg{55}}\CellText{0.55}{0}{0} & \SetCell{bg=\HeatCellBg{61}}\CellText{0.61}{0}{0} & \SetCell{bg=\HeatCellBg{51}}\CellText{0.51}{0}{0} & \SetCell{bg=\HeatCellBg{60}}\CellText{0.60}{0}{0} & \SetCell{bg=\HeatCellBg{55}}\CellText{0.55}{0}{0} & \SetCell{bg=\HeatCellBg{52}}\CellText{0.52}{0}{0} & \SetCell{bg=\HeatCellBg{47}}\CellText{0.47}{0}{0} & \SetCell{bg=\HeatCellBg{59}}\CellText{0.59}{0}{0} & \SetCell{bg=\HeatCellBg{55}}\CellText{0.55}{0}{0} & \SetCell{bg=\HeatCellBg{51}}\CellText{0.51}{0}{0} & \SetCell{bg=\HeatCellBg{53}}\CellText{0.53}{0}{0} & \SetCell{bg=\HeatCellBg{56}}\CellText{0.56}{0}{0} & \SetCell{bg=\HeatCellBg{57}}\CellText{0.57}{0}{0} & \SetCell{bg=\HeatCellBg{44}}\CellText{0.44}{0}{0} & \MeanCell{0.54}{0.10}{0} \\
        4 & \SetCell{bg=\HeatCellBg{94}}\CellText{0.94}{0}{1} & \SetCell{bg=\HeatCellBg{80}}\CellText{0.80}{0}{0} & \SetCell{bg=\HeatCellBg{69}}\CellText{0.69}{0}{0} & \SetCell{bg=\HeatCellBg{71}}\CellText{0.71}{0}{0} & \SetCell{bg=\HeatCellBg{71}}\CellText{0.71}{0}{0} & \SetCell{bg=\HeatCellBg{75}}\CellText{0.75}{1}{0} & \SetCell{bg=\HeatCellBg{74}}\CellText{0.74}{0}{0} & \SetCell{bg=\HeatCellBg{74}}\CellText{0.74}{1}{0} & \SetCell{bg=\HeatCellBg{75}}\CellText{0.75}{1}{0} & \SetCell{bg=\HeatCellBg{69}}\CellText{0.69}{0}{0} & \SetCell{bg=\HeatCellBg{67}}\CellText{0.67}{0}{0} & \SetCell{bg=\HeatCellBg{68}}\CellText{0.68}{0}{0} & \SetCell{bg=\HeatCellBg{71}}\CellText{0.71}{0}{0} & \SetCell{bg=\HeatCellBg{68}}\CellText{0.68}{0}{0} & \SetCell{bg=\HeatCellBg{68}}\CellText{0.68}{0}{0} & \SetCell{bg=\HeatCellBg{70}}\CellText{0.70}{0}{0} & \SetCell{bg=\HeatCellBg{69}}\CellText{0.69}{0}{0} & \SetCell{bg=\HeatCellBg{74}}\CellText{0.74}{1}{0} & \SetCell{bg=\HeatCellBg{69}}\CellText{0.69}{0}{0} & \MeanCell{0.72}{0.06}{0} \\
        5 & \SetCell{bg=\HeatCellBg{94}}\CellText{0.94}{0}{1} & \SetCell{bg=\HeatCellBg{80}}\CellText{0.80}{0}{0} & \SetCell{bg=\HeatCellBg{69}}\CellText{0.69}{0}{0} & \SetCell{bg=\HeatCellBg{71}}\CellText{0.71}{0}{0} & \SetCell{bg=\HeatCellBg{71}}\CellText{0.71}{0}{0} & \SetCell{bg=\HeatCellBg{75}}\CellText{0.75}{1}{0} & \SetCell{bg=\HeatCellBg{74}}\CellText{0.74}{0}{0} & \SetCell{bg=\HeatCellBg{74}}\CellText{0.74}{1}{0} & \SetCell{bg=\HeatCellBg{75}}\CellText{0.75}{1}{0} & \SetCell{bg=\HeatCellBg{69}}\CellText{0.69}{0}{0} & \SetCell{bg=\HeatCellBg{67}}\CellText{0.67}{0}{0} & \SetCell{bg=\HeatCellBg{68}}\CellText{0.68}{0}{0} & \SetCell{bg=\HeatCellBg{71}}\CellText{0.71}{0}{0} & \SetCell{bg=\HeatCellBg{68}}\CellText{0.68}{0}{0} & \SetCell{bg=\HeatCellBg{68}}\CellText{0.68}{0}{0} & \SetCell{bg=\HeatCellBg{70}}\CellText{0.70}{0}{0} & \SetCell{bg=\HeatCellBg{69}}\CellText{0.69}{0}{0} & \SetCell{bg=\HeatCellBg{74}}\CellText{0.74}{1}{0} & \SetCell{bg=\HeatCellBg{69}}\CellText{0.69}{0}{0} & \MeanCell{0.72}{0.06}{0} \\
        6a & \SetCell{bg=\HeatCellBg{96}}\CellText{0.96}{1}{1} & \SetCell{bg=\HeatCellBg{83}}\CellText{0.83}{1}{0} & \SetCell{bg=\HeatCellBg{77}}\CellText{0.77}{1}{0} & \SetCell{bg=\HeatCellBg{72}}\CellText{0.72}{1}{0} & \SetCell{bg=\HeatCellBg{73}}\CellText{0.73}{1}{0} & \SetCell{bg=\HeatCellBg{73}}\CellText{0.73}{0}{0} & \SetCell{bg=\HeatCellBg{77}}\CellText{0.77}{1}{0} & \SetCell{bg=\HeatCellBg{72}}\CellText{0.72}{0}{0} & \SetCell{bg=\HeatCellBg{74}}\CellText{0.74}{0}{0} & \SetCell{bg=\HeatCellBg{73}}\CellText{0.73}{1}{0} & \SetCell{bg=\HeatCellBg{69}}\CellText{0.69}{1}{0} & \SetCell{bg=\HeatCellBg{72}}\CellText{0.72}{1}{0} & \SetCell{bg=\HeatCellBg{76}}\CellText{0.76}{1}{0} & \SetCell{bg=\HeatCellBg{72}}\CellText{0.72}{1}{0} & \SetCell{bg=\HeatCellBg{70}}\CellText{0.70}{1}{0} & \SetCell{bg=\HeatCellBg{71}}\CellText{0.71}{1}{0} & \SetCell{bg=\HeatCellBg{70}}\CellText{0.70}{1}{0} & \SetCell{bg=\HeatCellBg{73}}\CellText{0.73}{0}{0} & \SetCell{bg=\HeatCellBg{72}}\CellText{0.72}{1}{0} & \MeanCell{0.74}{0.06}{1} \\
        Mean $\pm$ Std & \MeanCell{0.82}{0.19}{1} & \MeanCell{0.51}{0.31}{0} & \MeanCell{0.49}{0.26}{0} & \MeanCell{0.55}{0.21}{0} & \MeanCell{0.54}{0.21}{0} & \MeanCell{0.56}{0.21}{0} & \MeanCell{0.61}{0.15}{0} & \MeanCell{0.58}{0.16}{0} & \MeanCell{0.63}{0.14}{0} & \MeanCell{0.58}{0.14}{0} & \MeanCell{0.55}{0.15}{0} & \MeanCell{0.54}{0.16}{0} & \MeanCell{0.58}{0.17}{0} & \MeanCell{0.57}{0.15}{0} & \MeanCell{0.57}{0.13}{0} & \MeanCell{0.56}{0.16}{0} & \MeanCell{0.57}{0.14}{0} & \MeanCell{0.59}{0.17}{0} & \MeanCell{0.55}{0.16}{0} & \MeanCell{0.58}{0.19}{0} \\
}

\ClusterKTable{8}{\glsfmtshort{acu} --- \SpaceInvadersName.}{acu_space_invaders}{
        1 & \SetCell{bg=\HeatCellBg{-16}}\CellText{-0.16}{0}{0} & \SetCell{bg=\HeatCellBg{14}}\CellText{0.14}{0}{0} & \SetCell{bg=\HeatCellBg{4}}\CellText{0.04}{0}{0} & \SetCell{bg=\HeatCellBg{8}}\CellText{0.08}{0}{0} & \SetCell{bg=\HeatCellBg{-9}}\CellText{-0.09}{0}{0} & \SetCell{bg=\HeatCellBg{-7}}\CellText{-0.07}{0}{0} & \SetCell{bg=\HeatCellBg{21}}\CellText{0.21}{0}{1} & \SetCell{bg=\HeatCellBg{3}}\CellText{0.03}{0}{0} & \SetCell{bg=\HeatCellBg{3}}\CellText{0.03}{0}{0} & \SetCell{bg=\HeatCellBg{-3}}\CellText{-0.03}{0}{0} & \SetCell{bg=\HeatCellBg{-11}}\CellText{-0.11}{0}{0} & \SetCell{bg=\HeatCellBg{5}}\CellText{0.05}{0}{0} & \SetCell{bg=\HeatCellBg{12}}\CellText{0.12}{0}{0} & \SetCell{bg=\HeatCellBg{6}}\CellText{0.06}{0}{0} & \SetCell{bg=\HeatCellBg{-1}}\CellText{-0.01}{0}{0} & \SetCell{bg=\HeatCellBg{-7}}\CellText{-0.07}{0}{0} & \SetCell{bg=\HeatCellBg{15}}\CellText{0.15}{0}{0} & \SetCell{bg=\HeatCellBg{10}}\CellText{0.10}{0}{0} & \SetCell{bg=\HeatCellBg{12}}\CellText{0.12}{0}{0} & \MeanCell{0.03}{0.10}{0} \\
        2 & \SetCell{bg=\HeatCellBg{83}}\CellText{0.83}{1}{1} & \SetCell{bg=\HeatCellBg{19}}\CellText{0.19}{0}{0} & \SetCell{bg=\HeatCellBg{20}}\CellText{0.20}{0}{0} & \SetCell{bg=\HeatCellBg{41}}\CellText{0.41}{1}{0} & \SetCell{bg=\HeatCellBg{31}}\CellText{0.31}{0}{0} & \SetCell{bg=\HeatCellBg{22}}\CellText{0.22}{0}{0} & \SetCell{bg=\HeatCellBg{37}}\CellText{0.37}{0}{0} & \SetCell{bg=\HeatCellBg{24}}\CellText{0.24}{0}{0} & \SetCell{bg=\HeatCellBg{21}}\CellText{0.21}{0}{0} & \SetCell{bg=\HeatCellBg{17}}\CellText{0.17}{0}{0} & \SetCell{bg=\HeatCellBg{20}}\CellText{0.20}{0}{0} & \SetCell{bg=\HeatCellBg{5}}\CellText{0.05}{0}{0} & \SetCell{bg=\HeatCellBg{6}}\CellText{0.06}{0}{0} & \SetCell{bg=\HeatCellBg{31}}\CellText{0.31}{1}{0} & \SetCell{bg=\HeatCellBg{20}}\CellText{0.20}{0}{0} & \SetCell{bg=\HeatCellBg{18}}\CellText{0.18}{0}{0} & \SetCell{bg=\HeatCellBg{14}}\CellText{0.14}{0}{0} & \SetCell{bg=\HeatCellBg{26}}\CellText{0.26}{0}{0} & \SetCell{bg=\HeatCellBg{28}}\CellText{0.28}{0}{0} & \MeanCell{0.25}{0.16}{0} \\
        3 & \SetCell{bg=\HeatCellBg{82}}\CellText{0.82}{0}{1} & \SetCell{bg=\HeatCellBg{26}}\CellText{0.26}{0}{0} & \SetCell{bg=\HeatCellBg{44}}\CellText{0.44}{0}{0} & \SetCell{bg=\HeatCellBg{35}}\CellText{0.35}{0}{0} & \SetCell{bg=\HeatCellBg{22}}\CellText{0.22}{0}{0} & \SetCell{bg=\HeatCellBg{27}}\CellText{0.27}{0}{0} & \SetCell{bg=\HeatCellBg{29}}\CellText{0.29}{0}{0} & \SetCell{bg=\HeatCellBg{18}}\CellText{0.18}{0}{0} & \SetCell{bg=\HeatCellBg{23}}\CellText{0.23}{0}{0} & \SetCell{bg=\HeatCellBg{15}}\CellText{0.15}{0}{0} & \SetCell{bg=\HeatCellBg{14}}\CellText{0.14}{0}{0} & \SetCell{bg=\HeatCellBg{8}}\CellText{0.08}{0}{0} & \SetCell{bg=\HeatCellBg{33}}\CellText{0.33}{1}{0} & \SetCell{bg=\HeatCellBg{27}}\CellText{0.27}{0}{0} & \SetCell{bg=\HeatCellBg{22}}\CellText{0.22}{0}{0} & \SetCell{bg=\HeatCellBg{23}}\CellText{0.23}{0}{0} & \SetCell{bg=\HeatCellBg{33}}\CellText{0.33}{1}{0} & \SetCell{bg=\HeatCellBg{28}}\CellText{0.28}{0}{0} & \SetCell{bg=\HeatCellBg{20}}\CellText{0.20}{0}{0} & \MeanCell{0.28}{0.15}{0} \\
        4 & \SetCell{bg=\HeatCellBg{26}}\CellText{0.26}{0}{0} & \SetCell{bg=\HeatCellBg{53}}\CellText{0.53}{1}{0} & \SetCell{bg=\HeatCellBg{61}}\CellText{0.61}{1}{1} & \SetCell{bg=\HeatCellBg{40}}\CellText{0.40}{0}{0} & \SetCell{bg=\HeatCellBg{32}}\CellText{0.32}{1}{0} & \SetCell{bg=\HeatCellBg{32}}\CellText{0.32}{1}{0} & \SetCell{bg=\HeatCellBg{38}}\CellText{0.38}{1}{0} & \SetCell{bg=\HeatCellBg{36}}\CellText{0.36}{1}{0} & \SetCell{bg=\HeatCellBg{35}}\CellText{0.35}{1}{0} & \SetCell{bg=\HeatCellBg{35}}\CellText{0.35}{1}{0} & \SetCell{bg=\HeatCellBg{26}}\CellText{0.26}{1}{0} & \SetCell{bg=\HeatCellBg{20}}\CellText{0.20}{0}{0} & \SetCell{bg=\HeatCellBg{27}}\CellText{0.27}{0}{0} & \SetCell{bg=\HeatCellBg{31}}\CellText{0.31}{1}{0} & \SetCell{bg=\HeatCellBg{36}}\CellText{0.36}{1}{0} & \SetCell{bg=\HeatCellBg{33}}\CellText{0.33}{1}{0} & \SetCell{bg=\HeatCellBg{32}}\CellText{0.32}{0}{0} & \SetCell{bg=\HeatCellBg{43}}\CellText{0.43}{1}{0} & \SetCell{bg=\HeatCellBg{34}}\CellText{0.34}{1}{0} & \MeanCell{0.35}{0.09}{1} \\
        5 & \SetCell{bg=\HeatCellBg{26}}\CellText{0.26}{0}{0} & \SetCell{bg=\HeatCellBg{53}}\CellText{0.53}{1}{0} & \SetCell{bg=\HeatCellBg{61}}\CellText{0.61}{1}{1} & \SetCell{bg=\HeatCellBg{40}}\CellText{0.40}{0}{0} & \SetCell{bg=\HeatCellBg{32}}\CellText{0.32}{1}{0} & \SetCell{bg=\HeatCellBg{32}}\CellText{0.32}{1}{0} & \SetCell{bg=\HeatCellBg{38}}\CellText{0.38}{1}{0} & \SetCell{bg=\HeatCellBg{36}}\CellText{0.36}{1}{0} & \SetCell{bg=\HeatCellBg{35}}\CellText{0.35}{1}{0} & \SetCell{bg=\HeatCellBg{35}}\CellText{0.35}{1}{0} & \SetCell{bg=\HeatCellBg{26}}\CellText{0.26}{1}{0} & \SetCell{bg=\HeatCellBg{20}}\CellText{0.20}{0}{0} & \SetCell{bg=\HeatCellBg{27}}\CellText{0.27}{0}{0} & \SetCell{bg=\HeatCellBg{31}}\CellText{0.31}{1}{0} & \SetCell{bg=\HeatCellBg{36}}\CellText{0.36}{1}{0} & \SetCell{bg=\HeatCellBg{33}}\CellText{0.33}{1}{0} & \SetCell{bg=\HeatCellBg{32}}\CellText{0.32}{0}{0} & \SetCell{bg=\HeatCellBg{43}}\CellText{0.43}{1}{0} & \SetCell{bg=\HeatCellBg{34}}\CellText{0.34}{1}{0} & \MeanCell{0.35}{0.09}{1} \\
        6a & \SetCell{bg=\HeatCellBg{7}}\CellText{0.07}{0}{0} & \SetCell{bg=\HeatCellBg{39}}\CellText{0.39}{0}{1} & \SetCell{bg=\HeatCellBg{28}}\CellText{0.28}{0}{0} & \SetCell{bg=\HeatCellBg{24}}\CellText{0.24}{0}{0} & \SetCell{bg=\HeatCellBg{22}}\CellText{0.22}{0}{0} & \SetCell{bg=\HeatCellBg{26}}\CellText{0.26}{0}{0} & \SetCell{bg=\HeatCellBg{19}}\CellText{0.19}{0}{0} & \SetCell{bg=\HeatCellBg{23}}\CellText{0.23}{0}{0} & \SetCell{bg=\HeatCellBg{26}}\CellText{0.26}{0}{0} & \SetCell{bg=\HeatCellBg{23}}\CellText{0.23}{0}{0} & \SetCell{bg=\HeatCellBg{17}}\CellText{0.17}{0}{0} & \SetCell{bg=\HeatCellBg{24}}\CellText{0.24}{1}{0} & \SetCell{bg=\HeatCellBg{28}}\CellText{0.28}{0}{0} & \SetCell{bg=\HeatCellBg{25}}\CellText{0.25}{0}{0} & \SetCell{bg=\HeatCellBg{21}}\CellText{0.21}{0}{0} & \SetCell{bg=\HeatCellBg{26}}\CellText{0.26}{0}{0} & \SetCell{bg=\HeatCellBg{22}}\CellText{0.22}{0}{0} & \SetCell{bg=\HeatCellBg{27}}\CellText{0.27}{0}{0} & \SetCell{bg=\HeatCellBg{26}}\CellText{0.26}{0}{0} & \MeanCell{0.24}{0.06}{0} \\
        Mean $\pm$ Std & \MeanCell{0.35}{0.37}{0} & \MeanCell{0.34}{0.15}{0} & \MeanCell{0.36}{0.21}{1} & \MeanCell{0.31}{0.12}{0} & \MeanCell{0.22}{0.14}{0} & \MeanCell{0.22}{0.13}{0} & \MeanCell{0.30}{0.08}{0} & \MeanCell{0.23}{0.11}{0} & \MeanCell{0.24}{0.11}{0} & \MeanCell{0.20}{0.13}{0} & \MeanCell{0.15}{0.13}{0} & \MeanCell{0.14}{0.08}{0} & \MeanCell{0.22}{0.10}{0} & \MeanCell{0.25}{0.09}{0} & \MeanCell{0.22}{0.12}{0} & \MeanCell{0.21}{0.14}{0} & \MeanCell{0.25}{0.08}{0} & \MeanCell{0.29}{0.11}{0} & \MeanCell{0.26}{0.08}{0} & \MeanCell{0.25}{0.16}{0} \\
}

\ClusterKTable{8}[Observation baseline --- \SpaceInvadersName.]{Observation baseline ($\Similarity{obs}{}$) --- \SpaceInvadersName.}{obs_space_invaders}{
        1 & \SetCell{bg=\HeatCellBg{43}}\CellText{0.43}{1}{0} & \SetCell{bg=\HeatCellBg{82}}\CellText{0.82}{1}{0} & \SetCell{bg=\HeatCellBg{92}}\CellText{0.92}{1}{1} & \SetCell{bg=\HeatCellBg{87}}\CellText{0.87}{1}{0} & \SetCell{bg=\HeatCellBg{84}}\CellText{0.84}{1}{0} & \SetCell{bg=\HeatCellBg{76}}\CellText{0.76}{1}{0} & \SetCell{bg=\HeatCellBg{53}}\CellText{0.53}{1}{0} & \SetCell{bg=\HeatCellBg{56}}\CellText{0.56}{1}{0} & \SetCell{bg=\HeatCellBg{60}}\CellText{0.60}{1}{0} & \SetCell{bg=\HeatCellBg{58}}\CellText{0.58}{1}{0} & \SetCell{bg=\HeatCellBg{70}}\CellText{0.70}{1}{0} & \SetCell{bg=\HeatCellBg{59}}\CellText{0.59}{0}{0} & \SetCell{bg=\HeatCellBg{71}}\CellText{0.71}{1}{0} & \SetCell{bg=\HeatCellBg{70}}\CellText{0.70}{1}{0} & \SetCell{bg=\HeatCellBg{58}}\CellText{0.58}{1}{0} & \SetCell{bg=\HeatCellBg{63}}\CellText{0.63}{1}{0} & \SetCell{bg=\HeatCellBg{58}}\CellText{0.58}{1}{0} & \SetCell{bg=\HeatCellBg{57}}\CellText{0.57}{1}{0} & \SetCell{bg=\HeatCellBg{57}}\CellText{0.57}{1}{0} & \MeanCell{0.66}{0.13}{1} \\
        2 & \SetCell{bg=\HeatCellBg{17}}\CellText{0.17}{0}{0} & \SetCell{bg=\HeatCellBg{80}}\CellText{0.80}{0}{1} & \SetCell{bg=\HeatCellBg{78}}\CellText{0.78}{0}{0} & \SetCell{bg=\HeatCellBg{46}}\CellText{0.46}{0}{0} & \SetCell{bg=\HeatCellBg{41}}\CellText{0.41}{0}{0} & \SetCell{bg=\HeatCellBg{43}}\CellText{0.43}{0}{0} & \SetCell{bg=\HeatCellBg{41}}\CellText{0.41}{0}{0} & \SetCell{bg=\HeatCellBg{45}}\CellText{0.45}{0}{0} & \SetCell{bg=\HeatCellBg{41}}\CellText{0.41}{0}{0} & \SetCell{bg=\HeatCellBg{57}}\CellText{0.57}{0}{0} & \SetCell{bg=\HeatCellBg{54}}\CellText{0.54}{0}{0} & \SetCell{bg=\HeatCellBg{66}}\CellText{0.66}{1}{0} & \SetCell{bg=\HeatCellBg{47}}\CellText{0.47}{0}{0} & \SetCell{bg=\HeatCellBg{51}}\CellText{0.51}{0}{0} & \SetCell{bg=\HeatCellBg{46}}\CellText{0.46}{0}{0} & \SetCell{bg=\HeatCellBg{49}}\CellText{0.49}{0}{0} & \SetCell{bg=\HeatCellBg{48}}\CellText{0.48}{0}{0} & \SetCell{bg=\HeatCellBg{45}}\CellText{0.45}{0}{0} & \SetCell{bg=\HeatCellBg{45}}\CellText{0.45}{0}{0} & \MeanCell{0.49}{0.14}{0} \\
        3 & \SetCell{bg=\HeatCellBg{17}}\CellText{0.17}{0}{0} & \SetCell{bg=\HeatCellBg{73}}\CellText{0.73}{0}{1} & \SetCell{bg=\HeatCellBg{48}}\CellText{0.48}{0}{0} & \SetCell{bg=\HeatCellBg{43}}\CellText{0.43}{0}{0} & \SetCell{bg=\HeatCellBg{42}}\CellText{0.42}{0}{0} & \SetCell{bg=\HeatCellBg{39}}\CellText{0.39}{0}{0} & \SetCell{bg=\HeatCellBg{33}}\CellText{0.33}{0}{0} & \SetCell{bg=\HeatCellBg{35}}\CellText{0.35}{0}{0} & \SetCell{bg=\HeatCellBg{34}}\CellText{0.34}{0}{0} & \SetCell{bg=\HeatCellBg{37}}\CellText{0.37}{0}{0} & \SetCell{bg=\HeatCellBg{46}}\CellText{0.46}{0}{0} & \SetCell{bg=\HeatCellBg{50}}\CellText{0.50}{0}{0} & \SetCell{bg=\HeatCellBg{37}}\CellText{0.37}{0}{0} & \SetCell{bg=\HeatCellBg{41}}\CellText{0.41}{0}{0} & \SetCell{bg=\HeatCellBg{36}}\CellText{0.36}{0}{0} & \SetCell{bg=\HeatCellBg{37}}\CellText{0.37}{0}{0} & \SetCell{bg=\HeatCellBg{41}}\CellText{0.41}{0}{0} & \SetCell{bg=\HeatCellBg{37}}\CellText{0.37}{0}{0} & \SetCell{bg=\HeatCellBg{41}}\CellText{0.41}{0}{0} & \MeanCell{0.40}{0.10}{0} \\
        4 & \SetCell{bg=\HeatCellBg{6}}\CellText{0.06}{0}{0} & \SetCell{bg=\HeatCellBg{20}}\CellText{0.20}{0}{0} & \SetCell{bg=\HeatCellBg{31}}\CellText{0.31}{0}{0} & \SetCell{bg=\HeatCellBg{29}}\CellText{0.29}{0}{0} & \SetCell{bg=\HeatCellBg{28}}\CellText{0.28}{0}{0} & \SetCell{bg=\HeatCellBg{23}}\CellText{0.23}{0}{0} & \SetCell{bg=\HeatCellBg{21}}\CellText{0.21}{0}{0} & \SetCell{bg=\HeatCellBg{24}}\CellText{0.24}{0}{0} & \SetCell{bg=\HeatCellBg{25}}\CellText{0.25}{0}{0} & \SetCell{bg=\HeatCellBg{30}}\CellText{0.30}{0}{0} & \SetCell{bg=\HeatCellBg{33}}\CellText{0.33}{0}{1} & \SetCell{bg=\HeatCellBg{32}}\CellText{0.32}{0}{0} & \SetCell{bg=\HeatCellBg{28}}\CellText{0.28}{0}{0} & \SetCell{bg=\HeatCellBg{30}}\CellText{0.30}{0}{0} & \SetCell{bg=\HeatCellBg{32}}\CellText{0.32}{0}{0} & \SetCell{bg=\HeatCellBg{28}}\CellText{0.28}{0}{0} & \SetCell{bg=\HeatCellBg{31}}\CellText{0.31}{0}{0} & \SetCell{bg=\HeatCellBg{26}}\CellText{0.26}{0}{0} & \SetCell{bg=\HeatCellBg{30}}\CellText{0.30}{0}{0} & \MeanCell{0.27}{0.06}{0} \\
        5 & \SetCell{bg=\HeatCellBg{6}}\CellText{0.06}{0}{0} & \SetCell{bg=\HeatCellBg{20}}\CellText{0.20}{0}{0} & \SetCell{bg=\HeatCellBg{31}}\CellText{0.31}{0}{0} & \SetCell{bg=\HeatCellBg{29}}\CellText{0.29}{0}{0} & \SetCell{bg=\HeatCellBg{28}}\CellText{0.28}{0}{0} & \SetCell{bg=\HeatCellBg{23}}\CellText{0.23}{0}{0} & \SetCell{bg=\HeatCellBg{21}}\CellText{0.21}{0}{0} & \SetCell{bg=\HeatCellBg{24}}\CellText{0.24}{0}{0} & \SetCell{bg=\HeatCellBg{25}}\CellText{0.25}{0}{0} & \SetCell{bg=\HeatCellBg{30}}\CellText{0.30}{0}{0} & \SetCell{bg=\HeatCellBg{33}}\CellText{0.33}{0}{1} & \SetCell{bg=\HeatCellBg{32}}\CellText{0.32}{0}{0} & \SetCell{bg=\HeatCellBg{28}}\CellText{0.28}{0}{0} & \SetCell{bg=\HeatCellBg{30}}\CellText{0.30}{0}{0} & \SetCell{bg=\HeatCellBg{32}}\CellText{0.32}{0}{0} & \SetCell{bg=\HeatCellBg{28}}\CellText{0.28}{0}{0} & \SetCell{bg=\HeatCellBg{31}}\CellText{0.31}{0}{0} & \SetCell{bg=\HeatCellBg{26}}\CellText{0.26}{0}{0} & \SetCell{bg=\HeatCellBg{30}}\CellText{0.30}{0}{0} & \MeanCell{0.27}{0.06}{0} \\
        6a & \SetCell{bg=\HeatCellBg{1}}\CellText{0.01}{0}{0} & \SetCell{bg=\HeatCellBg{17}}\CellText{0.17}{0}{0} & \SetCell{bg=\HeatCellBg{23}}\CellText{0.23}{0}{0} & \SetCell{bg=\HeatCellBg{28}}\CellText{0.28}{0}{0} & \SetCell{bg=\HeatCellBg{25}}\CellText{0.25}{0}{0} & \SetCell{bg=\HeatCellBg{27}}\CellText{0.27}{0}{0} & \SetCell{bg=\HeatCellBg{20}}\CellText{0.20}{0}{0} & \SetCell{bg=\HeatCellBg{25}}\CellText{0.25}{0}{0} & \SetCell{bg=\HeatCellBg{25}}\CellText{0.25}{0}{0} & \SetCell{bg=\HeatCellBg{26}}\CellText{0.26}{0}{0} & \SetCell{bg=\HeatCellBg{31}}\CellText{0.31}{0}{1} & \SetCell{bg=\HeatCellBg{28}}\CellText{0.28}{0}{0} & \SetCell{bg=\HeatCellBg{24}}\CellText{0.24}{0}{0} & \SetCell{bg=\HeatCellBg{28}}\CellText{0.28}{0}{0} & \SetCell{bg=\HeatCellBg{30}}\CellText{0.30}{0}{0} & \SetCell{bg=\HeatCellBg{29}}\CellText{0.29}{0}{0} & \SetCell{bg=\HeatCellBg{30}}\CellText{0.30}{0}{0} & \SetCell{bg=\HeatCellBg{27}}\CellText{0.27}{0}{0} & \SetCell{bg=\HeatCellBg{28}}\CellText{0.28}{0}{0} & \MeanCell{0.25}{0.07}{0} \\
        Mean $\pm$ Std & \MeanCell{0.15}{0.14}{0} & \MeanCell{0.49}{0.30}{0} & \MeanCell{0.51}{0.26}{1} & \MeanCell{0.44}{0.21}{0} & \MeanCell{0.41}{0.20}{0} & \MeanCell{0.39}{0.18}{0} & \MeanCell{0.32}{0.12}{0} & \MeanCell{0.35}{0.12}{0} & \MeanCell{0.35}{0.13}{0} & \MeanCell{0.40}{0.13}{0} & \MeanCell{0.45}{0.14}{0} & \MeanCell{0.45}{0.15}{0} & \MeanCell{0.39}{0.16}{0} & \MeanCell{0.42}{0.15}{0} & \MeanCell{0.39}{0.10}{0} & \MeanCell{0.39}{0.13}{0} & \MeanCell{0.40}{0.10}{0} & \MeanCell{0.36}{0.12}{0} & \MeanCell{0.39}{0.10}{0} & \MeanCell{0.39}{0.18}{0} \\
}

\ClusterKTable{8}[Untrained-network baseline --- \SpaceInvadersName.]{Untrained-network baseline ($\Similarity{untr}{}$) --- \SpaceInvadersName.}{untr_space_invaders}{
        1 & \SetCell{bg=\HeatCellBg{52}}\CellText{0.52}{1}{0} & \SetCell{bg=\HeatCellBg{86}}\CellText{0.86}{1}{0} & \SetCell{bg=\HeatCellBg{89}}\CellText{0.89}{1}{1} & \SetCell{bg=\HeatCellBg{69}}\CellText{0.69}{1}{0} & \SetCell{bg=\HeatCellBg{80}}\CellText{0.80}{1}{0} & \SetCell{bg=\HeatCellBg{74}}\CellText{0.74}{1}{0} & \SetCell{bg=\HeatCellBg{59}}\CellText{0.59}{1}{0} & \SetCell{bg=\HeatCellBg{63}}\CellText{0.63}{1}{0} & \SetCell{bg=\HeatCellBg{61}}\CellText{0.61}{1}{0} & \SetCell{bg=\HeatCellBg{60}}\CellText{0.60}{1}{0} & \SetCell{bg=\HeatCellBg{72}}\CellText{0.72}{1}{0} & \SetCell{bg=\HeatCellBg{64}}\CellText{0.64}{1}{0} & \SetCell{bg=\HeatCellBg{69}}\CellText{0.69}{1}{0} & \SetCell{bg=\HeatCellBg{66}}\CellText{0.66}{1}{0} & \SetCell{bg=\HeatCellBg{65}}\CellText{0.65}{1}{0} & \SetCell{bg=\HeatCellBg{70}}\CellText{0.70}{1}{0} & \SetCell{bg=\HeatCellBg{68}}\CellText{0.68}{1}{0} & \SetCell{bg=\HeatCellBg{72}}\CellText{0.72}{1}{0} & \SetCell{bg=\HeatCellBg{71}}\CellText{0.71}{1}{0} & \MeanCell{0.69}{0.09}{1} \\
        2 & \SetCell{bg=\HeatCellBg{15}}\CellText{0.15}{0}{0} & \SetCell{bg=\HeatCellBg{76}}\CellText{0.76}{0}{1} & \SetCell{bg=\HeatCellBg{74}}\CellText{0.74}{0}{0} & \SetCell{bg=\HeatCellBg{51}}\CellText{0.51}{0}{0} & \SetCell{bg=\HeatCellBg{50}}\CellText{0.50}{0}{0} & \SetCell{bg=\HeatCellBg{63}}\CellText{0.63}{0}{0} & \SetCell{bg=\HeatCellBg{58}}\CellText{0.58}{0}{0} & \SetCell{bg=\HeatCellBg{61}}\CellText{0.61}{0}{0} & \SetCell{bg=\HeatCellBg{46}}\CellText{0.46}{0}{0} & \SetCell{bg=\HeatCellBg{54}}\CellText{0.54}{0}{0} & \SetCell{bg=\HeatCellBg{52}}\CellText{0.52}{0}{0} & \SetCell{bg=\HeatCellBg{49}}\CellText{0.49}{0}{0} & \SetCell{bg=\HeatCellBg{55}}\CellText{0.55}{0}{0} & \SetCell{bg=\HeatCellBg{47}}\CellText{0.47}{0}{0} & \SetCell{bg=\HeatCellBg{51}}\CellText{0.51}{0}{0} & \SetCell{bg=\HeatCellBg{51}}\CellText{0.51}{0}{0} & \SetCell{bg=\HeatCellBg{52}}\CellText{0.52}{0}{0} & \SetCell{bg=\HeatCellBg{52}}\CellText{0.52}{0}{0} & \SetCell{bg=\HeatCellBg{53}}\CellText{0.53}{0}{0} & \MeanCell{0.53}{0.12}{0} \\
        3 & \SetCell{bg=\HeatCellBg{17}}\CellText{0.17}{0}{0} & \SetCell{bg=\HeatCellBg{70}}\CellText{0.70}{0}{1} & \SetCell{bg=\HeatCellBg{42}}\CellText{0.42}{0}{0} & \SetCell{bg=\HeatCellBg{43}}\CellText{0.43}{0}{0} & \SetCell{bg=\HeatCellBg{52}}\CellText{0.52}{0}{0} & \SetCell{bg=\HeatCellBg{45}}\CellText{0.45}{0}{0} & \SetCell{bg=\HeatCellBg{39}}\CellText{0.39}{0}{0} & \SetCell{bg=\HeatCellBg{49}}\CellText{0.49}{0}{0} & \SetCell{bg=\HeatCellBg{40}}\CellText{0.40}{0}{0} & \SetCell{bg=\HeatCellBg{45}}\CellText{0.45}{0}{0} & \SetCell{bg=\HeatCellBg{41}}\CellText{0.41}{0}{0} & \SetCell{bg=\HeatCellBg{49}}\CellText{0.49}{0}{0} & \SetCell{bg=\HeatCellBg{41}}\CellText{0.41}{0}{0} & \SetCell{bg=\HeatCellBg{45}}\CellText{0.45}{0}{0} & \SetCell{bg=\HeatCellBg{49}}\CellText{0.49}{0}{0} & \SetCell{bg=\HeatCellBg{47}}\CellText{0.47}{0}{0} & \SetCell{bg=\HeatCellBg{44}}\CellText{0.44}{0}{0} & \SetCell{bg=\HeatCellBg{43}}\CellText{0.43}{0}{0} & \SetCell{bg=\HeatCellBg{56}}\CellText{0.56}{0}{0} & \MeanCell{0.45}{0.10}{0} \\
        4 & \SetCell{bg=\HeatCellBg{5}}\CellText{0.05}{0}{0} & \SetCell{bg=\HeatCellBg{16}}\CellText{0.16}{0}{0} & \SetCell{bg=\HeatCellBg{26}}\CellText{0.26}{0}{0} & \SetCell{bg=\HeatCellBg{24}}\CellText{0.24}{0}{0} & \SetCell{bg=\HeatCellBg{29}}\CellText{0.29}{0}{0} & \SetCell{bg=\HeatCellBg{25}}\CellText{0.25}{0}{0} & \SetCell{bg=\HeatCellBg{26}}\CellText{0.26}{0}{0} & \SetCell{bg=\HeatCellBg{25}}\CellText{0.25}{0}{0} & \SetCell{bg=\HeatCellBg{25}}\CellText{0.25}{0}{0} & \SetCell{bg=\HeatCellBg{31}}\CellText{0.31}{0}{1} & \SetCell{bg=\HeatCellBg{28}}\CellText{0.28}{0}{0} & \SetCell{bg=\HeatCellBg{28}}\CellText{0.28}{0}{0} & \SetCell{bg=\HeatCellBg{28}}\CellText{0.28}{0}{0} & \SetCell{bg=\HeatCellBg{29}}\CellText{0.29}{0}{0} & \SetCell{bg=\HeatCellBg{31}}\CellText{0.31}{0}{1} & \SetCell{bg=\HeatCellBg{30}}\CellText{0.30}{0}{0} & \SetCell{bg=\HeatCellBg{30}}\CellText{0.30}{0}{0} & \SetCell{bg=\HeatCellBg{25}}\CellText{0.25}{0}{0} & \SetCell{bg=\HeatCellBg{30}}\CellText{0.30}{0}{0} & \MeanCell{0.26}{0.06}{0} \\
        5 & \SetCell{bg=\HeatCellBg{5}}\CellText{0.05}{0}{0} & \SetCell{bg=\HeatCellBg{16}}\CellText{0.16}{0}{0} & \SetCell{bg=\HeatCellBg{26}}\CellText{0.26}{0}{0} & \SetCell{bg=\HeatCellBg{24}}\CellText{0.24}{0}{0} & \SetCell{bg=\HeatCellBg{29}}\CellText{0.29}{0}{0} & \SetCell{bg=\HeatCellBg{25}}\CellText{0.25}{0}{0} & \SetCell{bg=\HeatCellBg{26}}\CellText{0.26}{0}{0} & \SetCell{bg=\HeatCellBg{25}}\CellText{0.25}{0}{0} & \SetCell{bg=\HeatCellBg{25}}\CellText{0.25}{0}{0} & \SetCell{bg=\HeatCellBg{31}}\CellText{0.31}{0}{1} & \SetCell{bg=\HeatCellBg{28}}\CellText{0.28}{0}{0} & \SetCell{bg=\HeatCellBg{28}}\CellText{0.28}{0}{0} & \SetCell{bg=\HeatCellBg{28}}\CellText{0.28}{0}{0} & \SetCell{bg=\HeatCellBg{29}}\CellText{0.29}{0}{0} & \SetCell{bg=\HeatCellBg{31}}\CellText{0.31}{0}{1} & \SetCell{bg=\HeatCellBg{30}}\CellText{0.30}{0}{0} & \SetCell{bg=\HeatCellBg{30}}\CellText{0.30}{0}{0} & \SetCell{bg=\HeatCellBg{25}}\CellText{0.25}{0}{0} & \SetCell{bg=\HeatCellBg{30}}\CellText{0.30}{0}{0} & \MeanCell{0.26}{0.06}{0} \\
        6a & \SetCell{bg=\HeatCellBg{3}}\CellText{0.03}{0}{0} & \SetCell{bg=\HeatCellBg{12}}\CellText{0.12}{0}{0} & \SetCell{bg=\HeatCellBg{16}}\CellText{0.16}{0}{0} & \SetCell{bg=\HeatCellBg{27}}\CellText{0.27}{0}{0} & \SetCell{bg=\HeatCellBg{27}}\CellText{0.27}{0}{0} & \SetCell{bg=\HeatCellBg{26}}\CellText{0.26}{0}{0} & \SetCell{bg=\HeatCellBg{23}}\CellText{0.23}{0}{0} & \SetCell{bg=\HeatCellBg{28}}\CellText{0.28}{0}{1} & \SetCell{bg=\HeatCellBg{26}}\CellText{0.26}{0}{0} & \SetCell{bg=\HeatCellBg{27}}\CellText{0.27}{0}{0} & \SetCell{bg=\HeatCellBg{25}}\CellText{0.25}{0}{0} & \SetCell{bg=\HeatCellBg{21}}\CellText{0.21}{0}{0} & \SetCell{bg=\HeatCellBg{21}}\CellText{0.21}{0}{0} & \SetCell{bg=\HeatCellBg{23}}\CellText{0.23}{0}{0} & \SetCell{bg=\HeatCellBg{26}}\CellText{0.26}{0}{0} & \SetCell{bg=\HeatCellBg{24}}\CellText{0.24}{0}{0} & \SetCell{bg=\HeatCellBg{22}}\CellText{0.22}{0}{0} & \SetCell{bg=\HeatCellBg{23}}\CellText{0.23}{0}{0} & \SetCell{bg=\HeatCellBg{23}}\CellText{0.23}{0}{0} & \MeanCell{0.22}{0.06}{0} \\
        Mean $\pm$ Std & \MeanCell{0.16}{0.17}{0} & \MeanCell{0.46}{0.32}{1} & \MeanCell{0.46}{0.27}{0} & \MeanCell{0.40}{0.17}{0} & \MeanCell{0.45}{0.19}{0} & \MeanCell{0.43}{0.20}{0} & \MeanCell{0.39}{0.15}{0} & \MeanCell{0.42}{0.16}{0} & \MeanCell{0.37}{0.13}{0} & \MeanCell{0.41}{0.13}{0} & \MeanCell{0.41}{0.17}{0} & \MeanCell{0.40}{0.15}{0} & \MeanCell{0.40}{0.17}{0} & \MeanCell{0.40}{0.15}{0} & \MeanCell{0.42}{0.14}{0} & \MeanCell{0.42}{0.16}{0} & \MeanCell{0.41}{0.16}{0} & \MeanCell{0.40}{0.18}{0} & \MeanCell{0.44}{0.17}{0} & \MeanCell{0.40}{0.19}{0} \\
}

\ClusterKTable{8}[Action baseline --- \SpaceInvadersName.]{Action baseline ($\Similarity{act}{}$) --- \SpaceInvadersName.}{act_space_invaders}{
        1 & \SetCell{bg=\HeatCellBg{0}}\CellText{0.00}{1}{0} & \SetCell{bg=\HeatCellBg{0}}\CellText{0.00}{0}{0} & \SetCell{bg=\HeatCellBg{1}}\CellText{0.01}{1}{0} & \SetCell{bg=\HeatCellBg{1}}\CellText{0.01}{1}{0} & \SetCell{bg=\HeatCellBg{1}}\CellText{0.01}{1}{0} & \SetCell{bg=\HeatCellBg{1}}\CellText{0.01}{0}{0} & \SetCell{bg=\HeatCellBg{1}}\CellText{0.01}{0}{0} & \SetCell{bg=\HeatCellBg{1}}\CellText{0.01}{0}{0} & \SetCell{bg=\HeatCellBg{2}}\CellText{0.02}{1}{1} & \SetCell{bg=\HeatCellBg{1}}\CellText{0.01}{0}{0} & \SetCell{bg=\HeatCellBg{2}}\CellText{0.02}{1}{1} & \SetCell{bg=\HeatCellBg{1}}\CellText{0.01}{0}{0} & \SetCell{bg=\HeatCellBg{2}}\CellText{0.02}{0}{1} & \SetCell{bg=\HeatCellBg{2}}\CellText{0.02}{0}{1} & \SetCell{bg=\HeatCellBg{2}}\CellText{0.02}{0}{1} & \SetCell{bg=\HeatCellBg{2}}\CellText{0.02}{0}{1} & \SetCell{bg=\HeatCellBg{2}}\CellText{0.02}{0}{1} & \SetCell{bg=\HeatCellBg{2}}\CellText{0.02}{0}{1} & \SetCell{bg=\HeatCellBg{2}}\CellText{0.02}{0}{1} & \MeanCell{0.01}{0.01}{0} \\
        2 & \SetCell{bg=\HeatCellBg{0}}\CellText{0.00}{1}{0} & \SetCell{bg=\HeatCellBg{0}}\CellText{0.00}{0}{0} & \SetCell{bg=\HeatCellBg{1}}\CellText{0.01}{1}{0} & \SetCell{bg=\HeatCellBg{1}}\CellText{0.01}{1}{0} & \SetCell{bg=\HeatCellBg{1}}\CellText{0.01}{1}{0} & \SetCell{bg=\HeatCellBg{1}}\CellText{0.01}{0}{0} & \SetCell{bg=\HeatCellBg{1}}\CellText{0.01}{0}{0} & \SetCell{bg=\HeatCellBg{1}}\CellText{0.01}{0}{0} & \SetCell{bg=\HeatCellBg{2}}\CellText{0.02}{1}{0} & \SetCell{bg=\HeatCellBg{2}}\CellText{0.02}{1}{0} & \SetCell{bg=\HeatCellBg{2}}\CellText{0.02}{1}{0} & \SetCell{bg=\HeatCellBg{2}}\CellText{0.02}{0}{0} & \SetCell{bg=\HeatCellBg{2}}\CellText{0.02}{0}{0} & \SetCell{bg=\HeatCellBg{2}}\CellText{0.02}{0}{0} & \SetCell{bg=\HeatCellBg{2}}\CellText{0.02}{0}{0} & \SetCell{bg=\HeatCellBg{3}}\CellText{0.03}{1}{1} & \SetCell{bg=\HeatCellBg{3}}\CellText{0.03}{1}{1} & \SetCell{bg=\HeatCellBg{3}}\CellText{0.03}{1}{1} & \SetCell{bg=\HeatCellBg{3}}\CellText{0.03}{1}{1} & \MeanCell{0.02}{0.01}{0} \\
        3 & \SetCell{bg=\HeatCellBg{0}}\CellText{0.00}{1}{0} & \SetCell{bg=\HeatCellBg{0}}\CellText{0.00}{0}{0} & \SetCell{bg=\HeatCellBg{1}}\CellText{0.01}{1}{0} & \SetCell{bg=\HeatCellBg{1}}\CellText{0.01}{1}{0} & \SetCell{bg=\HeatCellBg{1}}\CellText{0.01}{1}{0} & \SetCell{bg=\HeatCellBg{1}}\CellText{0.01}{0}{0} & \SetCell{bg=\HeatCellBg{1}}\CellText{0.01}{0}{0} & \SetCell{bg=\HeatCellBg{1}}\CellText{0.01}{0}{0} & \SetCell{bg=\HeatCellBg{2}}\CellText{0.02}{1}{0} & \SetCell{bg=\HeatCellBg{2}}\CellText{0.02}{1}{0} & \SetCell{bg=\HeatCellBg{2}}\CellText{0.02}{1}{0} & \SetCell{bg=\HeatCellBg{2}}\CellText{0.02}{0}{0} & \SetCell{bg=\HeatCellBg{2}}\CellText{0.02}{0}{0} & \SetCell{bg=\HeatCellBg{2}}\CellText{0.02}{0}{0} & \SetCell{bg=\HeatCellBg{2}}\CellText{0.02}{0}{0} & \SetCell{bg=\HeatCellBg{3}}\CellText{0.03}{1}{1} & \SetCell{bg=\HeatCellBg{3}}\CellText{0.03}{1}{1} & \SetCell{bg=\HeatCellBg{3}}\CellText{0.03}{1}{1} & \SetCell{bg=\HeatCellBg{3}}\CellText{0.03}{1}{1} & \MeanCell{0.02}{0.01}{0} \\
        4 & \SetCell{bg=\HeatCellBg{0}}\CellText{0.00}{1}{0} & \SetCell{bg=\HeatCellBg{0}}\CellText{0.00}{0}{0} & \SetCell{bg=\HeatCellBg{1}}\CellText{0.01}{1}{0} & \SetCell{bg=\HeatCellBg{1}}\CellText{0.01}{1}{0} & \SetCell{bg=\HeatCellBg{1}}\CellText{0.01}{1}{0} & \SetCell{bg=\HeatCellBg{2}}\CellText{0.02}{1}{0} & \SetCell{bg=\HeatCellBg{1}}\CellText{0.01}{0}{0} & \SetCell{bg=\HeatCellBg{1}}\CellText{0.01}{0}{0} & \SetCell{bg=\HeatCellBg{2}}\CellText{0.02}{1}{0} & \SetCell{bg=\HeatCellBg{2}}\CellText{0.02}{1}{0} & \SetCell{bg=\HeatCellBg{2}}\CellText{0.02}{1}{0} & \SetCell{bg=\HeatCellBg{2}}\CellText{0.02}{0}{0} & \SetCell{bg=\HeatCellBg{3}}\CellText{0.03}{1}{1} & \SetCell{bg=\HeatCellBg{3}}\CellText{0.03}{1}{1} & \SetCell{bg=\HeatCellBg{3}}\CellText{0.03}{1}{1} & \SetCell{bg=\HeatCellBg{3}}\CellText{0.03}{1}{1} & \SetCell{bg=\HeatCellBg{3}}\CellText{0.03}{1}{1} & \SetCell{bg=\HeatCellBg{3}}\CellText{0.03}{1}{1} & \SetCell{bg=\HeatCellBg{3}}\CellText{0.03}{1}{1} & \MeanCell{0.02}{0.01}{0} \\
        5 & \SetCell{bg=\HeatCellBg{0}}\CellText{0.00}{1}{0} & \SetCell{bg=\HeatCellBg{0}}\CellText{0.00}{0}{0} & \SetCell{bg=\HeatCellBg{1}}\CellText{0.01}{1}{0} & \SetCell{bg=\HeatCellBg{1}}\CellText{0.01}{1}{0} & \SetCell{bg=\HeatCellBg{1}}\CellText{0.01}{1}{0} & \SetCell{bg=\HeatCellBg{2}}\CellText{0.02}{1}{0} & \SetCell{bg=\HeatCellBg{1}}\CellText{0.01}{0}{0} & \SetCell{bg=\HeatCellBg{1}}\CellText{0.01}{0}{0} & \SetCell{bg=\HeatCellBg{2}}\CellText{0.02}{1}{0} & \SetCell{bg=\HeatCellBg{2}}\CellText{0.02}{1}{0} & \SetCell{bg=\HeatCellBg{2}}\CellText{0.02}{1}{0} & \SetCell{bg=\HeatCellBg{2}}\CellText{0.02}{0}{0} & \SetCell{bg=\HeatCellBg{3}}\CellText{0.03}{1}{1} & \SetCell{bg=\HeatCellBg{3}}\CellText{0.03}{1}{1} & \SetCell{bg=\HeatCellBg{3}}\CellText{0.03}{1}{1} & \SetCell{bg=\HeatCellBg{3}}\CellText{0.03}{1}{1} & \SetCell{bg=\HeatCellBg{3}}\CellText{0.03}{1}{1} & \SetCell{bg=\HeatCellBg{3}}\CellText{0.03}{1}{1} & \SetCell{bg=\HeatCellBg{3}}\CellText{0.03}{1}{1} & \MeanCell{0.02}{0.01}{0} \\
        6a & \SetCell{bg=\HeatCellBg{0}}\CellText{-0.00}{1}{0} & \SetCell{bg=\HeatCellBg{1}}\CellText{0.01}{1}{0} & \SetCell{bg=\HeatCellBg{1}}\CellText{0.01}{1}{0} & \SetCell{bg=\HeatCellBg{1}}\CellText{0.01}{1}{0} & \SetCell{bg=\HeatCellBg{1}}\CellText{0.01}{1}{0} & \SetCell{bg=\HeatCellBg{2}}\CellText{0.02}{1}{0} & \SetCell{bg=\HeatCellBg{2}}\CellText{0.02}{1}{0} & \SetCell{bg=\HeatCellBg{2}}\CellText{0.02}{1}{0} & \SetCell{bg=\HeatCellBg{2}}\CellText{0.02}{1}{0} & \SetCell{bg=\HeatCellBg{2}}\CellText{0.02}{1}{0} & \SetCell{bg=\HeatCellBg{2}}\CellText{0.02}{1}{0} & \SetCell{bg=\HeatCellBg{3}}\CellText{0.03}{1}{1} & \SetCell{bg=\HeatCellBg{3}}\CellText{0.03}{1}{1} & \SetCell{bg=\HeatCellBg{2}}\CellText{0.02}{0}{0} & \SetCell{bg=\HeatCellBg{3}}\CellText{0.03}{1}{1} & \SetCell{bg=\HeatCellBg{3}}\CellText{0.03}{1}{1} & \SetCell{bg=\HeatCellBg{3}}\CellText{0.03}{1}{1} & \SetCell{bg=\HeatCellBg{3}}\CellText{0.03}{1}{1} & \SetCell{bg=\HeatCellBg{3}}\CellText{0.03}{1}{1} & \MeanCell{0.02}{0.01}{1} \\
        Mean $\pm$ Std & \MeanCell{0.00}{0.00}{0} & \MeanCell{0.00}{0.00}{0} & \MeanCell{0.01}{0.00}{0} & \MeanCell{0.01}{0.00}{0} & \MeanCell{0.01}{0.00}{0} & \MeanCell{0.01}{0.01}{0} & \MeanCell{0.01}{0.00}{0} & \MeanCell{0.01}{0.00}{0} & \MeanCell{0.02}{0.00}{0} & \MeanCell{0.02}{0.00}{0} & \MeanCell{0.02}{0.00}{0} & \MeanCell{0.02}{0.01}{0} & \MeanCell{0.03}{0.00}{0} & \MeanCell{0.02}{0.00}{0} & \MeanCell{0.03}{0.00}{0} & \MeanCell{0.03}{0.00}{1} & \MeanCell{0.03}{0.00}{1} & \MeanCell{0.03}{0.00}{1} & \MeanCell{0.03}{0.00}{1} & \MeanCell{0.02}{0.01}{0} \\
}


\clearpage
\section{\SurroundName}
\label{sec:appendix_additional_results_surround}
\ClusterKTable{10}{\gls{cu} --- \SurroundName.}{cu_surround}{
        1 & \SetCell{bg=\HeatCellBg{91}}\CellText{0.91}{1}{1} & \SetCell{bg=\HeatCellBg{41}}\CellText{0.41}{0}{0} & \SetCell{bg=\HeatCellBg{67}}\CellText{0.67}{0}{0} & \SetCell{bg=\HeatCellBg{84}}\CellText{0.84}{0}{0} & \SetCell{bg=\HeatCellBg{87}}\CellText{0.87}{1}{0} & \SetCell{bg=\HeatCellBg{78}}\CellText{0.78}{1}{0} & \SetCell{bg=\HeatCellBg{83}}\CellText{0.83}{1}{0} & \SetCell{bg=\HeatCellBg{73}}\CellText{0.73}{1}{0} & \SetCell{bg=\HeatCellBg{76}}\CellText{0.76}{1}{0} & \SetCell{bg=\HeatCellBg{79}}\CellText{0.79}{1}{0} & \SetCell{bg=\HeatCellBg{69}}\CellText{0.69}{1}{0} & \SetCell{bg=\HeatCellBg{76}}\CellText{0.76}{1}{0} & \SetCell{bg=\HeatCellBg{66}}\CellText{0.66}{1}{0} & \SetCell{bg=\HeatCellBg{53}}\CellText{0.53}{0}{0} & \SetCell{bg=\HeatCellBg{63}}\CellText{0.63}{1}{0} & \SetCell{bg=\HeatCellBg{56}}\CellText{0.56}{0}{0} & \SetCell{bg=\HeatCellBg{56}}\CellText{0.56}{0}{0} & \SetCell{bg=\HeatCellBg{68}}\CellText{0.68}{1}{0} & \SetCell{bg=\HeatCellBg{73}}\CellText{0.73}{1}{0} & \MeanCell{0.70}{0.12}{1} \\
        2 & \SetCell{bg=\HeatCellBg{57}}\CellText{0.57}{0}{0} & \SetCell{bg=\HeatCellBg{63}}\CellText{0.63}{0}{0} & \SetCell{bg=\HeatCellBg{86}}\CellText{0.86}{1}{1} & \SetCell{bg=\HeatCellBg{66}}\CellText{0.66}{0}{0} & \SetCell{bg=\HeatCellBg{65}}\CellText{0.65}{0}{0} & \SetCell{bg=\HeatCellBg{51}}\CellText{0.51}{0}{0} & \SetCell{bg=\HeatCellBg{48}}\CellText{0.48}{0}{0} & \SetCell{bg=\HeatCellBg{49}}\CellText{0.49}{0}{0} & \SetCell{bg=\HeatCellBg{54}}\CellText{0.54}{0}{0} & \SetCell{bg=\HeatCellBg{59}}\CellText{0.59}{0}{0} & \SetCell{bg=\HeatCellBg{53}}\CellText{0.53}{0}{0} & \SetCell{bg=\HeatCellBg{53}}\CellText{0.53}{0}{0} & \SetCell{bg=\HeatCellBg{46}}\CellText{0.46}{0}{0} & \SetCell{bg=\HeatCellBg{59}}\CellText{0.59}{0}{0} & \SetCell{bg=\HeatCellBg{59}}\CellText{0.59}{0}{0} & \SetCell{bg=\HeatCellBg{58}}\CellText{0.58}{0}{0} & \SetCell{bg=\HeatCellBg{48}}\CellText{0.48}{0}{0} & \SetCell{bg=\HeatCellBg{63}}\CellText{0.63}{0}{0} & \SetCell{bg=\HeatCellBg{62}}\CellText{0.62}{0}{0} & \MeanCell{0.58}{0.09}{0} \\
        3 & \SetCell{bg=\HeatCellBg{82}}\CellText{0.82}{0}{0} & \SetCell{bg=\HeatCellBg{61}}\CellText{0.61}{0}{0} & \SetCell{bg=\HeatCellBg{84}}\CellText{0.84}{0}{0} & \SetCell{bg=\HeatCellBg{85}}\CellText{0.85}{1}{1} & \SetCell{bg=\HeatCellBg{68}}\CellText{0.68}{0}{0} & \SetCell{bg=\HeatCellBg{72}}\CellText{0.72}{0}{0} & \SetCell{bg=\HeatCellBg{71}}\CellText{0.71}{0}{0} & \SetCell{bg=\HeatCellBg{59}}\CellText{0.59}{0}{0} & \SetCell{bg=\HeatCellBg{68}}\CellText{0.68}{0}{0} & \SetCell{bg=\HeatCellBg{68}}\CellText{0.68}{0}{0} & \SetCell{bg=\HeatCellBg{59}}\CellText{0.59}{0}{0} & \SetCell{bg=\HeatCellBg{61}}\CellText{0.61}{0}{0} & \SetCell{bg=\HeatCellBg{65}}\CellText{0.65}{0}{0} & \SetCell{bg=\HeatCellBg{56}}\CellText{0.56}{0}{0} & \SetCell{bg=\HeatCellBg{60}}\CellText{0.60}{0}{0} & \SetCell{bg=\HeatCellBg{56}}\CellText{0.56}{0}{0} & \SetCell{bg=\HeatCellBg{63}}\CellText{0.63}{1}{0} & \SetCell{bg=\HeatCellBg{57}}\CellText{0.57}{0}{0} & \SetCell{bg=\HeatCellBg{58}}\CellText{0.58}{0}{0} & \MeanCell{0.66}{0.09}{0} \\
        4 & \SetCell{bg=\HeatCellBg{85}}\CellText{0.85}{0}{1} & \SetCell{bg=\HeatCellBg{77}}\CellText{0.77}{1}{0} & \SetCell{bg=\HeatCellBg{80}}\CellText{0.80}{0}{0} & \SetCell{bg=\HeatCellBg{60}}\CellText{0.60}{0}{0} & \SetCell{bg=\HeatCellBg{69}}\CellText{0.69}{0}{0} & \SetCell{bg=\HeatCellBg{67}}\CellText{0.67}{0}{0} & \SetCell{bg=\HeatCellBg{75}}\CellText{0.75}{0}{0} & \SetCell{bg=\HeatCellBg{65}}\CellText{0.65}{0}{0} & \SetCell{bg=\HeatCellBg{73}}\CellText{0.73}{0}{0} & \SetCell{bg=\HeatCellBg{67}}\CellText{0.67}{0}{0} & \SetCell{bg=\HeatCellBg{63}}\CellText{0.63}{0}{0} & \SetCell{bg=\HeatCellBg{59}}\CellText{0.59}{0}{0} & \SetCell{bg=\HeatCellBg{62}}\CellText{0.62}{0}{0} & \SetCell{bg=\HeatCellBg{62}}\CellText{0.62}{1}{0} & \SetCell{bg=\HeatCellBg{57}}\CellText{0.57}{0}{0} & \SetCell{bg=\HeatCellBg{60}}\CellText{0.60}{1}{0} & \SetCell{bg=\HeatCellBg{53}}\CellText{0.53}{0}{0} & \SetCell{bg=\HeatCellBg{52}}\CellText{0.52}{0}{0} & \SetCell{bg=\HeatCellBg{53}}\CellText{0.53}{0}{0} & \MeanCell{0.65}{0.09}{0} \\
        5 & \SetCell{bg=\HeatCellBg{85}}\CellText{0.85}{0}{1} & \SetCell{bg=\HeatCellBg{77}}\CellText{0.77}{1}{0} & \SetCell{bg=\HeatCellBg{80}}\CellText{0.80}{0}{0} & \SetCell{bg=\HeatCellBg{60}}\CellText{0.60}{0}{0} & \SetCell{bg=\HeatCellBg{69}}\CellText{0.69}{0}{0} & \SetCell{bg=\HeatCellBg{67}}\CellText{0.67}{0}{0} & \SetCell{bg=\HeatCellBg{75}}\CellText{0.75}{0}{0} & \SetCell{bg=\HeatCellBg{65}}\CellText{0.65}{0}{0} & \SetCell{bg=\HeatCellBg{73}}\CellText{0.73}{0}{0} & \SetCell{bg=\HeatCellBg{67}}\CellText{0.67}{0}{0} & \SetCell{bg=\HeatCellBg{63}}\CellText{0.63}{0}{0} & \SetCell{bg=\HeatCellBg{59}}\CellText{0.59}{0}{0} & \SetCell{bg=\HeatCellBg{62}}\CellText{0.62}{0}{0} & \SetCell{bg=\HeatCellBg{62}}\CellText{0.62}{1}{0} & \SetCell{bg=\HeatCellBg{57}}\CellText{0.57}{0}{0} & \SetCell{bg=\HeatCellBg{60}}\CellText{0.60}{1}{0} & \SetCell{bg=\HeatCellBg{53}}\CellText{0.53}{0}{0} & \SetCell{bg=\HeatCellBg{52}}\CellText{0.52}{0}{0} & \SetCell{bg=\HeatCellBg{53}}\CellText{0.53}{0}{0} & \MeanCell{0.65}{0.09}{0} \\
        6b & \SetCell{bg=\HeatCellBg{83}}\CellText{0.83}{0}{1} & \SetCell{bg=\HeatCellBg{70}}\CellText{0.70}{0}{0} & \SetCell{bg=\HeatCellBg{74}}\CellText{0.74}{0}{0} & \SetCell{bg=\HeatCellBg{58}}\CellText{0.58}{0}{0} & \SetCell{bg=\HeatCellBg{60}}\CellText{0.60}{0}{0} & \SetCell{bg=\HeatCellBg{62}}\CellText{0.62}{0}{0} & \SetCell{bg=\HeatCellBg{49}}\CellText{0.49}{0}{0} & \SetCell{bg=\HeatCellBg{50}}\CellText{0.50}{0}{0} & \SetCell{bg=\HeatCellBg{61}}\CellText{0.61}{0}{0} & \SetCell{bg=\HeatCellBg{66}}\CellText{0.66}{0}{0} & \SetCell{bg=\HeatCellBg{57}}\CellText{0.57}{0}{0} & \SetCell{bg=\HeatCellBg{51}}\CellText{0.51}{0}{0} & \SetCell{bg=\HeatCellBg{50}}\CellText{0.50}{0}{0} & \SetCell{bg=\HeatCellBg{58}}\CellText{0.58}{0}{0} & \SetCell{bg=\HeatCellBg{57}}\CellText{0.57}{0}{0} & \SetCell{bg=\HeatCellBg{52}}\CellText{0.52}{0}{0} & \SetCell{bg=\HeatCellBg{56}}\CellText{0.56}{0}{0} & \SetCell{bg=\HeatCellBg{53}}\CellText{0.53}{0}{0} & \SetCell{bg=\HeatCellBg{53}}\CellText{0.53}{0}{0} & \MeanCell{0.59}{0.09}{0} \\
        7b & \SetCell{bg=\HeatCellBg{73}}\CellText{0.73}{0}{1} & \SetCell{bg=\HeatCellBg{63}}\CellText{0.63}{0}{0} & \SetCell{bg=\HeatCellBg{55}}\CellText{0.55}{0}{0} & \SetCell{bg=\HeatCellBg{57}}\CellText{0.57}{0}{0} & \SetCell{bg=\HeatCellBg{67}}\CellText{0.67}{0}{0} & \SetCell{bg=\HeatCellBg{53}}\CellText{0.53}{0}{0} & \SetCell{bg=\HeatCellBg{48}}\CellText{0.48}{0}{0} & \SetCell{bg=\HeatCellBg{46}}\CellText{0.46}{0}{0} & \SetCell{bg=\HeatCellBg{48}}\CellText{0.48}{0}{0} & \SetCell{bg=\HeatCellBg{42}}\CellText{0.42}{0}{0} & \SetCell{bg=\HeatCellBg{44}}\CellText{0.44}{0}{0} & \SetCell{bg=\HeatCellBg{45}}\CellText{0.45}{0}{0} & \SetCell{bg=\HeatCellBg{42}}\CellText{0.42}{0}{0} & \SetCell{bg=\HeatCellBg{43}}\CellText{0.43}{0}{0} & \SetCell{bg=\HeatCellBg{41}}\CellText{0.41}{0}{0} & \SetCell{bg=\HeatCellBg{41}}\CellText{0.41}{0}{0} & \SetCell{bg=\HeatCellBg{40}}\CellText{0.40}{0}{0} & \SetCell{bg=\HeatCellBg{39}}\CellText{0.39}{0}{0} & \SetCell{bg=\HeatCellBg{36}}\CellText{0.36}{0}{0} & \MeanCell{0.49}{0.10}{0} \\
        8b & \SetCell{bg=\HeatCellBg{77}}\CellText{0.77}{0}{1} & \SetCell{bg=\HeatCellBg{75}}\CellText{0.75}{0}{0} & \SetCell{bg=\HeatCellBg{67}}\CellText{0.67}{0}{0} & \SetCell{bg=\HeatCellBg{54}}\CellText{0.54}{0}{0} & \SetCell{bg=\HeatCellBg{52}}\CellText{0.52}{0}{0} & \SetCell{bg=\HeatCellBg{55}}\CellText{0.55}{0}{0} & \SetCell{bg=\HeatCellBg{52}}\CellText{0.52}{0}{0} & \SetCell{bg=\HeatCellBg{58}}\CellText{0.58}{0}{0} & \SetCell{bg=\HeatCellBg{61}}\CellText{0.61}{0}{0} & \SetCell{bg=\HeatCellBg{59}}\CellText{0.59}{0}{0} & \SetCell{bg=\HeatCellBg{51}}\CellText{0.51}{0}{0} & \SetCell{bg=\HeatCellBg{52}}\CellText{0.52}{0}{0} & \SetCell{bg=\HeatCellBg{50}}\CellText{0.50}{0}{0} & \SetCell{bg=\HeatCellBg{46}}\CellText{0.46}{0}{0} & \SetCell{bg=\HeatCellBg{50}}\CellText{0.50}{0}{0} & \SetCell{bg=\HeatCellBg{45}}\CellText{0.45}{0}{0} & \SetCell{bg=\HeatCellBg{43}}\CellText{0.43}{0}{0} & \SetCell{bg=\HeatCellBg{43}}\CellText{0.43}{0}{0} & \SetCell{bg=\HeatCellBg{43}}\CellText{0.43}{0}{0} & \MeanCell{0.54}{0.10}{0} \\
        Mean $\pm$ Std & \MeanCell{0.79}{0.10}{1} & \MeanCell{0.66}{0.11}{0} & \MeanCell{0.74}{0.10}{0} & \MeanCell{0.66}{0.11}{0} & \MeanCell{0.67}{0.09}{0} & \MeanCell{0.63}{0.09}{0} & \MeanCell{0.63}{0.14}{0} & \MeanCell{0.58}{0.09}{0} & \MeanCell{0.64}{0.09}{0} & \MeanCell{0.63}{0.10}{0} & \MeanCell{0.57}{0.07}{0} & \MeanCell{0.57}{0.09}{0} & \MeanCell{0.55}{0.09}{0} & \MeanCell{0.55}{0.07}{0} & \MeanCell{0.55}{0.06}{0} & \MeanCell{0.54}{0.07}{0} & \MeanCell{0.52}{0.07}{0} & \MeanCell{0.53}{0.09}{0} & \MeanCell{0.54}{0.11}{0} & \MeanCell{0.61}{0.12}{0} \\
}

\ClusterKTable{10}{\gls{as} --- \SurroundName.}{as_surround}{
        1 & \SetCell{bg=\HeatCellBg{10}}\CellText{0.10}{0}{0} & \SetCell{bg=\HeatCellBg{34}}\CellText{0.34}{0}{0} & \SetCell{bg=\HeatCellBg{27}}\CellText{0.27}{0}{0} & \SetCell{bg=\HeatCellBg{15}}\CellText{0.15}{0}{0} & \SetCell{bg=\HeatCellBg{15}}\CellText{0.15}{0}{0} & \SetCell{bg=\HeatCellBg{21}}\CellText{0.21}{0}{0} & \SetCell{bg=\HeatCellBg{20}}\CellText{0.20}{0}{0} & \SetCell{bg=\HeatCellBg{23}}\CellText{0.23}{0}{0} & \SetCell{bg=\HeatCellBg{26}}\CellText{0.26}{0}{0} & \SetCell{bg=\HeatCellBg{22}}\CellText{0.22}{0}{0} & \SetCell{bg=\HeatCellBg{34}}\CellText{0.34}{0}{0} & \SetCell{bg=\HeatCellBg{28}}\CellText{0.28}{0}{0} & \SetCell{bg=\HeatCellBg{26}}\CellText{0.26}{0}{0} & \SetCell{bg=\HeatCellBg{33}}\CellText{0.33}{0}{0} & \SetCell{bg=\HeatCellBg{30}}\CellText{0.30}{0}{0} & \SetCell{bg=\HeatCellBg{41}}\CellText{0.41}{0}{0} & \SetCell{bg=\HeatCellBg{35}}\CellText{0.35}{0}{0} & \SetCell{bg=\HeatCellBg{44}}\CellText{0.44}{0}{1} & \SetCell{bg=\HeatCellBg{34}}\CellText{0.34}{0}{0} & \MeanCell{0.27}{0.09}{0} \\
        2 & \SetCell{bg=\HeatCellBg{45}}\CellText{0.45}{0}{0} & \SetCell{bg=\HeatCellBg{29}}\CellText{0.29}{0}{0} & \SetCell{bg=\HeatCellBg{21}}\CellText{0.21}{0}{0} & \SetCell{bg=\HeatCellBg{32}}\CellText{0.32}{0}{0} & \SetCell{bg=\HeatCellBg{29}}\CellText{0.29}{0}{0} & \SetCell{bg=\HeatCellBg{46}}\CellText{0.46}{0}{0} & \SetCell{bg=\HeatCellBg{41}}\CellText{0.41}{0}{0} & \SetCell{bg=\HeatCellBg{49}}\CellText{0.49}{0}{0} & \SetCell{bg=\HeatCellBg{47}}\CellText{0.47}{0}{0} & \SetCell{bg=\HeatCellBg{56}}\CellText{0.56}{0}{0} & \SetCell{bg=\HeatCellBg{52}}\CellText{0.52}{0}{0} & \SetCell{bg=\HeatCellBg{57}}\CellText{0.57}{0}{1} & \SetCell{bg=\HeatCellBg{50}}\CellText{0.50}{0}{0} & \SetCell{bg=\HeatCellBg{52}}\CellText{0.52}{0}{0} & \SetCell{bg=\HeatCellBg{54}}\CellText{0.54}{0}{0} & \SetCell{bg=\HeatCellBg{56}}\CellText{0.56}{0}{0} & \SetCell{bg=\HeatCellBg{48}}\CellText{0.48}{0}{0} & \SetCell{bg=\HeatCellBg{48}}\CellText{0.48}{0}{0} & \SetCell{bg=\HeatCellBg{48}}\CellText{0.48}{0}{0} & \MeanCell{0.45}{0.10}{0} \\
        3 & \SetCell{bg=\HeatCellBg{74}}\CellText{0.74}{1}{1} & \SetCell{bg=\HeatCellBg{40}}\CellText{0.40}{0}{0} & \SetCell{bg=\HeatCellBg{31}}\CellText{0.31}{0}{0} & \SetCell{bg=\HeatCellBg{42}}\CellText{0.42}{0}{0} & \SetCell{bg=\HeatCellBg{47}}\CellText{0.47}{0}{0} & \SetCell{bg=\HeatCellBg{48}}\CellText{0.48}{0}{0} & \SetCell{bg=\HeatCellBg{57}}\CellText{0.57}{0}{0} & \SetCell{bg=\HeatCellBg{59}}\CellText{0.59}{0}{0} & \SetCell{bg=\HeatCellBg{60}}\CellText{0.60}{0}{0} & \SetCell{bg=\HeatCellBg{58}}\CellText{0.58}{0}{0} & \SetCell{bg=\HeatCellBg{61}}\CellText{0.61}{0}{0} & \SetCell{bg=\HeatCellBg{62}}\CellText{0.62}{0}{0} & \SetCell{bg=\HeatCellBg{64}}\CellText{0.64}{0}{0} & \SetCell{bg=\HeatCellBg{60}}\CellText{0.60}{0}{0} & \SetCell{bg=\HeatCellBg{61}}\CellText{0.61}{0}{0} & \SetCell{bg=\HeatCellBg{63}}\CellText{0.63}{0}{0} & \SetCell{bg=\HeatCellBg{56}}\CellText{0.56}{0}{0} & \SetCell{bg=\HeatCellBg{61}}\CellText{0.61}{0}{0} & \SetCell{bg=\HeatCellBg{62}}\CellText{0.62}{0}{0} & \MeanCell{0.56}{0.10}{0} \\
        4 & \SetCell{bg=\HeatCellBg{41}}\CellText{0.41}{0}{0} & \SetCell{bg=\HeatCellBg{36}}\CellText{0.36}{0}{0} & \SetCell{bg=\HeatCellBg{34}}\CellText{0.34}{0}{0} & \SetCell{bg=\HeatCellBg{46}}\CellText{0.46}{0}{0} & \SetCell{bg=\HeatCellBg{48}}\CellText{0.48}{0}{0} & \SetCell{bg=\HeatCellBg{50}}\CellText{0.50}{0}{0} & \SetCell{bg=\HeatCellBg{47}}\CellText{0.47}{0}{0} & \SetCell{bg=\HeatCellBg{46}}\CellText{0.46}{0}{0} & \SetCell{bg=\HeatCellBg{48}}\CellText{0.48}{0}{0} & \SetCell{bg=\HeatCellBg{50}}\CellText{0.50}{0}{0} & \SetCell{bg=\HeatCellBg{49}}\CellText{0.49}{0}{0} & \SetCell{bg=\HeatCellBg{53}}\CellText{0.53}{0}{0} & \SetCell{bg=\HeatCellBg{46}}\CellText{0.46}{0}{0} & \SetCell{bg=\HeatCellBg{47}}\CellText{0.47}{0}{0} & \SetCell{bg=\HeatCellBg{53}}\CellText{0.53}{0}{0} & \SetCell{bg=\HeatCellBg{58}}\CellText{0.58}{0}{0} & \SetCell{bg=\HeatCellBg{54}}\CellText{0.54}{0}{0} & \SetCell{bg=\HeatCellBg{57}}\CellText{0.57}{0}{0} & \SetCell{bg=\HeatCellBg{60}}\CellText{0.60}{0}{1} & \MeanCell{0.49}{0.07}{0} \\
        5 & \SetCell{bg=\HeatCellBg{41}}\CellText{0.41}{0}{0} & \SetCell{bg=\HeatCellBg{36}}\CellText{0.36}{0}{0} & \SetCell{bg=\HeatCellBg{34}}\CellText{0.34}{0}{0} & \SetCell{bg=\HeatCellBg{46}}\CellText{0.46}{0}{0} & \SetCell{bg=\HeatCellBg{48}}\CellText{0.48}{0}{0} & \SetCell{bg=\HeatCellBg{50}}\CellText{0.50}{0}{0} & \SetCell{bg=\HeatCellBg{47}}\CellText{0.47}{0}{0} & \SetCell{bg=\HeatCellBg{46}}\CellText{0.46}{0}{0} & \SetCell{bg=\HeatCellBg{48}}\CellText{0.48}{0}{0} & \SetCell{bg=\HeatCellBg{50}}\CellText{0.50}{0}{0} & \SetCell{bg=\HeatCellBg{49}}\CellText{0.49}{0}{0} & \SetCell{bg=\HeatCellBg{53}}\CellText{0.53}{0}{0} & \SetCell{bg=\HeatCellBg{46}}\CellText{0.46}{0}{0} & \SetCell{bg=\HeatCellBg{47}}\CellText{0.47}{0}{0} & \SetCell{bg=\HeatCellBg{53}}\CellText{0.53}{0}{0} & \SetCell{bg=\HeatCellBg{58}}\CellText{0.58}{0}{0} & \SetCell{bg=\HeatCellBg{54}}\CellText{0.54}{0}{0} & \SetCell{bg=\HeatCellBg{57}}\CellText{0.57}{0}{0} & \SetCell{bg=\HeatCellBg{60}}\CellText{0.60}{0}{1} & \MeanCell{0.49}{0.07}{0} \\
        6b & \SetCell{bg=\HeatCellBg{63}}\CellText{0.63}{0}{1} & \SetCell{bg=\HeatCellBg{49}}\CellText{0.49}{0}{0} & \SetCell{bg=\HeatCellBg{44}}\CellText{0.44}{0}{0} & \SetCell{bg=\HeatCellBg{60}}\CellText{0.60}{0}{0} & \SetCell{bg=\HeatCellBg{62}}\CellText{0.62}{0}{0} & \SetCell{bg=\HeatCellBg{61}}\CellText{0.61}{0}{0} & \SetCell{bg=\HeatCellBg{56}}\CellText{0.56}{0}{0} & \SetCell{bg=\HeatCellBg{60}}\CellText{0.60}{0}{0} & \SetCell{bg=\HeatCellBg{54}}\CellText{0.54}{0}{0} & \SetCell{bg=\HeatCellBg{54}}\CellText{0.54}{0}{0} & \SetCell{bg=\HeatCellBg{57}}\CellText{0.57}{0}{0} & \SetCell{bg=\HeatCellBg{58}}\CellText{0.58}{0}{0} & \SetCell{bg=\HeatCellBg{57}}\CellText{0.57}{0}{0} & \SetCell{bg=\HeatCellBg{55}}\CellText{0.55}{0}{0} & \SetCell{bg=\HeatCellBg{63}}\CellText{0.63}{0}{1} & \SetCell{bg=\HeatCellBg{59}}\CellText{0.59}{0}{0} & \SetCell{bg=\HeatCellBg{62}}\CellText{0.62}{0}{0} & \SetCell{bg=\HeatCellBg{53}}\CellText{0.53}{0}{0} & \SetCell{bg=\HeatCellBg{62}}\CellText{0.62}{0}{0} & \MeanCell{0.57}{0.05}{0} \\
        7b & \SetCell{bg=\HeatCellBg{66}}\CellText{0.66}{0}{0} & \SetCell{bg=\HeatCellBg{53}}\CellText{0.53}{1}{0} & \SetCell{bg=\HeatCellBg{70}}\CellText{0.70}{0}{0} & \SetCell{bg=\HeatCellBg{70}}\CellText{0.70}{1}{0} & \SetCell{bg=\HeatCellBg{71}}\CellText{0.71}{1}{0} & \SetCell{bg=\HeatCellBg{72}}\CellText{0.72}{1}{0} & \SetCell{bg=\HeatCellBg{71}}\CellText{0.71}{1}{0} & \SetCell{bg=\HeatCellBg{74}}\CellText{0.74}{1}{0} & \SetCell{bg=\HeatCellBg{71}}\CellText{0.71}{1}{0} & \SetCell{bg=\HeatCellBg{73}}\CellText{0.73}{1}{0} & \SetCell{bg=\HeatCellBg{75}}\CellText{0.75}{1}{0} & \SetCell{bg=\HeatCellBg{76}}\CellText{0.76}{1}{1} & \SetCell{bg=\HeatCellBg{75}}\CellText{0.75}{1}{0} & \SetCell{bg=\HeatCellBg{74}}\CellText{0.74}{1}{0} & \SetCell{bg=\HeatCellBg{76}}\CellText{0.76}{1}{1} & \SetCell{bg=\HeatCellBg{75}}\CellText{0.75}{1}{0} & \SetCell{bg=\HeatCellBg{75}}\CellText{0.75}{0}{0} & \SetCell{bg=\HeatCellBg{73}}\CellText{0.73}{0}{0} & \SetCell{bg=\HeatCellBg{76}}\CellText{0.76}{1}{1} & \MeanCell{0.72}{0.05}{1} \\
        8b & \SetCell{bg=\HeatCellBg{66}}\CellText{0.66}{0}{0} & \SetCell{bg=\HeatCellBg{53}}\CellText{0.53}{1}{0} & \SetCell{bg=\HeatCellBg{71}}\CellText{0.71}{1}{0} & \SetCell{bg=\HeatCellBg{70}}\CellText{0.70}{1}{0} & \SetCell{bg=\HeatCellBg{71}}\CellText{0.71}{1}{0} & \SetCell{bg=\HeatCellBg{70}}\CellText{0.70}{0}{0} & \SetCell{bg=\HeatCellBg{71}}\CellText{0.71}{1}{0} & \SetCell{bg=\HeatCellBg{69}}\CellText{0.69}{0}{0} & \SetCell{bg=\HeatCellBg{68}}\CellText{0.68}{0}{0} & \SetCell{bg=\HeatCellBg{71}}\CellText{0.71}{0}{0} & \SetCell{bg=\HeatCellBg{74}}\CellText{0.74}{0}{0} & \SetCell{bg=\HeatCellBg{74}}\CellText{0.74}{0}{0} & \SetCell{bg=\HeatCellBg{73}}\CellText{0.73}{0}{0} & \SetCell{bg=\HeatCellBg{74}}\CellText{0.74}{1}{0} & \SetCell{bg=\HeatCellBg{73}}\CellText{0.73}{0}{0} & \SetCell{bg=\HeatCellBg{75}}\CellText{0.75}{1}{0} & \SetCell{bg=\HeatCellBg{76}}\CellText{0.76}{1}{1} & \SetCell{bg=\HeatCellBg{75}}\CellText{0.75}{1}{0} & \SetCell{bg=\HeatCellBg{76}}\CellText{0.76}{1}{1} & \MeanCell{0.71}{0.05}{0} \\
        Mean $\pm$ Std & \MeanCell{0.51}{0.19}{0} & \MeanCell{0.41}{0.09}{0} & \MeanCell{0.41}{0.18}{0} & \MeanCell{0.48}{0.18}{0} & \MeanCell{0.49}{0.18}{0} & \MeanCell{0.52}{0.15}{0} & \MeanCell{0.51}{0.16}{0} & \MeanCell{0.53}{0.15}{0} & \MeanCell{0.53}{0.13}{0} & \MeanCell{0.54}{0.15}{0} & \MeanCell{0.56}{0.13}{0} & \MeanCell{0.58}{0.14}{0} & \MeanCell{0.55}{0.15}{0} & \MeanCell{0.55}{0.13}{0} & \MeanCell{0.58}{0.13}{0} & \MeanCell{0.61}{0.10}{1} & \MeanCell{0.57}{0.13}{0} & \MeanCell{0.58}{0.10}{0} & \MeanCell{0.60}{0.13}{0} & \MeanCell{0.53}{0.15}{0} \\
}

\ClusterKTable{10}{\gls{acu} --- \SurroundName.}{acu_surround}{
        1 & \SetCell{bg=\HeatCellBg{1}}\CellText{0.01}{0}{0} & \SetCell{bg=\HeatCellBg{-25}}\CellText{-0.25}{0}{0} & \SetCell{bg=\HeatCellBg{-6}}\CellText{-0.06}{0}{0} & \SetCell{bg=\HeatCellBg{-1}}\CellText{-0.01}{0}{0} & \SetCell{bg=\HeatCellBg{2}}\CellText{0.02}{0}{0} & \SetCell{bg=\HeatCellBg{-1}}\CellText{-0.01}{0}{0} & \SetCell{bg=\HeatCellBg{2}}\CellText{0.02}{0}{0} & \SetCell{bg=\HeatCellBg{-4}}\CellText{-0.04}{0}{0} & \SetCell{bg=\HeatCellBg{1}}\CellText{0.01}{0}{0} & \SetCell{bg=\HeatCellBg{1}}\CellText{0.01}{0}{0} & \SetCell{bg=\HeatCellBg{3}}\CellText{0.03}{0}{0} & \SetCell{bg=\HeatCellBg{4}}\CellText{0.04}{0}{0} & \SetCell{bg=\HeatCellBg{-8}}\CellText{-0.08}{0}{0} & \SetCell{bg=\HeatCellBg{-14}}\CellText{-0.14}{0}{0} & \SetCell{bg=\HeatCellBg{-8}}\CellText{-0.08}{0}{0} & \SetCell{bg=\HeatCellBg{-3}}\CellText{-0.03}{0}{0} & \SetCell{bg=\HeatCellBg{-9}}\CellText{-0.09}{0}{0} & \SetCell{bg=\HeatCellBg{12}}\CellText{0.12}{0}{1} & \SetCell{bg=\HeatCellBg{7}}\CellText{0.07}{0}{0} & \MeanCell{-0.02}{0.08}{0} \\
        2 & \SetCell{bg=\HeatCellBg{2}}\CellText{0.02}{0}{0} & \SetCell{bg=\HeatCellBg{-8}}\CellText{-0.08}{0}{0} & \SetCell{bg=\HeatCellBg{7}}\CellText{0.07}{0}{0} & \SetCell{bg=\HeatCellBg{-2}}\CellText{-0.02}{0}{0} & \SetCell{bg=\HeatCellBg{-5}}\CellText{-0.05}{0}{0} & \SetCell{bg=\HeatCellBg{-3}}\CellText{-0.03}{0}{0} & \SetCell{bg=\HeatCellBg{-11}}\CellText{-0.11}{0}{0} & \SetCell{bg=\HeatCellBg{-2}}\CellText{-0.02}{0}{0} & \SetCell{bg=\HeatCellBg{1}}\CellText{0.01}{0}{0} & \SetCell{bg=\HeatCellBg{15}}\CellText{0.15}{0}{1} & \SetCell{bg=\HeatCellBg{5}}\CellText{0.05}{0}{0} & \SetCell{bg=\HeatCellBg{10}}\CellText{0.10}{0}{0} & \SetCell{bg=\HeatCellBg{-5}}\CellText{-0.05}{0}{0} & \SetCell{bg=\HeatCellBg{11}}\CellText{0.11}{0}{0} & \SetCell{bg=\HeatCellBg{13}}\CellText{0.13}{0}{0} & \SetCell{bg=\HeatCellBg{14}}\CellText{0.14}{0}{0} & \SetCell{bg=\HeatCellBg{-4}}\CellText{-0.04}{0}{0} & \SetCell{bg=\HeatCellBg{12}}\CellText{0.12}{0}{0} & \SetCell{bg=\HeatCellBg{10}}\CellText{0.10}{0}{0} & \MeanCell{0.03}{0.08}{0} \\
        3 & \SetCell{bg=\HeatCellBg{55}}\CellText{0.55}{1}{1} & \SetCell{bg=\HeatCellBg{1}}\CellText{0.01}{0}{0} & \SetCell{bg=\HeatCellBg{15}}\CellText{0.15}{0}{0} & \SetCell{bg=\HeatCellBg{27}}\CellText{0.27}{1}{0} & \SetCell{bg=\HeatCellBg{15}}\CellText{0.15}{0}{0} & \SetCell{bg=\HeatCellBg{19}}\CellText{0.19}{0}{0} & \SetCell{bg=\HeatCellBg{29}}\CellText{0.29}{1}{0} & \SetCell{bg=\HeatCellBg{17}}\CellText{0.17}{0}{0} & \SetCell{bg=\HeatCellBg{28}}\CellText{0.28}{0}{0} & \SetCell{bg=\HeatCellBg{26}}\CellText{0.26}{0}{0} & \SetCell{bg=\HeatCellBg{20}}\CellText{0.20}{0}{0} & \SetCell{bg=\HeatCellBg{22}}\CellText{0.22}{0}{0} & \SetCell{bg=\HeatCellBg{29}}\CellText{0.29}{1}{0} & \SetCell{bg=\HeatCellBg{17}}\CellText{0.17}{0}{0} & \SetCell{bg=\HeatCellBg{21}}\CellText{0.21}{0}{0} & \SetCell{bg=\HeatCellBg{19}}\CellText{0.19}{0}{0} & \SetCell{bg=\HeatCellBg{19}}\CellText{0.19}{1}{0} & \SetCell{bg=\HeatCellBg{18}}\CellText{0.18}{1}{0} & \SetCell{bg=\HeatCellBg{20}}\CellText{0.20}{1}{0} & \MeanCell{0.22}{0.10}{0} \\
        4 & \SetCell{bg=\HeatCellBg{26}}\CellText{0.26}{0}{1} & \SetCell{bg=\HeatCellBg{12}}\CellText{0.12}{0}{0} & \SetCell{bg=\HeatCellBg{15}}\CellText{0.15}{0}{0} & \SetCell{bg=\HeatCellBg{5}}\CellText{0.05}{0}{0} & \SetCell{bg=\HeatCellBg{18}}\CellText{0.18}{0}{0} & \SetCell{bg=\HeatCellBg{17}}\CellText{0.17}{0}{0} & \SetCell{bg=\HeatCellBg{22}}\CellText{0.22}{0}{0} & \SetCell{bg=\HeatCellBg{11}}\CellText{0.11}{0}{0} & \SetCell{bg=\HeatCellBg{20}}\CellText{0.20}{0}{0} & \SetCell{bg=\HeatCellBg{17}}\CellText{0.17}{0}{0} & \SetCell{bg=\HeatCellBg{12}}\CellText{0.12}{0}{0} & \SetCell{bg=\HeatCellBg{12}}\CellText{0.12}{0}{0} & \SetCell{bg=\HeatCellBg{7}}\CellText{0.07}{0}{0} & \SetCell{bg=\HeatCellBg{9}}\CellText{0.09}{0}{0} & \SetCell{bg=\HeatCellBg{10}}\CellText{0.10}{0}{0} & \SetCell{bg=\HeatCellBg{18}}\CellText{0.18}{0}{0} & \SetCell{bg=\HeatCellBg{7}}\CellText{0.07}{0}{0} & \SetCell{bg=\HeatCellBg{9}}\CellText{0.09}{0}{0} & \SetCell{bg=\HeatCellBg{13}}\CellText{0.13}{0}{0} & \MeanCell{0.14}{0.05}{0} \\
        5 & \SetCell{bg=\HeatCellBg{26}}\CellText{0.26}{0}{1} & \SetCell{bg=\HeatCellBg{12}}\CellText{0.12}{0}{0} & \SetCell{bg=\HeatCellBg{15}}\CellText{0.15}{0}{0} & \SetCell{bg=\HeatCellBg{5}}\CellText{0.05}{0}{0} & \SetCell{bg=\HeatCellBg{18}}\CellText{0.18}{0}{0} & \SetCell{bg=\HeatCellBg{17}}\CellText{0.17}{0}{0} & \SetCell{bg=\HeatCellBg{22}}\CellText{0.22}{0}{0} & \SetCell{bg=\HeatCellBg{11}}\CellText{0.11}{0}{0} & \SetCell{bg=\HeatCellBg{20}}\CellText{0.20}{0}{0} & \SetCell{bg=\HeatCellBg{17}}\CellText{0.17}{0}{0} & \SetCell{bg=\HeatCellBg{12}}\CellText{0.12}{0}{0} & \SetCell{bg=\HeatCellBg{12}}\CellText{0.12}{0}{0} & \SetCell{bg=\HeatCellBg{7}}\CellText{0.07}{0}{0} & \SetCell{bg=\HeatCellBg{9}}\CellText{0.09}{0}{0} & \SetCell{bg=\HeatCellBg{10}}\CellText{0.10}{0}{0} & \SetCell{bg=\HeatCellBg{18}}\CellText{0.18}{0}{0} & \SetCell{bg=\HeatCellBg{7}}\CellText{0.07}{0}{0} & \SetCell{bg=\HeatCellBg{9}}\CellText{0.09}{0}{0} & \SetCell{bg=\HeatCellBg{13}}\CellText{0.13}{0}{0} & \MeanCell{0.14}{0.05}{0} \\
        6b & \SetCell{bg=\HeatCellBg{46}}\CellText{0.46}{0}{1} & \SetCell{bg=\HeatCellBg{19}}\CellText{0.19}{0}{0} & \SetCell{bg=\HeatCellBg{18}}\CellText{0.18}{0}{0} & \SetCell{bg=\HeatCellBg{18}}\CellText{0.18}{0}{0} & \SetCell{bg=\HeatCellBg{22}}\CellText{0.22}{0}{0} & \SetCell{bg=\HeatCellBg{22}}\CellText{0.22}{0}{0} & \SetCell{bg=\HeatCellBg{5}}\CellText{0.05}{0}{0} & \SetCell{bg=\HeatCellBg{10}}\CellText{0.10}{0}{0} & \SetCell{bg=\HeatCellBg{15}}\CellText{0.15}{0}{0} & \SetCell{bg=\HeatCellBg{19}}\CellText{0.19}{0}{0} & \SetCell{bg=\HeatCellBg{14}}\CellText{0.14}{0}{0} & \SetCell{bg=\HeatCellBg{10}}\CellText{0.10}{0}{0} & \SetCell{bg=\HeatCellBg{7}}\CellText{0.07}{0}{0} & \SetCell{bg=\HeatCellBg{13}}\CellText{0.13}{0}{0} & \SetCell{bg=\HeatCellBg{19}}\CellText{0.19}{0}{0} & \SetCell{bg=\HeatCellBg{11}}\CellText{0.11}{0}{0} & \SetCell{bg=\HeatCellBg{18}}\CellText{0.18}{0}{0} & \SetCell{bg=\HeatCellBg{6}}\CellText{0.06}{0}{0} & \SetCell{bg=\HeatCellBg{15}}\CellText{0.15}{0}{0} & \MeanCell{0.16}{0.09}{0} \\
        7b & \SetCell{bg=\HeatCellBg{39}}\CellText{0.39}{0}{1} & \SetCell{bg=\HeatCellBg{16}}\CellText{0.16}{0}{0} & \SetCell{bg=\HeatCellBg{25}}\CellText{0.25}{0}{0} & \SetCell{bg=\HeatCellBg{27}}\CellText{0.27}{1}{0} & \SetCell{bg=\HeatCellBg{38}}\CellText{0.38}{1}{0} & \SetCell{bg=\HeatCellBg{25}}\CellText{0.25}{1}{0} & \SetCell{bg=\HeatCellBg{19}}\CellText{0.19}{0}{0} & \SetCell{bg=\HeatCellBg{19}}\CellText{0.19}{0}{0} & \SetCell{bg=\HeatCellBg{19}}\CellText{0.19}{0}{0} & \SetCell{bg=\HeatCellBg{14}}\CellText{0.14}{0}{0} & \SetCell{bg=\HeatCellBg{19}}\CellText{0.19}{0}{0} & \SetCell{bg=\HeatCellBg{20}}\CellText{0.20}{0}{0} & \SetCell{bg=\HeatCellBg{17}}\CellText{0.17}{0}{0} & \SetCell{bg=\HeatCellBg{16}}\CellText{0.16}{0}{0} & \SetCell{bg=\HeatCellBg{17}}\CellText{0.17}{0}{0} & \SetCell{bg=\HeatCellBg{15}}\CellText{0.15}{0}{0} & \SetCell{bg=\HeatCellBg{15}}\CellText{0.15}{0}{0} & \SetCell{bg=\HeatCellBg{12}}\CellText{0.12}{0}{0} & \SetCell{bg=\HeatCellBg{11}}\CellText{0.11}{0}{0} & \MeanCell{0.20}{0.07}{0} \\
        8b & \SetCell{bg=\HeatCellBg{42}}\CellText{0.42}{0}{1} & \SetCell{bg=\HeatCellBg{28}}\CellText{0.28}{1}{0} & \SetCell{bg=\HeatCellBg{38}}\CellText{0.38}{1}{0} & \SetCell{bg=\HeatCellBg{24}}\CellText{0.24}{0}{0} & \SetCell{bg=\HeatCellBg{23}}\CellText{0.23}{0}{0} & \SetCell{bg=\HeatCellBg{24}}\CellText{0.24}{0}{0} & \SetCell{bg=\HeatCellBg{23}}\CellText{0.23}{0}{0} & \SetCell{bg=\HeatCellBg{27}}\CellText{0.27}{1}{0} & \SetCell{bg=\HeatCellBg{29}}\CellText{0.29}{1}{0} & \SetCell{bg=\HeatCellBg{30}}\CellText{0.30}{1}{0} & \SetCell{bg=\HeatCellBg{24}}\CellText{0.24}{1}{0} & \SetCell{bg=\HeatCellBg{26}}\CellText{0.26}{1}{0} & \SetCell{bg=\HeatCellBg{23}}\CellText{0.23}{0}{0} & \SetCell{bg=\HeatCellBg{20}}\CellText{0.20}{1}{0} & \SetCell{bg=\HeatCellBg{23}}\CellText{0.23}{1}{0} & \SetCell{bg=\HeatCellBg{20}}\CellText{0.20}{1}{0} & \SetCell{bg=\HeatCellBg{18}}\CellText{0.18}{0}{0} & \SetCell{bg=\HeatCellBg{18}}\CellText{0.18}{1}{0} & \SetCell{bg=\HeatCellBg{19}}\CellText{0.19}{0}{0} & \MeanCell{0.25}{0.06}{1} \\
        Mean $\pm$ Std & \MeanCell{0.30}{0.19}{1} & \MeanCell{0.07}{0.16}{0} & \MeanCell{0.16}{0.12}{0} & \MeanCell{0.13}{0.12}{0} & \MeanCell{0.16}{0.12}{0} & \MeanCell{0.15}{0.10}{0} & \MeanCell{0.14}{0.13}{0} & \MeanCell{0.11}{0.10}{0} & \MeanCell{0.17}{0.10}{0} & \MeanCell{0.17}{0.08}{0} & \MeanCell{0.14}{0.07}{0} & \MeanCell{0.14}{0.07}{0} & \MeanCell{0.10}{0.12}{0} & \MeanCell{0.10}{0.10}{0} & \MeanCell{0.13}{0.09}{0} & \MeanCell{0.14}{0.07}{0} & \MeanCell{0.09}{0.10}{0} & \MeanCell{0.12}{0.04}{0} & \MeanCell{0.14}{0.04}{0} & \MeanCell{0.14}{0.12}{0} \\
}

\ClusterKTable{10}{Observation baseline ($\Similarity{obs}{}$) --- \SurroundName.}{obs_surround}{
        1 & \SetCell{bg=\HeatCellBg{5}}\CellText{0.05}{0}{0} & \SetCell{bg=\HeatCellBg{33}}\CellText{0.33}{0}{0} & \SetCell{bg=\HeatCellBg{73}}\CellText{0.73}{0}{0} & \SetCell{bg=\HeatCellBg{85}}\CellText{0.85}{1}{1} & \SetCell{bg=\HeatCellBg{70}}\CellText{0.70}{1}{0} & \SetCell{bg=\HeatCellBg{78}}\CellText{0.78}{1}{0} & \SetCell{bg=\HeatCellBg{58}}\CellText{0.58}{1}{0} & \SetCell{bg=\HeatCellBg{58}}\CellText{0.58}{1}{0} & \SetCell{bg=\HeatCellBg{74}}\CellText{0.74}{1}{0} & \SetCell{bg=\HeatCellBg{59}}\CellText{0.59}{1}{0} & \SetCell{bg=\HeatCellBg{55}}\CellText{0.55}{1}{0} & \SetCell{bg=\HeatCellBg{61}}\CellText{0.61}{1}{0} & \SetCell{bg=\HeatCellBg{70}}\CellText{0.70}{1}{0} & \SetCell{bg=\HeatCellBg{48}}\CellText{0.48}{0}{0} & \SetCell{bg=\HeatCellBg{64}}\CellText{0.64}{1}{0} & \SetCell{bg=\HeatCellBg{56}}\CellText{0.56}{1}{0} & \SetCell{bg=\HeatCellBg{54}}\CellText{0.54}{1}{0} & \SetCell{bg=\HeatCellBg{52}}\CellText{0.52}{1}{0} & \SetCell{bg=\HeatCellBg{49}}\CellText{0.49}{1}{0} & \MeanCell{0.58}{0.17}{1} \\
        2 & \SetCell{bg=\HeatCellBg{20}}\CellText{0.20}{0}{0} & \SetCell{bg=\HeatCellBg{65}}\CellText{0.65}{1}{0} & \SetCell{bg=\HeatCellBg{76}}\CellText{0.76}{1}{1} & \SetCell{bg=\HeatCellBg{66}}\CellText{0.66}{0}{0} & \SetCell{bg=\HeatCellBg{70}}\CellText{0.70}{1}{0} & \SetCell{bg=\HeatCellBg{54}}\CellText{0.54}{0}{0} & \SetCell{bg=\HeatCellBg{52}}\CellText{0.52}{0}{0} & \SetCell{bg=\HeatCellBg{51}}\CellText{0.51}{0}{0} & \SetCell{bg=\HeatCellBg{51}}\CellText{0.51}{0}{0} & \SetCell{bg=\HeatCellBg{44}}\CellText{0.44}{0}{0} & \SetCell{bg=\HeatCellBg{43}}\CellText{0.43}{0}{0} & \SetCell{bg=\HeatCellBg{43}}\CellText{0.43}{0}{0} & \SetCell{bg=\HeatCellBg{47}}\CellText{0.47}{0}{0} & \SetCell{bg=\HeatCellBg{48}}\CellText{0.48}{0}{0} & \SetCell{bg=\HeatCellBg{40}}\CellText{0.40}{0}{0} & \SetCell{bg=\HeatCellBg{42}}\CellText{0.42}{0}{0} & \SetCell{bg=\HeatCellBg{45}}\CellText{0.45}{0}{0} & \SetCell{bg=\HeatCellBg{49}}\CellText{0.49}{0}{0} & \SetCell{bg=\HeatCellBg{46}}\CellText{0.46}{0}{0} & \MeanCell{0.50}{0.12}{0} \\
        3 & \SetCell{bg=\HeatCellBg{26}}\CellText{0.26}{0}{0} & \SetCell{bg=\HeatCellBg{60}}\CellText{0.60}{0}{0} & \SetCell{bg=\HeatCellBg{69}}\CellText{0.69}{0}{1} & \SetCell{bg=\HeatCellBg{58}}\CellText{0.58}{0}{0} & \SetCell{bg=\HeatCellBg{53}}\CellText{0.53}{0}{0} & \SetCell{bg=\HeatCellBg{52}}\CellText{0.52}{0}{0} & \SetCell{bg=\HeatCellBg{41}}\CellText{0.41}{0}{0} & \SetCell{bg=\HeatCellBg{39}}\CellText{0.39}{0}{0} & \SetCell{bg=\HeatCellBg{39}}\CellText{0.39}{0}{0} & \SetCell{bg=\HeatCellBg{36}}\CellText{0.36}{0}{0} & \SetCell{bg=\HeatCellBg{35}}\CellText{0.35}{0}{0} & \SetCell{bg=\HeatCellBg{36}}\CellText{0.36}{0}{0} & \SetCell{bg=\HeatCellBg{36}}\CellText{0.36}{0}{0} & \SetCell{bg=\HeatCellBg{39}}\CellText{0.39}{0}{0} & \SetCell{bg=\HeatCellBg{37}}\CellText{0.37}{0}{0} & \SetCell{bg=\HeatCellBg{34}}\CellText{0.34}{0}{0} & \SetCell{bg=\HeatCellBg{36}}\CellText{0.36}{0}{0} & \SetCell{bg=\HeatCellBg{38}}\CellText{0.38}{0}{0} & \SetCell{bg=\HeatCellBg{35}}\CellText{0.35}{0}{0} & \MeanCell{0.42}{0.11}{0} \\
        4 & \SetCell{bg=\HeatCellBg{59}}\CellText{0.59}{1}{0} & \SetCell{bg=\HeatCellBg{64}}\CellText{0.64}{0}{0} & \SetCell{bg=\HeatCellBg{66}}\CellText{0.66}{0}{1} & \SetCell{bg=\HeatCellBg{54}}\CellText{0.54}{0}{0} & \SetCell{bg=\HeatCellBg{52}}\CellText{0.52}{0}{0} & \SetCell{bg=\HeatCellBg{50}}\CellText{0.50}{0}{0} & \SetCell{bg=\HeatCellBg{53}}\CellText{0.53}{0}{0} & \SetCell{bg=\HeatCellBg{49}}\CellText{0.49}{0}{0} & \SetCell{bg=\HeatCellBg{52}}\CellText{0.52}{0}{0} & \SetCell{bg=\HeatCellBg{47}}\CellText{0.47}{0}{0} & \SetCell{bg=\HeatCellBg{45}}\CellText{0.45}{0}{0} & \SetCell{bg=\HeatCellBg{43}}\CellText{0.43}{0}{0} & \SetCell{bg=\HeatCellBg{49}}\CellText{0.49}{0}{0} & \SetCell{bg=\HeatCellBg{53}}\CellText{0.53}{1}{0} & \SetCell{bg=\HeatCellBg{45}}\CellText{0.45}{0}{0} & \SetCell{bg=\HeatCellBg{42}}\CellText{0.42}{0}{0} & \SetCell{bg=\HeatCellBg{45}}\CellText{0.45}{0}{0} & \SetCell{bg=\HeatCellBg{43}}\CellText{0.43}{0}{0} & \SetCell{bg=\HeatCellBg{37}}\CellText{0.37}{0}{0} & \MeanCell{0.50}{0.07}{0} \\
        5 & \SetCell{bg=\HeatCellBg{59}}\CellText{0.59}{1}{0} & \SetCell{bg=\HeatCellBg{64}}\CellText{0.64}{0}{0} & \SetCell{bg=\HeatCellBg{66}}\CellText{0.66}{0}{1} & \SetCell{bg=\HeatCellBg{54}}\CellText{0.54}{0}{0} & \SetCell{bg=\HeatCellBg{52}}\CellText{0.52}{0}{0} & \SetCell{bg=\HeatCellBg{50}}\CellText{0.50}{0}{0} & \SetCell{bg=\HeatCellBg{53}}\CellText{0.53}{0}{0} & \SetCell{bg=\HeatCellBg{49}}\CellText{0.49}{0}{0} & \SetCell{bg=\HeatCellBg{52}}\CellText{0.52}{0}{0} & \SetCell{bg=\HeatCellBg{47}}\CellText{0.47}{0}{0} & \SetCell{bg=\HeatCellBg{45}}\CellText{0.45}{0}{0} & \SetCell{bg=\HeatCellBg{43}}\CellText{0.43}{0}{0} & \SetCell{bg=\HeatCellBg{49}}\CellText{0.49}{0}{0} & \SetCell{bg=\HeatCellBg{53}}\CellText{0.53}{1}{0} & \SetCell{bg=\HeatCellBg{45}}\CellText{0.45}{0}{0} & \SetCell{bg=\HeatCellBg{42}}\CellText{0.42}{0}{0} & \SetCell{bg=\HeatCellBg{45}}\CellText{0.45}{0}{0} & \SetCell{bg=\HeatCellBg{43}}\CellText{0.43}{0}{0} & \SetCell{bg=\HeatCellBg{37}}\CellText{0.37}{0}{0} & \MeanCell{0.50}{0.07}{0} \\
        6b & \SetCell{bg=\HeatCellBg{37}}\CellText{0.37}{0}{0} & \SetCell{bg=\HeatCellBg{51}}\CellText{0.51}{0}{0} & \SetCell{bg=\HeatCellBg{56}}\CellText{0.56}{0}{1} & \SetCell{bg=\HeatCellBg{40}}\CellText{0.40}{0}{0} & \SetCell{bg=\HeatCellBg{38}}\CellText{0.38}{0}{0} & \SetCell{bg=\HeatCellBg{39}}\CellText{0.39}{0}{0} & \SetCell{bg=\HeatCellBg{44}}\CellText{0.44}{0}{0} & \SetCell{bg=\HeatCellBg{40}}\CellText{0.40}{0}{0} & \SetCell{bg=\HeatCellBg{42}}\CellText{0.42}{0}{0} & \SetCell{bg=\HeatCellBg{46}}\CellText{0.46}{0}{0} & \SetCell{bg=\HeatCellBg{41}}\CellText{0.41}{0}{0} & \SetCell{bg=\HeatCellBg{39}}\CellText{0.39}{0}{0} & \SetCell{bg=\HeatCellBg{43}}\CellText{0.43}{0}{0} & \SetCell{bg=\HeatCellBg{45}}\CellText{0.45}{0}{0} & \SetCell{bg=\HeatCellBg{35}}\CellText{0.35}{0}{0} & \SetCell{bg=\HeatCellBg{41}}\CellText{0.41}{0}{0} & \SetCell{bg=\HeatCellBg{37}}\CellText{0.37}{0}{0} & \SetCell{bg=\HeatCellBg{45}}\CellText{0.45}{0}{0} & \SetCell{bg=\HeatCellBg{37}}\CellText{0.37}{0}{0} & \MeanCell{0.42}{0.05}{0} \\
        7b & \SetCell{bg=\HeatCellBg{34}}\CellText{0.34}{0}{0} & \SetCell{bg=\HeatCellBg{47}}\CellText{0.47}{0}{1} & \SetCell{bg=\HeatCellBg{30}}\CellText{0.30}{0}{0} & \SetCell{bg=\HeatCellBg{30}}\CellText{0.30}{0}{0} & \SetCell{bg=\HeatCellBg{29}}\CellText{0.29}{0}{0} & \SetCell{bg=\HeatCellBg{28}}\CellText{0.28}{0}{0} & \SetCell{bg=\HeatCellBg{29}}\CellText{0.29}{0}{0} & \SetCell{bg=\HeatCellBg{24}}\CellText{0.24}{0}{0} & \SetCell{bg=\HeatCellBg{24}}\CellText{0.24}{0}{0} & \SetCell{bg=\HeatCellBg{27}}\CellText{0.27}{0}{0} & \SetCell{bg=\HeatCellBg{24}}\CellText{0.24}{0}{0} & \SetCell{bg=\HeatCellBg{24}}\CellText{0.24}{0}{0} & \SetCell{bg=\HeatCellBg{25}}\CellText{0.25}{0}{0} & \SetCell{bg=\HeatCellBg{26}}\CellText{0.26}{0}{0} & \SetCell{bg=\HeatCellBg{23}}\CellText{0.23}{0}{0} & \SetCell{bg=\HeatCellBg{25}}\CellText{0.25}{0}{0} & \SetCell{bg=\HeatCellBg{25}}\CellText{0.25}{0}{0} & \SetCell{bg=\HeatCellBg{27}}\CellText{0.27}{0}{0} & \SetCell{bg=\HeatCellBg{24}}\CellText{0.24}{0}{0} & \MeanCell{0.28}{0.05}{0} \\
        8b & \SetCell{bg=\HeatCellBg{34}}\CellText{0.34}{0}{0} & \SetCell{bg=\HeatCellBg{47}}\CellText{0.47}{0}{1} & \SetCell{bg=\HeatCellBg{29}}\CellText{0.29}{0}{0} & \SetCell{bg=\HeatCellBg{25}}\CellText{0.25}{0}{0} & \SetCell{bg=\HeatCellBg{24}}\CellText{0.24}{0}{0} & \SetCell{bg=\HeatCellBg{25}}\CellText{0.25}{0}{0} & \SetCell{bg=\HeatCellBg{23}}\CellText{0.23}{0}{0} & \SetCell{bg=\HeatCellBg{20}}\CellText{0.20}{0}{0} & \SetCell{bg=\HeatCellBg{20}}\CellText{0.20}{0}{0} & \SetCell{bg=\HeatCellBg{21}}\CellText{0.21}{0}{0} & \SetCell{bg=\HeatCellBg{18}}\CellText{0.18}{0}{0} & \SetCell{bg=\HeatCellBg{19}}\CellText{0.19}{0}{0} & \SetCell{bg=\HeatCellBg{20}}\CellText{0.20}{0}{0} & \SetCell{bg=\HeatCellBg{19}}\CellText{0.19}{0}{0} & \SetCell{bg=\HeatCellBg{19}}\CellText{0.19}{0}{0} & \SetCell{bg=\HeatCellBg{17}}\CellText{0.17}{0}{0} & \SetCell{bg=\HeatCellBg{18}}\CellText{0.18}{0}{0} & \SetCell{bg=\HeatCellBg{22}}\CellText{0.22}{0}{0} & \SetCell{bg=\HeatCellBg{19}}\CellText{0.19}{0}{0} & \MeanCell{0.23}{0.07}{0} \\
        Mean $\pm$ Std & \MeanCell{0.34}{0.17}{0} & \MeanCell{0.54}{0.11}{0} & \MeanCell{0.58}{0.17}{1} & \MeanCell{0.52}{0.18}{0} & \MeanCell{0.48}{0.16}{0} & \MeanCell{0.47}{0.16}{0} & \MeanCell{0.44}{0.12}{0} & \MeanCell{0.41}{0.13}{0} & \MeanCell{0.44}{0.16}{0} & \MeanCell{0.41}{0.11}{0} & \MeanCell{0.38}{0.11}{0} & \MeanCell{0.39}{0.12}{0} & \MeanCell{0.42}{0.15}{0} & \MeanCell{0.41}{0.12}{0} & \MeanCell{0.39}{0.13}{0} & \MeanCell{0.37}{0.11}{0} & \MeanCell{0.38}{0.11}{0} & \MeanCell{0.40}{0.10}{0} & \MeanCell{0.35}{0.09}{0} & \MeanCell{0.43}{0.15}{0} \\
}

\ClusterKTable{10}{Untrained-network (latent) baseline ($\Similarity{untr}{}$) --- \SurroundName.}{untr_surround}{
        1 & \SetCell{bg=\HeatCellBg{90}}\CellText{0.90}{1}{1} & \SetCell{bg=\HeatCellBg{66}}\CellText{0.66}{1}{0} & \SetCell{bg=\HeatCellBg{70}}\CellText{0.70}{0}{0} & \SetCell{bg=\HeatCellBg{73}}\CellText{0.73}{1}{0} & \SetCell{bg=\HeatCellBg{85}}\CellText{0.85}{1}{0} & \SetCell{bg=\HeatCellBg{79}}\CellText{0.79}{1}{0} & \SetCell{bg=\HeatCellBg{80}}\CellText{0.80}{1}{0} & \SetCell{bg=\HeatCellBg{77}}\CellText{0.77}{1}{0} & \SetCell{bg=\HeatCellBg{70}}\CellText{0.70}{1}{0} & \SetCell{bg=\HeatCellBg{78}}\CellText{0.78}{1}{0} & \SetCell{bg=\HeatCellBg{66}}\CellText{0.66}{1}{0} & \SetCell{bg=\HeatCellBg{72}}\CellText{0.72}{1}{0} & \SetCell{bg=\HeatCellBg{60}}\CellText{0.60}{1}{0} & \SetCell{bg=\HeatCellBg{67}}\CellText{0.67}{1}{0} & \SetCell{bg=\HeatCellBg{70}}\CellText{0.70}{1}{0} & \SetCell{bg=\HeatCellBg{59}}\CellText{0.59}{1}{0} & \SetCell{bg=\HeatCellBg{65}}\CellText{0.65}{1}{0} & \SetCell{bg=\HeatCellBg{54}}\CellText{0.54}{1}{0} & \SetCell{bg=\HeatCellBg{66}}\CellText{0.66}{1}{0} & \MeanCell{0.71}{0.09}{1} \\
        2 & \SetCell{bg=\HeatCellBg{55}}\CellText{0.55}{0}{0} & \SetCell{bg=\HeatCellBg{46}}\CellText{0.46}{0}{0} & \SetCell{bg=\HeatCellBg{79}}\CellText{0.79}{1}{1} & \SetCell{bg=\HeatCellBg{68}}\CellText{0.68}{0}{0} & \SetCell{bg=\HeatCellBg{70}}\CellText{0.70}{0}{0} & \SetCell{bg=\HeatCellBg{51}}\CellText{0.51}{0}{0} & \SetCell{bg=\HeatCellBg{55}}\CellText{0.55}{0}{0} & \SetCell{bg=\HeatCellBg{46}}\CellText{0.46}{0}{0} & \SetCell{bg=\HeatCellBg{53}}\CellText{0.53}{0}{0} & \SetCell{bg=\HeatCellBg{40}}\CellText{0.40}{0}{0} & \SetCell{bg=\HeatCellBg{47}}\CellText{0.47}{0}{0} & \SetCell{bg=\HeatCellBg{41}}\CellText{0.41}{0}{0} & \SetCell{bg=\HeatCellBg{50}}\CellText{0.50}{0}{0} & \SetCell{bg=\HeatCellBg{44}}\CellText{0.44}{0}{0} & \SetCell{bg=\HeatCellBg{46}}\CellText{0.46}{0}{0} & \SetCell{bg=\HeatCellBg{44}}\CellText{0.44}{0}{0} & \SetCell{bg=\HeatCellBg{52}}\CellText{0.52}{0}{0} & \SetCell{bg=\HeatCellBg{52}}\CellText{0.52}{0}{0} & \SetCell{bg=\HeatCellBg{52}}\CellText{0.52}{0}{0} & \MeanCell{0.52}{0.10}{0} \\
        3 & \SetCell{bg=\HeatCellBg{3}}\CellText{0.03}{0}{0} & \SetCell{bg=\HeatCellBg{43}}\CellText{0.43}{0}{0} & \SetCell{bg=\HeatCellBg{68}}\CellText{0.68}{0}{1} & \SetCell{bg=\HeatCellBg{58}}\CellText{0.58}{0}{0} & \SetCell{bg=\HeatCellBg{51}}\CellText{0.51}{0}{0} & \SetCell{bg=\HeatCellBg{50}}\CellText{0.50}{0}{0} & \SetCell{bg=\HeatCellBg{43}}\CellText{0.43}{0}{0} & \SetCell{bg=\HeatCellBg{41}}\CellText{0.41}{0}{0} & \SetCell{bg=\HeatCellBg{39}}\CellText{0.39}{0}{0} & \SetCell{bg=\HeatCellBg{42}}\CellText{0.42}{0}{0} & \SetCell{bg=\HeatCellBg{39}}\CellText{0.39}{0}{0} & \SetCell{bg=\HeatCellBg{38}}\CellText{0.38}{0}{0} & \SetCell{bg=\HeatCellBg{36}}\CellText{0.36}{0}{0} & \SetCell{bg=\HeatCellBg{38}}\CellText{0.38}{0}{0} & \SetCell{bg=\HeatCellBg{39}}\CellText{0.39}{0}{0} & \SetCell{bg=\HeatCellBg{37}}\CellText{0.37}{0}{0} & \SetCell{bg=\HeatCellBg{44}}\CellText{0.44}{0}{0} & \SetCell{bg=\HeatCellBg{38}}\CellText{0.38}{0}{0} & \SetCell{bg=\HeatCellBg{38}}\CellText{0.38}{0}{0} & \MeanCell{0.41}{0.12}{0} \\
        4 & \SetCell{bg=\HeatCellBg{43}}\CellText{0.43}{0}{0} & \SetCell{bg=\HeatCellBg{38}}\CellText{0.38}{0}{0} & \SetCell{bg=\HeatCellBg{54}}\CellText{0.54}{0}{1} & \SetCell{bg=\HeatCellBg{48}}\CellText{0.48}{0}{0} & \SetCell{bg=\HeatCellBg{48}}\CellText{0.48}{0}{0} & \SetCell{bg=\HeatCellBg{45}}\CellText{0.45}{0}{0} & \SetCell{bg=\HeatCellBg{52}}\CellText{0.52}{0}{0} & \SetCell{bg=\HeatCellBg{53}}\CellText{0.53}{0}{0} & \SetCell{bg=\HeatCellBg{52}}\CellText{0.52}{0}{0} & \SetCell{bg=\HeatCellBg{50}}\CellText{0.50}{0}{0} & \SetCell{bg=\HeatCellBg{51}}\CellText{0.51}{0}{0} & \SetCell{bg=\HeatCellBg{45}}\CellText{0.45}{0}{0} & \SetCell{bg=\HeatCellBg{53}}\CellText{0.53}{0}{0} & \SetCell{bg=\HeatCellBg{52}}\CellText{0.52}{0}{0} & \SetCell{bg=\HeatCellBg{46}}\CellText{0.46}{0}{0} & \SetCell{bg=\HeatCellBg{41}}\CellText{0.41}{0}{0} & \SetCell{bg=\HeatCellBg{46}}\CellText{0.46}{0}{0} & \SetCell{bg=\HeatCellBg{38}}\CellText{0.38}{0}{0} & \SetCell{bg=\HeatCellBg{40}}\CellText{0.40}{0}{0} & \MeanCell{0.47}{0.05}{0} \\
        5 & \SetCell{bg=\HeatCellBg{43}}\CellText{0.43}{0}{0} & \SetCell{bg=\HeatCellBg{38}}\CellText{0.38}{0}{0} & \SetCell{bg=\HeatCellBg{54}}\CellText{0.54}{0}{1} & \SetCell{bg=\HeatCellBg{48}}\CellText{0.48}{0}{0} & \SetCell{bg=\HeatCellBg{48}}\CellText{0.48}{0}{0} & \SetCell{bg=\HeatCellBg{45}}\CellText{0.45}{0}{0} & \SetCell{bg=\HeatCellBg{52}}\CellText{0.52}{0}{0} & \SetCell{bg=\HeatCellBg{53}}\CellText{0.53}{0}{0} & \SetCell{bg=\HeatCellBg{52}}\CellText{0.52}{0}{0} & \SetCell{bg=\HeatCellBg{50}}\CellText{0.50}{0}{0} & \SetCell{bg=\HeatCellBg{51}}\CellText{0.51}{0}{0} & \SetCell{bg=\HeatCellBg{45}}\CellText{0.45}{0}{0} & \SetCell{bg=\HeatCellBg{53}}\CellText{0.53}{0}{0} & \SetCell{bg=\HeatCellBg{52}}\CellText{0.52}{0}{0} & \SetCell{bg=\HeatCellBg{46}}\CellText{0.46}{0}{0} & \SetCell{bg=\HeatCellBg{41}}\CellText{0.41}{0}{0} & \SetCell{bg=\HeatCellBg{46}}\CellText{0.46}{0}{0} & \SetCell{bg=\HeatCellBg{38}}\CellText{0.38}{0}{0} & \SetCell{bg=\HeatCellBg{40}}\CellText{0.40}{0}{0} & \MeanCell{0.47}{0.05}{0} \\
        6b & \SetCell{bg=\HeatCellBg{30}}\CellText{0.30}{0}{0} & \SetCell{bg=\HeatCellBg{39}}\CellText{0.39}{0}{0} & \SetCell{bg=\HeatCellBg{34}}\CellText{0.34}{0}{0} & \SetCell{bg=\HeatCellBg{29}}\CellText{0.29}{0}{0} & \SetCell{bg=\HeatCellBg{32}}\CellText{0.32}{0}{0} & \SetCell{bg=\HeatCellBg{35}}\CellText{0.35}{0}{0} & \SetCell{bg=\HeatCellBg{36}}\CellText{0.36}{0}{0} & \SetCell{bg=\HeatCellBg{38}}\CellText{0.38}{0}{0} & \SetCell{bg=\HeatCellBg{46}}\CellText{0.46}{0}{1} & \SetCell{bg=\HeatCellBg{45}}\CellText{0.45}{0}{0} & \SetCell{bg=\HeatCellBg{43}}\CellText{0.43}{0}{0} & \SetCell{bg=\HeatCellBg{42}}\CellText{0.42}{0}{0} & \SetCell{bg=\HeatCellBg{42}}\CellText{0.42}{0}{0} & \SetCell{bg=\HeatCellBg{41}}\CellText{0.41}{0}{0} & \SetCell{bg=\HeatCellBg{37}}\CellText{0.37}{0}{0} & \SetCell{bg=\HeatCellBg{39}}\CellText{0.39}{0}{0} & \SetCell{bg=\HeatCellBg{37}}\CellText{0.37}{0}{0} & \SetCell{bg=\HeatCellBg{44}}\CellText{0.44}{0}{0} & \SetCell{bg=\HeatCellBg{38}}\CellText{0.38}{0}{0} & \MeanCell{0.38}{0.05}{0} \\
        7b & \SetCell{bg=\HeatCellBg{27}}\CellText{0.27}{0}{0} & \SetCell{bg=\HeatCellBg{35}}\CellText{0.35}{0}{1} & \SetCell{bg=\HeatCellBg{21}}\CellText{0.21}{0}{0} & \SetCell{bg=\HeatCellBg{22}}\CellText{0.22}{0}{0} & \SetCell{bg=\HeatCellBg{28}}\CellText{0.28}{0}{0} & \SetCell{bg=\HeatCellBg{26}}\CellText{0.26}{0}{0} & \SetCell{bg=\HeatCellBg{24}}\CellText{0.24}{0}{0} & \SetCell{bg=\HeatCellBg{22}}\CellText{0.22}{0}{0} & \SetCell{bg=\HeatCellBg{22}}\CellText{0.22}{0}{0} & \SetCell{bg=\HeatCellBg{22}}\CellText{0.22}{0}{0} & \SetCell{bg=\HeatCellBg{20}}\CellText{0.20}{0}{0} & \SetCell{bg=\HeatCellBg{24}}\CellText{0.24}{0}{0} & \SetCell{bg=\HeatCellBg{18}}\CellText{0.18}{0}{0} & \SetCell{bg=\HeatCellBg{20}}\CellText{0.20}{0}{0} & \SetCell{bg=\HeatCellBg{23}}\CellText{0.23}{0}{0} & \SetCell{bg=\HeatCellBg{22}}\CellText{0.22}{0}{0} & \SetCell{bg=\HeatCellBg{22}}\CellText{0.22}{0}{0} & \SetCell{bg=\HeatCellBg{18}}\CellText{0.18}{0}{0} & \SetCell{bg=\HeatCellBg{18}}\CellText{0.18}{0}{0} & \MeanCell{0.23}{0.04}{0} \\
        8b & \SetCell{bg=\HeatCellBg{21}}\CellText{0.21}{0}{0} & \SetCell{bg=\HeatCellBg{25}}\CellText{0.25}{0}{1} & \SetCell{bg=\HeatCellBg{18}}\CellText{0.18}{0}{0} & \SetCell{bg=\HeatCellBg{19}}\CellText{0.19}{0}{0} & \SetCell{bg=\HeatCellBg{18}}\CellText{0.18}{0}{0} & \SetCell{bg=\HeatCellBg{18}}\CellText{0.18}{0}{0} & \SetCell{bg=\HeatCellBg{16}}\CellText{0.16}{0}{0} & \SetCell{bg=\HeatCellBg{16}}\CellText{0.16}{0}{0} & \SetCell{bg=\HeatCellBg{15}}\CellText{0.15}{0}{0} & \SetCell{bg=\HeatCellBg{17}}\CellText{0.17}{0}{0} & \SetCell{bg=\HeatCellBg{15}}\CellText{0.15}{0}{0} & \SetCell{bg=\HeatCellBg{16}}\CellText{0.16}{0}{0} & \SetCell{bg=\HeatCellBg{17}}\CellText{0.17}{0}{0} & \SetCell{bg=\HeatCellBg{14}}\CellText{0.14}{0}{0} & \SetCell{bg=\HeatCellBg{18}}\CellText{0.18}{0}{0} & \SetCell{bg=\HeatCellBg{14}}\CellText{0.14}{0}{0} & \SetCell{bg=\HeatCellBg{15}}\CellText{0.15}{0}{0} & \SetCell{bg=\HeatCellBg{19}}\CellText{0.19}{0}{0} & \SetCell{bg=\HeatCellBg{17}}\CellText{0.17}{0}{0} & \MeanCell{0.17}{0.03}{0} \\
        Mean $\pm$ Std & \MeanCell{0.39}{0.24}{0} & \MeanCell{0.41}{0.11}{0} & \MeanCell{0.50}{0.22}{1} & \MeanCell{0.46}{0.19}{0} & \MeanCell{0.47}{0.21}{0} & \MeanCell{0.44}{0.17}{0} & \MeanCell{0.45}{0.19}{0} & \MeanCell{0.43}{0.18}{0} & \MeanCell{0.44}{0.17}{0} & \MeanCell{0.43}{0.18}{0} & \MeanCell{0.41}{0.16}{0} & \MeanCell{0.40}{0.15}{0} & \MeanCell{0.41}{0.15}{0} & \MeanCell{0.41}{0.16}{0} & \MeanCell{0.41}{0.15}{0} & \MeanCell{0.37}{0.13}{0} & \MeanCell{0.41}{0.15}{0} & \MeanCell{0.38}{0.12}{0} & \MeanCell{0.39}{0.15}{0} & \MeanCell{0.42}{0.17}{0} \\
}

\ClusterKTable{10}{Action baseline ($\Similarity{act}{}$) --- \SurroundName.}{act_surround}{
        1 & \SetCell{bg=\HeatCellBg{-1}}\CellText{-0.01}{0}{0} & \SetCell{bg=\HeatCellBg{-5}}\CellText{-0.05}{0}{0} & \SetCell{bg=\HeatCellBg{0}}\CellText{0.00}{0}{0} & \SetCell{bg=\HeatCellBg{1}}\CellText{0.01}{0}{0} & \SetCell{bg=\HeatCellBg{2}}\CellText{0.02}{0}{0} & \SetCell{bg=\HeatCellBg{1}}\CellText{0.01}{0}{0} & \SetCell{bg=\HeatCellBg{2}}\CellText{0.02}{0}{0} & \SetCell{bg=\HeatCellBg{2}}\CellText{0.02}{0}{0} & \SetCell{bg=\HeatCellBg{3}}\CellText{0.03}{0}{0} & \SetCell{bg=\HeatCellBg{3}}\CellText{0.03}{0}{0} & \SetCell{bg=\HeatCellBg{3}}\CellText{0.03}{0}{0} & \SetCell{bg=\HeatCellBg{2}}\CellText{0.02}{0}{0} & \SetCell{bg=\HeatCellBg{3}}\CellText{0.03}{0}{0} & \SetCell{bg=\HeatCellBg{3}}\CellText{0.03}{0}{0} & \SetCell{bg=\HeatCellBg{4}}\CellText{0.04}{0}{1} & \SetCell{bg=\HeatCellBg{3}}\CellText{0.03}{0}{0} & \SetCell{bg=\HeatCellBg{4}}\CellText{0.04}{0}{1} & \SetCell{bg=\HeatCellBg{3}}\CellText{0.03}{0}{0} & \SetCell{bg=\HeatCellBg{4}}\CellText{0.04}{0}{1} & \MeanCell{0.02}{0.02}{0} \\
        2 & \SetCell{bg=\HeatCellBg{9}}\CellText{0.09}{0}{1} & \SetCell{bg=\HeatCellBg{-2}}\CellText{-0.02}{0}{0} & \SetCell{bg=\HeatCellBg{0}}\CellText{0.00}{0}{0} & \SetCell{bg=\HeatCellBg{2}}\CellText{0.02}{0}{0} & \SetCell{bg=\HeatCellBg{3}}\CellText{0.03}{0}{0} & \SetCell{bg=\HeatCellBg{5}}\CellText{0.05}{0}{0} & \SetCell{bg=\HeatCellBg{4}}\CellText{0.04}{0}{0} & \SetCell{bg=\HeatCellBg{6}}\CellText{0.06}{0}{0} & \SetCell{bg=\HeatCellBg{4}}\CellText{0.04}{0}{0} & \SetCell{bg=\HeatCellBg{6}}\CellText{0.06}{0}{0} & \SetCell{bg=\HeatCellBg{5}}\CellText{0.05}{0}{0} & \SetCell{bg=\HeatCellBg{5}}\CellText{0.05}{0}{0} & \SetCell{bg=\HeatCellBg{5}}\CellText{0.05}{0}{0} & \SetCell{bg=\HeatCellBg{5}}\CellText{0.05}{0}{0} & \SetCell{bg=\HeatCellBg{4}}\CellText{0.04}{0}{0} & \SetCell{bg=\HeatCellBg{5}}\CellText{0.05}{0}{0} & \SetCell{bg=\HeatCellBg{5}}\CellText{0.05}{0}{0} & \SetCell{bg=\HeatCellBg{5}}\CellText{0.05}{0}{0} & \SetCell{bg=\HeatCellBg{5}}\CellText{0.05}{0}{0} & \MeanCell{0.04}{0.02}{0} \\
        3 & \SetCell{bg=\HeatCellBg{12}}\CellText{0.12}{0}{1} & \SetCell{bg=\HeatCellBg{2}}\CellText{0.02}{1}{0} & \SetCell{bg=\HeatCellBg{0}}\CellText{0.00}{0}{0} & \SetCell{bg=\HeatCellBg{7}}\CellText{0.07}{0}{0} & \SetCell{bg=\HeatCellBg{7}}\CellText{0.07}{0}{0} & \SetCell{bg=\HeatCellBg{6}}\CellText{0.06}{0}{0} & \SetCell{bg=\HeatCellBg{6}}\CellText{0.06}{0}{0} & \SetCell{bg=\HeatCellBg{7}}\CellText{0.07}{0}{0} & \SetCell{bg=\HeatCellBg{7}}\CellText{0.07}{0}{0} & \SetCell{bg=\HeatCellBg{7}}\CellText{0.07}{0}{0} & \SetCell{bg=\HeatCellBg{5}}\CellText{0.05}{0}{0} & \SetCell{bg=\HeatCellBg{5}}\CellText{0.05}{0}{0} & \SetCell{bg=\HeatCellBg{6}}\CellText{0.06}{0}{0} & \SetCell{bg=\HeatCellBg{6}}\CellText{0.06}{0}{0} & \SetCell{bg=\HeatCellBg{5}}\CellText{0.05}{0}{0} & \SetCell{bg=\HeatCellBg{5}}\CellText{0.05}{0}{0} & \SetCell{bg=\HeatCellBg{6}}\CellText{0.06}{0}{0} & \SetCell{bg=\HeatCellBg{5}}\CellText{0.05}{0}{0} & \SetCell{bg=\HeatCellBg{5}}\CellText{0.05}{0}{0} & \MeanCell{0.06}{0.02}{0} \\
        4 & \SetCell{bg=\HeatCellBg{4}}\CellText{0.04}{0}{0} & \SetCell{bg=\HeatCellBg{0}}\CellText{-0.00}{0}{0} & \SetCell{bg=\HeatCellBg{0}}\CellText{0.00}{0}{0} & \SetCell{bg=\HeatCellBg{4}}\CellText{0.04}{0}{0} & \SetCell{bg=\HeatCellBg{8}}\CellText{0.08}{0}{0} & \SetCell{bg=\HeatCellBg{9}}\CellText{0.09}{0}{1} & \SetCell{bg=\HeatCellBg{9}}\CellText{0.09}{0}{1} & \SetCell{bg=\HeatCellBg{8}}\CellText{0.08}{0}{0} & \SetCell{bg=\HeatCellBg{8}}\CellText{0.08}{0}{0} & \SetCell{bg=\HeatCellBg{8}}\CellText{0.08}{0}{0} & \SetCell{bg=\HeatCellBg{7}}\CellText{0.07}{0}{0} & \SetCell{bg=\HeatCellBg{6}}\CellText{0.06}{0}{0} & \SetCell{bg=\HeatCellBg{6}}\CellText{0.06}{0}{0} & \SetCell{bg=\HeatCellBg{6}}\CellText{0.06}{0}{0} & \SetCell{bg=\HeatCellBg{6}}\CellText{0.06}{0}{0} & \SetCell{bg=\HeatCellBg{6}}\CellText{0.06}{0}{0} & \SetCell{bg=\HeatCellBg{7}}\CellText{0.07}{0}{0} & \SetCell{bg=\HeatCellBg{5}}\CellText{0.05}{0}{0} & \SetCell{bg=\HeatCellBg{6}}\CellText{0.06}{0}{0} & \MeanCell{0.06}{0.02}{0} \\
        5 & \SetCell{bg=\HeatCellBg{4}}\CellText{0.04}{0}{0} & \SetCell{bg=\HeatCellBg{0}}\CellText{-0.00}{0}{0} & \SetCell{bg=\HeatCellBg{0}}\CellText{0.00}{0}{0} & \SetCell{bg=\HeatCellBg{4}}\CellText{0.04}{0}{0} & \SetCell{bg=\HeatCellBg{8}}\CellText{0.08}{0}{0} & \SetCell{bg=\HeatCellBg{9}}\CellText{0.09}{0}{1} & \SetCell{bg=\HeatCellBg{9}}\CellText{0.09}{0}{1} & \SetCell{bg=\HeatCellBg{8}}\CellText{0.08}{0}{0} & \SetCell{bg=\HeatCellBg{8}}\CellText{0.08}{0}{0} & \SetCell{bg=\HeatCellBg{8}}\CellText{0.08}{0}{0} & \SetCell{bg=\HeatCellBg{7}}\CellText{0.07}{0}{0} & \SetCell{bg=\HeatCellBg{6}}\CellText{0.06}{0}{0} & \SetCell{bg=\HeatCellBg{6}}\CellText{0.06}{0}{0} & \SetCell{bg=\HeatCellBg{6}}\CellText{0.06}{0}{0} & \SetCell{bg=\HeatCellBg{6}}\CellText{0.06}{0}{0} & \SetCell{bg=\HeatCellBg{6}}\CellText{0.06}{0}{0} & \SetCell{bg=\HeatCellBg{7}}\CellText{0.07}{0}{0} & \SetCell{bg=\HeatCellBg{5}}\CellText{0.05}{0}{0} & \SetCell{bg=\HeatCellBg{6}}\CellText{0.06}{0}{0} & \MeanCell{0.06}{0.02}{0} \\
        6b & \SetCell{bg=\HeatCellBg{14}}\CellText{0.14}{0}{1} & \SetCell{bg=\HeatCellBg{0}}\CellText{-0.00}{0}{0} & \SetCell{bg=\HeatCellBg{1}}\CellText{0.01}{0}{0} & \SetCell{bg=\HeatCellBg{13}}\CellText{0.13}{0}{0} & \SetCell{bg=\HeatCellBg{12}}\CellText{0.12}{0}{0} & \SetCell{bg=\HeatCellBg{10}}\CellText{0.10}{0}{0} & \SetCell{bg=\HeatCellBg{11}}\CellText{0.11}{0}{0} & \SetCell{bg=\HeatCellBg{12}}\CellText{0.12}{0}{0} & \SetCell{bg=\HeatCellBg{12}}\CellText{0.12}{0}{0} & \SetCell{bg=\HeatCellBg{12}}\CellText{0.12}{0}{0} & \SetCell{bg=\HeatCellBg{10}}\CellText{0.10}{0}{0} & \SetCell{bg=\HeatCellBg{11}}\CellText{0.11}{0}{0} & \SetCell{bg=\HeatCellBg{9}}\CellText{0.09}{0}{0} & \SetCell{bg=\HeatCellBg{11}}\CellText{0.11}{0}{0} & \SetCell{bg=\HeatCellBg{11}}\CellText{0.11}{0}{0} & \SetCell{bg=\HeatCellBg{11}}\CellText{0.11}{0}{0} & \SetCell{bg=\HeatCellBg{11}}\CellText{0.11}{0}{0} & \SetCell{bg=\HeatCellBg{10}}\CellText{0.10}{0}{0} & \SetCell{bg=\HeatCellBg{10}}\CellText{0.10}{0}{0} & \MeanCell{0.10}{0.03}{0} \\
        7b & \SetCell{bg=\HeatCellBg{11}}\CellText{0.11}{0}{0} & \SetCell{bg=\HeatCellBg{0}}\CellText{-0.00}{0}{0} & \SetCell{bg=\HeatCellBg{16}}\CellText{0.16}{0}{0} & \SetCell{bg=\HeatCellBg{16}}\CellText{0.16}{0}{0} & \SetCell{bg=\HeatCellBg{16}}\CellText{0.16}{0}{0} & \SetCell{bg=\HeatCellBg{16}}\CellText{0.16}{0}{0} & \SetCell{bg=\HeatCellBg{17}}\CellText{0.17}{0}{0} & \SetCell{bg=\HeatCellBg{20}}\CellText{0.20}{0}{0} & \SetCell{bg=\HeatCellBg{21}}\CellText{0.21}{0}{1} & \SetCell{bg=\HeatCellBg{18}}\CellText{0.18}{0}{0} & \SetCell{bg=\HeatCellBg{17}}\CellText{0.17}{0}{0} & \SetCell{bg=\HeatCellBg{17}}\CellText{0.17}{0}{0} & \SetCell{bg=\HeatCellBg{17}}\CellText{0.17}{0}{0} & \SetCell{bg=\HeatCellBg{17}}\CellText{0.17}{0}{0} & \SetCell{bg=\HeatCellBg{17}}\CellText{0.17}{0}{0} & \SetCell{bg=\HeatCellBg{16}}\CellText{0.16}{0}{0} & \SetCell{bg=\HeatCellBg{16}}\CellText{0.16}{0}{0} & \SetCell{bg=\HeatCellBg{16}}\CellText{0.16}{0}{0} & \SetCell{bg=\HeatCellBg{14}}\CellText{0.14}{0}{0} & \MeanCell{0.16}{0.04}{0} \\
        8b & \SetCell{bg=\HeatCellBg{15}}\CellText{0.15}{1}{0} & \SetCell{bg=\HeatCellBg{0}}\CellText{-0.00}{0}{0} & \SetCell{bg=\HeatCellBg{24}}\CellText{0.24}{1}{0} & \SetCell{bg=\HeatCellBg{28}}\CellText{0.28}{1}{0} & \SetCell{bg=\HeatCellBg{26}}\CellText{0.26}{1}{0} & \SetCell{bg=\HeatCellBg{26}}\CellText{0.26}{1}{0} & \SetCell{bg=\HeatCellBg{28}}\CellText{0.28}{1}{0} & \SetCell{bg=\HeatCellBg{31}}\CellText{0.31}{1}{0} & \SetCell{bg=\HeatCellBg{32}}\CellText{0.32}{1}{1} & \SetCell{bg=\HeatCellBg{29}}\CellText{0.29}{1}{0} & \SetCell{bg=\HeatCellBg{26}}\CellText{0.26}{1}{0} & \SetCell{bg=\HeatCellBg{26}}\CellText{0.26}{1}{0} & \SetCell{bg=\HeatCellBg{27}}\CellText{0.27}{1}{0} & \SetCell{bg=\HeatCellBg{26}}\CellText{0.26}{1}{0} & \SetCell{bg=\HeatCellBg{27}}\CellText{0.27}{1}{0} & \SetCell{bg=\HeatCellBg{25}}\CellText{0.25}{1}{0} & \SetCell{bg=\HeatCellBg{23}}\CellText{0.23}{1}{0} & \SetCell{bg=\HeatCellBg{23}}\CellText{0.23}{1}{0} & \SetCell{bg=\HeatCellBg{23}}\CellText{0.23}{1}{0} & \MeanCell{0.24}{0.07}{1} \\
        Mean $\pm$ Std & \MeanCell{0.09}{0.05}{0} & \MeanCell{-0.01}{0.02}{0} & \MeanCell{0.05}{0.09}{0} & \MeanCell{0.09}{0.09}{0} & \MeanCell{0.10}{0.07}{0} & \MeanCell{0.10}{0.07}{0} & \MeanCell{0.11}{0.08}{0} & \MeanCell{0.12}{0.09}{0} & \MeanCell{0.12}{0.09}{1} & \MeanCell{0.11}{0.08}{0} & \MeanCell{0.10}{0.07}{0} & \MeanCell{0.10}{0.08}{0} & \MeanCell{0.10}{0.08}{0} & \MeanCell{0.10}{0.07}{0} & \MeanCell{0.10}{0.08}{0} & \MeanCell{0.10}{0.07}{0} & \MeanCell{0.10}{0.06}{0} & \MeanCell{0.09}{0.07}{0} & \MeanCell{0.09}{0.06}{0} & \MeanCell{0.09}{0.08}{0} \\
}


\end{landscape}
